\documentclass[a4paper,10pt,oneside,onecolumn,1p,times]{elsarticle}

%\documentclass[letter,10pt,oneside,onecolumn,1p,times]{elsarticle}

% ---- 中文排版支持（xelatex 编译）----
\usepackage{xeCJK}
\IfFontExistsTF{SimSun}{%
  \setCJKmainfont[BoldFont=SimHei, ItalicFont=KaiTi]{SimSun}%
  \setCJKsansfont{SimHei}%
  \setCJKmonofont{FangSong}%
}{%
  % 本机无 SimSun：XeTeX 按家族名会命中 Noto TTC 的第 0 面（jp），
  % 故用 fontTools 从 TTC 抽出 SC 面存于 ~/.fonts，并按文件路径加载
  \setCJKmainfont[Path=/home/w/.fonts/, BoldFont={NotoSerifCJKsc-Bold.otf}]{NotoSerifCJKsc-Regular.otf}%
  \setCJKsansfont[Path=/home/w/.fonts/, BoldFont={NotoSansCJKsc-Bold.otf}]{NotoSansCJKsc-Regular.otf}%
  \setCJKmonofont[Path=/home/w/.fonts/]{NotoSansMonoCJKsc-Regular.otf}%
}
\renewcommand\figurename{图}
\renewcommand\tablename{表}
\renewcommand\refname{参考文献}
\renewcommand\contentsname{目\quad 录}
\abstracttitle{摘要}
% ------------------------------------

\newcommand{\mc}{\mathcal}
\newcommand{\ham}{\hat{H}}                          % Hamiltonian 
\newcommand{\Tr}{{\rm Tr}}
\newcommand{\tr}{{\rm Tr }}
\newcommand{\diff}{{\rm d}}
\newcommand{\eul}{{\rm e }}
\newcommand{\imag}{{\rm i }}
\newcommand{\obs}{{\rm obs}}
\newcommand{\fit}{{\rm fit}}
\newcommand{\bra}[1]{\mbox{$\langle #1  |$}} 
\newcommand{\ket}[1]{\mbox{$| #1  \rangle$}} 
\newcommand{\identity}{I}
\newcommand{\braket}[2]{\mbox{$\langle #1  | #2 \rangle$}} 	% <x|y>
\newcommand{\fat}[1]{\mbox{\boldmath$#1$\unboldmath}}	% fat also for Greek
\newcommand{\syslength}{L}
\newcommand{\Eq}[1]{式~(\ref{#1})}



\begin{document}
\title{矩阵乘积态时代的密度矩阵重正化群\\ The density-matrix renormalization group in the age of matrix product states}
\author{Ulrich Schollw\"{o}ck}
\ead{schollwoeck@lmu.de}
\address{Department of Physics, Arnold Sommerfeld Center for Theoretical Physics and Center for NanoScience, University of Munich, Theresienstrasse 37, 80333 Munich, Germany \\
Institute for Advanced Study Berlin, Wallotstrasse 19, 14159 Berlin, Germany}

\begin{abstract}
密度矩阵重正化群方法（DMRG）在过去十年中已确立了其作为模拟一维强关联量子格点系统静态与动力学的首要方法的地位。在该方法的进一步发展中，人们认识到 DMRG 所处理的是一类极有研究价值的量子态，即所谓矩阵乘积态（MPS），这使得我们对 DMRG 方法的内在结构、进一步潜力及其局限有了深入得多的理解。在本文中，我想用 MPS 语言对当前的 DMRG 思想作详细阐述，以便使 DMRG 算法族的推荐实现方式能够完全用 MPS 术语透明地表达出来。随后我将转而讨论一些可能富有成果的进一步算法发展方向：尽管 DMRG 如今已是非常成熟的方法，我仍然看到进一步改进的潜力，若干新近提出的算法即是明证。   
\end{abstract}

\maketitle 

\tableofcontents

\section{引言}
低维格点上的强关联量子系统至今仍在向现代量子多体物理提出一些最引人入胜的挑战。在凝聚态物理中，关联效应在量子自旋链与自旋梯子、一维与二维空间中的阻挫磁体以及高温超导体等众多令人关注的物理系统中占据主导地位。而近年来，光晶格中高度可控、可调谐的强相互作用超冷原子气体的出现，为这一领域开辟了一个全新的方向\cite{Bloch08}。 

无论从解析还是数值的角度，这些系统都难以研究：只有在极少数情形下才能得到精确的解析解，例如一维情形下借助 Bethe 拟设的求解\cite{Bethe31,Lieb68}。在强相互作用存在时，微扰理论会失效。其他方法，如场论方法，曾给出深刻的洞见，例如关于整数自旋反铁磁链的 Haldane 能隙物理\cite{Haldane83}，但它们要作可能相当严重的近似，这些近似最终必须靠数值方法加以检验控制。这类数值方法包括精确对角化、量子蒙特卡罗、级数展开或耦合簇方法。

自 1992 年 Steve White 发明以来\cite{White92,White93}，密度矩阵重正化群（DMRG）已牢固确立了其作为当前
研究一维量子格点的最强大数值方法的地位\cite{Schollwoeck05,Hallberg06}。在对诸如海森堡、$t$-$J$ 与 Hubbard 等强关联哈密顿量的低激发本征态（特别是基态）的静态性质（能量、序参量、$n$ 点关联函数）作了初步研究之后，该方法又被推广到研究
本征态的动力学性质，如动力学结构函数或频率依赖的电导率\cite{Hallberg95,Ramasesha97,Kuhner99,Jeckelmann02}。与此同时，将其推广用于分析二维经典 \cite{Nishino95} 与一维量子 \cite{Wang97,Shibata97} 转移矩阵，使人们得以获得经典二维系统与量子一维系统的高精度有限温度信息；更近一些时候，DMRG 的转移矩阵变体也被推广到有限温度下的动力学\cite{Sirker05}。它甚至被推广到数值上要求高得多的非厄米（赝）哈密顿量的研究，这类哈密顿量出现在对一维系统远离平衡时向经典稳态弛豫的分析中\cite{Hieida98,Carlon99,Carlon01,Henkel01}。 

在 DMRG 的许多应用中，结果的精度基本上只受机器精度的限制，即便使用的数值资源相当有限也是如此，而且与哈密顿量的具体性质几乎无关。因此毫不奇怪，除了把算法推广到越来越多的问题类别之外，人们也想知道 DMRG 优异表现的物理根源，以及这一成功故事能否扩展到实时动力学或二维系统的研究。 

事实上，这两个问题紧密相关：人们很快就意识到，DMRG 应用于二维格点时只能算中等成功：虽然相对较小的系统可以高精度地加以研究\cite{White96,White98a,White98b,White00,White03,Chernyshev03,Chernyshev05,White07}，但所需的数值资源量基本上随系统尺寸指数增长，使大的格点系统无法企及。事后来看，DMRG 在一维与二维中截然不同的行为，与多体态中量子纠缠在一维和二维的不同标度密切相关\cite{Vidal03,Latorre04}，后者由所谓面积定律所支配（近期综述见 \cite{Eisert10}）。 

在本文中，我将局限于一维物理的框架之内；虽然把 DMRG 推广到更高维度会自然地约化为一维 DMRG，但由此涌现的结构远比一维丰富，超出了本文的范围。  
  
在一项原本不相干的发展中，所谓{\em 矩阵乘积态}（MPS）作为适合解析研究的一类有趣量子态被发现了。事实上，这一结构如此简单却又如此有力，因此在过去五十多年里它们以各种各样新的名称被反复引入和使用（最著名的或许是 Baxter \cite{Baxter68}）也就不足为奇了。就本文的语境而言，最相关的史前事例可以说是标志性的一维 AKLT 态在此形式下的精确表达式\cite{Affleck87,Affleck88,Fannes89}，它引发了对 MPS 中平移不变子类——即所谓有限关联态\cite{Fannes92}——的大量研究。此后这一形式在各种语境中被用于解析变分计算，例如自旋-1 海森堡反铁磁体\cite{Klumper93,Kolezhuk96a,Kolezhuk96b,Kolezhuk98}与亚铁磁体\cite{Kolezhuk97,Kolezhuk99}。

MPS 与 DMRG 的联系分两步建立起来。第一步，Ostlund 与 Rommer \cite{Ostlund95} 意识到，所谓无限系统 DMRG 的块增长步骤可以按 MPS 中矩阵所取的形式表达为一个矩阵。由于在均匀系统中这一块增长步骤在热力学极限下导向一个不动点，他们便取该不动点矩阵作为平移不变 MPS 的构建单元。更进一步，人们认识到更重要的有限系统 DMRG 给出 MPS 形式的量子态，并在其上作变分优化\cite{Dukelsky98}。人们还认识到，在传统 DMRG 中，作变分优化的态类会随算法的推进而改变，因此如果要求某种意义上``完美''的变分性质，就需要对算法作一点小改动，然而发现这会加剧（可解决的）亚稳性问题\cite{Takasaki99,White05}。 

一个耐人寻味的历史事实是，直到 2004 年前后 Cirac、Verstraete、Vidal 及其合作者开始非常系统地发掘 MPS 的威力之前，只有少数 DMRG 实践者认真对待这一进展。虽然人们认为它对概念性目的有用，但令人惊讶的是，很少有人去思考如何纯粹用 MPS 语言重新思考和重新表述实际使用的 DMRG 实现；可以说，这是因为绝大多数传统 DMRG 应用（即具有开边界条件的量子链的基态问题）几乎从中无利可图。而被忽视的是：它轻而易举地开辟了通往 DMRG 强大扩展的道路，这些扩展在传统 DMRG 语言中既难以看出也难以表述。 

一个非穷尽的扩展清单可以列出：实时演化\cite{Vidal03a,Vidal04b,Daley04,WhiteFeiguin04,VerstraeteRipoll04,Prosen09,Hartmann09,Banuls09}（也包括有限温度下的\cite{VerstraeteRipoll04,Feiguin05a}）、周期边界条件的高效使用\cite{VerstraetePorras04,Pippan10,Pirvu10}、可靠的单格点 DMRG\cite{White05}、数值重正化群（NRG）应用\cite{Weichselbaum09}、无限系统算法\cite{Vidal07,Orus08,McCulloch08}、连续系统\cite{Verstraete10}，更不必说从 \cite{VerstraeteCirac04} 开始、利用 MPS 态类的推广\cite{Nishino01} 而在更高维度取得的进展了。 

本文的目标不可能是提供一部从 2010 年视角回顾 1992 年以来 DMRG 的完整综述，尤其是鉴于已有综述\cite{Schollwoeck05}试图对截至 2004 年夏的 DMRG 作出相当全面的论述。我更愿意把自己限定在较新的进展上，并完全用矩阵乘积态的语言来表述它们，重点放在方法的具体构造上，而不是展示大量应用。我希望这篇综述能让本领域的新手迅速写出自己的代码，并牢牢掌握 MPS 算法的关键构件。它与 \cite{VerstraeteMurg08,McCulloch07} 的引论部分在所介绍的关键方法上有重叠，但侧重于不同的扩展（其中一些产生于这些文章之后），并且在许多地方力求更加显明。它采取的观点不同于 \cite{Peschel98}——那是 1998 年对 DMRG 的首次全面阐述，因为在那时与 MPS 的联系（虽然已知）尤其是与量子信息的联系在许多方面尚未被开发利用，而后者正是本文的焦点。尽管如此，作为第一个``历史性''步骤，我想提醒读者回顾引入 DMRG 的原始方式，它并不借助矩阵乘积态的思想。这应当能让读者更容易进入较早的文献，不过如果对此不感兴趣，也可以直接跳到第 \ref{sec:matrixproductstates} 节。 
 
第二步，我将证明任何量子态都可以精确地写成一种非常特定的形式，也就是前面已经提到的矩阵乘积态所给出的形式。事实上，之所以限于一维，是因为只有在这种情形下 MPS 才在数值上可处理。我将着重指出 MPS 的特殊规范形式，并建立它们与作为数学工具的奇异值分解（SVD）以及作为量子态紧凑表示的 Schmidt 分解之间的联系。在此之后，我将解释 MPS 如何成为一维状态抽取方案（如 DMRG 与 Wilson 的 NRG 中出现的那类方案）的自然框架。作为一个简单而非平凡的例子，我将显式讨论其 MPS 形式下的 AKLT 态。
然后我们转向显式讨论对 MPS 的各种操作：重叠、归一化、算符矩阵元、期望值以及 MPS 相加。这些是对任何量子态都会做的操作；更具 MPS 特色的是把它们化为便于计算的规范形式的方法，以及用一个维数更小的 MPS 来逼近另一个 MPS 的方法。在结束这篇关于 MPS 的阐述时，我将讨论我所偏爱的 MPS 记法、Vidal 提出的另一种记法以及 DMRG 书写态的方式之间的关系与相互转换；这一相对技术性的章节应有助于读者更顺利地阅读文献。

MPS 的思想可以从态推广到算符的表示，因此我接着讨论矩阵乘积算符（MPO）的使用\cite{VerstraeteRipoll04,McCulloch07,VerstraetePirvu08,Crosswhite08,Frowis10}。据我所见，所有我们感兴趣的算符（从局域算符、键演化算符一直到完整哈密顿量）都能用 MPO 得到非常紧凑而透明的表述。这使 DMRG 相关算法的表述大为简洁，有时甚至在数值上更为精确，但它们的使用尚未广泛传播。

不可否认，读到这里，读者的耐心可能已经被消磨了不少，因为迄今为止还没有一个现实可用的算法（例如基态搜索或时间演化）用 MPS 语言表述出来；但将会显而易见的是，DMRG 算法中大量繁琐的数值细节已经被干净利落地隐藏在 MPS 与 MPO 结构之中。 

我将讨论基态算法，讨论单中心格点或双中心格点的 DMRG 与完全基于 MPS 的算法之间的等价性与差异，包括避免陷入亚稳态的改进。我将聚焦于开边界条件的有限系统，这些方法在此表现卓越。

在此之后，我转向动力学的时间依赖方法，涵盖纯态与混合态。在讨论了基本算法及其细微差别之后，我将聚焦于延长此类模拟时间范围这一关键问题：能够对复杂量子多体态计算高精度的实时与虚时演化，令许多人尤为振奋，部分也因为它来得正是时候，恰逢高度可调谐的超冷原子系统的研究兴起。虽然这一发展已经带来了大量有趣的洞见与应用，但人们很快意识到，含时 DMRG 及相关方法的时间范围受制于非平衡量子态中纠缠的增长，使得长时间物理无法企及。在这一背景下，近期在设法超越这一限制方面已经取得了有趣的进展。

本综述以两条进一步的发展轴线作结。我将首先讨论 DMRG 与 Wilson 的 NRG 之间的联系，展示如何以非常简洁的方式表述 NRG，以及如何在多个方向上对它加以改进。这闭合了一个有趣的历史循环，因为 NRG 在均匀一维量子格点上的彻底失败——与之相对的是映射到特殊非均匀一维量子格点上的量子杂质模型——正是 White 发明 DMRG 的出发点\cite{WhiteNoack92}。

接下来我考察直接在热力学极限下工作的基于 MPS 的无限尺寸算法，其中一种基于时间演化（iTEBD）\cite{Vidal07}。另一种（iDMRG）\cite{McCulloch08}是无限系统 DMRG 算法的扩展，它有一段有趣的历史：在许多关于 DMRG 的早期论述中，它被当作 DMRG 的关键所在，而有限系统 DMRG 只被当作实践者为进一步提高数值精度而附加的手段。后来人们认识到，在许多情形下即使只求定性正确，也必须运用有限系统 DMRG，而无限系统 DMRG 只被视为一种热身步骤。直到最近，McCulloch\cite{McCulloch08} 才指出一种方法，把无限系统 DMRG 变成生成均匀系统热力学极限态的高效工具。 

最后但同样重要的是，我将对本文未能涵盖的 MPS 进一步应用作一展望。

\section{密度矩阵重正化群（DMRG）}
\subsection{无限系统与有限系统算法}
作为玩具模型，让我们考虑一维空间中长度为 $L$、外磁场为 $h$ 的（各向异性）$S=\frac{1}{2}$ 海森堡反铁磁（$J=1$）自旋链， 
\begin{equation}
\hat{H} = \sum_{i=1}^{L-1} \frac{J}{2} (\hat{S}^+_{i}  \hat{S}^-_{i+1} + \hat{S}^-_{i} \hat{S}^+_{i+1}) + J^z \hat{S}^z_{i} \hat{S}^z_{i+1} - \sum_{i=1}^L h \hat{S}^z_i .
\label{eq:basicHam}
\end{equation}
我们考虑{\em 开边界条件}（图~\ref{fig:toymodel}），这一点非常值得强调：从解析的角度看，周期边界条件通常更方便；许多数值方法其实并不强烈依赖边界条件，有些方法（如精确对角化）在周期边界条件下甚至更为高效。而 DMRG 则偏爱开边界条件。

\begin{figure}
\centering\includegraphics[width=0.8\textwidth]{PyMPS_Toymodel.eps}
\caption{我们的玩具模型：一条长度为 $L$、两端开放的自旋链，每个格点上有一个自旋-$\frac{1}{2}$，并与最近邻相互作用。}
\label{fig:toymodel}
\end{figure}

还应强调的是，DMRG 作为在某类态上作变分的方法，并不受诸如费米子符号问题之类的困扰，可以同样应用于玻色系统与费米系统。 

DMRG 的出发点是求 $\hat{H}$ 的基态与基态能量。我们可以对热力学极限 $L\rightarrow\infty$ 提出这个问题，也可以更谦逊一些对有限 $L$ 提出。在前一种情形，答案由{\em 无限系统} DMRG 给出，尽管精度相当有限；在后一种情形，答案虽然可以从无限系统 DMRG 读出，但更合适的做法是运行一个两步流程，先运行无限系统 DMRG，再继续运行{\em 有限系统} DMRG。

无论如何，数值上的拦路虎来自希尔伯特空间维数的指数增长，在我们的例子中即 $d^L$，其中 $d=2$ 是自旋-$\frac{1}{2}$ 的局域态空间维数。 

\subsection{无限系统 DMRG}

无限系统 DMRG 处理这一问题的办法是考虑长度渐增的链，通常为 $L=2$、$4$、$6$、$\ldots$，并舍弃足够多的态以使希尔伯特空间的大小保持在可处理的范围内。这一{\em 状态抽取}步骤是算法成功的关键：我们假设存在一个能够描述相关物理的约化态空间，并且我们能够设计出一个把它识别出来的流程。第一个假设是一切变分方法的典型假设，而我们将会看到，在一维情形我们确实很幸运：一维中所有短程哈密顿量都存在这样一个包含相关物理的约化态空间！

如何找到它？在无限系统 DMRG 中（图~\ref{fig:DMRGCombined}），构建过程如下：我们引入左、右两个{\em 块} A 与 B，第一步中它们可以各由一个自旋（或格点）组成，使总链长为 2。更长的链随后从左右两端迭代地构建，在两块之间插入成对的自旋，使链长依次增长到 4、6 等等；每一步中，之前的自旋被并入左右两块，使块尺寸按 1、2、3 等等增长，导致整块态空间维数按 $2^\ell$ 指数增长，其中 $\ell$ 是当前块尺寸。于是我们的链始终具有块-格点-格点-块的结构，即 A$\bullet\bullet$B。

让我们假定我们的数值资源足以处理维数为 $D$ 的约化块态空间（实践中 $D$ 大约取 $O(100)$ 到 $O(1000)$ 之间的值），并且对长度为 $\ell$ 的块 A，我们在一个 $D$ 维约化希尔伯特空间中拥有块 A 的有效描述，其正交归一基为 $\{ \ket{a_\ell}_A \}$。对于非常小的块，基的维数可能小于 $D$；对较大的块，则必定已发生过某种截断，这一点我们暂且视作理所当然。这里不妨立即提到，在文献中，$D$——它将被证明是一个关键数字——有着各种各样新的名称：在传统 DMRG 文献中，它通常被记作 $m$（或 $M$）；在较新的矩阵乘积态文献中，则 $D$ 与 $\chi$ 两个记号都在使用。

在这一框架内，我们现在可以问：(i) 当前长度为 $2\ell +2$ 的链（也称为{\em 超块}）的基态是什么；(ii) 如何为新块 A$\bullet$ 与 $\bullet$B 找到维数为 $D$ 的约化希尔伯特空间。 

超块 A$\bullet\bullet$B 的任意态都可以描述为
\begin{equation}
\ket{\psi}= \sum_{a_A \sigma_A \sigma_B a_B} \psi_{a_A \sigma_A \sigma_B a_B} \ket{a}_A \ket{\sigma}_A \ket{\sigma}_B \ket{a}_B \equiv \sum_{i_A j_B} \psi_{i_A j_B} \ket{i}_A \ket{j}_B,
\label{US:eq:DMRGstate}
\end{equation}
其中紧邻 A 的格点的态位于局域态空间维数为 $d$ 的集合 $\{ \ket{\sigma}_A \}$ 中，紧邻 B 的格点的态亦然。通过数值对角化，我们求出使能量
\begin{equation}
E = \frac{\bra{\psi} \hat{H}_{{\rm A}\bullet\bullet{\rm B}} \ket{\psi}}{\braket{\psi}{\psi}}
\end{equation}
相对于超块的哈密顿量取极小的 $\ket{\psi}$，这就回答了 (i)。为此，我们需要某种迭代法稀疏矩阵本征值求解器，例如 Lanczos 方法或 Jacobi-Davidson 方法所提供的那些。鉴于我们的哈密顿量 [式~(\ref{eq:basicHam})] 是在全张量积空间中给出的，这一极小化当然预设了它能够方便地在超块基中表达出来。这个问题我们先搁置一下，因为它与块的生长及其态空间的状态抽取紧密相关。由于矩阵维数为 $d^2 D^2$，对典型的 $D$ 而言，即使假定利用了量子对称性，该本征值问题也大得无法直接求解。短程哈密顿量的大多数矩阵元为零，矩阵因而是稀疏的，这样迭代法稀疏矩阵本征值求解器所需的基本矩阵—矢量乘法就可以非常高效地实现。我们在下文还会进一步讨论这个问题。

如果我们现在取态 $\{ \ket{i}_A \}$ 作为下一个更大的左块 A$\bullet$ 的基，基维数就增长到 $dD$。为避免指数增长，我们按如下步骤把这个基截断回 $D$ 个态（对块 A$\bullet$ 这样做，对 $\bullet$B 同理）：我们考虑 A$\bullet$ 的 {\em 约化密度算符}，即
\begin{equation}
\hat{\rho}_{A\bullet} = \Tr_{\bullet B} \ket{\psi}\bra{\psi} 	\quad\quad (\rho_{A\bullet})_{ii'} = \sum_j \psi_{ij}\psi^*_{i'j} .
\end{equation}
$\hat{\rho}_{A\bullet}$ 的本征系统由精确对角化确定；做法是保留关联本征值 $w$ 最大的那 $D$ 个正交归一本征态作为约化基。若把它们记为 $\ket{b}_A$，则矢量的分量不过是在先前的块—格点基中的展开系数 $_A \braket{a \sigma}{b}_A$。

在把 A$\bullet$ 上所有需要的算符（近似地）变换到新基之后，系统尺寸就可以再次增大，直到达到想要的最终尺寸。B 同时生长，对反射对称的系统只需简单地镜像即可。

截断步骤的动机有二。第一个理由基础较弱：我们关心的是 A$\bullet$ 中对基态贡献最大的那些态，这里的基态是指 A$\bullet$ 嵌入在最终的、大得多的乃至无限大的系统中时的基态。我们以目前所知，用 A$\bullet\bullet$B——即我们能高效构造的最大超块——的基态来近似这个尚不知晓的最终系统基态。在这个意义上，我们是在做自举（bootstrap），而正如我们将看到的，有限系统 DMRG 会处理由此可能带来的大近似。
这里我要立即指出，在无限系统 DMRG 的 MPS 表述中，一幅清晰得多的图景将会浮现。选取态的第二个动机则根基牢固得多：上述方案既可以由统计物理的论证得出，也可以由如下要求导出：当前基态 $\ket{\psi}$ 与其到维数为 $D$ 的截断块基上的投影 $\ket{\psi}_{{\rm trunc}}$ 之间的 2-范数距离 $\| \ket{\psi} - \ket{\psi}_{{\rm trunc}} \|_2$ 为最小。最透明的证明来自对矩阵 $\Psi$ 的奇异值分解，其矩阵元 $\Psi_{(a_A \sigma_A), (\sigma_B a_B)}$ 由波函数系数 $\psi_{a_A \sigma_A\sigma_B a_B}$ 构成。DMRG 的整体成功建立在如下观察之上：即使对于中等的 $D$（常常只有几百），被截断的本征值的总权重——若假定 $w_a$ 按降序排列，则由 {\em 截断误差} $\epsilon=1-\sum_{a>D} w_a$ 给出——也极其接近 0，比如说 $10^{-10}$ 或更小。

在算法上，我们当然也可以固定一个我们愿意接受的（小的）$\epsilon$，并在每次截断时选取足够大的 $D$ 来满足这一要求。这样计算资源的利用更为高效（典型地，链两端附近量子涨落减弱，那里的 $D$ 会稍小一些），当然，为实现这种灵活性所需的编程工作量也要更高一些。

提升任何量子算法性能的一个重要方面是利用量子对称性：从离散对称性（如当哈密顿量在从左到右的镜像变换下不变时的 $Z_2$ 镜像对称），到阿贝尔对称性（如粒子数（``电荷''）守恒或磁化强度守恒的 $U(1)$ 对称性），再到非阿贝尔对称性（如旋转不变性的 {\em SU}(2) 对称性）\cite{McCulloch07,SierraNishino97,McCulloch00,McCulloch01,McCulloch02}。大量这类对称性已在众多应用中得到成功使用，最常见的是电荷守恒与磁化强度守恒这两个 $U(1)$ 对称性，但也有 {\em SU}(2) 对称性\cite{McCulloch08a,Smerat09,Langer09,Holzner09}。让我们聚焦于磁化强度，并假定总磁化强度 $\hat{M} = \sum_i \hat{S}^z_i$ 与哈密顿量对易，$[\hat{H},\hat{M}]=0$，于是 $\hat{H}$ 的本征态可以选成 $\hat{M}$ 的本征态。特别地，我们假定基态的好量子数为 $M=0$。若块态与格点态都是磁化强度的本征态，则仅当 $M (\ket{a}_A) +  M (\ket{\sigma}_A) + M (\ket{\sigma}_B) + M (\ket{a}_B) = 0$ 时才有 $\psi_{a_A \sigma_A \sigma_B a_B} \neq 0$；
这里 $M$ 是相应块和格点磁化强度的简写，前提是这些态以磁化强度为好量子数。这一约束使我们能把大量系数从计算中排除，从而带来计算上的巨大加速以及（较次要的）内存节省。

决定性的一点是：如果态 $\ket{a}_A$ 和 $\ket{\sigma}_A$ 是磁化强度的本征态，那么约化密度算符 $\hat{\rho}_{A\bullet}$ 的本征态也将是，而这些本征态正是增大后的块的态 $\ket{a}_A$。由于局域格点态可以选成这样的本征态，且生长过程中的第一个块仅由一个局域格点组成，本征态性质便由此贯穿整个算法。为证明这一论断，我们考虑 $(\rho_{A\bullet})_{ii'} = \sum_j \psi_{ij} \psi^*_{i'j}$。态 $\ket{i}_A$ 和 $\ket{j}_B$ 按构造就是磁化强度的本征态，故
$M (\ket{i}_A) + M(\ket{j}_B) = 0 = M (\ket{i'}_A) + M (\ket{j}_B)$，即 $M (\ket{i}_A) = M (\ket{i'}_A)$。因此密度矩阵分解为若干块，每块由磁化强度相同的态构成，可以在这些集合内逐块对角化，于是其本征态也是磁化强度的本征态。

在实际实现中，使用这类好量子数的做法是：把算符的矩阵表示排列成以好量子数标记的块结构，再将它们组合起来以满足局域与全局的约束。这在概念上容易，但编码会更为复杂，此处不打算具体展开；在 MPS 各节中会给出提示。

到此为止，我们一直搁置着如何把作用在块上的算符表达在当前块基中的问题，而这是构造哈密顿量和所关心的可观测量所必需的。考虑作用在格点 $\ell$ 上的算符 $\hat{O}$，其矩阵元为 $O^{\sigma_\ell,\sigma'_\ell}=\bra{\sigma_\ell} \hat{O}_{\ell} \ket{\sigma'_\ell}$。不失一般性，设格点 $\ell$ 在左侧并被并入块 A。那么当块 A 从 $\ell-1 \rightarrow \ell$ 生长时，它的构造初始化为
\begin{equation}
\bra{a_\ell} \hat{O} \ket{a'_\ell} = \sum_{a_{\ell-1},\sigma_\ell,\sigma'_\ell}
\braket{a_\ell}{a_{\ell-1}\sigma_\ell} \bra{\sigma_\ell} \hat{O} \ket{\sigma'_\ell} \braket{a_{\ell-1}\sigma'_\ell}{a'_\ell} .
\end{equation}
这里 $\ket{a_\ell}$ 和 $\ket{a_{\ell-1}}$ 分别是长度为 $\ell$ 和 $\ell-1$ 的块的有效基态。当然，在块 A 进一步的生长步骤中还需要更新，例如
\begin{equation}
\bra{a_{\ell+1}} \hat{O} \ket{a'_{\ell+1}} = \sum_{a_{\ell},a'_{\ell}\sigma_{\ell+1}}
\braket{a_{\ell+1}}{a_{\ell}\sigma_{\ell+1}} \bra{a_\ell} \hat{O} \ket{a'_\ell} \braket{a'_{\ell}\sigma_{\ell+1}}{a'_{\ell+1}} .
\label{eq:update}
\end{equation}
重要的是要意识到，式~(\ref{eq:update}) 中的求和必须拆分为
\begin{equation}
\sum_{a_{\ell}\sigma_{\ell+1}} \braket{a_{\ell+1}}{a_{\ell}\sigma_{\ell+1}} \left( \sum_{a'_{\ell}}
 \bra{a_\ell} \hat{O} \ket{a'_\ell} \braket{a'_{\ell}\sigma_{\ell+1}}{a'_{\ell+1}} \right),
\end{equation}
从而把计算量从 $O(D^4 d)$ 降到 $2O(D^3 d)$。

在哈密顿量中会出现算符乘积 $\hat{O}\hat{P}$。把量子力学的恒等式插入当前块基并写成下式虽然诱人，但由于截断步骤众多而极不精确：
\begin{equation}
\bra{a_\ell} \hat{O}\hat{P} \ket{a'_\ell} = \sum_{\tilde{a}_\ell} \bra{a_\ell} \hat{O} \ket{\tilde{a}_\ell} \bra{\tilde{a}_\ell} \hat{P} \ket{a'_\ell} . \quad\quad {\rm NO!}
\end{equation}
正确的做法是：不断更新第一个算符（对左块 A 从左往右数），直到到达第二个算符所在的格点，再把它并入如下：
\begin{equation}
\bra{a_\ell} \hat{O} \hat{P} \ket{a'_\ell} = \sum_{a_{\ell-1},a'_{\ell-1},\sigma_\ell,\sigma'_\ell}
\braket{a_\ell}{a_{\ell-1}\sigma_\ell}  \bra{a_{\ell-1}} \hat{O} \ket{a'_{\ell-1}} \bra{\sigma_\ell} \hat{P} \ket{\sigma'_\ell} \braket{a'_{\ell-1}\sigma'_\ell}{a'_\ell} .
\label{eq:incorporate}
\end{equation}
其后的更新与单个算符的情形相同。显然，我们看到的是一串直截了当的（约化）基变换，只是记号略显繁琐；但这些公式在 MPS 语言中会显得简单多少，将会是一件有趣的事，尽管内容完全相同。



期望值的最终计算在生长过程结束时由下式给出
\begin{equation}
\bra{\psi} \hat{O} \ket{\psi} = \sum_{a_A a'_A \sigma_A \sigma_B a_B} \braket{\psi}{a_A \sigma_A \sigma_B a_B} \bra{a_A} \hat{O} \ket{a'_A} \braket{a'_A \sigma_A \sigma_B a_B}{\psi},
\end{equation}
其中通过适当的加括号方式，这成为 $O(D^3 d^2)$ 阶的运算。一个重要的特例是：当我们要计算作用在最终的块—格点构型中某个自由格点 $\bullet$ 上的局域算符的期望值时。此时
\begin{equation}
\bra{\psi} \hat{O} \ket{\psi} = \sum_{a_A \sigma_A \sigma'_A \sigma_B a_B} \braket{\psi}{a_A \sigma_A \sigma_B a_B} \bra{\sigma_A} \hat{O} \ket{\sigma'_A} \braket{a_A \sigma'_A \sigma_B a_B}{\psi},
\end{equation}
这是 $O(D^2 d^3)$ 阶的表达式（对许多算符而言甚至只有 $O(D^2 d^2)$）。鉴于在实际应用中总有 $D \gg d$，这类计算在计算上是极为有利的。

在结束这些较为技术性的讨论时我要指出：要从这类算符高效地构造 $\hat{H}$，关键在于 {\em 绝不}把它造成一个完整矩阵，而是利用其特定的块—格点—格点—块结构。例如，设某项含一个作用在 A 上的算符和一个作用在 B 上的算符（这实际上是最复杂的情形），$\hat{h}=\hat{O}_A \hat{O}_B$。那么
\begin{equation}
\bra{a_A \sigma_A \sigma_B a_B} \hat{h}\ket{\psi} = \sum_{a'_A}
\bra{a_A}\hat{O}_A\ket{a'_A} \left( \sum_{a'_B} \bra{a_B}\hat{O}_B\ket{a'_B}  \braket{a'_A \sigma_A \sigma_B a'_B}{\psi} \right),
\label{eq:hamfragment}
\end{equation}
这是两次 $O(D^3 d^2)$ 乘法的序列（而非一次朴素的 $O( D^4 d^2)$ 计算），对所有系数 $\bra{a_A \sigma_A \sigma_B a_B} \hat{h}\ket{\psi}$ 皆如此，其他各项也类似。

如果我们在某个超块尺寸 $L$ 处停止无限系统 DMRG，可以用两种方式解读最终的波函数。我们可以把它视为尺寸为 $L$ 的超块的精确态的近似，并计算期望值。其精度不仅受截断的限制，还受如下事实的影响：最初的几次截断是在极小的超块上进行的：为短块挑选的相关态，很可能与把这些短块嵌入长度为 $L$ 的最终系统时所选出的态相差不小。

另一种做法是忽略这些边界效应，聚焦于中央格点，把它们当作无限大系统行为的近似，前提是 $L$ 足够大，并且对结果向 $L\rightarrow\infty$ 作了仔细的外推。这要做到非常可控并不简单，但正如我们将看到的，无限系统 DMRG 可用来构造高度可控的、平移不变的（在相差比如说一个长度为 2 的单胞的意义下）热力学极限态。

\subsection{有限系统 DMRG}

一旦无限系统 DMRG 达到了想要的最终系统尺寸，除了最平凡的应用之外，重要的是接着进行所谓的有限系统 DMRG 步骤。这不只会给我们的结果带来一些微小的定量改进，还可能使其彻底改观：参见 \cite{Schollwoeck03} 中的一个例子，其中连核心的物理结论都变了：对于一个梯子上的 $t$-$J$-$V$-$V'$ 模型、中等空穴掺杂 $\delta=0.1$，无限系统 DMRG 计算表明在方格（plaquette）上存在交替环流，由梯子左端的一个无穷小电流所触发，这是一种所谓 $d$-密度波态的信号。只有在应用有限系统算法之后才变得明显：这个电流实际上向体内指数衰减，从而排除了这种类型的序（图~\ref{fig:PlaquetteCurrent}）。

\begin{figure}
\centering\includegraphics[width=\textwidth]{PyMPS_DMRGCombined.eps}
\caption{图的左半部分与右半部分分别展示无限系统与有限系统 DMRG 步骤中的迭代过程。在两种情形下，新块都是通过把一个格点并入块而形成的，并按照 DMRG 的密度矩阵方案作态空间截断。在无限系统版本中，这种生长发生在链的两侧，导致链变长；而在有限系统算法中，生长只发生在一侧，以牺牲另一侧为代价，链长保持不变。}
\label{fig:DMRGCombined}
\end{figure}

有限系统算法所修正的，是在某个超块的背景下为约化基所作的选择——该超块并非所关心的系统（最终长度 $L$），而只是它的某种较小的替身。

\begin{figure}
\centering\includegraphics[width=250pt]{PyMPS_PlaquetteCurrent.eps}
\caption{$t$-$J$-$V$-$V'$ 梯子上的方格（plaquette）电流，$J=0.4t$，$V=3t$，$V'=t$（$V$ 为最近邻、$V'$ 为次近邻相互作用），空穴掺杂 $\delta=0.1$，系统尺寸 $2\times 60$，由横杆 $1$ 上的有限边界电流诱导。图中显示的是电流的绝对强度；无限系统 DMRG 与第一次扫描都表明产生了长程图样，而完全收敛的计算（此处为 6 到 7 次扫描之后）则揭示出向体内的指数衰减。取自文献~\cite{Schollwoeck03}。}
\label{fig:PlaquetteCurrent}
\end{figure}

有限系统算法所做的事情如下（Figure \ref{fig:DMRGCombined}）：它按照与之前相同的方案继续（比如说）块 B 的生长过程：求超块系统的基态，确定约化密度算符，求本征系统，为下一个更大的块保留权重最高的 $D$ 个本征态。但这是以牺牲块 A 为代价的，A 收缩（即复用旧的较短的块 A）。如此继续，直到 A 小到具有完整的希尔伯特空间，即维数不超过 $D$（也可以一直继续到 A 只剩一个格点长；结果不受影响）。然后生长方向反转：A 以牺牲 B 为代价生长，包括为 A 重新确定基态和选取基，直到 B 小到具有完整的希尔伯特空间，这又导致生长方向再次反转。

这种贯穿系统的 {\em 扫描}一直继续，直到能量（或者更准确地说，波函数）收敛。这个（实践中极为成功的）步骤的直观动机是：每完成一次扫描，块 A 或 B 都是在一个不断改善的嵌入环境中确定的。

在实践中，这个算法涉及大量簿记工作，因为我们需要的所有算符都必须在当前的有效基中维护，而这些基会一步步改变。这意味着每一步之后都要执行所确定的截断基变换；所有基中的算符表示都必须存储下来，因为收缩中的块还会再次用到它们。

另一个重要特点是：为每个 A$\bullet\bullet$B 构型求基态 $\ket{\psi}$ 时，使用的是某种迭代法大型稀疏矩阵本征值求解器，它基于把 $\hat{H}$ 逐次作用于某个初始起始矢量。为加速算法中这个最耗时的部分，非常希望有一个好的起始矢量预测，即尽可能接近最终解。这可以通过（近似地）把上一步的结果变换到移动后的 A$\bullet\bullet$B 构型 \cite{White96} 来实现，办法是施加两次基变换：例如向右扫描时作 A$\bullet\rightarrow$A 和 B$\rightarrow\bullet$B。显式公式（见 \cite{Schollwoeck05,White96}）可以这样导出：写出
\begin{equation}
\ket{\psi} = \sum_{a_\ell\sigma_{\ell+1}\sigma_{\ell+2}b_{\ell+2}} \psi_{a_\ell\sigma_{\ell+1}\sigma_{\ell+2}b_{\ell+2}} \ket{a_\ell}_A \ket{\sigma_{\ell+1}} \ket{\sigma_{\ell+2}} \ket{b_{\ell+2}}_B,
\end{equation}
其中 $\ket{a_\ell}_A$ 和 $\ket{b_{\ell+2}}_B$ 分别是包含格点 1 到 $\ell$ 的块 A 与包含格点 $\ell+3$ 到 $L$ 的块 B 的块态（块态的标记取自其末端所截的键的标记；见 图~\ref{fig:labelling}），并两次插入近似恒等式，即 $\hat{I} = \sum_{a_{\ell+1}} \ket{a_{\ell+1}}_A {\,}_A\bra{a_{\ell+1}}$ 与
$\hat{I} = \sum_{\sigma_{\ell+3} b_{\ell+3}} \ket{\sigma_{\ell+3}}\ket{b_{\ell+3}}_B {\, }_B \bra{b_{\ell+3}} \bra{\sigma_{\ell+3}}$。
于是得到
\begin{equation}
\ket{\psi} = \sum_{a_{\ell+1}\sigma_{\ell+2}\sigma_{\ell+3}b_{\ell+3}} \psi_{a_{\ell+1}\sigma_{\ell+2}\sigma_{\ell+3}b_{\ell+3}} \ket{a_{\ell+1}}_A \ket{\sigma_{\ell+2}} \ket{\sigma_{\ell+3}} \ket{b_{\ell+3}}_B,
\end{equation}
其中
\begin{equation}
\psi_{a_{\ell+1}\sigma_{\ell+2}\sigma_{\ell+3}b_{\ell+3}} =
\sum_{a_\ell \sigma_{\ell+1} b_{\ell+2}}
\psi_{a_\ell\sigma_{\ell+1}\sigma_{\ell+2}b_{\ell+2}}
\braket{a_{\ell+1}}{a_\ell \sigma_{\ell+1}} \braket{b_{\ell+3}\sigma_{\ell+3}}{b_{\ell+2}} .
\end{equation}
上面最后一个等式所需的基变换都可以从 DMRG 步骤中先前的步骤获得。类似的操作也可以用于向左的扫描。正如我们将看到的，这一巧妙步骤——它曾使 DMRG 的性能大幅提升——在把 DMRG 改写为 MPS 语言时已经隐含其中，所以我们不在此详细讨论。
一个重要的观察是，无限系统与有限系统算法也都可以只在其中插入单个显式格点 $\bullet$ 来执行，因此人们可以研究形如 A$\bullet$B 的超块，只需对增长步骤稍作调整。这种方法的优点是大型稀疏本征值求解器可以加速大约 $d$ 倍；例如，把 $\hat{h}$ 作用到式~(\ref{eq:hamfragment}) 中的态上，将只需 $O(D^3 d)$ 次运算。
在无限系统算法中，一个明显的缺点是超块长度会在奇偶之间来回振荡；而在有限系统算法中，二者相对优劣的问题要有趣得多，将在第 \ref{subsec:conventionalDMRGinMPS} 节详细讨论。

显然，当 $D\rightarrow\infty$ 时不会发生任何截断，DMRG 变为精确方法；增大 $D$ 会减少截断，从而单调地改善观测量结果，人们把结果在 $D\rightarrow\infty$ 下外推（更好的做法是在截断误差 $\epsilon\rightarrow 0$ 下外推，因为局域观测量对 $\epsilon$ 的误差依赖往往实际上是线性的），以获得最佳结果。

关于 DMRG 及其应用的更多细节，我推荐参阅 \cite{Schollwoeck05}。
  
\section{DMRG 与纠缠：DMRG 为何成功、为何失败}
DMRG 算法相当自然地把人们引向对二分量子系统的考察，其两部分为 A$\bullet$ 和 $\bullet$B。对任意一个二分，$\ket{\psi} = \sum_{ij} \psi_{ij} \ket{i}_A \ket{j}_B$，其中态 $\ket{i}_A$ 和 $\ket{j}_B$ 分别构成维数为 $N_A$ 和 $N_B$ 的正交归一基。把 $\psi_{ij}$ 看作矩形矩阵 $\Psi$（维数 $N_A \times N_B$）的矩阵元，则约化密度矩阵 $\rho_A$ 与 $\rho_B$ 取如下形式
\begin{equation}
\rho_A = \Psi \Psi^\dagger \quad\quad \rho_B = \Psi^\dagger \Psi.
\label{eq:ABdensityoperators}
\end{equation}
假设我们精确地知道 $\ket{\psi}$，但在 DMRG 中每块最多只能用 $D$ 个态来近似它，那么最优的 DMRG 近似就是保留属于最大的 $D$ 个本征值的本征态作为块态。如果我们碰巧知道 $\ket{\psi}$ 的约化密度算符的本征谱，就能容易地评估 DMRG 近似所能达到的质量；它只取决于本征值 $w_a$ 下降得多快。事实上，这类分析已在若干一维和二维的严格可解系统中进行过 \cite{Peschel99,Okunishi99,Chung00,Chung01,Chan02}。结果表明：在一维，对有能隙的系统，本征值 $w_a$ 通常呈指数式快速衰减（大致如 $\eul^{- c \ln^2 a}$），这解释了 DMRG 的成功；但在尺寸为 $L \times W$、$L\gg W$ 的二维条带几何中，$c \propto 1/W$，于是随着宽度 $W$ 增大（二维性增强），本征谱衰减如此之慢，以至于 DMRG 变得低效。

通常我们对本征值谱并没有清晰的认识；但事实证明，在这种情况下
纠缠熵可以作为"代理"量，即 von Neumann 纠缠或纠缠熵。它由 $\rho_A$ 的本征值谱的非零部分给出（如我们将在下文讨论的，它与 $\rho_B$ 的相同）：
\begin{equation}
S_{A|B}  = - \Tr\, \rho_A \log_2 \rho_A = - \sum_\alpha w_a \log_2 w_a .
\end{equation}
表面上看我们一无所获，因为我们并不知道 $w_a$，但关于纠缠标度的普遍定律是现成的。考虑一个二分 A$|$B，其中 AB 处于热力学极限，A 的大小为 $L^{\cal D}$，${\cal D}$ 为空间维数，那么所谓的{\em 面积定律}
\cite{Eisert10,Bekenstein73,Srednicki93,Callan94,Plenio05} 预言：对具有激发能隙的短程哈密顿量的基态而言，纠缠熵不是广延量，而是正比于表面积，即 $S(A|B) \sim L^{{\cal D}-1}$，与热熵相反。这意味着在一维 $S \sim $ 常数，在二维 $S \sim L$。在临界点处，会出现远为丰富的结构：一维中，
$S = \frac{c+\overline{c}}{6} \log_2 L + k$，其中 $c$ 和 $\overline{c}$ 是来自共形场论的（非）全息中心荷\cite{Vidal03,Latorre04}；二维中，玻色系统似乎对临界性不敏感（即 $S \propto L$）\cite{Srednicki93,Barthel06}，而费米系统在一维费米面情形会得到对数修正 $S \propto L \log_2 L$（前因子正比于费米面大小），但当费米面由点构成时，似乎只以低于对数的速度增长 \cite{Barthel06,Gioev06}。值得强调的是，基态的这些性质是非常不寻常的：在热力学极限下，从希尔伯特空间中随机抽取的一个态确实会以概率 1 展现广延的纠缠熵。

现在可以用一种在数学上不严格的方式，把 DMRG 与量子纠缠的面积定律联系起来：对 A 与 B 各自的 $D$ 维态空间而言，最大纠缠为 $\log_2 D$，出现在 $\rho_A$ 的所有本征值都相同且为 $D^{-1}$ 的情形（此时 $\rho_A$ 为最大混合态）；这意味着要恰当地编码纠缠 $S$，需要一个维数为 $2^S$ 以上的态。这表明，对一维有能隙系统，系统尺寸的增大不会导致 $D$ 的急剧增大；而在二维中 $D \sim 2^L$，由于资源必须指数增长，DMRG 即使对相对较小的系统尺寸也会失效（但这并不排除对小型二维团簇或相当大的条带得到非常精确的结果）。一维临界系统是边缘情形：$D \sim L^{\frac{c+\overline{c}}{6}}$；这意味着热力学极限不可企及，但 $D$ 的增长足够缓慢（由于中心荷的典型取值，幂次通常较弱，比如 $1/3$ 或 $1/6$），因而可以达到很大的系统尺寸（$L \sim O(10^3)$）；这使人们能够进行非常精确的有限尺寸外推。

显然，这一论证隐含了一个过于随意的假设，即本征值谱接近平坦，从而给出最大纠缠，这样才能对 $D$ 作出近似估计。实际上，谱由具体问题决定，而且确实远离平坦：如我们已看到的，它通常呈指数衰减。但数值上，对标准问题而言，资源 $D$ 的标度在定性层面上还是能被正确预言的。

更进一步，在数学上严格的分析中，von Neumann 纠缠{\em 并不能}对资源用量作出普遍预言：这是因为人们可以构造出人工的本征值谱，它们或者允许、或者禁止高效模拟，而若按上述论证，其 von Neumann 纠缠却会给出相反的暗示 \cite{Schuch08}。然而，人们感兴趣的典型多体态并不具有这种"病态"谱。事实上，Renyi 纠缠熵——von Neumann 纠缠熵的一种推广——确实允许建立数学上严格的联系 \cite{Schuch08}，但通常难以计算，多亏共形场论，一维临界情形是个例外。

\section{矩阵乘积态（MPS）}
\label{sec:matrixproductstates}

考虑我们的范式问题——一维海森堡反铁磁体，关键困难在于希尔伯特空间似乎是指数般巨大的（$d^L = 2^L$）。寻找基态因此可能像大海捞针。我们的主张是：至少对基态与第一激发态之间存在能隙的局域哈密顿量而言，这片"干草堆"并不大，与整个希尔伯特空间的大小相比实际上小到可以忽略——这一点我们从纠缠标度的奇特性质中已有领略。更重要的是，希尔伯特空间的这个相关角落可以被高效地参数化，即只需适度的数值资源，并能被高效地操作，而且确实存在 DMRG 类型的高效算法来解决量子物理中的重要问题。这一参数化就由{\em 矩阵乘积态（MPS）}给出。

上面讲解的两种 DMRG 算法也许看起来实现起来非常繁琐。但事实证明，一旦我们在矩阵乘积态所提供的受限态类中做量子力学，DMRG 及其他方法几乎是自己主动浮现出来的。操纵矩阵乘积态起初似乎非常复杂，但实际上可以被优美地形式化，还配上一种图示记法，几乎能让人自动生成允许的操作；正如任何好的形式体系（比如 bra 与 ket 记号）一样，它本质上就保证了正确性。

\subsection{矩阵乘积态的引入}

\subsubsection{奇异值分解（SVD）与 Schmidt 分解}
在本文余下的部分中，我们将大量使用线性代数中最多能的工具之一——所谓的{\em 奇异值分解}（SVD），它是一种非常紧凑的量子态表示的基础，适用于生活在二分宇宙 AB 中的量子态，即{\em Schmidt 分解}。让我们简要回顾一下它们的内容。

对任意维数为 $(N_A \times N_B)$ 的（矩形）矩阵 $M$，SVD 保证存在如下分解
\begin{equation}
M =   U S V^\dagger ,
\end{equation}
其中 
\begin{itemize}
\item $U$ 的维数为 $(N_A \times \min (N_A,N_B))$，各列正交归一（即{\em 左奇异向量}），即 $U^\dagger U = I$；若 $N_A \leq N_B$，这意味着它是幺正的，并且还有 $UU^\dagger = I$。 
\item $S$ 的维数为 $(\min (N_A,N_B) \times \min (N_A,N_B))$，是对角矩阵，其对角元为非负的 $S_{aa} \equiv s_a$。它们就是所谓的{\em 奇异值}。非零奇异值的个数 $r$ 称为 $M$ 的{\em (Schmidt) 秩}。下面我们假定奇异值按降序排列：$s_1 \geq s_2 \geq \ldots \geq s_r > 0.$
\item $V^\dagger$ 的维数为 $(\min (N_A,N_B) \times N_B)$，各行正交归一（即{\em 右奇异向量}），即 $V^\dagger V = I$。若 $N_A \geq N_B$，这意味着它是幺正的，并且还有 $VV^\dagger  = I$。
\end{itemize}
这一情形示意性地画在图~\ref{fig:SVD} 中。奇异值与奇异向量有许多极为有趣的性质。其中对下文具有实用价值的一条是：$M$（秩 $r$）可以被矩阵 $M'$（秩 $r'<r$）在 Frobenius 范数 $\| M \|_F^2 = \sum_{ij} |M_{ij}|^2$ 下最优地近似，该范数由内积 $\braket{M}{N}= \tr\, M^\dagger N$ 诱导。最优近似为 
\begin{equation}
M' = U S' V^\dagger \quad\quad S'={\rm diag} (s_1, s_2, \ldots, s_{r'}, 0,\ldots ),
\end{equation} 
即把除前 $r'$ 个之外的所有奇异值都置为零（在数值实践中，还会相应地把 $U$ 的列维数和 $V^\dagger$ 的行维数缩减到 $r'$）。

\begin{figure}
\centering\includegraphics[width=\textwidth]{PyMPS_SVD.eps}
\caption{奇异值分解（SVD）所导致的矩阵形状，对应于可能出现的两种矩形形状。奇异值对角线提醒我们，在 $M=USV^{\ \dagger}$ 中 $S$ 是纯非负对角矩阵。}
\label{fig:SVD}
\end{figure}

作为 SVD 的第一个应用，我们用它来导出一般量子态的{\em Schmidt 分解}。AB 上的任意纯态 
$\ket{\psi}$
都可以写成
\begin{equation}
\ket{\psi} = \sum_{ij} \Psi_{ij} \ket{i}_A \ket{j}_B,
\label{eq:generalquantumstate}
\end{equation}
其中 $\{ \ket{i}_A \}$ 与 $\{ \ket{j}_B \}$ 分别是 A 和 B 的维数为 $N_A$ 与 $N_B$ 的正交归一基；我们把系数读作矩阵 $\Psi$ 的矩阵元。由这一表示出发，可以导出约化密度算符 $\hat{\rho}_A = \tr_B \ket{\psi}\bra{\psi}$ 与 $\hat{\rho}_B = \tr_A \ket{\psi}\bra{\psi}$，它们在块基下取矩阵形式 
\begin{equation} 
\rho_A = \Psi \Psi^\dagger \quad\quad \rho_B = \Psi^\dagger\Psi .
\end{equation}

如果对式~(\ref{eq:generalquantumstate}) 中的矩阵 $\Psi$ 作 SVD，我们得到
\begin{eqnarray}
\ket{\psi} &=& \sum_{ij} \sum_{a=1}^{\min(N_A,N_B)} U_{ia} S_{aa} V^*_{ja} \ket{i}_A \ket{j}_B \nonumber \\
&=& \sum_{a=1}^{\min(N_A,N_B)} \left( \sum_i U_{ia} \ket{i}_A \right) s_a \left( \sum_j V_{ja}^* \ket{j} \right) \nonumber \\
&=& \sum_{a=1}^{\min(N_A,N_B)} s_a \ket{a}_A \ket{a}_B .
\end{eqnarray}
由于 $U$ 与 $V^\dagger$ 的正交归一性，集合 $\{ \ket{a}_A \}$ 与 $\{ \ket{a}_B \}$ 是正交归一的，并可扩展为 A 和 B 的正交归一基。如果把求和限制为只跑遍 $r \leq \min(N_A,N_B)$ 个正的非零奇异值，我们就得到{\em Schmidt 分解}
\begin{equation}
\ket{\psi} = \sum_{a=1}^r s_a \ket{a}_A \ket{a}_B .
\end{equation}
显而易见，$r=1$ 对应（经典）直积态，而 $r>1$ 对应纠缠的（量子）态。

Schmidt 分解使我们能非常方便地读出前面引入的 A 和 B 的约化密度算符：完成偏迹运算，可得
\begin{equation}
\hat{\rho}_A = \sum_{a=1}^r s^2_a \ket{a}_A  \phantom\rangle_A \bra{a} \quad\quad \hat{\rho}_B = \sum_{a=1}^r s^2_a \ket{a}_B \phantom\rangle_B \bra{a} ,
\end{equation}
这表明二者共享谱的非零部分，但本征态不同。本征值是奇异值的平方，$w_a = s^2_a$，相应的本征向量就是左、右奇异向量。因此，von Neumann 纠缠熵可以直接从 SVD 读出，
\begin{equation}
S_{A|B}(\ket{\psi}) = - \tr\, \hat{\rho}_A \log_2 \hat{\rho}_A = - \sum_{a=1}^r s_a^2 \log_2 s_a^2 .
\end{equation}

考虑到希尔伯特空间的巨大尺寸，用一个只在 A 和 B 的维数仅为 $r'$ 的态空间上张成的 $\ket{\tilde{\psi}}$ 来近似 $\ket{\psi}$ 也是自然的。这个问题可以与 SVD 联系起来，因为 $\ket{\psi}$ 的 2-范数与矩阵 $\Psi$ 的 Frobenius 范数相等，
\begin{equation}
\| \ket{\psi} \|_2^2 = \sum_{ij} | \Psi_{ij} |^2 = \| \Psi \|_F^2 ,
\end{equation}
当且仅当集合 $\{ \ket{i} \}$ 与 $\{ \ket{j} \}$ 正交归一（此处正是如此）。
因此，最优近似在 2-范数意义下就是 $\Psi$ 在 Frobenius 范数意义下被 $\tilde{\Psi}$ 的最优近似，其中 $\tilde{\Psi}$ 是秩为 $r'$ 的矩阵。如上所述，$\tilde{\Psi} = US' V^\dagger$，其中 $S'={\rm diag}(s_1, \ldots, s_{r'}, 0, \ldots)$，由 $\Psi$ 的最大的那些奇异值构成。于是，近似态的 Schmidt 分解为
\begin{equation}
\ket{\tilde{\psi}} = \sum_{a=1}^{r'} s_a \ket{a}_A \ket{a}_B ,
\label{eq:Schmidtapproximation}
\end{equation}
若要求归一化，则 $s_a$ 还需重新标度。

\subsubsection{QR 分解}
虽然后面会看到 SVD 能覆盖我们的一切需求，但有时它有点大材小用：在 $M=USV^\dagger$ 的许多应用场景中，我们只关心性质 $U^\dagger U = I$ 以及乘积 $SV^\dagger$，例如凡是奇异值的实际数值不会被显式使用的时候。此时有一种数值上更廉价的技术，即 {\em QR 分解}：对任意维数为 $(N_A \times N_B)$ 的矩阵 $M$，它给出分解
\begin{equation}
M = QR ,
\label{eq:QRdecomposition}
\end{equation}
此名即由此而来。其中 $Q$ 的维数为 $(N_A \times N_A)$ 且是幺正的，$Q^\dagger Q = I = Q Q^\dagger$；$R$ 的维数为 $(N_A \times N_B)$ 且是上三角的，即当 $i>j$ 时 $R_{ij}=0$。这一{\em 完全} QR 分解可以约化为{\em 精简} QR 分解：设 $N_A > N_B$，则 R 底部的 $N_A-N_B$ 行为零，于是可以写成
\begin{equation}
M=Q \cdot \left[ \begin{array}{c} R_1 \\ 0 \end{array} \right] = \left[ \begin{array}{cc} Q_1 & Q_2 \end{array} \right] \left[ \begin{array}{c} R_1 \\ 0 \end{array} \right] = Q_1 R_1 ,
\label{eq:thinQRdecomposition}
\end{equation}
其中 $Q_1$ 的维数为 $(N_A \times N_B)$，$R_1$ 的维数为 $(N_B \times N_B)$；虽然 $Q_1^\dagger Q_1 = I$，但一般而言 $Q_1 Q_1^\dagger \neq I$。在下面的叙述中，凡我提到 QR 分解，指的都是精简版。还应当清楚的是，矩阵 $Q$（或 $Q_1$）与 SVD 中的 $U$ 具有共同的性质，但一般并不相同；不过正如我们将看到的，态的 MPS 表示本来就是非唯一的，所以这无关紧要。

\subsubsection{任意量子态分解为 MPS}
\label{subsubsec:MPSdecomposition}
考虑一个由 $L$ 个格点组成的晶格，各格点 $i=1,\ldots,L$ 上有 $d$ 维局域态空间 $\{ \sigma_i \}$。事实上，虽然我们自然会想到一维晶格，但以下论述对任意维度的晶格同样成立，只要格点已被编号；然而，由高维晶格上的态生成的 MPS 在数值实践中将难以处理。晶格上最一般的纯量子态为
\begin{equation}
\ket{\psi} = \sum_{\sigma_1,\ldots,\sigma_L} c_{\sigma_1 \ldots \sigma_L} \ket{\sigma_1,\ldots,\sigma_L},
\end{equation}
其中系数 $c_{\sigma_1 \ldots \sigma_L}$ 的数目呈指数增长，而在典型的量子多体问题中它们的内容相当晦涩。让我们假定该态已归一化。能否找到一种记法，使这个态具有更局域的面貌（同时保持态的量子非定域性）？确实，SVD 恰好能做到这一点。其结果初看可能相当令人生畏，但我们将看到它与量子物理中熟悉的概念有着深刻的联系。与我们相关的做法有三种。

{\em (i) 左规范矩阵乘积态。}第一步，我们把具有 $d^L$ 个分量的态矢量{\em 重排}为维数 $(d \times d^{L-1})$ 的矩阵 $\Psi$，其中系数的关系为 
\begin{equation}
\Psi_{\sigma_1, (\sigma_2 \ldots \sigma_L)} = c_{\sigma_1 \ldots \sigma_L}  .
\end{equation}
对 $\Psi$ 作 SVD 给出
\begin{equation}
c_{\sigma_1 \ldots \sigma_L}  = \Psi_{\sigma_1, (\sigma_2 \ldots \sigma_L)} = \sum_{a_1}^{r_1} U_{\sigma_1,a_1} S_{a_1,a_1} (V^\dagger)_{a_1, (\sigma_2 \ldots \sigma_L)} \equiv \sum_{a_1}^{r_1}  U_{\sigma_1,a_1} c_{a_1\sigma_2\ldots\sigma_L} ,
\end{equation}
其中在最后一个等式中 $S$ 与 $V^\dagger$ 已相乘，所得矩阵又被重排回一个矢量。秩为 $r_1 \leq d$。现在我们把矩阵 $U$ 分解为 $d$ 个行矢量 $A^{\sigma_1}$ 的集合，其矩阵元为 $A^{\sigma_1}_{a_1}=U_{\sigma_1,a_1}$。与此同时，我们把 $c_{a_1\sigma_2\ldots\sigma_L}$ 重排为维数 $(r_1 d \times d^{L-2})$ 的矩阵 $\Psi_{(a_1\sigma_2),(\sigma_3 \ldots \sigma_L)}$，得到
\begin{equation}
c_{\sigma_1 \ldots \sigma_L}  = \sum_{a_1}^{r_1} A^{\sigma_1}_{a_1} \Psi_{(a_1\sigma_2), (\sigma_3 \ldots \sigma_L)} .
\end{equation}
对 $\Psi$ 作 SVD，我们有 
\begin{equation}
c_{\sigma_1 \ldots \sigma_L}  = \sum_{a_1}^{r_1} \sum_{a_2}^{r_2} A^{\sigma_1}_{a_1} U_{(a_1\sigma_2),a_2} S_{a_2,a_2} (V^\dagger)_{a_2, (\sigma_3 \ldots \sigma_L)} =  
\sum_{a_1}^{r_1} \sum_{a_2}^{r_2} A^{\sigma_1}_{a_1} A^{\sigma_2}_{a_1,a_2} \Psi_{(a_2\sigma_3), (\sigma_4 \ldots \sigma_L)} ,
\end{equation}
其中我们把 $U$ 换成了 $d$ 个维数为 $(r_1 \times r_2)$ 的矩阵 $A^{\sigma_2}$ 的集合，其矩阵元为 $A^{\sigma_2}_{a_1,a_2} = U_{(a_1\sigma_2),a_2}$，并把 $S$ 与 $V^\dagger$ 相乘，再重排为维数 $(r_2 d \times d ^{L-3})$ 的矩阵 $\Psi$，其中 $r_2 \leq r_1 d \leq d^2$。继续作更多的 SVD，我们得到
\begin{equation}
c_{\sigma_1 \ldots \sigma_L} = \sum_{a_1, \ldots, a_{L-1}} A^{\sigma_1}_{a_1} A^{\sigma_2}_{a_1,a_2} \ldots A^{\sigma_{L_1}}_{a_{L-2},a_{L-1}} A^{\sigma_{L}}_{a_{L-1}}
\end{equation}
或者更紧凑地写作
\begin{equation}
c_{\sigma_1 \ldots \sigma_L} = A^{\sigma_1} A^{\sigma_2} \ldots A^{\sigma_{L-1}} A^{\sigma_L} ,
\end{equation}
其中我们已经认出对 $a_1$、$a_2$ 等等的求和正是矩阵乘法。最后一组矩阵 $A^{\sigma_L}$ 实际上由列矢量组成。如果愿意，可以在第一个和最后一个 $A$ 中引入哑指标 1，把它们也变成矩阵。无论如何，这个（任意的）量子态现在已被精确地表示为{\em 矩阵乘积态：}的形式
\begin{equation}
\ket{\psi} = \sum_{\sigma_1,\ldots,\sigma_L} A^{\sigma_1} A^{\sigma_2} \ldots A^{\sigma_{L-1}} A^{\sigma_L} \ket{\sigma_1,\ldots,\sigma_L} .
\end{equation}

让我们研究 $A$ 矩阵的性质。矩阵维数取到最大值的条件是：每次所作的 SVD 中非零奇异值的个数都等于上界（即被分解矩阵两维中较小者）。计数表明，从第一个格点到最后一个格点，维数最大可为 $(1 \times d), (d \times d^2), \ldots, (d^{L/2-1} \times d^{L/2}), (d^{L/2} \times d^{L/2-1}), \ldots, (d^2 \times d), (d \times 1)$（这里为简单起见我假定 $L$ 为偶数）。这表明在实际计算中通常不可能显式地完成这种精确分解，因为矩阵维数会指数式地膨胀。

但事情还不止于此。由于每次 SVD 都满足 $U^\dagger U = I$，把 $U$ 换成一组 $A^\sigma$ 就蕴含如下关系：
\begin{eqnarray*}
\delta_{a_\ell,a'_\ell} &=& \sum_{a_{\ell-1}\sigma_\ell} (U^\dagger)_{a_\ell, (a_{\ell-1}\sigma_\ell)} U_{(a_{\ell-1}\sigma_\ell), a'_\ell} \\
&=& \sum_{a_{\ell-1}\sigma_\ell} (A^{\sigma_\ell\dagger})_{a_\ell, a_{\ell-1}} A^{\sigma_\ell}_{a_{\ell-1}, a'_\ell} \\
&=& \sum_{\sigma_\ell} ( A^{\sigma_\ell\dagger} A^{\sigma_\ell})_{a_\ell,a'_\ell} 
\end{eqnarray*}
或者更简洁地写作，
\begin{equation}
\sum_{\sigma_\ell} A^{\sigma_\ell \dagger} A^{\sigma_\ell} = I . 
\end{equation}
满足这一条件的矩阵我们将称为{\em 左归一化}的，而仅由左归一化矩阵构成的矩阵乘积态我们将称为{\em 左规范}的。事实上，细看会发现最后一个格点上该条件在技术上可能不成立，但正如我们在计算 MPS 的范数时将会看到的，这对应于原来的态未归一化到 1。让我们暂且忽略这个微妙之处。

鉴于 DMRG 把整个宇宙分解为 A、B 两块，把晶格分成 A、B 两部分是很有启发性的，其中 A 包含格点 $1$ 至 $\ell$，B 包含格点 $\ell+1$ 至 $L$。于是我们可以引入态
\begin{eqnarray}
\ket{a_\ell}_A &=& \sum_{\sigma_1,\ldots,\sigma_\ell} (A^{\sigma_1} A^{\sigma_2} \ldots A^{\sigma_\ell})_{1,a_\ell} \ket{\sigma_1,\ldots,\sigma_\ell} \\
\ket{a_\ell}_B &=& \sum_{\sigma_{\ell+1},\ldots,\sigma_L} (A^{\sigma_{\ell+1}} A^{\sigma_{\ell+2}} \ldots A^{\sigma_L})_{a_\ell,1} \ket{\sigma_{\ell+1},\ldots,\sigma_L}
\end{eqnarray}
于是 MPS 可以写为
\begin{equation}
\ket{\psi} = \sum_{a_\ell} \ket{a_\ell}_A \ket{a_\ell}_B . 
\end{equation}
这种态的配对看起来与 $\ket{\psi}$ 的 Schmidt 分解极其接近，但事实并非如此。原因在于，虽然 $\{ \ket{a_\ell}_A \}$ 构成正交归一集合，$\{ \ket{a_\ell}_B \}$ 一般却不然。这是 $A$ 矩阵左归一性的直接推论。对 A 部分我们求得 
\begin{eqnarray*}
\phantom\langle_A \braket{a'_\ell}{a_\ell}_A &=& \sum_{\sigma_1,\ldots,\sigma_\ell} (A^{\sigma_1} \ldots A^{\sigma_\ell})^*_{1,a'_\ell} (A^{\sigma_1} \ldots A^{\sigma_\ell})_{1,a_\ell} \\
&=& \sum_{\sigma_1,\ldots,\sigma_\ell} (A^{\sigma_1} \ldots A^{\sigma_\ell})^\dagger_{a'_\ell,1} (A^{\sigma_1} \ldots A^{\sigma_\ell})_{1,a_\ell} \\
&=& \sum_{\sigma_1,\ldots,\sigma_\ell} (A^{\sigma_\ell\dagger} \ldots A^{\sigma_1\dagger}A^{\sigma_1} \ldots A^{\sigma_\ell})_{a'_\ell,a_\ell} \\
&=& \delta_{a'_\ell,a_\ell} ,
\end{eqnarray*}
其中我们迭代地完成了从 $\sigma_1$ 到 $\sigma_\ell$ 的求和并利用了左归一性。另一方面，对 B 部分作同样的计算得到
\begin{eqnarray*}
\phantom\langle_B \braket{a'_\ell}{a_\ell}_B &=& \sum_{\sigma_{\ell+1},\ldots,\sigma_L} (A^{\sigma_{\ell+1}} \ldots A^{\sigma_L})^*_{a'_\ell,1} (A^{\sigma_{\ell+1}} \ldots A^{\sigma_L})_{a_\ell,1} \\
&=& \sum_{\sigma_{\ell+1},\ldots,\sigma_L} (A^{\sigma_L\dagger} \ldots A^{\sigma_{\ell+1}\dagger})_{1,a'_\ell} (A^{\sigma_{\ell+1}} \ldots A^{\sigma_L})_{a_\ell,1} \\
&=& \sum_{\sigma_{\ell+1},\ldots,\sigma_L} (A^{\sigma_{\ell+1}} \ldots A^{\sigma_L}A^{\sigma_L\dagger} \ldots A^{\sigma_{\ell+1}\dagger})_{a'_\ell,a_\ell},
\end{eqnarray*}
它无法进一步化简，因为一般来说 $\sum_\sigma A^\sigma A^{\sigma\dagger} \neq I$。

态系数表示方式的改变也可以用图形来表示（图~\ref{fig:MPSBySVD}）。我们把系数 $c_{\sigma_1 \ldots \sigma_L}$ 画成一个黑箱（边角圆润），物理指标 $\sigma_1$ 至 $\sigma_L$ 垂直伸出。第一次分解之后的结果如第二行所示：左边是一个代表 $A^{\sigma_1}_{a_1}$ 的对象，右边是 $c_{a_1 \sigma_2 \ldots \sigma_L}$。辅助自由度（$a_1$）用水平线表示，规则是{\em 相连的线要求和。}第二步就显而易见了：我们有 
$A^{\sigma_1}_{a_1}$，然后是 $A^{\sigma_2}_{a_1,a_2}$，右边是 $c_{a_2 \sigma_3 \ldots \sigma_L}$，所有相连的线都要求和。最后，我们得到 $L$ 个相乘并带物理指标标记的 $A$ 矩阵（图的最后一行）。

$A$ 矩阵的图形规则——它们在第一个和最后一个格点上分别是行矢量和列矢量——总结于图~\ref{fig:ABulkEdge}：格点 $\ell$ 用实心圆表示，物理指标 $\sigma_\ell$ 用一条竖线表示，两个矩阵指标用水平线表示。

\begin{figure}
\centering\includegraphics[scale=0.6]{PyMPS_MPSBySVD.eps}
\caption{通过一系列奇异值分解迭代构造任意量子态的精确 MPS 表示的图形表示。}
\label{fig:MPSBySVD}
\end{figure}

\begin{figure}
\centering\includegraphics[scale=0.6]{PyMPS_ABulkEdge.eps}
\caption{链两端及链内部 $A$ 矩阵的图形表示：左图表示 $A\ ^{\ \sigma_{1}}_{\ 1,a_1}$，即左端的行矢量，右图表示 $A\ ^{\sigma_{L}}_{a_L,1}$，即右端的列矢量。中间则是 $A\ ^{\sigma_{\ell}}_{a_{\ell-1},a_\ell}$。}
\label{fig:ABulkEdge}
\end{figure}

在本小节的最后，让我展示如何通过一系列 QR 分解生成左规范矩阵乘积态。我们从下式开始
\begin{equation}
c_{\sigma_1\ldots\sigma_L} = \Psi_{\sigma_1,(\sigma_2\ldots\sigma_L)} =
\sum_{a_1} Q_{\sigma_1,a_1} R_{a_1,(\sigma_2\ldots\sigma_L)} =
\sum_{a_1} A^{\sigma_1}_{1,a_1} \Psi_{(a_1\sigma_2),(\sigma_3\ldots\sigma_L)},
\end{equation}
其中与 SVD 步骤类似，我们把 $Q\rightarrow A$、$R\rightarrow \Psi$ 重排。下一次 QR 分解给出
\begin{equation}
c_{\sigma_1\ldots\sigma_L} = \sum_{a_1,a_2} A^{\sigma_1}_{1,a_1} Q_{(a_1\sigma_2),a_2} R_{a_2,(\sigma_3\ldots\sigma_L)} = \sum_{a_1,a_2} A^{\sigma_1}_{1,a_1} A^{\sigma_2}_{a_1,a_2} \Psi_{(a_2\sigma_3),(\sigma_4\ldots\sigma_L)}
\end{equation}
依此类推（对维数的分析表明，在链的右半部分需要用 thin QR）。$Q^\dagger Q=I$ 蕴含 $A$ 矩阵所期望的左归一性。在数值上可行的情况下，这比 SVD 更快。我们失去的是看不到奇异值的谱；而且除非使用更先进的秩显式（rank-revealing）QR 分解，我们也无法像 SVD 那样确定秩 $r_1, r_2, \ldots$。这意味着这种分解充分利用了 $A$ 矩阵的最大维数。

{\em (ii) 右规范矩阵乘积态。}显然，从左边（即格点 1）开始的分解没有任何特殊之处。类似地，我们也可以从右边开始，以便得到
\begin{eqnarray*}
c_{\sigma_1 \ldots \sigma_L} &=& \Psi_{(\sigma_1 \ldots \sigma_{L-1}), \sigma_L} \\
&=& \sum_{a_{L-1}} U_{(\sigma_1 \ldots \sigma_{L-1}), a_{L-1}} S_{a_{L-1},a_{L-1}} (V^\dagger)_{a_{L-1},\sigma_L} \\
&=& \sum_{a_{L-1}} \Psi_{(\sigma_1 \ldots \sigma_{L-2}),(\sigma_{L-1} a_{L-1})} B^{\sigma_L}_{a_{L-1}} \\
&=& \sum_{a_{L-2}, a_{L-1}} U_{(\sigma_1 \ldots \sigma_{L-2}), a_{L-2}} S_{a_{L-2},a_{L-2}} (V^\dagger)_{a_{L-2},(\sigma_{L-1}a_{L-1})} B^{\sigma_L}_{a_{L-1}} \\
&=& \sum_{a_{L-2}, a_{L-1}} \Psi_{(\sigma_1 \ldots \sigma_{L-3}), (\sigma_{L-2} a_{L-2})} B^{\sigma_{L-1}}_{a_{L-2},a_{L-1}} B^{\sigma_L}_{a_{L-1}} = \ldots \\
&=& \sum_{a_{1},\ldots, a_{L-1}} B^{\sigma_1}_{a_{1}} B^{\sigma_{2}}_{a_{1},a_{2}} \ldots B^{\sigma_{L-1}}_{a_{L-2},a_{L-1}} B^{\sigma_L}_{a_{L-1}}.
\end{eqnarray*}

这里，我们把 $(V^\dagger)_{a_{L-1},\sigma_L}$ 重排成 $d$ 个列向量 $B^{\sigma_L}_{a_{L-1}}$，把 $(V^\dagger)_{(a_{L-2}\sigma_{L-1}), a_{L-1}}$ 重排成 $d$ 个矩阵 $B^{\sigma_{L-1}}_{a_{L-2},a_{L-1}}$，依此类推，并且在每一步重排为 $\Psi$ 之前先将 $U$ 与 $S$ 相乘。显而易见的图形表示见图~\ref{fig:rightMPSbySVD}。为使记号简单，我们在图形表示中不对 $A$ 矩阵和 $B$ 矩阵加以区分。

\begin{figure}
\centering\includegraphics[scale=0.6]{PyMPS_RightMPSBySVD.eps}
\caption{通过一列奇异值分解迭代构造任意量子态的精确 MPS 表示的图形表示，这次从右侧开始。}
\label{fig:rightMPSbySVD}
\end{figure}


我们得到形如
\begin{equation}
\ket{\psi} = \sum_{\sigma_1,\ldots,\sigma_L} B^{\sigma_1} B^{\sigma_2} \ldots B^{\sigma_{L-1}} B^{\sigma_L} \ket{\sigma_1,\ldots,\sigma_L}
\end{equation}
其中可以证明 $B$ 矩阵具有与 $A$ 矩阵相同的矩阵维数上界，并且由 $V^\dagger V =I$ 还可证明它们满足
\begin{equation}
\sum_{\sigma_\ell} B^{\sigma_\ell} B^{\sigma_\ell\dagger} = I ,
\end{equation}
因此我们称它们为{\em 右规范化}矩阵。完全由这类矩阵构造而成的 MPS 我们称之为{\em 右规范}的。

同样，我们可以把格子分成 A 和 B 两部分，即格点 1 到 $\ell$ 和 $\ell+1$ 到 $L$，并引入态
\begin{eqnarray}
\ket{a_\ell}_A &=& \sum_{\sigma_1,\ldots,\sigma_\ell} (B^{\sigma_1} B^{\sigma_2} \ldots B^{\sigma_\ell})_{1,a_\ell} \ket{\sigma_1,\ldots,\sigma_\ell} \\
\ket{a_\ell}_B &=& \sum_{\sigma_{\ell+1},\ldots,\sigma_L} (B^{\sigma_{\ell+1}} B^{\sigma_{\ell+2}} \ldots B^{\sigma_L})_{a_\ell,1} \ket{\sigma_{\ell+1},\ldots,\sigma_L}
\end{eqnarray}
这样 MPS 就可以写成
\begin{equation}
\ket{\psi} = \sum_{a_\ell} \ket{a_\ell}_A \ket{a_\ell}_B . 
\end{equation}
这种态的配对看起来又一次诱人地接近 $\ket{\psi}$ 的 Schmidt 分解，但同样并非如此。原因在于，虽然这次 $\{ \ket{a_\ell}_B \}$ 构成正交归一集合，但 $\{ \ket{a_\ell}_A \}$ 一般并非如此，这可以由 $B$ 矩阵的右规范性得到证明。

同样，右规范化形式也可以通过一列 QR 分解得到。与左规范化形式的区别在于，我们不是对 $\Psi=QR$ 作 QR 分解，而是对 $\Psi^\dagger=QR$ 作 QR 分解，于是 $\Psi = R^\dagger Q^\dagger$。如果我们用 $Q^\dagger$ 来构造这些矩阵，这就直接给出 $B$ 矩阵的右规范化性质。让我把头两步明确写出；我们从
\begin{equation}
c_{\sigma_1\ldots\sigma_L} = \Psi_{(\sigma_1\ldots\sigma_{L-1}),\sigma_L} =
\sum_{a_{L-1}} R^\dagger_{(\sigma_1\ldots\sigma_{L-1}),a_{L-1}} Q^\dagger_{a_{L-1},\sigma_L} =
\sum_{a_{L-1}} \Psi_{(\sigma_1\ldots\sigma_{L-2}),(\sigma_{L-1}a_{L-1})} B^{\sigma_L}_{a_{L-1},1},
\end{equation}
把 $R^\dagger$ 重排为 $\Psi$、$Q^\dagger$ 重排为 $B$，然后继续对 $\Psi^\dagger$ 作 QR 分解，得
\begin{equation}
c_{\sigma_1\ldots\sigma_L} = \sum_{a_{L-1},a_{L-2}} R^\dagger_{(\sigma_1\ldots\sigma_{L-2}),a_{L-2}} Q^\dagger_{a_{L-2},(\sigma_{L-1}a_{L-1})}  B^{\sigma_L}_{a_{L-1},1} 
= \sum_{a_{L-1},a_{L-2}} \Psi_{(\sigma_1\ldots\sigma_{L-3}),(\sigma_{L-2}a_{L-2})}
B^{\sigma_{L-1}}_{a_{L-2},a_{L-1}} B^{\sigma_L}_{a_{L-1},1} .
\end{equation}

至此我们已经得到了 $\ket{\psi}$ 的多种不同的精确 MPS 形式表示，这已经表明一个态的 MPS 表示{\em 并不唯一}，我们后面将利用这一事实。

{\em (iii) 混合规范矩阵乘积态。}我们也可以把从左和从右对态的分解混合起来。假设我们已经从左分解到格点 $\ell$，使得
\begin{equation}
c_{\sigma_1\ldots\sigma_L} = \sum_{a_\ell} (A^{\sigma_1} \ldots A^{\sigma_\ell})_{a_\ell} S_{a_\ell,a_\ell} (V^\dagger)_{a_\ell, (\sigma_{\ell+1} \ldots \sigma_L)} .
\end{equation}
我们把 $V^\dagger$ 重排为 $\Psi_{(a_\ell \sigma_{\ell+1} \ldots \sigma_{L-1}), \sigma_L}$，并像原先从右分解那样从右进行逐次 SVD，直到并包括格点 $\sigma_{\ell+2}$；在最后一次 SVD 中剩下 $U_{(a_{\ell}\sigma_{\ell+1}),a_{\ell+1}} S_{a_{\ell+1},a_{\ell+1}}$，我们把它重排为 $B^{\sigma_{\ell+1}}_{a_{\ell}a_{\ell+1}}$。于是我们
得到
\begin{equation}
(V^\dagger)_{a_\ell, (\sigma_{\ell+1} \ldots \sigma_L)} = \sum_{a_{\ell+1},\ldots,a_{L-1}} B^{\sigma_{\ell+1}}_{a_\ell, a_{\ell+1}} \ldots B^{\sigma_L}_{a_{L-1}} .
\end{equation}
所有 $B$ 矩阵都是右规范化的。对格点 $\ell+2$ 到 $L$ 而言，这只是 SVD 的直接结果；在格点 $\ell+1$ 上，则由性质 $V^\dagger V = I$ 得到：
\begin{eqnarray*}
\delta_{a_\ell,a'_\ell} &=&
\sum_{\sigma_{\ell+1},\ldots} (V^\dagger)_{a_\ell, (\sigma_{\ell+1} \ldots \sigma_L)} V_{( \sigma_{\ell+1} \ldots \sigma_L),a'_\ell} \\
&=& ( \sum_{\sigma_{\ell+1},\ldots} B^{\sigma_{\ell+1}} \ldots B^{\sigma_L} B^{\sigma_L\dagger} \ldots B^{\sigma_{\ell+1}\dagger} )_{a_\ell,a'_\ell} \\
&=& ( \sum_{\sigma_{\ell+1}} B^{\sigma_{\ell+1}} B^{\sigma_{\ell+1}\dagger} )_{a_\ell,a'_\ell} ,
\end{eqnarray*}
其中最后一行利用了格点 $\ell+2,\ldots,L$ 上所有 $B$ 矩阵的右规范化性质，从而得到所要的结果。

于是我们最终得到分解
\begin{equation}
c_{\sigma_1 \ldots \sigma_L} = A^{\sigma_1} \ldots A^{\sigma_\ell} S B^{\sigma_{\ell+1}} \ldots B^{\sigma_L} ,
\end{equation}
它包含键 $(\ell,\ell+1)$ 上的奇异值，可以像图~\ref{fig:bothMPSbySVD} 那样作图形表示。

\begin{figure}
\centering\includegraphics[scale=0.6]{PyMPS_BothMPSBySVD.eps}
\caption{由一列奇异值分解从左和从右出发得到的精确 MPS 的图形表示。菱形表示对角的奇异值矩阵。左边的矩阵是左规范化的，右边的矩阵是右规范化的。}
\label{fig:bothMPSbySVD}
\end{figure}

不仅如此，分成 A 与 B 的 Schmidt 分解——其中 A 覆盖格点 1 到 $\ell$，B 覆盖格点 $\ell+1$ 到 $L$——现在可以立即读出。如果我们引入矢量
\begin{eqnarray}
\ket{a_\ell}_A &=& \sum_{\sigma_1,\ldots,\sigma_\ell}
(A^{\sigma_1} \ldots A^{\sigma_\ell})_{1,a_\ell}
 \ket{\sigma_1,\ldots,\sigma_\ell} \\
\ket{a_\ell}_B &=&  \sum_{\sigma_{\ell+1},\ldots,\sigma_L} 
(B^{\sigma_{\ell+1}} \ldots B^{\sigma_L})_{a_\ell,1}
\ket{\sigma_{\ell+1},\ldots,\sigma_L}
\end{eqnarray}
则态便取形式（$s_a = S_{a,a}$）
\begin{equation}
\ket{\psi} = \sum_{a_\ell} s_a \ket{a_\ell}_A \ket{a_\ell}_B ,
\end{equation}
这就是 Schmidt 分解，{\em 只要} A {\em 和} B {\em 上的态各自正交归一。}而按构造这确实成立。

(iv) {\em 规范自由度。}至此，我们已有把任意量子态写成 MPS 的三种不同方式，它们各有优劣。虽然这三种可以说是书写 MPS 最重要的方式，但必须认识到非唯一性的程度远不止于此：MPS 不唯一，其意义在于存在规范自由度。考虑两个相邻的矩阵集合 $M^{\sigma_i}$ 和 $M^{\sigma_{i+1}}$，其共享的列/行维数为 $D$。那么对任意维数为 $(D\times D)$ 的可逆矩阵 $X$，MPS 在变换
\begin{equation}
M^{\sigma_i} \rightarrow M^{\sigma_i} X, \quad \quad M^{\sigma_{i+1}} \rightarrow X^{-1}M^{\sigma_{i+1}} .
\end{equation}
之下不变。这一规范自由度可用来极大地简化各种操作，我们的三种构造只是它的特例。
 
由此产生几个问题。这种记法与多体物理中人们更熟悉的概念之间有联系吗？确实，它与重正化群方案中出现的迭代状态抽取过程有深刻的联系，我们将在第~\ref{subsubsec:MPSsinglesitedecimation} 节讨论这一点。

这些矩阵的规模可能大到指数级，在计算机上我们必须把它们的大小限制在某个 $D$ 以内。这能否做到而不使对态的描述过于失真？在一维这确实可能：如果我们考虑混合规范表示，就会看到，对于约化密度算符指数衰减的本征值谱（因而指数衰减的奇异值 $s_a$），可以按照式~(\ref{eq:Schmidtapproximation}) 在最大的 $D$ 个奇异值处截断谱（就 2-范数意义下的最优近似而言），而没有可察觉的精度损失。这一论证可以从单次截断造成的近似推广到 $L-1$ 次截断（每个键一次）造成的近似，从而表明误差最坏为 \cite{Verstraete06}
\begin{equation}
\| \ket{\psi} - \ket{\psi_{{\rm trunc}}} \|_2^2 \leq 2 \sum_{i=1}^L \epsilon_i(D) ,
\end{equation}
其中 $\epsilon_i(D)$ 是键 $i$ 上截断到最大的 $D$ 个奇异值所造成的截断误差（被丢弃的奇异值平方之和）。因此，与 DMRG 中一样，可近似性问题与约化密度算符的本征值谱相关联，这预示它在二维会失效，并且（论证稍弱一些）与面积定律的存在相关联。

\subsubsection{一维中的 MPS 与单格点状态抽取}
\label{subsubsec:MPSsinglesitedecimation}

为了把 MPS 与更传统的概念联系起来，让我们设想为自旋链建立一个迭代生长过程 $\ell \rightarrow \ell + 1$，如图~\ref{fig:BlockGrowth} 所示，使得相应的态空间每步增大 $d$ 倍。为了避免指数增长，我们现在要求态空间维数有一个上限 $D$。一旦态空间维数超过 $D$，就必须用某种{\em 目前尚未定义}的过程把态空间截断。

假设我们以某种方式得到了长度为 $\ell-1$ 的系统（用 DMRG 的语言即左块 A）的这样一个 $D$ 维有效基 $\{ \ket{a_{\ell-1}}_A \}$。如果截断后长度为 $\ell$ 的（左）块 A 的 $D$ 个基态是 $\{ \ket{a_{\ell}}_A \}$，而新增格点的局域态是 $\{ \ket{\sigma_\ell} \}$，则必须有
\begin{equation}
\ket{a_\ell}_A = \sum_{a_{\ell-1}\sigma_\ell} \phantom\langle _A \braket{a_{\ell-1}\sigma_\ell}{a_\ell}_A \,\, \ket{a_{\ell-1}}_A\ket{\sigma_\ell}
\label{eq:basistrafoblocksite}
\end{equation}
对这些态成立，其中 $\phantom\langle _A \braket{a_{\ell-1}\sigma_\ell}{a_\ell}_A$ 目前尚未指定。现在我们在格点 $\ell$ 引入 $d$ 个维数各为 $(D \times D)$ 的矩阵 $A^{[\ell]\sigma_\ell}$，每个可能的局域态 $\ket{\sigma_\ell}$ 对应一个。于是可以把式~(\ref{eq:basistrafoblocksite}) 改写为
\begin{equation}
\ket{a_\ell}_A = \sum_{a_{\ell-1}\sigma_\ell} A^{[\ell]\sigma_\ell}_{a_{\ell-1}, a_\ell} \ket{a_{\ell-1}}_A  \ket{\sigma_\ell} 
\label{eq:matrixbyRG}
\end{equation}
其中矩阵 $A^{[\ell]\sigma_\ell}$ 的矩阵元由下式给出（见图~\ref{fig:ABulk}）
\begin{equation}
A^{[\ell]\sigma_\ell}_{a_{\ell-1}, a_\ell} \equiv \phantom\langle_A \braket{a_{\ell-1}\sigma_\ell}{a_\ell}_A. 
\label{eq:Amatrixelements}
\end{equation}
这里让我们对{\em 记法}作一简短说明：在 $A^{[\ell]\sigma_\ell}$ 中，$[\ell]$ 标明考虑的是哪一组 $A$ 矩阵，而 $\sigma_\ell$ 标明具体是哪一个 $A$ 矩阵。在目前的情形，这种记法有些多余，因为局域态 $\ket{\sigma_\ell}$ 就取自引入这些矩阵的格点。在这类情形，我们略去两个 $\ell$ 中的一个，通常是 $[\ell]$：
\begin{equation}
A^{[\ell]\sigma_\ell} \rightarrow A^{\sigma_\ell} .
\end{equation}
不过，我们也会遇到矩阵 $A$ 由{\em 并非}其引入格点上的局域态来选取的情形。在这种情形下，显然必须恢复完整的记号！

类似地，当态处于块 A 上这一事实无关紧要或完全显而易见时，我们将把 $\ket{a_\ell}_A \rightarrow \ket{a_\ell}$ 加以简写。

\begin{figure}
\centering\includegraphics[scale=0.8]{PyMPS_BlockGrowth.eps}
\caption{长度为 $\ell-1$ 的块通过添加一个格点 $\ell$ 向右生长为长度为 $\ell$ 的块。}
\label{fig:BlockGrowth}
\end{figure}

\begin{figure}
\centering\includegraphics[scale=0.8]{PyMPS_ABulk.eps}
\caption{$A$ 矩阵的图形表示：左图表示 $A\ ^{\sigma_{\ell}}_{a_{\ell-1},a_\ell}$，右图表示{\em 共轭}的 $A\ ^{\sigma_{\ell}*}_{a_{\ell-1},a_\ell}$。实心圆代表格点，竖线代表物理指标，横线代表矩阵指标。}
\label{fig:ABulk}
\end{figure}

矩阵记法的优点——其中的状态抽取步骤尚未具体给出——在于它允许从长度为 $\ell$ 的块出发，进行简单的{\em 递归}，一直到最小的、即消失的块。以这种方式得到的量子态具有非常特殊的形式：
\begin{eqnarray}
\ket{a_{\ell}}_A &=& \sum_{a_{\ell-1}} \sum_{\sigma_\ell} A^{\sigma_\ell}_{a_{\ell-1}, a_{\ell}} \ket{a_{\ell-1}}_A \ket{\sigma_{\ell}} \nonumber \\
&=& \sum_{a_{\ell-1},a_{\ell-2}} \sum_{\sigma_{\ell-1},\sigma_\ell} 
A^{\sigma_{\ell-1}}_{a_{\ell-2}, a_{\ell-1}} A^{\sigma_\ell}_{a_{\ell-1}, a_{\ell}} \ket{a_{\ell-2}}_A \ket{\sigma_{\ell-1}}\ket{\sigma_{\ell}} = \ldots \nonumber \\
&=& \sum_{a_1,a_2,\ldots,a_{\ell-1}} \sum_{\sigma_1,\sigma_2,\ldots,\sigma_\ell} A^{\sigma_1}_{1,a_1} A^{\sigma_2}_{a_1,a_2} \ldots A^{\sigma_\ell}_{a_{\ell-1},a_\ell} \ket{\sigma_1}\ket{\sigma_2} \ldots \ket{\sigma_\ell} \nonumber \\
&=& \sum_{\sigma_i \in {\rm A}}  (A^{\sigma_1}A^{\sigma_2}\ldots A^{\sigma_\ell})_{1,a_\ell}  \ket{\sigma_1}\ket{\sigma_2} \ldots \ket{\sigma_\ell} ,
\end{eqnarray}
其中 $i$ 遍历块 A 的所有格点。指标 'A' 表明我们所考虑的是正在构建的链左侧的态。在第一个格点上，为了与矩阵记法保持一致，我们引入了一个哑行指标 1：块长为 1 的态由格点 1 上的局域态与长度为 0 的“块”的块态构成，对后者我们引入一个哑态及指标 1。这意味着 $A^{\sigma_1}$ 实际上是一个（行）矢量（见图~\ref{fig:ABulkEdge}）。我们还看到，$A$ 的左指标即行指标，对应于由右指标即列指标所标记的那些态“左侧”的态。相当一般地，我们可以如图~\ref{fig:BlockByMPS} 那样展示这一构造，只要——如前所述——我们引入规则：所有相连的腿都要求和（收缩）。矩阵记法的优点在于可以把求和隐藏在矩阵乘法之中。

\begin{figure}
\centering\includegraphics[scale=0.8]{PyMPS_BlockByMPS.eps}
\caption{通过 $A$ 矩阵的收缩（相乘）递归构造态 $\ket{a{\,}_\ell}$ 的图形表示。收缩遍历所有相连的腿。}
\label{fig:BlockByMPS}
\end{figure}

类似地，我们也可以让块向左而不是向右生长（图~\ref{fig:BlockGrowth_Right}）：我们有
\begin{equation}
\ket{a_\ell}_B = \sum_{a_{\ell+1}\sigma_{\ell+1}} \phantom\langle_B \braket{a_{\ell+1} \sigma_{\ell+1}}{a_\ell}_B \,\, \ket{a_{\ell+1}}_B \ket{\sigma_{\ell+1}} 
\end{equation}
或
\begin{equation}
\ket{a_\ell}_B = \sum_{a_{\ell+1}\sigma_{\ell+1}}B^{[\ell+1]\sigma_{\ell+1}}_{a_\ell, a_{\ell+1}}  \ket{a_{\ell+1}}_B \ket{\sigma_{\ell+1}} 
\end{equation}
其中
\begin{equation}
B^{[\ell+1]\sigma_{\ell+1}}_{a_\ell, a_{\ell+1}} = \phantom\langle_B \braket{a_{\ell+1} \sigma_{\ell+1}}{a_\ell}_B .
\end{equation}
我们把这些矩阵记为 $B$，以表明它们源自向左的生长过程；用 DMRG 的语言来说，这对应于块 B。递归给出
\begin{equation}
\ket{a_{\ell}}_B = \sum_{\sigma_i \in {\rm B}}  (B^{\sigma_{\ell+1}}B^{\sigma_{\ell+2}}\ldots B^{\sigma_L})_{a_{\ell+1},1}  \ket{\sigma_{\ell+1}}\ket{\sigma_{\ell+2}} \ldots \ket{\sigma_L}, 
\end{equation}   
其中 $i$ 从 $\ell+1$ 遍历到 $L$，即块 B 的格点。与位置 1 处类似，在位置 $L$ 处也引入了一个类似的哑指标，此处 $B$ 矩阵是一个（列）矢量。

注意与 $A$ 矩阵相比记法上的细微不对称：为了日后能把块 A 与块 B 拼接起来，我们按照块态所终止的键来标记它们：键 $\ell$ 连接格点 $\ell$ 与 $\ell+1$，于是得到如图~\ref{fig:labelling} 所示的标记方式。

\begin{figure}
\centering\includegraphics[width=\textwidth]{PyMPS_BlockGrowth_Right.eps}
\caption{长度为 $L-\ell-1$ 的块 B 通过添加格点 $\ell+1$ 向左生长为长度为 $L-\ell$ 的块 B。}
\label{fig:BlockGrowth_Right}
\end{figure}

\begin{figure}
\centering\includegraphics[scale=1.0]{PyMPS_Labelling.eps}
\caption{块 A（格点 1 到 $\ell$）与块 B（格点 $\ell+1$ 到 $L$）在键 $\ell$ 处相接。态标记为 $\ket{a\ _{\ell}}_A$ 与 $\ket{a'_\ell}_B$。}
\label{fig:labelling}
\end{figure}

如果以这种方式引入 $A$ 矩阵与 $B$ 矩阵，可以看到它们具有非常特殊的性质。考虑从左边的生长，即 $A$ 矩阵，并合理地要求所选的态对每个块长都彼此正交归一，则利用 式~(\ref{eq:matrixbyRG}) 可得
\begin{eqnarray}
\delta_{a'_\ell, a_\ell} = \phantom\langle_A \braket{a'_\ell}{a_\ell}_A &=& \sum_{\sigma'_\ell,\sigma_\ell} \sum_{a'_{\ell-1},a_{\ell-1}} A^{\sigma'_\ell*}_{a'_{\ell-1}, a'_\ell} A^{\sigma_\ell}_{a_{\ell-1},a_\ell} 
 \phantom\langle_A\braket{a'_{\ell-1} \sigma'_\ell}{a_{\ell-1} \sigma_\ell}_A \\
&=& \sum_{\sigma_\ell} \sum_{a_{\ell-1}} A^{\sigma_\ell \dagger}_{a'_\ell,a_{\ell-1}} A^{\sigma_\ell}_{a_{\ell-1},a_\ell} = \sum_{\sigma_\ell}
(A^{\sigma_\ell \dagger}A^{\sigma_\ell})_{a'_{\ell},a_{\ell}} .
\end{eqnarray}
总结起来，我们发现 $A$ 矩阵是{\em 左规范}的：
\begin{equation}
\sum_{\sigma} A^{\sigma\dagger} A^\sigma = I .
\end{equation}
图~\ref{fig:LeftNormalization} 给出了图形表示：该乘法也可以解读为 $A$ 与 $A^*$ 在 $\sigma$ 及其左指标上的收缩。

类似地，对于从右边构建的块 B 的 $B$ 矩阵，可以推导出{\em 右规范}恒等式
\begin{equation}
\sum_{\sigma} B^\sigma B^{\sigma\dagger} = I
\end{equation}
成立（通常用 $A$ 与 $B$ 来区分这两种情形）。见 图~\ref{fig:RightNormalization}。这意味着正交归一的态总可以按 MPS 的意义分解为左规范或右规范的矩阵，而且凡由左规范或右规范矩阵构造出来的态都构成正交归一集，只要规范化的类型与生长的方向相匹配。

\begin{figure}
\centering\includegraphics[scale=1.0]{PyMPS_LeftNormalization.eps}
\caption{若把两个{\em 左}规范的 $A$ 矩阵沿其{\em 左}指标与物理指标收缩，就会得到一条 $\delta\ _{a'_\ell,a_\ell}$ 线。}
\label{fig:LeftNormalization}
\end{figure}

\begin{figure}
\centering\includegraphics[scale=1.0]{PyMPS_RightNormalization.eps}
\caption{若把两个{\em 右}规范的 $B$ 矩阵沿其{\em 右}指标与物理指标收缩，就会得到一条 $\delta\ _{a'_\ell,a_\ell}$ 线。}
\label{fig:RightNormalization}
\end{figure}

让我们更仔细地考察矩阵的维数。从左边生长时，矩阵维数依次为 $(1 \times d)$、$(d \times d^2)$、$(d^2 \times d^3)$、$(d^3 \times D)$，这里我假定了 $d^4 > D$。此后维数保持在 $(D \times D)$。在右端，维数将依次为 $(D \times D)$、$(D \times d^3)$、$(d^3 \times d^2)$、$(d^2 \times d)$ 和 $(d \times 1)$。

现在我们又可以写下一个矩阵乘积态了。把一条长度为 $L$ 的链由长度为 $\ell$ 的（左）块 A（格点 $1$ 到 $\ell$）与长度为 $L-\ell$ 的（右）块 B（格点 $\ell+1$ 到 $L$）拼合起来，就可以构成一般的叠加
\begin{equation}
\ket{\psi} = \sum_{a_\ell, a'_\ell} \Psi_{a_\ell, a'_\ell} \ket{a_\ell}_A \ket{a'_\ell}_B .
\end{equation}
把态显式地代入，我们发现
\begin{equation}
\ket{\psi} = \sum_{\fat{\sigma}} (A^{\sigma_1} \ldots A^{\sigma_\ell})_{1,a_\ell} \Psi_{a_\ell, a'_\ell} 
 (B^{\sigma_{\ell+1}} \ldots B^{\sigma_L})_{a'_\ell,1} \ket{\fat{\sigma}} .
\end{equation}
粗体的 $\fat{\sigma}$ 代表所有局域态指标，$\ket{\fat{\sigma}} =\ket{\sigma_1,\sigma_2,\ldots,\sigma_L}$。这一记法提示我们把 $\Psi$ 解释为一个矩阵；于是记法简化为
\begin{equation}
\ket{\psi} = \sum_{\fat{\sigma}} A^{\sigma_1} \ldots A^{\sigma_\ell}\Psi
 B^{\sigma_{\ell+1}} \ldots B^{\sigma_L} \ket{\fat{\sigma}} .
 \label{eq:MPSwithPsi}
\end{equation}


如果我们允许一般矩阵，且不担心左规范、右规范或无归一化，只需把 $\Psi$-矩阵乘入相邻的某个 $A$ 或 $B$ 矩阵，就得到开边界条件下的普遍 MPS 形式（见图~\ref{fig:MPS_OBC}）：
\begin{equation}
\ket{\psi} = \sum_{\fat{\sigma}} M^{\sigma_1} \ldots M^{\sigma_L} \ket{\fat{\sigma}}  \quad{\rm (MPS \ for \ OBC)},
\label{eq:MPSOBC}
\end{equation}
其中对归一化并{\em 不}作任何假设（正因如此我把这些矩阵记作 $M$）。由于第一个和最后一个矩阵具有矢量性质，乘积最终给出一个标量。这正是上一节已经讨论过的 MPS 形式。

\begin{figure}
\centering\includegraphics[scale=1.0]{PyMPS_MPS_OBC.eps}
\caption{开边界条件 MPS 的图示。}
\label{fig:MPS_OBC}
\end{figure}

至此很容易看出矩阵乘积态如何利用好量子数。让我们聚焦于磁化强度，并假设整体态的磁化强度为 $M$。这个阿贝尔量子数是可加的，$M = \sum_i M_i$。我们选取局域基 $\{ \sigma_i \}$，其状态是局域磁化强度的本征态。现在考虑从左边的生长过程。如果我们选 $\ket{a_1}$ 为局域磁化强度的本征态（例如直接取 $\ket{\sigma_1}$），那么式~(\ref{eq:basistrafoblocksite}) 使我们能够通过归纳法构造出磁化强度的本征态 $\ket{a_\ell}$，只要矩阵 $A^{\sigma_\ell}_{a_{\ell-1}, a_\ell}$ 具有分块结构，使得对每个非零矩阵元
\begin{equation}
M(\ket{a_{\ell-1}}) + M(\ket{\sigma_\ell}) = M(\ket{a_\ell})
\end{equation}
成立。这一点很容易在图示中表达出来：给图形表示（图~\ref{fig:MPS_OBC_quantumnumbers}）中的线加上方向，即画出射入和射出的箭头。规则很简单：射入线上磁化强度之和等于射出线上磁化强度之和。为了强制某个整体磁化强度 $M$，我们只需在第一个格点之前和最后一个格点之后的虚拟键上，分别给射入和射出的键赋予磁化强度值 0 和 $M$。我们可以设想 MPS 矩阵的指标是针对给定磁化强度的多重指标，允许简并，从而得到一种优雅的编码表示。键箭头的反转则直接对应于从右边生长所得 $B$-矩阵的结构，不过只要记账得当，箭头的取法有很大的自由度：反转只不过意味着符号要变号。

要在实际中利用好量子数，它们必须在我们对矩阵乘积态执行的典型操作下得以保持。事实证明，所有并非显而易见地无碍、且保持好量子数的操作都可以用 SVD 表达。SVD 会作用在形如 $A_{(a_{i-1}\sigma_i),a_i}$ 的矩阵上。如果我们按好量子数对态 $\ket{a_{i-1}\sigma_i}$ 和 $\ket{a_i}$ 进行分组，$A$ 就会由若干块组成；只要适当地重排标记，就可以写出 $A = A_1 \oplus A_2 \oplus \ldots = U_1 S_1 V^\dagger_1 \oplus U_2 S_2 V^\dagger_2 \oplus \ldots$，或者 $A=USV^\dagger$，其中 $U= U_1 \oplus U_2 \oplus \ldots$，依此类推。但这意味着经由 $U$ 从 $\ket{a_{i-1}\sigma_i}$ 生成的新态也具有好量子数。当需要截断时，这个性质对保留的态当然仍然成立。如果在可行之处用 QR 分解代替 SVD 并对各个块分别进行，$A_i =Q_i R_i$，那么幺正矩阵 $Q_i$ 只在具有相同量子数的态集合内部变换，因此好量子数得以保持。

\begin{figure}
\centering\includegraphics[scale=0.85]{PyMPS_MPS_OBC_QuantumNumbers.eps}
\caption{具有好（可加）量子数的开边界条件 MPS 的图示。物理态和键都变成有向的，使得射入线上的量子数等于射出线上的量子数。对于第一个格点之前和最后一个格点之后的虚拟键，我们取适当的值以固定整体好量子数。}
\label{fig:MPS_OBC_quantumnumbers}
\end{figure}

现在假设我们的格点系统满足周期边界条件。在态系数 $c_{\sigma_1\ldots\sigma_L}$ 的层面上并不存在边界条件的概念，因此我们的 MPS 标准形式完全能够描述一个体现周期边界条件的态。从这个意义上说，断言式~(\ref{eq:MPSOBC}) 只对开边界条件成立实际上是错误的。它的正确之处仅在于：第一个和最后一个格点上矩阵的异常结构对周期边界条件并不方便；事实上，跨越键 $(L,1)$ 的纠缠必须编码为贯穿整条链的拉伸。这会导致数值效率极低的 MPS 表示。

对于周期边界条件，MPS 形式的自然推广是让所有矩阵都具有相同的维度 $(D \times D)$；由于格点 $L$ 又连回格点 $1$，我们通过取迹使 MPS 在所有键上都与矩阵乘法自洽（见图~\ref{fig:MPS_PBC}）：
\begin{equation}
\ket{\psi} = \sum_{\fat{\sigma}} \tr (M^{\sigma_1} \ldots M^{\sigma_L}) \ket{\fat{\sigma}}  \quad{\rm (MPS \ for \ PBC)}.
\label{eq:MPSforPBC}
\end{equation}
虽然{\em 先验地}看它并不比另一种形式更精确，但它要合适得多，在计算上也高效得多。
\begin{figure}
\centering\includegraphics[scale=1.0]{PyMPS_MPS_PBC.eps}
\caption{周期边界条件 MPS 的图示；底部的长线对应于求迹操作。}
\label{fig:MPS_PBC}
\end{figure}

本节的重点在于量子态的近似表示，而非通常无法实现的精确表示。虽然我们还没有给出如何构造这些近似表示的具体方案，但有几点评注是有必要的。

即使是近似的 MPS 也仍然是希尔伯特空间中所有态的线性组合，没有任何乘积基态被舍弃。限制其实在于线性组合的形式：不再是 $d^L$ 个系数，而是 $dL$ 个维度为 $(D\times D)$ 的矩阵，加上一个给出 $LD^2$ 个标量约束的矩阵型归一化约束，因此只有 $(d-1)LD^2$ 个独立参数，这使得态的各个系数之间产生了相互依赖。
 
对于给定的矩阵维度，任意量子态的最优近似的质量会随 $D$ 单调改善：取 $D<D'$，则 $D$ 所能达到的最佳近似可以写成一个键维数为 $D'$ 的 MPS，即在 $(D' \times D')$ 矩阵中嵌入 $(D \times D)$ 子矩阵，而所有新增的行和列均为零。它们为进一步改善态的近似提供了额外的参数。

乘积态（对任何 Schmidt 分解 Schmidt 秩均为 1）可以用 $D=1$ 的 MPS 精确写出。包含纠缠态的真实量子物理从 $D=2$ 开始。鉴于一个量子态中系数的数目呈指数增长，可能令人惊讶的是，即便在这个最简单的非平庸情形下，竟能精确地完成有趣的量子物理！但确实存在一些重要的量子态，它们在这个新格式下能找到紧凑的精确表达式。

\subsubsection{作为矩阵乘积态的 AKLT 态}
为了让 MPS 框架不那么抽象，让我们构造一个非平庸量子态的 MPS 表示。关联物理中最有趣的量子态之一是 1987 年提出的{\em Affleck-Kennedy-Lieb-Tasaki 态}，它是 AKLT 哈密顿量\cite{Affleck87,Affleck88} 的基态
\begin{equation}
\ham = \sum_{i} {\bf S}_i \cdot {\bf S}_{i+1} + 
\frac{1}{3} ({\bf S}_i \cdot {\bf S}_{i+1})^2 , 
\label{eq:aklt_hamiltonian}
\end{equation}
其中本文破例处理 $S=1$ 的自旋。可以证明，这个哈密顿量的基态可以按图~\ref{fig:akltstate} 所示的方式构造。每个自旋-1 都被一对完全对称化的自旋-$\frac{1}{2}$ 所取代，也就是说，在四个可能的态中我们只考虑自然地被识别为 $S=1$ 态的三个三重态：
\begin{eqnarray}
\ket{+}&=&\ket{\uparrow\uparrow} \nonumber \\
\ket{0}&=&\frac{\ket{\uparrow\downarrow}+\ket{\downarrow\uparrow}}{\sqrt{2}} 
\label{eq:tripletmapping} \\
\ket{-}&=&\ket{\downarrow\downarrow} \nonumber
\end{eqnarray}
在相邻格点上，相邻的自旋-$\frac{1}{2}$ 对以单态相连接
\begin{equation}
\frac{\ket{\uparrow\downarrow}-\ket{\downarrow\uparrow}}{\sqrt{2}} .
\end{equation}

事实证明，这个态可以用最低非平庸维度 $D=2$ 的矩阵乘积态编码，而且已经蕴含了大量令人振奋的物理 \cite{Affleck87,Affleck88,Fannes89}。在长度为 $\syslength$ 的链上 $2\syslength$ 个辅助自旋-$\frac{1}{2}$ 态的语言中，任意态都写成
\begin{equation}
    \ket{\psi} = \sum_{{\bf a}}\sum_{{\bf b}} c_{{\bf ab}} \ket{{\bf 
    ab}}
\end{equation}
其中 $\ket{{\bf a}}=\ket{a_1,\ldots,a_\syslength}$ 和 $\ket{{\bf b}}=\ket{b_1,\ldots,b_\syslength}$ 分别表示每个格点上的第一个和第二个自旋-$\frac{1}{2}$。现在我们把连接格点 $i$ 和 $i+1$ 的键 $i$ 上的单态键 $\Sigma_i$ 编码为
\begin{equation}
\ket{\Sigma^{[i]}} = \sum_{b_i a_{i+1}} \Sigma_{ba} \ket{b_i}\ket{a_{i+1}}
\end{equation}
并引入一个 $2 \times 2$ 矩阵 $\Sigma$
\begin{equation}
\Sigma = \left[ 
\begin{array}{cc} 0 & \frac{1}{\sqrt{2}} \\ -\frac{1}{\sqrt{2}} & 0 \end{array}
\right] .
\end{equation}
于是所有键上都有单态的态为
\begin{equation}
\ket{\psi_\Sigma} = \sum_{{\bf a}}\sum_{{\bf b}} \Sigma_{b_1 a_2} \Sigma_{b_2 a_3} \ldots \Sigma_{b_{\syslength-1} a_\syslength}\Sigma_{b_\syslength a_1} \ket{{\bf 
    ab}} 
\end{equation}
这对应周期边界条件。如果考虑开边界条件，则略去 $\Sigma_{b_\syslength a_1}$，第一个和最后一个自旋-$\frac{1}{2}$ 保持为单态。
\begin{figure}
\centering\includegraphics[scale=1.0]{PyMPS_AKLTState.eps}
\caption{Affleck-Kennedy-Lieb-Tasaki（AKLT）态的构造方式是：把局域自旋-1 表示成两个完全对称化的自旋-$\frac{1}{2}$ 粒子，它们通过单态跨格点相连。图中展示的是 PBC 的情形；对于 OBC，``长''键被切断，两端各出现一个单独的自旋-$\frac{1}{2}$。}
\label{fig:akltstate}
\end{figure}

注意，这个态是一个乘积态：只要把任意格点 $i$ 拆成它的两个组分，它就会因子化。现在我们通过引入一个{\em 映射}来编码辅助自旋的对称化态与物理自旋的对应关系：从
两个辅助自旋-$\frac{1}{2}$ 的态 $\ket{a_i}\ket{b_i} \in \{
\ket{\uparrow},\ket{\downarrow}\}^{\otimes 2}$ 映射到
物理自旋-1 的态 $\ket{\sigma_i} \in \{ \ket{+},\ket{0},\ket{-} \}$。为了表示 \Eq{eq:tripletmapping}，我们引入 $M_{ab}^{\sigma}
\ket{\sigma} \bra{ab}$，其中 $\ket{ab}$ 和 $\ket{\sigma}$ 分别表示格点 $i$ 上的辅助自旋和物理自旋。把 $M_{ab}^{\sigma}$ 写成三个 $2\times 2$ 矩阵，每个对应 $\ket{\sigma}$ 的一个取值，行与列的指标分别代表 $\ket{a}$ 和 $\ket{b}$ 的取值，我们得到
\begin{equation}
    M^+ = \left[ \begin{array}{cc} 1 & 0 \\ 0 & 0 \end{array} \right] \quad
    M^0 = \left[ \begin{array}{cc} 0 & \frac{1}{\sqrt{2}} \\ 
    \frac{1}{\sqrt{2}} & 0 \end{array} \right] \quad
    M^- = \left[ \begin{array}{cc} 0 & 0 \\ 0 & 1 \end{array} \right] .
\end{equation}
自旋-1 链希尔伯特空间 $\{ \ket{\fat{\sigma}} \}$ 上的这个映射于是读作
\begin{equation} 
\sum_{\fat{\sigma}} \sum_{{\bf a}{\bf b}} M^{\sigma_1}_{a_1 b_1} M^{\sigma_2}_{a_2 b_2} \ldots M^{\sigma_\syslength}_{a_\syslength b_\syslength}
\ket{\fat{\sigma}} \bra{{\bf a}{\bf b}} .
\end{equation}
因此 $\ket{\psi_\Sigma}$ 被映射为
\begin{equation}
\sum_{\fat{\sigma}} \sum_{{\bf a}{\bf b}} M^{\sigma_1}_{a_1 b_1} \Sigma_{b_1 a_2}
M^{\sigma_2}_{a_2 b_2} \Sigma_{b_2 a_3} \ldots \Sigma_{b_{\syslength-1} a_\syslength} 
M^{\sigma_\syslength}_{a_\syslength b_\syslength} \Sigma_{b_\syslength a_1} 
\ket{\fat{\sigma}}
\end{equation}
或者
\begin{equation}
    \ket{\psi} = \sum_{\fat{\sigma}} \tr (M^{\sigma_{1}} \Sigma
    M^{\sigma_{2}} \Sigma \ldots M^{\sigma_{\syslength}} \Sigma] ) \ket{\fat{\sigma}} ,
\end{equation}
这里使用了矩阵记法。为了进一步简化，我们引入 $\tilde{A}^\sigma= M^\sigma \Sigma$，于是
\begin{equation}
    \tilde{A}^+ = \left[ \begin{array}{cc} 0 & \frac{1}{\sqrt{2}} \\ 0 & 0 \end{array} \right] \quad
    \tilde{A}^0 = \left[ \begin{array}{cc} -\frac{1}{2} & 0 \\ 
    0 & +\frac{1}{2} \end{array} \right] \quad
    \tilde{A}^- = \left[ \begin{array}{cc} 0 & 0 \\ -\frac{1}{\sqrt{2}} & 0 \end{array} \right] .
    \label{eq:unnormalizeda}
\end{equation}
AKLT 态现在取如下形式
\begin{equation}
    \ket{\psi} = \sum_{\fat{\sigma}} \tr (\tilde{A}^{\sigma_{1}} 
    \tilde{A}^{\sigma_{2}} \ldots \tilde{A}^{\sigma_{\syslength}}) \ket{\fat{\sigma}} .
    \label{eq:akltstate}
\end{equation}
让我们把 $\tilde{A}^\sigma$ 左规范化。由 $\sum_\sigma \tilde{A}^{\sigma\dagger}\tilde{A}^\sigma = \frac{3}{4}I$ 可知，矩阵 $\tilde{A}$  应当
按 $\frac{2}{\sqrt{3}}$ 重新标度，从而得到归一化的矩阵 $A$，
\begin{equation}
    A^+ = \left[ \begin{array}{cc} 0 & \sqrt{\frac{2}{3}} \\ 0 & 0 \end{array} \right] \quad
    A^0 = \left[ \begin{array}{cc} -\frac{1}{\sqrt{3}} & 0 \\ 
    0 & \frac{1}{\sqrt{3}} \end{array} \right] \quad
    A^- = \left[ \begin{array}{cc} 0 & 0 \\ -\sqrt{\frac{2}{3}} & 0 \end{array} \right] .
    \label{eq:normalizeda}
\end{equation}
这在热力学极限下把态归一化：我们有
\begin{eqnarray*}
\braket{\psi}{\psi} &=& \sum_{\fat{\sigma}} \tr (A^{\sigma_1} \ldots A^{\sigma_\syslength})^* \tr (A^{\sigma_1} \ldots A^{\sigma_\syslength}) \\
&=&  \tr \left( \sum_{\sigma_1} A^{\sigma_1*} \otimes A^{\sigma_1}\right) \ldots \left( \sum_{\sigma_L} A^{\sigma_\syslength*}\otimes A^{\sigma_\syslength} \right) \\
&=& \tr E^L = \sum_{i=1}^4 \lambda_i^L .
\end{eqnarray*}
在这个表达式中，$\lambda_i$ 是
\begin{equation}
E = \sum_\sigma A^{\sigma*} \otimes A^{\sigma} = \left[ \begin{array}{cccc} 
\frac{1}{4} & 0 & 0 & \frac{1}{2} \\
0 & -\frac{1}{4} & 0 & 0 \\
0 & 0 & -\frac{1}{4} & 0 \\
\frac{1}{2} & 0 & 0 & \frac{1}{4}
\end{array} \right] ,
\end{equation}
的 4 个本征值，即 $1,-\frac{1}{3},-\frac{1}{3},-\frac{1}{3}$。于是当 $L\rightarrow\infty$ 时 $\braket{\psi}{\psi} = 1 + 3(-1/3)^L \rightarrow 1$。

下一节的方法现在可以用来解析地求出 AKLT 态的关联子：反铁磁关联函数呈指数衰减，$\langle S^z_i S^z_j \rangle \propto (-1/3)^{i-j}$，而弦关联函数在 $j-i>2$ 时为 $\langle S^z_i \eul^{\imag \pi\sum_{i<k<j} S^z_k}S^z_j \rangle =-4/9$，这表明存在隐藏序。

总而言之，AKLT 态可以被表达为一个 $D=2$ 的矩阵乘积态，即最简单的非平凡 MPS！事实上，把由最大纠缠态（这里是单态）连接起来的辅助自旋的较大态空间投影到较小的物理态空间，这一做法也可以作为引入 MPS 及其高维推广的出发点\cite{VerstraetePorras04,VerstraeteCirac04}。

\subsection{重叠、期望值与矩阵元}

现在我们转向对 MPS 的各种操作，从期望值的计算开始。
期望值显然是广义矩阵元的特例，此时态 $\ket{\psi}$ 与 $\ket{\phi}$ 相同。为保持一般性，让我们考虑态 $\ket{\psi}$ 与 $\ket{\phi}$ 之间的重叠，它们分别由矩阵 $M$ 和 $\tilde{M}$ 描述，并集中于开边界条件的情形。

取 $\ket{\phi}$ 的伴随，并注意到波函数系数是标量，重叠为
\begin{equation}
\braket{\phi}{\psi} = \sum_{\fat{\sigma}} \tilde{M}^{\sigma_1*} \ldots \tilde{M}^{\sigma_L*} M^{\sigma_1} \ldots M^{\sigma_L} .
\end{equation}
把由 $\tilde{M} \ldots \tilde{M}$ 构成的标量转置（这是恒等操作），我们得到顺序颠倒的伴随：
\begin{equation}
\braket{\phi}{\psi} = \sum_{\fat{\sigma}} \tilde{M}^{\sigma_L\dagger} \ldots \tilde{M}^{\sigma_1\dagger} M^{\sigma_1} \ldots M^{\sigma_L} .
\label{eq:finaloverlap}
\end{equation}
在图形表示中（图~\ref{fig:overlap}），如果我们遵循所有键指标都要求和这一规则，此计算会变得简单得多。
\begin{figure}
\centering\includegraphics[width=\textwidth]{PyMPS_Overlap.eps}
\caption{态 $\ket{\phi}$ 与 $\ket{\psi}$ 之间的重叠。所有对相同指标的收缩（求和）用箭头标出。}
\label{fig:overlap}
\end{figure}

\subsubsection{收缩的高效求值}

详细求值表达式（\ref{eq:finaloverlap}）表明，在矩阵网络或更一般的张量网络中，找到{\em 正确的（最优的）收缩次序}十分重要。我们要做的收缩，既有矩阵乘法中隐含的对矩阵指标的收缩，也有对物理指标的收缩。如果决定先收缩矩阵指标、再收缩物理指标，就必须对 $d^L$ 组矩阵乘法串求和，其代价是指数级的。但我们可以如下重组这些求和：
\begin{equation}
\braket{\phi}{\psi} = \sum_{\sigma_L} 
\tilde{M}^{\sigma_L\dagger} \left( 
\ldots \left( 
\sum_{\sigma_2}  \tilde{M}^{\sigma_2 \dagger} \left(
\sum_{\sigma_1} \tilde{M}^{\sigma_1\dagger} M^{\sigma_1}\right) 
M^{\sigma_2}\right) 
 \ldots  \right) 
 M^{\sigma_L} .
\end{equation}
这意味着，在（特殊的）第一步中，我们把列向量与行向量 $\tilde{M}^{\sigma_1\dagger}$ 和 $M^{\sigma_1}$ 相乘得到一个矩阵，并对（第一个）物理指标求和。下一步，我们对第二个物理指标收缩一个三矩阵乘积，依此类推（图~\ref{fig:overlaporder}）。重要的观察是，从第二步起复杂度不再增长。此外，把矩阵乘积 $ABC$ 分解为 $(AB)C$ 或 $A(BC)$ 当然也是最高效的做法。这样我们共进行 $(2L-1)d$ 次乘法，每次复杂度为 $O(D^3)$，这里暂且为简单起见忽略矩阵一般不是方阵的事实。决定性的一点在于：复杂度从指数级降为弱多项式级，总运算次数为 $O(LD^3 d)$。

\begin{figure}
\centering\includegraphics[width=\textwidth]{PyMPS_OverlapOrder.eps}
\caption{态 $\ket{\phi}$ 与 $\ket{\psi}$ 之间的重叠，并标明了最优收缩次序，它像拉链一样沿链推进。}
\label{fig:overlaporder}
\end{figure}


同样立即可见的是，对于范数计算 $\braket{\psi}{\psi}$，在开边界条件下，若态处于左规范或右规范形式，则立即意味着其范数为 1。从上面的计算可以看出，对于左规范矩阵 $A$，最内层的求和正是左规范化条件，给出 $I$，因而该项消去，接着下一个左规范化条件显现，如此继续直到走完整条链（图~\ref{fig:overlaplnormalized}）：
\begin{eqnarray*}
\braket{\psi}{\psi} &=& \sum_{\sigma_L} A^{\sigma_L\dagger} \left( \ldots \left( \sum_{\sigma_2}  A^{\sigma_2 \dagger} \left(\sum_{\sigma_1} A^{\sigma_1\dagger} A^{\sigma_1}\right) A^{\sigma_2}\right)  \ldots  \right) A^{\sigma_L} \\
&=&  \sum_{\sigma_L} A^{\sigma_L\dagger} \left( \ldots \left( \sum_{\sigma_2}  A^{\sigma_2 \dagger}  A^{\sigma_2}\right)  \ldots  \right) A^{\sigma_L} = \ldots \\
&=& \sum_{\sigma_L} A^{\sigma_L\dagger} A^{\sigma_L} = 1.
\end{eqnarray*}

\begin{figure}
\centering\includegraphics[width=\textwidth]{PyMPS_OverlapLNormalized.eps}
\caption{对{\em 左规范}态 $\ket{\psi}$ 计算范数的各步骤：依次应用左规范 $A$-矩阵的收缩规则。}
\label{fig:overlaplnormalized}
\end{figure}

为计算一般矩阵元，我们考虑 $\bra{\phi} \hat{O}^{[i]} \hat{O}^{[j]} \ldots \ket{\psi}$，即作用于格点 $i$ 和 $j$ 的张量积算符。这类算符的矩阵元取自
\begin{equation}
 \hat{O}^{[\ell]} = \sum_{\sigma_\ell,\sigma'_\ell} O^{\sigma_\ell,\sigma'_\ell} \ket{\sigma_\ell} \bra{\sigma'_\ell} .
\end{equation}
让我们把它推广为每个格点上都放一个算符，实际上几乎所有格点上都是恒等算符，例如对于局域期望值或两格点关联子就是如此。因此我们考虑的是算符矩阵元 $O^{\sigma_1,\sigma'_1} O^{\sigma_2,\sigma'_2} \cdots  O^{\sigma_L,\sigma'_L}$。
在解析表达式中，我们再次把（现在是双重的）对局域态的求和转置并分配开（$M$-矩阵的矩阵乘法与之前相同）：
\begin{eqnarray*}
& & \sum_{\fat{\sigma},\fat{\sigma}'} \tilde{M}^{\sigma_1*} \ldots \tilde{M}^{\sigma_L*} 
O^{\sigma_1,\sigma'_1} O^{\sigma_2,\sigma'_2} \cdots  O^{\sigma_L,\sigma'_L}
M^{\sigma'_1} \ldots M^{\sigma'_L} \\
&=& \sum_{\sigma_L,\sigma'_L} O^{\sigma_L,\sigma'_L} \tilde{M}^{\sigma_L\dagger}  \left( \ldots \left( \sum_{\sigma_2,\sigma'_2}  O^{\sigma_2,\sigma'_2} \tilde{M}^{\sigma_2 \dagger} \left(\sum_{\sigma_1,\sigma'_1} O^{\sigma_1,\sigma'_1} \tilde{M}^{\sigma_1\dagger} M^{\sigma'_1}\right) M^{\sigma'_2}\right)  \ldots  \right) M^{\sigma'_L}
\end{eqnarray*}
这与重叠的计算完全相同，唯一的例外是形式上对物理指标的单重求和变成了双重求和（图~\ref{fig:matrixelement}）。对于典型的关联子，双重求和在大多数格点上会平凡地约化为单重求和，因为大多数格点上作用的只是恒等算符 $\hat{O}^{[i]}=\hat{I}$；在少数非平凡格点上，多达 $d^2$ 个矩阵元中，对常规算符而言大多数为零，这强烈限制了项数，因此本质上运算次数仍为 $O(LD^3 d)$。

\begin{figure}
\centering\includegraphics[scale=1.0]{PyMPS_MatrixElement.eps}
\caption{态 $\ket{\phi}$ 与 $\ket{\psi}$ 之间的矩阵元按与重叠相同的方式计算，只需在正确位置插入算符，如箭头所示，在这些位置产生物理指标的双重求和。}
\label{fig:matrixelement}
\end{figure}

\begin{figure}
\centering\includegraphics[scale=1.0]{PyMPS_LocalOverlap.eps}
\caption{$\bra{\psi} \hat{O}^{[\ell]} \ket{\psi}$，其中态 $\ket{\psi}$ 在格点 $\ell$ 左侧的矩阵为左规范、右侧的矩阵为右规范。}
\label{fig:localoverlap}
\end{figure}

对期望值 $\bra{\psi} \hat{O}^{[\ell]} \ket{\psi}$ 可以做重要简化，应尽可能加以利用：假设我们考察局域算符 $\hat{O}^{[\ell]}$，且规范化使得格点 $\ell$ 左侧的所有矩阵都是左规范的、格点 $\ell$ 右侧的所有矩阵都是右规范的；格点 $\ell$ 本身的状态任意。于是可以利用左规范和右规范，像 图~\ref{fig:overlaplnormalized} 中那样无需显式计算地收缩网络，使得只剩下两个矩阵（图~\ref{fig:localoverlap}）。剩下的计算就是
\begin{equation}
\bra{\psi} \hat{O}^{[\ell]} \ket{\psi} = \sum_{\sigma_\ell \sigma'_\ell} O^{\sigma_\ell \sigma'_\ell} \tr
(M^{\sigma_\ell \dagger} M^{\sigma'_\ell}),
\end{equation} 
这是一个 $O(D^2 d^2)$ 量级的运算，在计算时间上节省了一个 $D$ 的量级。由于我们将会遇到态处于这种混合规范表示的算法，“在扫描中”计算可观测量是有意义的。这与 DMRG 中在显式格点上计算期望值完全相同。
    
计算重叠、期望值和矩阵元的一个常用记法，是把括号的层级结构解读为对一个矩阵的迭代更新，最终给出标量结果。现在我们引入矩阵 $C^{[\ell]}$，其中 $C^{[0]}$ 是哑矩阵，即标量 1。于是重叠 $\braket{\phi}{\psi}$ 可以通过让 $\ell$ 从 1 取到 $L$ 迭代地进行：
\begin{equation}
C^{[\ell]} = \sum_{\sigma_\ell} \tilde{M}^{\sigma_{\ell}\dagger} C^{[\ell-1]} M^{\sigma_{\ell}} ,
\end{equation}
其中 $C^{[L]}$ 又是一个标量，包含着结果。对于插入两态之间的算符，这一方法的自然推广是
\begin{equation}
C^{[\ell]} = \sum_{\sigma_\ell,\sigma'_\ell} O^{\sigma_\ell,\sigma'_\ell} \tilde{M}^{\sigma_{\ell}\dagger} C^{[\ell-1]} M^{\sigma'_{\ell}} .
\end{equation}
同样，求值矩阵乘积的正确次序对效率影响巨大：
\begin{equation}
C^{[\ell]}_{a_\ell,a'_\ell} = \sum_{\sigma_\ell,a_{\ell-1}}  \tilde{M}^{\sigma_{\ell}*}_{a_{\ell-1},a_\ell} \left( \sum_{\sigma'_\ell} O^{\sigma_\ell,\sigma'_\ell} 
\left( \sum_{a'_{\ell-1}} C^{[\ell-1]}_{a_{\ell-1},a'_{\ell-1}} 
M^{\sigma'_{\ell}}_{a'_{\ell-1},a'_\ell} \right) \right) 
\end{equation}
把一个 $O(D^4 d^2)$ 的运算约化为 $O(D^3 d) + O(D^2d^2) + O(D^3 d)$。

当然，我们也可以从右端出发，引入矩阵 $D^{[\ell]}$，从标量
$D^{[L]}=1$ 开始。为此，我们交换 $\braket{\phi}{\psi}$ 中标量的次序，
\begin{equation}
\braket{\phi}{\psi} = \sum_{\fat{\sigma}}M^{\sigma_1} \ldots M^{\sigma_L}  \tilde{M}^{\sigma_1*} \ldots \tilde{M}^{\sigma_L*}  ,
\end{equation}
并再次转置 $\tilde{M}$-矩阵，得到一个带括号求和的层级结构，其中对 $\sigma_L$ 的求和位于最内层。从 $\ell=L$ 跑到 $\ell=1$ 的迭代于是为：
\begin{equation}
D^{[\ell-1]} = \sum_{\sigma_\ell}  M^{\sigma_{\ell}} D^{[\ell]} \tilde{M}^{\sigma_{\ell}\dagger} ,
\end{equation}
可以把它推广用来计算矩阵元：
\begin{equation}
D^{[\ell-1]} = \sum_{\sigma_\ell,\sigma'_\ell}  O^{\sigma_\ell,\sigma'_\ell} M^{\sigma'_{\ell}} D^{[\ell]} \tilde{M}^{\sigma_{\ell}\dagger} .
\end{equation}
结果就包含在 $D^{[0]}$ 中。

这种方法对簿记非常有用，因为在 DMRG 中我们需要左块和右块的算符矩阵元，而这正是块 A 和块 B 的 $C$-矩阵与 $D$-矩阵所给出的内容。随着块迭代式地增长，上述矩阵序列也会随着块的增长方便地一同生成。

\subsubsection{转移算符与关联结构}
\label{subsubsec:transferoperator}
让我们把上一节中 $C^{[\ell]}$-矩阵的迭代构建再稍作形式化，因为引入一个转移（超）算符 $\hat{E}^{[\ell]}$ 有助于理解 MPS 中关联的本质，它是从定义在长度 $\ell-1$ 的块 A 上的算符到定义在长度 $\ell$ 的块 A 上的算符的一个映射，
\begin{equation}
\{ \ket{a_{\ell-1}} \bra{a'_{\ell-1}} \} \rightarrow \{ \ket{a_{\ell}} \bra{a'_{\ell}} \},
\end{equation}
其定义为
\begin{equation}
\hat{E}^{[\ell]} = \sum_{a_{\ell-1},a'_{\ell-1}} \sum_{a_\ell,a'_\ell} \left( \sum_{\sigma_\ell} M^{[\ell]\sigma_\ell *} \otimes M^{[\ell]\sigma_\ell} \right)_{(a_{\ell-1}a'_{\ell-1}),(a_\ell,a'_\ell)} (\ket{a_{\ell-1}} \bra{a'_{\ell-1}}) (\ket{a_{\ell}} \bra{a'_{\ell}}),
\end{equation}
其中括号中的表达式应读作 $E^{[\ell]}$ 的矩阵元，其维数为 $(D^2_{\ell-1} \times D^2_{\ell})$，即相应键上 $M$-矩阵的维数。它可以推广为在格点 $\ell$ 处插入一个算符后的收缩：
\begin{equation}
E^{[\ell]}_O = \sum_{\sigma_\ell,\sigma'_\ell} O^{\sigma_\ell,\sigma'_\ell} M^{[\ell]\sigma_\ell *} \otimes M^{[\ell]\sigma'_\ell} .
\end{equation}
$\hat{E}^{[\ell]} $ 是如何作用的？由显式写法
\begin{equation}
E^{[\ell]}_{(a_{\ell-1}a'_{\ell-1}),(a_{\ell}a'_{\ell})} = \sum_{\sigma_\ell} M^{[\ell]\sigma_\ell *}_{a_{\ell-1},a_\ell} \cdot M^{[\ell]\sigma_\ell}_{a'_{\ell-1},a'_\ell}
\end{equation}
我们可以把 $\hat{E}^{[\ell]}[ \hat{O}^{[\ell-1]}]$ 读作对矩阵 $O^{[\ell-1]}_{a,a'}$ 的一种操作，
\begin{equation}
E^{[\ell]} [O^{[\ell-1]}]  = \sum_{\sigma_\ell} M^{\sigma_\ell\dagger} O^{[\ell-1]} M^{\sigma_\ell}
\end{equation}
或者是对长度为 $D_{\ell-1}^2$、系数为 $v_{aa'} = O^{[\ell-1]}_{a,a'}$ 的行向量从左边乘的一种操作，
\begin{equation}
\sum_{\sigma_\ell} \sum_{a_{\ell-1},a'_{\ell-1}} v_{a_{\ell-1}a'_{\ell-1}} M^{\sigma_\ell*}_{a_{\ell-1},a_\ell} M^{\sigma_\ell}_{a'_{\ell-1},a'_\ell} .
\end{equation}
于是上一节的 $C$-矩阵之间有如下关系
\begin{equation}
C^{[\ell]} = E^{[\ell]} [C^{[\ell-1]}]  ,
\end{equation}
但我们现在要把这一结果用于数值方便之外的目的：若 $D_{\ell-1}=D_\ell$，我们还可以等价地谈论本征值、本征矩阵以及（左或右）本征矢。在此背景下我们得到 $E$ 最重要的性质：若它由左规范矩阵 $A$ 或右规范矩阵 $B$ 构成，则所有本征值满足 $|\lambda_k| \leq 1$。

事实上，当 $\lambda_1=1$ 且 $A$-矩阵为左规范时，相应的左本征矢为 $v_{aa'}=\delta_{aa'}$，这可以通过直接计算看出，或者平凡地把它翻译为单位矩阵：
\begin{equation}
E[I] = \sum_\sigma A^{\sigma\dagger} \cdot I \cdot A^{\sigma} = 1 \cdot I .
\end{equation}
由左规范 $A$ 矩阵构造的 $E$ 的右本征矢是非平凡的，但我们暂且忽略它。对于右规范的 $B$ 矩阵，情形正好反过来：显式计算表明 $v_{aa'}=\delta_{aa'}$ 此时是本征值 $\lambda_1=1$ 的右本征矢，而左本征矢是非平凡的。

为了证明 1 是最大本征值\cite{Weichselbaum05}，我们考虑 $C' = E[C]$。思路是：若 $E$ 由左规范或右规范矩阵构成，则可以证明 $C'$ 与 $C$ 的最大奇异值满足 $s'_1 \leq s_1$。这立即意味着 $E$ 的所有本征值满足 $|\lambda_i| \leq 1$：$C'=\lambda_i C$ 意味着 $s'_1 = |\lambda_i| s_1$，因此 $|\lambda_i|>1$ 将与前面的论断矛盾。无法排除存在多个 $|\lambda_i|=1$ 的情形。
证明进行如下（此处针对左规范矩阵）：考虑 SVD $C = U^\dagger S V$。$C$ 是方阵，故 $U^\dagger U = UU^\dagger = V^\dagger V = VV^\dagger = I$。于是我们可以写
\begin{equation}
C' = \sum_\sigma A^{\sigma\dagger} U^\dagger S V A^{\sigma} =
\left[ \begin{array}{ccc} (UA^1)^\dagger & \ldots  & (UA^d)^\dagger \end{array} \right]
\left[ \begin{array}{ccc} S & & \\ & S & \\ & & S \end{array} \right]
\left[ \begin{array}{c} VA^1 \\ \vdots \\ VA^d \end{array} \right] = P^\dagger \left[ \begin{array}{ccc} S & & \\ & S & \\ & & S \end{array} \right] Q .
\end{equation}
若 $A^\sigma$ 是左规范的，我们有 $P^\dagger P = I$ 和 $Q^\dagger Q = I$（但 $PP^\dagger \neq I$，$QQ^\dagger \neq I$）：$P^\dagger P =\sum_\sigma A^{\sigma\dagger} U^\dagger U A^\sigma = \sum_\sigma A^{\sigma\dagger} A^\sigma = I$，对 $Q$ 亦然；因此它们是到正交归一子空间的约化基底变换，故 $C'$ 的最大奇异值必小于等于 $S$ 的最大奇异值，也就是 $s_1$。

与归一化方式无关，重叠的计算变为
\begin{equation}
\braket{\psi}{\psi} = E^{[1]} E^{[2]} E^{[3]} \ldots E^{[L-2]} E^{[L-1]} E^{[L]}  ,
\end{equation}
而在用 $\braket{\psi}{\psi}$ 正确归一化之前的期望值计算则应读作
\begin{equation}
\bra{\psi} \hat{O}^{[1]} \hat{O}^{[2]} \ldots \hat{O}^{[L-1]} \hat{O}^{[L]}\ket{\psi} = E^{[1]}_{O_1} E^{[2]}_{O_2} E^{[3]}_{O_3} \ldots E^{[L-2]}_{O_{L-2}} E^{[L-1]}_{O_{L-1}} E^{[L]}_{O_L}  .
\end{equation}
在数值上，这种记法若按字面直接使用并不太有用，因为它意味着 $O(D^6)$ 的运算量；当然，若计及其内部的乘积结构，我们就回到前面讨论过的 $O(D^3)$ 运算量。但在解析上，它揭示了非常有趣的洞见。让我们考虑一个具有周期边界条件、由左规范且与格点无关的 $A$-矩阵构成的平移不变态（因而 $E$ 也与格点无关）。于是在极限 $L\rightarrow\infty$ 下我们得到
\begin{eqnarray*}
\bra{\psi} \hat{O}^{[i]}\hat{O}^{[j]} \ket{\psi} &=& \tr E^{[1]} \ldots E^{[i-1]} E^{[i]}_O E^{[i+1]} \ldots E^{[j-1]}
E^{[j]}_O E^{[j+1]} \ldots E^{[L]} \\
&=& \tr E^{[i]}_O E^{j-i-1} E^{[j]}_O E^{L-j+i-1} \\
&=& \sum_{l,k} \bra{l} E^{[i]}_O \ket{k} \lambda_k^{j-i-1} \bra{k} E^{[j]}_O \ket{l} \lambda_l^{L-j+i-1}\\
&=& \sum_k \bra{1} E^{[i]}_O \ket{k} \lambda_k^{j-i-1} \bra{k} E^{[j]}_O \ket{1} \quad (L\rightarrow\infty)
\end{eqnarray*}
其中 $\lambda_k$ 是 $E$ 的本征值。我们用到了由归一化矩阵构成的 $E$ 满足 $|\lambda_k|\leq 1$，以及 $\lambda_1=1$ 是唯一的模为 1 的本征值；但放宽后一（并不必然成立的）假设只会带来微小的修改。$\ket{k}$ 和 $\bra{k}$ 是（非厄米的）$E$ 对本征值 $\lambda_k$ 的右本征矢和左本征矢。$\bra{1}$ 对应于由左规范 $A$ 构成的 $E$ 的本征算符 $I$。

决定性的观察是：关联子既可以是长程的（若矩阵元 $ \bra{1} E^{[i]}_O \ket{1}$ 有限），也可以是衰减长度为 $\xi_k = -1/ \ln \lambda_k$ 的若干指数衰减的叠加，于是 MPS 的两点关联子取如下的一般形式
\begin{equation}
\frac{\bra{\psi} \hat{O}^{[i]} \hat{O}^{[j]} \ket{\psi}}{\braket{\psi}{\psi}} = c_1 + \sum_{k=2}^{D^2} c_k \eul^{-r/\xi_k} ,
\end{equation}
其中 $r=|j-i-1|$，且当 $i<j$ 时 $c_k = \bra{1} E^{[i]}_O \ket{k} \bra{k} E^{[j]}_O \ket{1}$。

一个简单的例子可以从上一节的 AKLT 态中提取。$E$ 的本征值此前已求得为 $1,-1/3,-1/3,-1/3$。对自旋算符而言，对应长程序的矩阵元为零，因此当 $j>i$ 时关联为 $\langle \hat{S}_i^z \hat{S}_j^z \rangle = (12/9)(-1)^{j-i} \eul^{-(j-i)\ln 3}$，是关联长度 $\xi = 1/\ln 3= 0.9102$ 的纯指数衰减。另一方面，对弦关联子而言，长程矩阵元是有限的，长程序在弦关联子中出现。

MPS 关联子的形式有重要的后果：在一维量子系统的热力学极限下，关联子通常取的形式要么是 Ornstein-Zernike 形式
\begin{equation}
\langle O_0 O_x \rangle \sim \frac{\eul^{-x/\xi}}{\sqrt{x}}
\end{equation}
要么是临界幂律形式（可能带有对数修正），
\begin{equation}
\langle O_0 O_x \rangle \sim x^{-\alpha} .
\end{equation}
AKLT 态属于一个非常特殊的态类，其关联函数模仿空间维数更低的量子系统（所谓的维数约化），从而去掉了 $\sqrt{x}$ 项；AKLT 态恰好位于所谓的无序线上，这类现象正是在那里发生 \cite{Garel96}。

因此，任何有限维的 MPS 都只能用指数衰减的叠加来逼近真实的关联子。事实证明，这在短距离上效果非常好，甚至对幂律也是如此。数值上观察到的是：随着 $D$ 的增大，真实关联函数会在越来越长的长度尺度上被准确表示。最终，最慢的指数衰减将存活下来，使关联变成带 $\xi =  -1/\ln \lambda$ 的纯指数衰减，其中 $\lambda$ 是对关联子有贡献的 $E$ 的最大本征值。因此，比较不同 $D$ 值的曲线是评估关联函数收敛性以及收敛已达成的长度尺度的一个极佳工具。

\subsubsection{MPS 与约化密度算符}
正如我们在 DMRG 中已经看到的，约化密度算符的概念在多个方面都很重要。让我们用 MPS 记法来表达它们。我们有
\begin{equation}
\ket{\psi}\bra{\psi} = \sum_{\fat{\sigma},\fat{\sigma}'} A^{\sigma_1} \ldots A^{\sigma_L} A^{\sigma'_1*} \ldots A^{\sigma'_L*} \ket{\fat{\sigma}} \bra{\fat{\sigma}'} =
\sum_{\fat{\sigma},\fat{\sigma}'} A^{\sigma_1} \ldots A^{\sigma_L} A^{\sigma'_L\dagger} \ldots A^{\sigma'_1\dagger} \ket{\fat{\sigma}} \bra{\fat{\sigma}'} .
\label{eq:fullstateprojector}
\end{equation}
我们现在把它二分为 A 和 B，其中 A 包含格点 1 至 $\ell$，B 包含格点 $\ell+1$ 至 $L$。在上式中对 B 的自由度求迹，我们得到
\begin{equation}
\hat{\rho}_A^{[\ell]} = \tr_B \ket{\psi}\bra{\psi} = \sum_{\fat{\sigma},\fat{\sigma}'\in A}
A^{\sigma_1} \ldots A^{\sigma_\ell} \rho_A^{[\ell]} A^{\sigma'_\ell\dagger} \ldots A^{\sigma'_1\dagger} \ket{\fat{\sigma}} \bra{\fat{\sigma}'} ,
\end{equation}
其中
\begin{equation}
\rho_A^{[\ell]} =  \sum_{\fat{\sigma}\in B} A^{\sigma_{\ell+1}} \ldots A^{\sigma_L} A^{\sigma_L\dagger} \ldots A^{\sigma_{\ell+1}\dagger} .
\end{equation}
这个方程立即给出不同约化密度矩阵之间的一个递推关系，即
\begin{equation}
\rho_A^{[\ell-1]} = \sum_{\sigma_\ell} A^{\sigma_\ell} \rho_A^{[\ell]} A^{\sigma_\ell\dagger} .
\label{eq:rhoarecursion}
\end{equation}
在平移不变系统的热力学极限 $L\rightarrow\infty$、$\ell\rightarrow\infty$ 下，我们可以由此追问是否满足如下不动点关系
\begin{equation}
\rho_A^f = \sum_{\sigma} A^{\sigma} \rho_A^{f} A^{\sigma\dagger}
\end{equation}
。

即使 MPS 的矩阵并非左规范，所有这些等式依然成立。若它们是左规范的，我们就可以直接在由 $A$ 矩阵生成的正交归一基中表达密度算符，即
\begin{equation}
\hat{\rho}_A^{[\ell]} = \sum_{a_\ell, a'_\ell} (\rho_A^{[\ell]})_{a_\ell, a'_\ell} \ket{a_\ell}_A \phantom{{}}_A\bra{a'_\ell} .
\end{equation}

对 B 的约化密度算符也有类似的关系式，其中（现在使用 $B$ 矩阵）我们得到
\begin{equation}
\hat{\rho}_B^{[\ell]} = \tr_A \ket{\psi}\bra{\psi} = \sum_{\fat{\sigma},\fat{\sigma}'\in B} 
B^{\sigma'_L\dagger} \ldots B^{\sigma'_{\ell+1}\dagger} \rho_B^{[\ell]} B^{\sigma_{\ell+1}} \ldots B^{\sigma_L} \ket{\fat{\sigma}} \bra{\fat{\sigma}'} ,
\end{equation}
其中
\begin{equation}
\rho_B^{[\ell]} =  \sum_{\fat{\sigma}\in A} B^{\sigma_{\ell}\dagger} \ldots B^{\sigma_1\dagger} B^{\sigma_1} \ldots B^{\sigma_{\ell} }
\end{equation}
以及递推关系
\begin{equation}
\rho_B^{[\ell]} = \sum_{\sigma_\ell} B^{\sigma_\ell\dagger} \rho_B^{[\ell-1]} B^{\sigma_\ell} .
\label{eq:rhobrecursion}
\end{equation}
由此产生一个潜在的不动点关系
\begin{equation}
\rho_B^f = \sum_{\sigma} B^{\sigma\dagger} \rho_B^{f} B^{\sigma} .
\end{equation}
同样，所有这些关系对任意的 MPS 矩阵都成立，但如果它们是右规范的，我们又会得到一个在正交归一基中的表达式，此时该基由 $B$ 矩阵生成，
\begin{equation}
\hat{\rho}_B^{[\ell]} = \sum_{a_\ell, a'_\ell} (\rho_B^{[\ell]})_{a_\ell, a'_\ell} \ket{a_\ell}_B \phantom{{}}_B\bra{a'_\ell} .
\end{equation}
对于混合规范形式的态$\ket{\psi}=\sum_{\fat{\sigma}} A^{\sigma_1}\ldots A^{\sigma_\ell}\Psi B^{\sigma_{\ell+1}} \ldots B^{\sigma_L} \ket{\fat{\sigma}}$，重新进行一遍计算表明
\begin{equation}
\hat{\rho}^{[\ell]}_A = \Psi \Psi^\dagger
\end{equation}
而
\begin{equation}
\hat{\rho}^{[\ell]}_B = \Psi^\dagger\Psi ,
\end{equation}
均在正交归一基中表达。

\subsection{两个矩阵乘积态的相加}

在实践中相对较少用到的一种操作是两个 MPS 的相加。我们先考虑较容易的 PBC 情形。取两个不作归一化假设的 MPS，
\begin{equation}
\ket{\psi} = \sum_{\fat{\sigma}} \tr (M^{\sigma_1} \ldots M^{\sigma_L}) \ket{\fat{\sigma}} 
\quad\quad
\ket{\phi} = \sum_{\fat{\sigma}} \tr (\tilde{M}^{\sigma_1} \ldots \tilde{M}^{\sigma_L}) \ket{\fat{\sigma}} 
\end{equation}
我们可以写出
\begin{equation}
\ket{\psi} + \ket{\phi} = \sum_{\fat{\sigma}} \tr (N^{\sigma_1} \ldots N^{\sigma_L}) \ket{\fat{\sigma}}
\end{equation}
其中
\begin{equation}
N^{\sigma_i} =  M^{\sigma_i} \oplus \tilde{M}^{\sigma_i} .
\end{equation} 
这意味着我们只是把 $M$ 和 $\tilde{M}$ 取为矩阵 $N$ 的对角块。对角块结构意味着，将各 $N$ 矩阵相乘后结果仍是对角的，第一块中是 $MMMMM\ldots$，第二块中是 $\tilde{M}\tilde{M}\tilde{M}\tilde{M}\tilde{M}\ldots$。于是迹可以拆开，我们又回到了原来的两个态：
\begin{equation}
\tr (NNNNNN) = \tr \left( \begin{array}{cc} MMMMMM & 0 \\ 0 & \tilde{M}\tilde{M}\tilde{M}\tilde{M}\tilde{M}\tilde{M} \end{array} \right) = \tr (MMMMMM) + \tr (\tilde{M}\tilde{M}\tilde{M}\tilde{M}\tilde{M}\tilde{M}).
\end{equation} 

在 OBC 情形下，我们可以完全照搬同样的做法。在第一个和最后一个格点上必须有特殊处理：若不加斟酌，首末两处的维数会升到 2，标量性质随之丧失。但在物理上这些指标本来就是无足轻重的哑指标。因此我们要做的（一个简单的计算表明这确实可行）是：利用原来两个态的行向量与列向量，在首末格点处构造出行向量 $[ M\ \tilde{M}]$ 和列向量 $[M\ \tilde{M}]^T$。

因此，MPS 相加会得到维数为 $D_N = D_M + D_{\tilde{M}}$ 的新矩阵，也就是说具有一定维数的 MPS 在加法下并不封闭。同样显而易见的是，在许多情况下这种描述新态的方式是不经济的：极端情形是相加 $\ket{\psi}+\ket{\psi}$，所得的态其实与原态相同，只是带一个前因子 2，矩阵维数本不应增大。所以在相加之后，值得考虑把 MPS 重新{\em 压缩}到较低的维数；视被加的态而定，这一步可能（也可能不，如上例那样）造成信息损失。

\subsection{把矩阵乘积态化为规范形式}

对于一般的矩阵乘积态，除了维数必须恰当匹配之外，对矩阵 $M^{\sigma_i}$ 并无特别的要求。有几类矩阵是值得优先采用的，即左规范和右规范矩阵，它们分别给出左规范和右规范的 MPS：某些收缩变得平庸，正交归一的约化基自动生成。

为了把任意 MPS 化为规范形式，我们利用这样一个事实：SVD 产生的或者是幺正矩阵，或者是行、列正交归一的矩阵，可以证明后者满足左规范或右规范条件。

\subsubsection{左规范 MPS 的生成}
从一个不作归一化假设的一般 MPS 出发，把收缩显式地写出来，
\begin{equation}
\ket{\psi}= \sum_{\fat{\sigma}} \sum_{a_1,\ldots} M^{\sigma_ 1}_{1,a_1} M^{\sigma_2}_{a_1,a_2} M^{\sigma_3}_{a_2,a_3} \ldots \ket{\fat{\sigma}}
\end{equation}
我们将物理指标与左（行）指标合并，对 $M^{\sigma_ 1}_{1,a_1}$ 进行重排，然后对新矩阵 $M$ 作 SVD，得到 $M=ASV^\dagger$：
\begin{eqnarray}
& & \sum_{\fat{\sigma}} \sum_{a_1,\ldots} M_{(\sigma_ 1,1),a_1} M^{\sigma_2}_{a_1,a_2} M^{\sigma_3}_{a_2,a_3} \ldots \ket{\fat{\sigma}} \nonumber \\
&=&  \sum_{\fat{\sigma}} \sum_{a_1,\ldots} \sum_{s_1} A_{(\sigma_ 1,1),s_1} S_{s_1,s_1} V^\dagger_{s_1,a_1}  M^{\sigma_2}_{a_1,a_2}  \ldots \ket{\fat{\sigma}} \nonumber \\
&=& \sum_{\fat{\sigma}} \sum_{a_2,\ldots} \sum_{s_1} A^{\sigma_1}_{1,s_1} \left( \sum_{a_1} S_{s_1,s_1} V^\dagger_{s_1,a_1}  M^{\sigma_2}_{a_1,a_2} \right) M^{\sigma_3}_{a_2,a_3} \ldots \ket{\fat{\sigma}} \nonumber \\
&=& \sum_{\fat{\sigma}} \sum_{a_2,\ldots} \sum_{s_1} A^{\sigma_1}_{1,s_1} \tilde{M}^{\sigma_2}_{s_1,a_2} M^{\sigma_3}_{a_2,a_3} \ldots \ket{\fat{\sigma}} .
\end{eqnarray}
由于 SVD 保证 $A^\dagger A = I$，重排为 $A^{\sigma_1}$ 之后，$A^{\sigma_1}$ 满足左规范条件。SVD 其余的两个矩阵则被乘入 $M^{\sigma_2}$，于是生成一个新的 MPS，其中
$\tilde{M}^{\sigma_2}_{s_1,a_2} = \sum_{a_1} S_{s_1,s_1} V^\dagger_{s_1,a_1}  M^{\sigma_2}_{a_1,a_2}$ is generated. 

现在这一步骤可以迭代进行：把 $\tilde{M}^{\sigma_2}_{s_1,a_2}$ 重排为 $\tilde{M}_{(\sigma_2,s_1),a_2}$（图~\ref{fig:Mreshaping}），作奇异值分解 $ASV^\dagger$，生成 $A_{(\sigma_2,s_1),s_2}$，再重排为左规范的 $A^{\sigma_2}_{s_1,s_2}$。SVD 右边的两个矩阵又被乘入下一个待分解的矩阵，如此继续。最后一步之后，所有格点上都驻留着左规范矩阵 $A^{\sigma_i}_{s_{i-1},s_i}$。$S_{1,1} (V^\dagger)_{1,1}$——由于 $A^{\sigma_L}$ 是列向量，它是一个标量——留在最后一个格点上，但这个标量不过是 $\ket{\psi}$ 的范数。如果想使用非归一化的态，可以把它单独保存。

\begin{figure}
\centering\includegraphics[scale=0.6]{PyMPS_MatrixReshaping.eps}
\caption{为实现规范化，把同一格点上的矩阵集合合并为一个矩阵。}
\label{fig:Mreshaping}
\end{figure}
    
这一过程可以轻而易举地推广到归一化：我们找出态的前因子，但不存储它，而是简单地把它设为 1。由于我们并不显式地使用奇异值，上述过程可以按照第 \ref{subsubsec:MPSdecomposition} 节的思路，用 QR 分解轻松地重新表述。然而标准 QR 分解无法告诉我们所用矩阵是否大于必要的尺寸，即是否含有为零的奇异值，从而可以把矩阵裁剪到更小的尺寸而不损失精度；要做到这一点只能使用秩显现（rank revealing）QR 分解；不过对许多 MPS 而言，从其背后的物理可以清楚地知道奇异值谱具有长尾，因此这个问题不会出现。同样的论证也适用于右规范 MPS 的生成，下面就来讨论。
    
\subsubsection{右规范 MPS 的生成}
同样的步骤也可用来得到由右规范矩阵构成的态：对重排后的矩阵 $\{ M^\sigma \} \rightarrow M=USB$ 从右端开始依次作一列 SVD，把 $B$ 拆分成右规范的矩阵 $B^\sigma$（由于 $BB^\dagger = I$），并把 $U$ 和 $S$ 乘到左边，构成下一个待作奇异值分解的矩阵：

\begin{eqnarray}
& & \sum_{\fat{\sigma}} \sum_{\ldots,a_{L-1}}  \ldots M^{\sigma_{L-2}}_{a_{L-3},a_{L-2}} M^{\sigma_{L-1}}_{a_{L-2},a_{L-1}} M^{\sigma_{L}}_{a_{L-1},1}  \ket{\fat{\sigma}} \nonumber \\
&=& \sum_{\fat{\sigma}} \sum_{\ldots,a_{L-1}}  \ldots M^{\sigma_{L-2}}_{a_{L-3},a_{L-2}} M^{\sigma_{L-1}}_{a_{L-2},a_{L-1}} M_{a_{L-1},(\sigma_{L},1)}  \ket{\fat{\sigma}} \nonumber \\
&=&  \sum_{\fat{\sigma}} \sum_{\ldots a_{L-1}} \sum_{s_{L-1}}  \ldots M^{\sigma_{L-2}}_{a_{L-3},a_{L-2}} M^{\sigma_{L-1}}_{a_{L-2},a_{L-1}} 
U_{a_{L-1},s_{L-1}} S_{s_{L-1},s_{L-1}} B_{s_{L-1},(\sigma_L,1)}   \ket{\fat{\sigma}} \nonumber \\
&=& \sum_{\fat{\sigma}} \sum_{\ldots a_{L-2}} \sum_{s_{L-1}} 
\ldots M^{\sigma_{L-2}}_{a_{L-3},a_{L-2}} \left( \sum_{a_{L-1}} M^{\sigma_{L-1}}_{a_{L-2},a_{L-1}} U_{a_{L-1},s_{L-1}} S_{s_{L-1},s_{L-1}}\right) 
B_{s_{L-1},(\sigma_L,1)}   \ket{\fat{\sigma}} \nonumber \\
&=& \sum_{\fat{\sigma}} \sum_{\ldots a_{L-2}} \sum_{s_{L-1}}
\ldots M^{\sigma_{L-2}}_{a_{L-3},a_{L-2}} \tilde{M}^{\sigma_{L-1}}_{a_{L-2},s_{L-1}} B^{\sigma_L}_{s_{L-1},1}   \ket{\fat{\sigma}} ,
\end{eqnarray}
按照前面的方式继续进行，仅有的差别是：(i) 方向相反；(ii) 重排时现在把物理指标与列指标（而不是行指标）合并：$\tilde{M}^{\sigma_i}_{a_{i-1},s_{i}} \rightarrow \tilde{M}_{a_{i-1},(\sigma_i s_i)} \rightarrow B_{s_{i-1},(\sigma_i s_i)} \rightarrow B^{\sigma_i}_{s_{i-1},s_i}$。
    
    
\subsection{MPS 的近似压缩}

MPS 的（少见的）相加以及各种能够用 MPS 表述的算法，会使结果的矩阵维数大幅增加。因此，一个反复出现的问题就是：如何把一个矩阵维数为 $(D'_i \times D'_{i+1})$ 的给定 MPS，
尽可能好地用另一个矩阵维数为 $(D_i \times D_{i+1})$（其中 $D_i < D'_i$）的 MPS 来近似。

原则上存在两种做法：{\em SVD 压缩}和{\em 变分压缩}。二者各有优劣：当压缩程度不大（$D \sim D'$）时，SVD 很快，但绝非最优；若 $D' \gg D$，它会变得非常慢。变分压缩是最优的，但若起点随机选取则很慢；不过，如果提供一个好的试探态（例如来自 SVD 方法），它可以大大加速。一般说来，变分拟设中可能出现陷入非最优压缩的问题。
 
根据待压缩态的具体性质，可以把步骤进一步优化，例如当待压缩的 MPS 是若干 MPS 之和，或者是把矩阵乘积算符（MPO；第~\ref{sec:MPO}~节）作用到某个 MPS 上的结果时。

\subsubsection{用 SVD 压缩矩阵乘积态}
考虑一个处于混合规范表示中的 MPS，
\begin{equation}
\ket{\psi} = \sum_{\fat{\sigma}} A^{\sigma_1} A^{\sigma_2} \ldots A^{\sigma_{\ell}} \Lambda^{[\ell]} B^{\sigma_{\ell+1}} \ldots B^{\sigma_{L-1}} B^{\sigma_L}  \ket{\fat{\sigma}} ,
\label{eq:compressionstart}
\end{equation}
从中读出 Schmidt 分解 $\ket{\psi} = \sum_{a_\ell=1}^{D'} s_{a_\ell} \ket{a_\ell}_A \ket{a_\ell}_B$，其中 A 与 B 上的态各自构成正交归一集合；这由规范构造直接得出。设此分解双方各有 $D'$ 个态。我们现在寻找在 2-范数意义下最好地逼近 $\ket{\psi}$、且在 A 与 B 中各能由 $D$ 个态张成的态 $\ket{\tilde{\psi}}$。前面已经证明，SVD 通过保留 $D$ 个最大的奇异值给出结果，压缩后的态简单地为 $\ket{\psi} = \sum_{a_\ell=1}^{D} s_{a_\ell} \ket{a_\ell}_A \ket{a_\ell}_B$，由此得到一个简单的截断处方：保留 $A^{\sigma_{\ell}}$ 的前 $D$ 列、$B^{\sigma_{\ell+1}}$ 的前 $D$ 行，以及 $\Lambda^{[\ell]}$ 的前 $D$ 行和列。若需要归一化，保留的奇异值必须重新标度。

这一步骤依赖于 A 与 B 上态的正交归一性，因此只能在单个键上执行。为了缩小所有矩阵，我们必须遍历所有混合规范表示——比如从右到左——进行截断，并借助规范化技术把左规范与右规范矩阵之间的边界向左移动一个格点。

对一个左规范的态做右规范化的第一步之后，它变为：
\begin{equation}
\ket{\psi^{(L-1)}} = \sum_{\fat{\sigma}} A^{\sigma_1} \ldots A^{\sigma_{L-1}} U S B^{\sigma_L} \ket{\fat{\sigma}} ,
\end{equation}
其中我已经把 $B$ 重排成形，它是右规范的，并保证形如 $\ket{a_{L-1}}_B = \sum_{\sigma_L} (B^{\sigma_L})_{a_{L-1},1} \ket{\sigma_L}$ 的态是正交归一的。但如下这些态也是正交归一的
\begin{equation}
\ket{a_{L-1}}_A = \sum_{\sigma_1 \ldots \sigma_{L-1}} (A^{\sigma_1} \ldots A^{\sigma_{L-1}} U)_{1,a_{L-1}} \ket{\sigma_1 \ldots \sigma_{L-1}},
\end{equation}
因为 SVD 保证 $U^\dagger U = 1$：我们只是在由左规范的 $A^{\sigma_i}$ 所构造的正交归一基集合之内做一次基变换。于是我们得到一个正确的 Schmidt 分解
\begin{equation}
\ket{\psi^{(L-1)}} = \sum_{a_{L-1}} s_{a_{L-1}} \ket{a_{L-1}}_A \ket{a_{L-1}}_B .
\end{equation}
与右规范化的区别现在在于截断：矩阵 $U$、$S$ 和 $B^{\sigma_L}$ 按照前面的说明被截断（奇异值可能重新归一化）为 $\tilde{U}$、$\tilde{S}$ 和 $\tilde{B}^{\sigma_L}$：保留 $D$ 个最大的奇异值。$\tilde{B}^{\sigma_L}$ 仍是右规范的。左边下一个 $A^{\sigma_{L-1}}$ 与 $\tilde{U}$、$\tilde{S}$ 相乘构成 $M^{\sigma_{L-1}}$。通过如下的重排、SVD、再重排
\begin{equation}
M^{\sigma}_{ij} = M_{i, (\sigma j)} = \sum_k U_{ik} S_{kk} B_{k,(\sigma j)} = \sum_k U_{ik} S_{kk} B^{\sigma}_{kj} 
\end{equation} 
我们得到右规范的 $B^{\sigma_{L-1}}$，把 $U$、$S$ 和 $B^{\sigma_{L-1}}$ 截断为 $\tilde{U}$、$\tilde{S}$ 和 $\tilde{B}^{\sigma_{L-1}}$，然后步骤继续：
\begin{eqnarray*}
\ket{\psi} &=& \sum_{\fat{\sigma}}  A^{\sigma_1} \ldots A^{\sigma_{L-3}}A^{\sigma_{L-2}} \left( A^{\sigma_{L-1}} U S \right)B^{\sigma_L}   \ket{\fat{\sigma}} \\
&\rightarrow& \sum_{\fat{\sigma}} A^{\sigma_1} \ldots A^{\sigma_{L-3}}A^{\sigma_{L-2}}\left( A^{\sigma_{L-1}} \tilde{U} \tilde{S} \right) \tilde{B}^{\sigma_L}   \ket{\fat{\sigma}} \\
&=&  \sum_{\fat{\sigma}} A^{\sigma_1} \ldots A^{\sigma_{L-3}}A^{\sigma_{L-2}} M^{\sigma_{L-1}} \tilde{B}^{\sigma_L}   \ket{\fat{\sigma}} \\
&=& \sum_{\fat{\sigma}} A^{\sigma_1} \ldots A^{\sigma_{L-3}} \left( A^{\sigma_{L-2}} U S \right) B^{\sigma_{L-1}} \tilde{B}^{\sigma_L}   \ket{\fat{\sigma}} \\
&\rightarrow& \sum_{\fat{\sigma}} A^{\sigma_1} \ldots A^{\sigma_{L-3}} \left( A^{\sigma_{L-2}} \tilde{U} \tilde{S} \right) \tilde{B}^{\sigma_{L-1}} \tilde{B}^{\sigma_L}   \ket{\fat{\sigma}} \\
&=& \ldots
\end{eqnarray*}
最终，压缩后的 MPS $\ket{\tilde{\psi}}$ 是右规范的，由 $\tilde{B}$ 矩阵给出。后面会看到，这一压缩步骤正是（含时）DMRG 或 TEBD 扫描经过链时所执行的截断。这两种方法在每个键的左右两侧都拥有正确归一化的矩阵（即正交归一态），执行切割后继续前进。

该方法的缺点是会出现截断之间的单侧相互依赖：矩阵 $M$ 总是包含来自上一步的截断后的 $\tilde{U}$，因此依赖于截断。一般而言，各次截断不可能彼此独立，因为对每个分解 (\ref{eq:compressionstart})，$A$ 矩阵和 $B$ 矩阵的截断都会影响正交归一系；但这里的依赖是单侧且``不平衡''的：越靠左的截断依赖于越靠右的截断，反之则不然。如果截断很小——在含时 DMRG 中对小时间步长通常如此——引入的额外不精确性是次要的；但在可能出现大截断的情形中，这种依赖可能变得过强，截断也就远非最优。

第二个问题关乎效率：对维数为 $(m\times n)$、$m\ge n$ 的矩阵，SVD 的代价是 $O(m n^2)$。这意味着当 $D' \le dD$ 时 SVD 代价为 $O((D')^2 dD)$，否则为 $O(D' d^2 D^2)$；矩阵乘法的代价是 $O(dD(D')^2)$。在许多应用中 $D'\gg D$；此时该方法变得相当慢。若原始态不在规范形式、必须先通过一串阶为 $O((D')^3)$ 的 SVD 化为规范形式，情况还要更糟。

最后以一点评论结束本节：当然，我们也可以通过施加某个 $\epsilon$ 来压缩——在每次截断时把它接受为原态与压缩态之间 2-范数距离的最大值（由被丢弃奇异值的平方和给出），从而隐式地定义 $D$。

\subsubsection{迭代压缩矩阵乘积态}
最优的做法是从一个具有期望的约化维数的拟设 MPS 出发，迭代地极小化它与待近似 MPS 的距离，也就是说，通过迭代地改动它的 $\tilde{M}^\sigma$ 矩阵。这些矩阵扮演变分参数的角色。

把 $\ket{\psi}$ 从维数 $D'$ 最优压缩为维数 $D$ 的 $\ket{\tilde{\psi}}$，其数学上的精确表述是极小化 $\| \ket{\psi} - \ket{\tilde{\psi}} \|_2^2$，也就是说，我们要关于 $\ket{\tilde{\psi}}$ 极小化 $\braket{\psi}{\psi} - \braket{\tilde{\psi}}{\psi} -\braket{\psi}{\tilde{\psi}} + \braket{\tilde{\psi}}{\tilde{\psi}}$。把这两组矩阵分别记为 $M$ 和 $\tilde{M}$，以强调我们对归一化不作任何假设。用底层的矩阵 $\tilde{M}$ 表出时，这是一个高度非线性的优化问题。

但这可以按如下方式迭代进行。先给 $\ket{\tilde{\psi}}$ 一个初始猜测，它可以取 $\ket{\psi}$ 的 SVD 压缩，虽未必最优，却是一个不错的起点。然后我们逐格点扫描 $\tilde{M}^{\sigma_i}$ 集合：保持所有其他矩阵固定，选取新的 $\tilde{M}^{\sigma_i}$ 使距离极小。（通常可以指望的是）把这样的矩阵扫描重复若干次，就会收敛到最优近似。

新的 $\tilde{M}^{\sigma_i}$ 由对 $\tilde{M}^{\sigma_i*}_{a_{i-1},a_{i}}$ 取极值得到，它只出现在 $ - \braket{\tilde{\psi}}{\psi} + \braket{\tilde{\psi}}{\tilde{\psi}}$ 中。我们求得
\begin{eqnarray*}
& & \frac{\partial}{\partial \tilde{M}^{\sigma_i*}_{a_{i-1},a_i}} ( \braket{\tilde{\psi}}{\tilde{\psi}}
- \braket{\tilde{\psi}}{\psi} ) = \\
& & \sum_{\fat{\sigma}*} (\tilde{M}^{\sigma_1*} \ldots \tilde{M}^{\sigma_{i-1}*})_{1, a_{i-1}}(\tilde{M}^{\sigma_{i+1}*} \ldots \tilde{M}^{\sigma_L*})_{a_i,1} \tilde{M}^{\sigma_1} \ldots \tilde{M}^{\sigma_i} \ldots \tilde{M}^{\sigma_L} -  \\
& & \sum_{\fat{\sigma}*}
(\tilde{M}^{\sigma_1*} \ldots \tilde{M}^{\sigma_{i-1}*})_{1, a_{i-1}}(\tilde{M}^{\sigma_{i+1}*} \ldots \tilde{M}^{\sigma_L*})_{a_i,1} M^{\sigma_1} \ldots M^{\sigma_i} \ldots M^{\sigma_L}= 0.
\end{eqnarray*}
对 $\fat{\sigma}*$ 的求和跑遍除 $i$ 以外的所有物理格点。这个方程组看起来复杂，实际上相当容易。把待求的矩阵 $\tilde{M}^{\sigma_i}$ 显式地分离出来，可以把这个方程改写为
\begin{equation}
\sum_{a'_{i-1}a'_i} \tilde{O}_{a_{i-1}a_i,a'_{i-1}a'_i} \tilde{M}^{\sigma_i}_{a'_{i-1}a'_i} = O^{\sigma_i}_{a_{i-1}a_i}.
\end{equation}
如果对每个 $\sigma_i$，我们把矩阵 $\tilde{M}^{\sigma_i}$ 解释为长度 $D^2$ 的向量 $v$，把 $\tilde{O}$ 解释为维数 $(D^2 \times D^2)$ 的矩阵 $P$，把 $O^{\sigma_i}$ 解释为长度 $D^2$ 的向量 $b$，就得到一个线性方程组
\begin{equation}
P v = b. 
\end{equation}
所得的 $v$ 即可取作我们要找的矩阵。由于这个方程组通常太大而无法直接求解，必须使用迭代求解器，例如共轭梯度法。该方程组是厄米的，从 $P$ 的构造可以看出：非共轭与共轭的 $\tilde{M}$ 在转置下恰好互换角色。图示表示再次最为简单（图~\ref{fig:iterativecompression}）。

\begin{figure}
\centering\includegraphics[scale=0.6]{PyMPS_MPS_IterativeCompression.eps}
\caption{MPS 迭代压缩时需求解的线性方程组。较粗的线对应待压缩的态，较细的线对应压缩后的态。未知矩阵被圈出。}
\label{fig:iterativecompression}
\end{figure}
    

正如后面变分求基态的情形一样，必须认识到共轭梯度法中矩阵--向量乘法的代价并不是按维数天真估计的 $O(D^4)$。存在一个明显的因式分解，它在图示中格外一目了然，可以把代价降到 $O(D^3)$：
\begin{equation}
\tilde{O}_{a_{i-1}a_i,a'_{i-1}a'_i} = \tilde{L}_{a_{i-1},a'_{i-1}} \cdot \tilde{R}_{a_{i},a'_{i}}
\end{equation}
其中
\begin{equation}
\tilde{L}_{a_{i-1},a'_{i-1}} = \left( \sum_{\sigma_{i-1}}  \tilde{M}^{\sigma_{i-1}\dagger} \left( \ldots \left( \sum_{\sigma_1}  
\tilde{M}^{\sigma_{1}\dagger} \tilde{M}^{\sigma_{1}} \right) \ldots \right) \tilde{M}^{\sigma_{i-1}} \right)_{a_{i-1},a'_{i-1}}
\end{equation}
对 $\tilde{R}$ 亦类似。在图示表示（图~\ref{fig:normalizediterativecompressionLR}）中，它们不过是待解的带圈 $\tilde{M}$ 矩阵左右两边已缩并的对象。

于是
\begin{equation}
\sum_{a'_{i-1}a'_i} \tilde{O}_{a_{i-1}a_i,a'_{i-1}a'_i} \tilde{M}^{\sigma_i}_{a'_{i-1}a'_i} =
\sum_{a'_{i-1}} \tilde{L}_{a_{i-1},a'_{i-1}} \left( \sum_{a'_i} \tilde{R}_{a_{i},a'_{i}}
 \tilde{M}^{\sigma_i}_{a'_{i-1}a'_i} \right),
\end{equation}
两次代价为 $O(D^3)$ 的操作。

一个类似的分解简化了向量 $b$ 的计算，它由 $O^{\sigma_i}_{a_{i-1}a_i}$ 按下式构成：
\begin{equation}
\sum_{a'_{i-1}a'_i} L_{a_{i-1},a'_{i-1}} M^{\sigma_i}_{a'_{i-1}a'_i} R_{a_{i},a'_{i}} , 
\end{equation}
其中 $L$ 和 $R$ 如图~\ref{fig:normalizediterativecompression} 所示。
事实上，计算 $L$ 和 $R$ 无非就是从左端或右端开始执行重叠（overlap）计算的头几步，其结果正是该计算在中间步骤产生的 $C$ 矩阵。如果从左到右再折返地扫描整个系统，就可以利用前面各步迭代地构建 $L$ 和 $R$，这是最有效率的方式。

然而，如果我们利用规范形式，就可以{\em 大幅}简化压缩过程！假设 $\ket{\tilde{\psi}}$ 处于混合规范形式 $\tilde{A}^{\sigma_1} \ldots \tilde{A}^{\sigma_{i-1}} \tilde{M}^{\sigma_i} \tilde{B}^{\sigma_{i+1}}  \ldots \tilde{B}^{\sigma_L}$，我们想要更新 $\tilde{M}^{\sigma_i}$：在待更新矩阵的左边，一切都已左规范化，右边一切都已右规范化，而 $\tilde{M}^{\sigma_i}$ 的形式无关紧要，因为它反正会被重新计算。

于是，由于左规范化，$\tilde{L}_{a_{i-1},a'_{i-1}} = \delta_{a_{i-1},a'_{i-1}}$；类似地，由于右规范化，$\tilde{R}_{a_{i},a'_{i}} = \delta_{a_{i},a'_{i}}$，因此 
$\tilde{O}_{a_{i-1}a_i,a'_{i-1}a'_i} = \delta_{a_{i-1},a'_{i-1}} \delta_{a_{i},a'_{i}}$。在线性方程组中这意味着 $P=I$，于是我们平凡地得到
\begin{equation}
v = b,
\end{equation}
所以无需求解大型方程组（图~\ref{fig:normalizediterativecompression}）。

\begin{figure}
\centering\includegraphics[scale=0.6]{PyMPS_MPS_NormalizedIterativeCompression.eps}
\caption{对经过适当归一化的态作 MPS 迭代压缩的方程。较粗的线对应待压缩的态，较细的线对应压缩后的态。}
\label{fig:normalizediterativecompression}
\end{figure}

\begin{figure}
\centering\includegraphics[scale=0.6]{PyMPS_MPS_IterativeCompressionLRObjects.eps}
\caption{压缩中迭代构造的对象 $\tilde{L}$ 和 $\tilde{R}$。}
\label{fig:normalizediterativecompressionLR}
\end{figure}


要让这一做法适用于整条链，我们必须在沿链移动时移动左规范化矩阵与右规范化矩阵之间的分界。假设开始时所有矩阵都是左规范化的。然后我们从右到左穿过链，先在最后一个格点上求解 $\tilde{M}^{\sigma_L}$，像前面一样通过 SVD（或者，由于我们不需要奇异值，用更廉价的 QR）将其右规范化，得到 $\tilde{A}^{\sigma_1} \ldots \tilde{A}^{\sigma_{L-2}} \tilde{M}^{\sigma_{L-1}} \tilde{B}^{\sigma_{L}}$，其中 $\tilde{M}^{\sigma_{L-1}}$ 一般而言未归一化。它现在按如下方式优化
\begin{equation}
\tilde{M}^{\sigma_i}_{a_{i-1},a_i} =  \sum_{a'_{i-1}} L_{a_{i-1},a'_{i-1}} \left( \sum_{a'_i} R_{a_{i},a'_{i}}
M^{\sigma_i}_{a'_{i-1},a'_i} \right), 
\label{eq:compressionrhs}
\end{equation}
其中 
\begin{equation}
L_{a_{i-1},a'_{i-1}} = \left( \sum_{\sigma_{i-1}}  \tilde{M}^{\sigma_{i-1}\dagger} \left( \ldots \left( \sum_{\sigma_1}  
\tilde{M}^{\sigma_{1}\dagger} M^{\sigma_{1}} \right) \ldots \right) M^{\sigma_{i-1}} \right)_{a_{i-1},a'_{i-1}}
\end{equation}
类似地处理 $R$，结果是右规范化的，如此这般沿链继续。到最后，所有矩阵都右规范化了，我们再从左边重新开始。

为了评估收敛性，我们可以在每一步监测 $\| \ket{\psi} - \ket{\tilde{\psi}} \|^2$ 并观察该值的收敛情况；必要时必须增大 $D$。这个计算看似代价高昂，实则不然。如果我们让 $\ket{\tilde{\psi}}$ 保持恰当的混合归一化，并利用式~(\ref{eq:compressionrhs}) 化简重叠 $\braket{\psi}{\tilde{\psi}}$，就会发现
\begin{equation}
\| \ket{\psi} - \ket{\tilde{\psi}} \|^2 = 1 - \sum_{\sigma_i} \tr (\tilde{M}^{\sigma_i\dagger}\tilde{M}^{\sigma_i}) , 
\end{equation}
该式易于计算。被减去的求和项正是 $\braket{\tilde{\psi}}{\tilde{\psi}}$；最后，这使我们能够通过简单的重新标度来归一化态 $\ket{\tilde{\psi}}$。

正如单格点 DMRG 中已暗示的那样——我们将在第~\ref{sec:groundstates} 节详细讨论这一问题——这种变分拟设存在陷入非全局极小值的危险，即压缩态与原态之间距离的非全局极小。通常（但不总是）仿照两格点 DMRG，同时考虑两个格点进行优化会有所帮助。运算量统计表明这会稍慢一些。假设压缩后的 $\ket{\tilde{\psi}}$ 处于混合规范表示
\begin{equation}
\ket{\tilde{\psi}} = \sum_{{\fat \sigma}} \tilde{A}^{\sigma_1} \ldots \tilde{A}^{\sigma_{\ell-1}} \tilde{M}^{\sigma_{\ell}\sigma_{\ell+1}} \tilde{B}^{\sigma_{\ell+2}} \ldots \tilde{B}^{\sigma_L}  \ket{{\fat \sigma}}.
\end{equation}
重复前面同样的论证，对 $\tilde{M}^{\sigma_{\ell}\sigma_{\ell+1}*}_{a_{\ell-1},a_{\ell+1}}$ 作优化，便得到如图~\ref{fig:twositecompression} 所示关于 $\tilde{M}^{\sigma_{\ell}\sigma_{\ell+1}}_{a_{\ell-1},a_{\ell+1}}$ 的方程。现在发生主要变化：我们把新的 $\tilde{M}^{\sigma_{\ell}\sigma_{\ell+1}}$ 重塑为 $\tilde{M}_{(a_{\ell-1}\sigma_{\ell}),(\sigma_{\ell+1}a_{\ell+1})}$，作一次 SVD 得到 
\begin{equation}
\sum_{a_\ell} \tilde{U}_{(a_{\ell-1}\sigma_\ell), a_\ell} S_{a_\ell} (V^\dagger)_{a_\ell, (\sigma_{\ell+1}a_{\ell+1})} = \sum_{a_\ell} \tilde{M}^{\sigma_\ell}_{a_{\ell-1}, a_\ell} B^{\sigma_{\ell+1}}_{a_\ell, a_{\ell+1}} ,
\end{equation}
其中的 $\tilde{M}$ 由重塑 $\tilde{U} S$ 得到。实际上它被丢弃了，因为我们向左移动一个格点，去寻找 $\tilde{M}^{\sigma_{\ell-1}\sigma_{\ell}}$。我们也可以改用对 $\tilde{M}^\dagger$ 作更廉价的 QR 分解，从 $Q^\dagger$ 得到 $B$。

\begin{figure}
\centering\includegraphics[scale=0.6]{PyMPS_MPS_TwoSiteCompression.eps}
\caption{两格点方案下 MPS 迭代压缩的方程。}
\label{fig:twositecompression}
\end{figure}


在结束本节关于压缩的讨论之前，让我比较一下各种方法的相对优劣。如果压缩量不大，SVD 方法的相互依赖性并不太要紧。尽管如此，变分拟设仍更优越；其唯一的弱点是由于其迭代性质，必须为压缩后的态提供初始猜测。若随机选取，该方法会把大量时间浪费在仅仅进入正确的邻域上。因此，聪明的做法是把 SVD 压缩得到的态作为迭代方法的最初输入。我们如何才能至少部分地避免由 $D'$ 带来的潜在高成本？

实践中，压缩主要出现在两种情形：（i）MPS 已相加，因而矩阵维数也相加了；（ii）矩阵乘积算符（MPO）已作用于某个 MPS；我们将看到这会导致 MPS 与 MPO 的矩阵维数相乘。

在第一种情形，可以利用 $\ket{\psi} = \ket{\phi_1} + \ket{\phi_2} + \ldots \ket{\phi_n}$ 这一事实来加速变分压缩。于是我们可以把变分方程改写为
\begin{equation}
 \frac{\partial}{\partial \tilde{M}^{\sigma_i*}_{a_{i-1}a_i}} ( \braket{\tilde{\psi}}{\tilde{\psi}}
- \braket{\tilde{\psi}}{\psi} ) =  \frac{\partial}{\partial \tilde{M}^{\sigma_i*}_{a_{i-1}a_i}} ( \braket{\tilde{\psi}}{\tilde{\psi}} - \braket{\tilde{\psi}}{\phi_1}
- \braket{\tilde{\psi}}{\phi_2} - \ldots - \braket{\tilde{\psi}}{\phi_n}) = 0 .
\end{equation}
如果把方程展开（假设混合规范表示），右边现在由 $n$ 个涉及 $D$ 维矩阵的重叠之和构成，而不是一个涉及 $D$ 维与 $nD$ 维矩阵的重叠（该重叠中出现的两类矩阵乘法分别代价高达 $O(n^2D^3)$ 与 $O(nD^3)$）；现在我们有 $n$ 个重叠，每个代价为 $O(D^3)$。当 $n$ 很大时，``分解''方法应当最多快 $n$ 倍。

第二种情形的细节我们推迟到讨论过 MPO 之后。想法是进行 SVD 压缩，但省去事先保证正确归一化这一代价特别高的步骤；如果由于某种原因块态仍然几乎正交归一，那么结果应当相当不错（并且可以转化为规范形式，这在压缩之后是很廉价的），或者至少可以作为变分方法的合理输入 \cite{Stoudenmire10}。

\subsection{记法与转换}
迄今为止，我们探讨的 MPS 记法基于每个格点一组矩阵；我们利用了这些矩阵的特殊归一化性质，从而得到具有诱人附加特性的 MPS（比如生成正交归一的态集合，或编码一个 Schmidt 分解）。考虑我们的晶格，其格点编号为 $1$ 到 $L$，就 DMRG 构造而言，若能方便地访问用单次切割可得到的系统 AB 的全部 $L-1$ 种可能二分，将很有用处。

这种记法由 Vidal\cite{Vidal03a} 引入，形式如下：
\begin{equation}
\ket{\psi} = \sum_{\sigma_1,\ldots,\sigma_L} \Gamma^{\sigma_1} \Lambda^{[1]}  \Gamma^{\sigma_2} \Lambda^{[2]}  \Gamma^{\sigma_3} \Lambda^{[3]} \ldots  \Gamma^{\sigma_{L-1}} \Lambda^{[L-1]}  \Gamma^{\sigma_L} \ket{\sigma_1,\ldots,\sigma_L} ,
\label{eq:canonicalform}
\end{equation} 
其中我们在每个格点 $\ell$ 上引入一组 $d$ 个矩阵 $\Gamma^{\sigma_\ell}$，并在每条键 $\ell$ 上引入一个对角矩阵 $\Lambda^{[\ell]}$。这些矩阵由如下要求确定：对任意 $1\leq \ell < L$，我们都能读出 Schmidt 分解
\begin{equation}
\ket{\psi} = \sum_{a_\ell} s_{a_\ell} \ket{a_\ell}_A \ket{a_\ell}_B
\end{equation}
其中 Schmidt 系数是 $\Lambda^{[\ell]}$ 的对角元，$s_{a_\ell} = \Lambda^{[\ell]}_{a_\ell,a_\ell}$，而 A 和 B 上的态由下式给出
\begin{eqnarray}
\ket{a_\ell}_A &=& \sum_{\sigma_1,\ldots,\sigma_\ell}
 (\Gamma^{\sigma_1}\Lambda^{[1]}\Gamma^{\sigma_2} \ldots \Lambda^{[\ell-1]} \Gamma^{\sigma_\ell})_{a_{\ell}}
\ket{\sigma_1,\ldots,\sigma_\ell}, \\
\ket{a_\ell}_B &=& \sum_{\sigma_{\ell+1},\ldots,\sigma_L} 
(\Gamma^{\sigma_{\ell+1}}\Lambda^{[\ell+1]}\Gamma^{\sigma_{\ell+2}} \ldots \Lambda^{[L-1]} \Gamma^{\sigma_L})_{a_\ell}
\ket{\sigma_{\ell+1},\ldots,\sigma_L},
\end{eqnarray}
其中 A 上与 B 上的态各自正交归一，这让人想起由 $A$-矩阵和 $B$-矩阵出发的类似构造。图形上，新记法可表示成图~\ref{fig:vidalmps} 那样。它显然是 $A$-矩阵记法的一个更显式的版本，其优点是显式保留了奇异值、约化密度矩阵本征值与纠缠：切开键 $\ell$，约化密度算符 $\hat{\rho}_A$ 与 $\hat{\rho}_B$ 在本征基表示下读作
\begin{equation}
\rho_A^{[\ell]} = \rho_B^{[\ell]} = (\Lambda^{[\ell]})^2 ,
\end{equation}
更精确地（但对实且对角的 $\Lambda^{[\ell]}$ 无关紧要）$\rho_A^{[\ell]} = \Lambda^{[\ell]}\Lambda^{[\ell]\dagger}$ 与 $\rho_B^{[\ell]} = \Lambda^{[\ell]\dagger}\Lambda^{[\ell]}$，
其中 $\rho_A^{[\ell]}$ 与 $\rho_B^{[\ell]}$ 的本征态分别由 $\{ \ket{a_\ell}_A \}$ 与 $\{ \ket{a_\ell}_B \}$ 给出。纠缠的 von Neumann 熵可直接从 $\Lambda^{[\ell]}$ 读出：$S_{A|B} = -\tr (\Lambda^{[\ell]})^2 \log_2 (\Lambda^{[\ell]})^2$。

在探讨与其他记法的联系之前，让我们首先证明：任何量子态确实都能通过一个与前述把 $\ket{\psi}$ 分解为 $A$-矩阵乘积（$B$-矩阵亦然）的过程极为类似的程序化成该形式。从系数 $c_{\sigma_1\ldots\sigma_L}$ 出发，我们将其重塑为 $\Psi_{\sigma_1, (\sigma_2 \ldots \sigma_L)}$，然后对其进行迭代 SVD。我们把奇异值矩阵记为 $\Lambda^{[i]}$。第一次 SVD 之后，我们把 $A^{\sigma_1}$ 改名为 $\Gamma^{\sigma_1}$。在随后的各次 SVD 中，与前面一样，我们把 $\Lambda$ 与 $V^\dagger$ 乘入 $\Psi$ 以构成下一个待 SVD 的矩阵，并作重塑使得行指标总是一个 $a$-指标和一个 $\sigma$-指标。利用已用过的重塑 $U_{(a_{\ell-1}\sigma_\ell), a_{\ell}} \rightarrow A^{\sigma_\ell}_{a_{\ell-1},a_\ell}$，我们得到
\begin{eqnarray*}
c_{\sigma_1\ldots\sigma_L} &=& \Psi_{\sigma_1, (\sigma_2 \ldots \sigma_L)} \\
&=& \sum_{a_1} A^{\sigma_1}_{a_1} \underbrace{\Lambda^{[1]}_{a_1,a_1} (V^\dagger)_{a_1, (\sigma_2 \ldots \sigma_L)}} \\
&=& \sum_{a_1} \Gamma^{\sigma_1}_{a_1} \Psi_{(a_1\sigma_2), (\sigma_3 \ldots \sigma_L)} \\
&=& \sum_{a_1,a_2} \Gamma^{\sigma_1}_{a_1} A^{\sigma_2}_{a_1,a_2} \underbrace{\Lambda^{[2]}_{a_2,a_2} (V^\dagger)_{a_2, (\sigma_3 \ldots \sigma_L)}} \\
&=& \sum_{a_1,a_2} \Gamma^{\sigma_1}_{a_1} \overbrace{\Lambda^{[1]}_{a_1,a_1} \Gamma^{\sigma_2}_{a_1,a_2}} \Psi_{(a_2\sigma_3), (\sigma_4 \ldots \sigma_L)} \\
&=& \sum_{a_1,a_2,a_3} \Gamma^{\sigma_1}_{a_1} \Lambda^{[1]}_{a_1,a_1} \Gamma^{\sigma_2}_{a_1,a_2} A^{\sigma_3}_{a_2,a_3}  \underbrace{\Lambda^{[3]}_{a_3,a_3} (V^\dagger)_{a_3, (\sigma_4 \ldots \sigma_L)}} \\ 
&=& \sum_{a_1,a_2,a_3} \Gamma^{\sigma_1}_{a_1} \Lambda^{[1]}_{a_1,a_1} \Gamma^{\sigma_2}_{a_1,a_2} \overbrace{\Lambda^{[2]}_{a_2,a_2} \Gamma^{\sigma_3}_{a_2,a_3}}  \Psi_{(a_3\sigma_4), (\sigma_5 \ldots \sigma_L)} 
\end{eqnarray*}
如此继续下去。与分解为 $A$-矩阵的关键区别在于：利用上一步得到的 $\Lambda^{[\ell-1]}$ 的信息，每个 $A$ 被分解为 
\begin{equation}
A^{\sigma_\ell}_{a_{\ell-1},a_\ell} = \Lambda^{[\ell-1]}_{a_{\ell-1},a_{\ell-1}} \Gamma^{\sigma_\ell}_{a_{\ell-1},a_\ell} ,
\label{eq:agammalambda}
\end{equation}
这意味着要除以 $\Lambda^{[\ell-1]}$ 的对角元。如果在 SVD 中我们只保留非零奇异值，这在数学上是合法的运算，尽管在数值上可能相当棘手。在这个概念性演示中忽略这一问题，我们确实得到了所需形式的分解；为了证明它确实正确，我们必须证明在每一步迭代中我们确实得到一个 Schmidt 分解。但这很容易看出：任何 $\Lambda^{[\ell]}$ 左边的矩阵都可以归并成（或者更确切地说，是由其生成的）左规范化的 $A$-矩阵，它们在晶格从格点 1 到 $\ell$ 的部分上生成一组正交归一的态。在任何 $\Lambda^{[\ell]}$ 的右边，有一个各行正交归一的矩阵 $V^\dagger$，这意味着态 $\ket{a_\ell}_B = \sum_{\sigma_{\ell+1},\ldots}
(V^\dagger)_{a_\ell, \sigma_{\ell+1} \ldots} \ket{\sigma_{\ell+1} \ldots}$ 也是正交归一的。因此，给出 $\Lambda^{[\ell]}$ 的 SVD 确实导致一个合法的 Schmidt 分解。

得到这种记法的另一种途径是：先做一次标准的左规范分解，把生成的（以及此前被丢弃的）所有奇异值矩阵存为 $\Lambda^{[i]}$，然后再在所有相邻的 $A$-矩阵 $A^{\sigma_i}$ 与 $A^{\sigma_{i+1}}$ 之间插入恒等式 $\Lambda^{[i]}(\Lambda^{[i]})^{-1}$。随后利用式~(\ref{eq:agammalambda}) 便得到同样的结果。

\begin{figure}
\centering\includegraphics[width=\textwidth]{PyMPS_VidalMPS.eps}
\caption{用 Vidal 记法表示的 MPS。奇异值在键上保持显式（菱形）。$\Lambda$ 位于键上，$\Gamma$ 位于格点上。按构造，相邻的 $\Lambda$ 与 $\Gamma$ 可缩并成左规范或右规范的 $A$ 或 $B$ 矩阵。该态可以平凡地分组为一串 $A$（给出正交归一的块态）、一个奇异值矩阵、以及一串 $B$（给出正交归一的块态）。}
\label{fig:vidalmps}
\end{figure}

类似地，若从右端出发、利用 $B$-矩阵的右规范进行分解，则同一态可用如下分组得到
\begin{equation}
B^{\sigma_\ell}_{a_{\ell-1},a_{\ell}} = \Gamma^{\sigma_\ell}_{a_{\ell-1},a_\ell} \Lambda^{[\ell]}_{a_\ell,a_\ell} ,
\label{eq:bgammalambda}
\end{equation}
其中，为简化本式以及 $A$-矩阵对应方程（式~(\ref{eq:agammalambda})）的记法，引入哑元 $\Lambda^{[0]}$ 与 $\Lambda^{[L]}$ 会很有用，二者均为标量 1。

$A$ 与 $B$-矩阵的这种分组使得我们能够用 $\Gamma\Lambda$ 语言重新表述左规范与右规范条件：左规范条件为
\begin{equation}
I = \sum_{\sigma_i} A^{\sigma_i\dagger} A^{\sigma_i} = \sum_{\sigma_i} \Gamma^{\sigma_i\dagger} \Lambda^{[i-1]\dagger} \Lambda^{[i-1]} \Gamma^{\sigma_i}
\end{equation}
或者更紧凑地写为
\begin{equation}
\sum_{\sigma_i}  \Gamma^{\sigma_i\dagger} \rho_B^{[i-1]} \Gamma^{\sigma_i} = I .
\label{eq:leftnormalizevidal}
\end{equation}
右规范条件为
\begin{equation}
\sum_{\sigma_i} \Gamma^{\sigma_i} \rho_A^{[i]} \Gamma^{\sigma_i\dagger} = I . 
\label{eq:rightnormalizevidal}
\end{equation}
有趣的是，若利用式~(\ref{eq:agammalambda}) 与式~(\ref{eq:bgammalambda}) 来翻译密度算符的递推关系（式~(\ref{eq:rhoarecursion}) 与式~(\ref{eq:rhobrecursion})），同样会得到式~(\ref{eq:leftnormalizevidal}) 与式~(\ref{eq:rightnormalizevidal})。
形如式~(\ref{eq:canonicalform}) 且满足约束条件式~(\ref{eq:leftnormalizevidal}) 与式~(\ref{eq:rightnormalizevidal}) 的矩阵乘积态称为{\em 规范}的。

在 $AB$ 记法、$\Gamma\Lambda$ 记法以及 DMRG 的块-格点记法之间都可以相互转换，只是其中暗藏一些数值上的陷阱。

{\em 转换} $\Gamma\Lambda \rightarrow A,B$：从 $\Gamma\Lambda \rightarrow A,B$ 的转换很容易。若在 $\ket{\psi}$ 的最左端引入一个额外的哑标量 $\Lambda^{[0]} = 1$ 充当``矩阵''，就可以利用上面的定义式~(\ref{eq:agammalambda}) 分组得到
\begin{equation}
(\Lambda^{[0]} \Gamma^{\sigma_1})
(\Lambda^{[1]} \Gamma^{\sigma_2})(\Lambda^{[2]} \Gamma^{\sigma_3})
 \ldots \rightarrow A^{\sigma_1}A^{\sigma_2}A^{\sigma_3}\ldots 
 \end{equation}
或者利用式~(\ref{eq:bgammalambda}) 得到
\begin{equation}
(\Gamma^{\sigma_1}\Lambda^{[1]})
(\Gamma^{\sigma_2}\Lambda^{[2]} )(\Gamma^{\sigma_3}\Lambda^{[3]})
 \ldots \rightarrow B^{\sigma_1}B^{\sigma_2}B^{\sigma_3}\ldots 
 \end{equation}

考虑到 DMRG 及其他 MPS 方法的实际做法，考察{\em 混合}转换是很有意思的。考虑格点 $\ell$ 与 $\ell+1$ 之间的键 $\ell$。我们可以从左端开始将 $\Lambda \Gamma \rightarrow A$ 缩并，得到左规范矩阵；从右端将 $\Gamma\Lambda \rightarrow B$ 缩并，得到右规范矩阵；只在中心留下 $\Lambda^{[\ell]}$ 不动（图~\ref{fig:dmrg0mps}）：
\begin{figure}
\centering\includegraphics[width=\textwidth]{PyMPS_DMRG0MPS.eps}
\caption{Vidal 的 MPS 记法、$A$、$B$-矩阵 MPS 记法以及 DMRG 块记法。$A$-矩阵生成左块态，$B$ 矩阵生成右块态。矩阵 $\Lambda^{[\ell]}$ 通过奇异值将二者连接起来。}
\label{fig:dmrg0mps}
\end{figure}

\begin{figure}
\centering\includegraphics[width=\textwidth]{PyMPS_DMRG1MPS.eps}
\caption{单格点 DMRG 中一个态的表示：把 Vidal 的 MPS 记法及 $A,B$-矩阵 MPS 记法翻译成 DMRG 块记法。$A$-矩阵生成左块态，$B$ 矩阵生成右块态。$\Psi^{\sigma_\ell}$ 的矩阵元正是 DMRG 态的系数。}
\label{fig:dmrg1mps}
\end{figure}


\begin{equation}
(\Lambda^{[0]}\Gamma)(\Lambda^{[1]}\Gamma)\ldots(\Lambda^{[\ell-2]}\Gamma)(\Lambda^{[\ell-1]}\Gamma) \Lambda^{[\ell]} (\Gamma\Lambda^{[\ell+1]})(\Gamma\Lambda^{[\ell+2]}) \ldots (\Gamma\Lambda^{[L]}) .
\end{equation}
由于键 $\ell$ 左侧的加括号方式生成左规范的 $A$-矩阵，右侧生成右规范的 $B$ 矩阵，我们可以像上一节的递推那样把它们乘开，得到 A 与 B 的正交归一块基，从而得到 Schmidt 分解
\begin{equation}
\ket{\psi} = \sum_{a_{\ell}} \ket{a_\ell}_A s_{a_\ell} \ket{a_\ell}_B .
\end{equation}
更进一步，我们还可以取一个格点（$\ell$）或两个格点（$\ell,\ell+1$），把那里的所有矩阵乘成一个矩阵（图~\ref{fig:dmrg1mps} 与图~\ref{fig:dmrg2mps}）：
\begin{figure}
\centering\includegraphics[width=\textwidth]{PyMPS_DMRG2MPS.eps}
\caption{双格点 DMRG 中一个态的表示：把 Vidal 的 MPS 记法及 $A,B$-矩阵 MPS 记法翻译成 DMRG 块记法。$A$-矩阵生成左块态，$B$ 矩阵生成右块态。矩阵 $\Psi^{\sigma_{\ell}\sigma_{\ell+1}}$ 的元素正是 DMRG 态的系数。}
\label{fig:dmrg2mps}
\end{figure}


\begin{equation}
(\Lambda^{[0]}\Gamma)(\Lambda^{[1]}\Gamma)\ldots(\Lambda^{[\ell-2]}\Gamma)(\Lambda^{[\ell-1]}\Gamma \Lambda^{[\ell]}) (\Gamma\Lambda^{[\ell+1]})(\Gamma\Lambda^{[\ell+2]}) \ldots (\Gamma\Lambda^{[L]}) .
\end{equation}
把中心矩阵记为 $\Psi^{\sigma_\ell}= \Lambda^{[\ell-1]}\Gamma^{\sigma_\ell} \Lambda^{[\ell]}$，我们就可以写
\begin{equation}
\ket{\psi} = \sum_{\fat{\sigma}} A^{\sigma_1} \ldots A^{\sigma_{\ell-1}} \Psi^{\sigma_\ell} B^{\sigma_{\ell+1}} \ldots B^{\sigma_L}
\ket{\fat{\sigma}}
\end{equation}
或者，再次构建块基，
\begin{equation}
\ket{\psi} = \sum_{a_{\ell-1},a_{\ell},\sigma_\ell} \ket{a_{\ell-1}}_A \Psi^{\sigma_\ell}_{a_{\ell-1}a_\ell} \ket{a_\ell}_B .
\end{equation}
如果我们把两个格点也一并分组，则有
\begin{equation}
(\Lambda^{[0]}\Gamma)(\Lambda^{[1]}\Gamma)\ldots(\Lambda^{[\ell-2]}\Gamma)(\Lambda^{[\ell-1]}\Gamma \Lambda^{[\ell]} \Gamma\Lambda^{[\ell+1]})(\Gamma\Lambda^{[\ell+2]}) \ldots (\Gamma\Lambda^{[L]}) 
\end{equation}
或者，取中心矩阵 $\Psi^{\sigma_\ell \sigma_{\ell+1}}=\Lambda^{[\ell-1]}\Gamma^{\sigma_\ell} \Lambda^{[\ell]} \Gamma^{\sigma_{\ell+1}}\Lambda^{[\ell+1]}$，
\begin{equation}
\ket{\psi} = \sum_{\fat{\sigma}} A^{\sigma_1} \ldots A^{\sigma_{\ell-1}} \Psi^{\sigma_\ell \sigma_{\ell+1}} B^{\sigma_{\ell+2}} \ldots B^{\sigma_L} \ket{\fat{\sigma}}
\end{equation}
或者，用块基表示，
\begin{equation}
\ket{\psi} = \sum_{a_{\ell-1},a_{\ell+1},\sigma_\ell,\sigma_{\ell+1}} \ket{a_{\ell-1}}_A \Psi^{\sigma_\ell \sigma_{\ell+1}}_{a_{\ell-1}a_{\ell+1}} \ket{a_{\ell+1}}_B .
\end{equation}
这正是``单格点''与最初的``双格点''DMRG 所考虑的态，它们在两个块之间显式地保留一个或两个格点。

{\em 转换} $A,B \rightarrow \Gamma\Lambda$：反方向的转换 $A,B \rightarrow \Gamma\Lambda$ 要更复杂一些。其想法是逐次迭代地求出该态的 Schmidt 分解（从而得到奇异值），也就是对角矩阵 $\Lambda^{[\ell]}$，再由此算出 $\Gamma^{\sigma}$-矩阵。这一做法实际上与证明任意量子态都存在规范 $\Gamma\Lambda$ 表示所用的方法非常相似。

设 $\ket{\psi}$ 是右规范的，则态系数形如 $B^{\sigma_1}B^{\sigma_2}B^{\sigma_3}B^{\sigma_4}\ldots$。于是我们可以通过一连串 SVD 逐步进行
\begin{eqnarray*}
& & B^{\sigma_1}B^{\sigma_2}B^{\sigma_3}B^{\sigma_4} \ldots \\
&=& (A^{\sigma_1} \Lambda^{[1]} V^\dagger) B^{\sigma_2}B^{\sigma_3}B^{\sigma_4} \ldots \\
&=& \Gamma^{\sigma_1} M^{\sigma_2} B^{\sigma_3}B^{\sigma_4} \ldots \\
&=& \Gamma^{\sigma_1} (A^{\sigma_2}  \Lambda^{[2]} V^\dagger) B^{\sigma_3}B^{\sigma_4} \ldots \\
&=& \Gamma^{\sigma_1} \Lambda^{[1]} \Gamma^{\sigma_2} M^{\sigma_3} B^{\sigma_4} \ldots \\
\end{eqnarray*} 
如此继续下去。这里，待作 SVD 的矩阵为 $M^{\sigma_\ell} = \Lambda^{[\ell-1]} V^\dagger B^{\sigma_\ell}$。$\Gamma^{\sigma_\ell}$-矩阵则由 $A^{\sigma_\ell}$-矩阵借助记下的 $\Lambda^{[\ell-1]}$ 并利用式~(\ref{eq:agammalambda}) 得到，这意味着要除以奇异值。

除以奇异值在数值上是个令人头疼的问题，因为奇异值可能而且常常会非常小，特别是在尝试高精度计算时，连非常小的奇异值也会被一路带着参与运算。数值上明智的做法是仿照计算近乎奇异矩阵的（伪）逆时的处理方式，把所有 $s_a < \epsilon$（比如取 $\epsilon=10^{-8}$）置为 0，并把它们从所有求和中排除（例如在 Schmidt 分解中）。由于我们按大小对 $s_a$ 排序，这意味着相应地收缩矩阵 $U$ 和 $V^\dagger$。这些小奇异值在约化密度算符（即它们的平方）中所占权重很小，因此与数值上的隐患相比，态描述精度的损失非常小。事实上，问题在于算法中有若干处我们隐含地依赖（正交）归一性假设，而在一次``野蛮''的除法之后这些假设可能不再成立。

在本节（篇幅颇长）结束之际，让我用一幅图（图~\ref{fig:conversions}）总结各种转换与规范化过程，不过应当记住，有些转换只在理论上可行，在数值实践中并不可行。

\begin{figure}
\centering\includegraphics[scale=0.5]{PyMPS_MPSConversions.eps}
\caption{各种表示之间的转换：$c$ 是由指数多的态系数给出的显式表示；$A$、$B$ 和 $\Gamma\Lambda$ 分别代表左规范、右规范和规范 MPS；$M$ 代表任意 MPS。实线表示计算上可行的转换，虚线表示较为假设性的转换。}
\label{fig:conversions}
\end{figure}

\section{矩阵乘积算符（MPO）}
\label{sec:MPO}
如果我们考虑 MPS 的单个系数 $\braket{\fat{\sigma}}{\psi}$，
\[
\braket{\fat{\sigma}}{\psi} = M^{\sigma_1}M^{\sigma_2}\ldots M^{\sigma_{L-1}}M^{\sigma_L},
\]
一个自然的推广是尝试把算符的系数 $\bra{\fat{\sigma}} \hat{O} \ket{\fat{\sigma}'}$ 写成 \cite{VerstraeteRipoll04,McCulloch07,VerstraetePirvu08,Crosswhite08,Frowis10}
\begin{equation}
\bra{\fat{\sigma}} \hat{O} \ket{\fat{\sigma}'} = W^{\sigma_1\sigma'_1}W^{\sigma_2\sigma'_2} \ldots
W^{\sigma_{L-1}\sigma'_{L-1}}W^{\sigma_L\sigma'_L}
\end{equation}
其中 $W^{\sigma\sigma'}$ 是与 $M^{\sigma}$ 类似的矩阵，唯一区别在于作为算符的表示，它们同时需要出射和入射的物理态：
\begin{equation}
\hat{O} = \sum_{\fat{\sigma},\fat{\sigma}'} W^{\sigma_1\sigma'_1}W^{\sigma_2\sigma'_2} \ldots
W^{\sigma_{L-1}\sigma'_{L-1}}W^{\sigma_L\sigma'_L} \ket{\fat{\sigma}}\bra{\fat{\sigma}'} ,
\label{eq:mpoform}
\end{equation}
与 MPS 一样，可以照同样方式推广到周期边界条件。为 MPS 引入的图示表示可以径直推广：在 $M$ 的表示中原来有一条代表物理态的竖线，现在 $W$ 中则有两条竖线，一向下、一向上，分别代表入射和出射的物理态（图~\ref{fig:singleMPO}）。完整的 MPO 本身则如图~\ref{fig:MPO} 所示。如果想利用好量子数，MPS 的那套方法可以直接搬过来：从顶部引入一个入射的局域态量子数，向底部引出一个出射的量子数，并从左侧引入一个入射量子数、向右侧引出一个出射量子数。规则与 MPS 相同：入射与出射量子数之总和必须相等，即 $M(\ket{\sigma_i}) + M(\ket{b_{i-1}}) = M(\ket{\sigma'_i}) + M(\ket{b_{i}})$，这里为记法方便我把键上的标记理解为态。与 MPS 中一样，我们也可以设想第一个格点之前和最后一个格点之后的哑指标，它们反映算符以何种（确定的！）方式改变总量子数。对于与相应算符对易的哈密顿量，这一改变量为零，可以忽略哑指标。我们要构建的 MPO 都可以证明在键上具有好量子数，因为它们要么源自 SVD（例如用于时间演化），要么源自涉及量子数改变明确的算符的构造规则（例如哈密顿量的 MPO）。

事实上，{\em 任何}算符都可以化成式~(\ref{eq:mpoform}) 的形式，因为它可以写成
\begin{eqnarray}
\hat{O} &=& \sum_{\sigma_1,\ldots,\sigma_L,\sigma'_1,\ldots,\sigma'_L} c_{(\sigma_1\ldots\sigma_L)(\sigma'_1\ldots\sigma'_L)} \ket{\sigma_1,\ldots,\sigma_L} \bra{\sigma'_1,\ldots,\sigma'_L}  \nonumber \\ 
&=& \sum_{\sigma_1,\ldots,\sigma_L,\sigma'_1,\ldots,\sigma'_L} c_{(\sigma_1\sigma'_1) \ldots (\sigma_L \sigma'_L)} \ket{\sigma_1,\ldots,\sigma_L} \bra{\sigma'_1,\ldots,\sigma'_L}
\end{eqnarray}
然后可以像对 MPS 那样对它作分解，让双指标 $\sigma_i \sigma'_i$ 扮演 MPS 中指标 $\sigma_i$ 的角色。

\begin{figure}
\centering\includegraphics[width=250pt]{PyMPS_MPOElements.eps}
\caption{矩阵乘积算符的元素：(i) 链左端的角矩阵算符 $W^{[1]\sigma_1\sigma'_1}_{1,b_1}$；(ii) 体矩阵算符 $W^{[\ell]\sigma_\ell\sigma'_\ell}_{b_{\ell-1},b_{\ell}}$；(iii) 链右端的角算符 $W^{[L]\sigma_L\sigma'_L}_{b_{L-1},1}$：物理指标朝上和朝下，矩阵指标用水平横线表示。}
\label{fig:singleMPO}
\end{figure}

\begin{figure}
\centering\includegraphics[width=250pt]{PyMPS_MPO.eps}
\caption{作用在整条链上的矩阵乘积算符：水平的矩阵指标已缩并，MPO 只需简单缩并竖直（物理）指标即可作用于一个 MPS。}
\label{fig:MPO}
\end{figure}

与 MPS 一样，我们必须追问如何用它们进行操作以及它们在实践中如何构造，因为朴素的分解可能是指数复杂的。事实证明，大多数操作与 MPS 情形完全类似地运行。

\subsection{将 MPO 作用于 MPS}
矩阵乘积算符作用于矩阵乘积态的过程如下
\begin{eqnarray*}
\hat{O}\ket{\psi} &=& \sum_{\fat{\sigma},\fat{\sigma}'} (W^{\sigma_1,\sigma'_1} W^{\sigma_2,\sigma'_2}\ldots)(M^{\sigma'_1} M^{\sigma'_2}\ldots) \ket{\fat{\sigma}} \\
&=&  \sum_{\fat{\sigma},\fat{\sigma}'}\sum_{\fat{a},\fat{b}} (W^{\sigma_1,\sigma'_1}_{1,b_1} W^{\sigma_2,\sigma'_2}_{b_1,b_2}\ldots)(M^{\sigma'_1}_{1,a_1} M^{\sigma'_2}_{a_1,a_2} \ldots) \ket{\fat{\sigma}} \\
&=& \sum_{\fat{\sigma},\fat{\sigma}'}\sum_{\fat{a},\fat{b}} (W^{\sigma_1,\sigma'_1}_{1,b_1} M^{\sigma'_1}_{1,a_1})(W^{\sigma_2,\sigma'_2}_{b_1,b_2}M^{\sigma'_2}_{a_1,a_2}) \ldots \ket{\fat{\sigma}} \\
&=& \sum_{\fat{\sigma}}\sum_{\fat{a},\fat{b}} N^{\sigma_1}_{(1,1),(b_1,a_1)} N^{\sigma_2}_{(b_1,a_1),(b_2,a_2)} \ldots \ket{\fat{\sigma}} \\
&=& \sum_{\fat{\sigma}} N^{\sigma_1} N^{\sigma_2} \ldots \ket{\fat{\sigma}}
\end{eqnarray*}

MPO 的美妙之处在于它保持 MPS 的形式不变，代价是矩阵尺寸增大：新 MPS 的维数是原 MPS 维数与 MPO 维数的乘积（图~\ref{fig:MPOtimesMPS}）。

\begin{figure}
\centering\includegraphics[width=250pt]{PyMPS_MPOtimesMPS.eps}
\caption{矩阵乘积算符作用于矩阵乘积态：相匹配的物理（竖直）指标被缩并，产生一个新的矩阵乘积态，其矩阵维数相乘并保持乘积结构。}
\label{fig:MPOtimesMPS}
\end{figure}

结果可以概括为 $\ket{\phi}=\hat{O}\ket{\psi}$，其中 $\ket{\phi}$ 是由如下矩阵 $N^{\sigma_i}$ 构成的 MPS：
\begin{equation}
N^{\sigma_i}_{(b_{i-1},a_{i-1}),(b_i,a_i)} = \sum_{\sigma_i'} W^{\sigma_i\sigma_i'}_{b_{i-1}b_i} M^{\sigma_i'}_{a_{i-1}a_i} .
\end{equation}
如果使用（可加的）好量子数，则由每个张量处的求和规则可以证明，它们在入射和出射的水平键上是可加的。

又一次，一个看似指数复杂的操作（对指数多个 $\fat{\sigma}$ 求和）被约化为低代价的操作：运算量为
$Ld^2 D_W^2 D^2$ 的量级，其中 $D_W$ 是 MPO 的维数。

\subsection{MPO 的加法与乘法}

对 MPO 的操作大体上沿袭 MPS 的思路。考虑两个算符 $\hat{O}$ 与 $\hat{P}$ 的加法，它们分别有 MPO 表示 $W^{\sigma_i\sigma'_i}$ 与
$\tilde{W}^{\sigma_i\sigma'_i}$，则所得 MPO 的构造与 MPS 情形完全一样：对所有格点 $1 < i < L$ 取直和 $W^{\sigma_i\sigma'_i} \oplus \tilde{W}^{\sigma_i\sigma'_i}$，格点 $1$ 与 $L$ 则适用同样的特殊规则。本质上，我们（又一次）只需把 $\sigma_i$ 与 $\sigma'_i$ 看成一个``大''的物理指标。

两个算符的乘积 $\hat{P} \hat{O}$（更确切地说，随后的相继作用）也不难理解：把 $\hat{O}$ 作用到某个态 $\ket{\psi}$ 上得到一个新的 MPS，其
矩阵为 $N^{\sigma_i}_{(b_{i-1}a_{i-1}),(b_i a_i)} = \sum_{\sigma'_i} W^{\sigma_i\sigma'_i}_{b_{i-1}b_i} M^{\sigma'_i}_{a_{i-1}a_i}$。随后再作用 $\hat{P}$，又给出一个新的 MPS，其 $K^{\sigma_i}_{(\tilde{b}_{i-1}b_{i-1}a_{i-1}),(\tilde{b}_i b_i a_i)} =  \sum_{\sigma'_i} \tilde{W}^{\sigma_i\sigma'_i}_{\tilde{b}_{i-1} \tilde{b}_i} N^{\sigma'_i}_{(b_{i-1}a_{i-1}),(b_i a_i)}$。但由此我们可以立即读出（从两个 MPO 相继作用于一个态的图示中也一目了然），新的 MPO（矩阵为 $V^{\sigma_i\sigma'_i}$）由下式给出
\begin{equation}
V^{\sigma_i\sigma'_i}_{(\tilde{b}_{i-1}b_{i-1}),(\tilde{b}_{i}b_{i})} = \sum_{\sigma''_i} 
\tilde{W}^{\sigma_i\sigma''_i}_{\tilde{b}_{i-1} \tilde{b}_i} W^{\sigma''_i\sigma'_i}_{b_{i-1}b_i} .
\end{equation}
因此，MPO 的维数就像张量那样简单地相乘。如果把 MPS 看成在一个物理方向上带哑指标的 MPO，那么把 MPO 作用于 MPS 的规则便作为特例而得到。
 
\subsection{压缩 MPO 与 MPO-MPS 乘积}

与 MPS 一样，可能会出现压缩 MPO 的问题。这一点从上一节应当看得很清楚，那里 MPO 的维数被求和或相乘。如果无法把压缩问题推给 MPO 向 MPS 的作用（那时维数为 $D_W D$，是压缩的自然候选对象），我们就必须压缩 MPO 本身。一个典型例子是较程长的哈密顿量的 MPO 形式表示，它会很快导致很大的维数。

但我们可以采用与压缩 MPS 相同的技巧，既可以用 SVD 也可以迭代进行，用一个 $D_W$ 更小的 MPO 来逼近精确的 MPO。唯一的改变是：现在不是单一物理指标 $\sigma$，而是两个物理指标 $\sigma,\sigma'$，我们可以把它们当作一个单一指标 $(\sigma,\sigma')$。逼近随后在 Frobenius 范数下进行，它自然地推广了我们用于 MPS 逼近的向量 2-范数。

在此值得一提的是，文献 \cite{Stoudenmire10} 提出，在压缩 MPO-MPS 乘积这一特殊情形下，可以相对标准方法取得重要加速：如果必须先做归一化（代价为 $O(D_W^3 D^3)$），SVD 可能非常慢，而变分方法又亟需一个好的起点。不过，由 SVD 压缩给出的这个建议解也许并不差，只要块态接近正交归一；而在 MPO-MPS 乘积的情形，若 MPO 与 MPS 都处于规范形式（MPO 再次通过把双指标看成一个大指标而构成），这一点似乎基本成立，而且这种规范形式可以以低得多的代价（$O(d D^3)$ 与 $O(d^2 D_W^3)$，通常 $D_W \ll D$，对比 $O(d D^3 D_W^3)$）实现。即使所建议的压缩并不太好，它仍会为变分压缩提供一个好得多的起点。于是流程为：(i) 把 MPO 与 MPS 都化到同一规范形式；(ii) 做 SVD 压缩，当然只随算随乘 MPO 与 MPS 的矩阵；(iii) 如果对结果不太放心，就把它用作变分输入。


 
\section{用 MPS 作基态计算}
\label{sec:groundstates}
假设我们想求某个哈密顿量 $\hat{H}$ 的基态。为了找到对它的最优逼近，我们必须找出某个维数 $D$ 的、使下式极小的 MPS $\ket{\psi}$：
\begin{equation}
E = \frac{\bra{\psi} \hat{H} \ket{\psi}}{\braket{\psi}{\psi}} .
\end{equation} 
做到这一点最有效率的方式（特别是与从某个随机态出发作虚时演化这一同样可行的做法相比）是在 MPS 空间中作变分搜索。为了让这个算法清晰透明，我们先把 $\hat{H}$ 表示成 MPO。

\subsection{哈密顿量的 MPO 表示}
由于 MPO 表示中固有的乘积结构，尽管其存在性有保证，要为如下这样的哈密顿量显式构造出紧致的 MPO 表示似乎仍是一项无望的任务
\begin{equation}
\hat{H} = \sum_{i=1}^{L-1} \frac{J}{2}\hat{S}^+_i \hat{S}^-_{i+1} + \frac{J}{2}\hat{S}^-_i \hat{S}^+_{i+1} + J^z \hat{S}^z_i \hat{S}^z_{i+1} - h\sum_{i} \hat{S}^z_i .
\end{equation}
这种常见记法当然是算符张量积之和的缩写：
\begin{eqnarray*}
\hat{H} &=& J^z\hat{S}^z_1 \otimes \hat{S}^z_2 \otimes \hat{I} \otimes \hat{I} \otimes \hat{I} \ldots + \\
 & & \hat{I} \otimes J^z\hat{S}^z_2 \otimes \hat{S}^z_3 \otimes \hat{I} \otimes \hat{I} \ldots + \\
 & & \ldots
 \end{eqnarray*}
然而，把这个张量积之和写成 MPO 形式出奇地容易 \cite{McCulloch07} -- 为此，方便的做法是重新考虑基本构件 $W^{\sigma\sigma'}_{bb'}$，把它与其相伴的投影算符 $\ket{\sigma}\bra{\sigma'}$ 合并成一个算符值矩阵 $\hat{W}_{bb'} = \sum_{\sigma\sigma'} W^{\sigma\sigma'}_{bb'} \ket{\sigma}\bra{\sigma'}$，于是 MPO 取如下简单形式
\begin{equation}
\hat{O} = \hat{W}^{[1]}\hat{W}^{[2]} \ldots \hat{W}^{[L]} .
\end{equation}
每个 $\hat{W}^{[i]}$ 作用在格点 $i$ 处不同的局域希尔伯特空间上，这些空间的张量积给出全局希尔伯特空间。把这样的算符值矩阵相乘就得到算符张量积之和，因此把 $\hat{H}$ 写成紧致形式看来是可行的。

为了理解这一构造，我们遍历 $\hat{H}$ 中出现的任意一个算符串：从右端开始，串中只含单位算符，直到某处遇到（在我们的例子中）4 个非平凡算符之一。对于场算符，其左边更远的串部分只能含单位算符；对于相互作用算符，互补算符必须紧随其后以补全相互作用项，再往左则继续由单位算符接续。为便于记账，我们在某个给定键上引入串的 5 个相应状态：状态 1：右边只有单位算符，
状态 2,3,4：右边紧邻处有一个 $\hat{S}^+$、$\hat{S}^-$、$\hat{S}^z$，状态 5：右边某处已有补全的相互作用项或场项 $-h\hat{S}^z$。比较一个格点左右两侧串的状态，只允许少数几种跃迁：$1 \rightarrow 1$ 由单位算符 $\hat{I}$ 实现，$1 \rightarrow 2$ 由 $\hat{S}^+$ 实现，$1 \rightarrow 3$ 由 $\hat{S}^-$ 实现，$1 \rightarrow 4$ 由 $\hat{S}^z$ 实现，$1 \rightarrow 5$ 由 $-h\hat{S}^z$ 实现。此外，$2\rightarrow 5$ 由 $(J/2)\hat{S}^-$ 实现，$3 \rightarrow 5$ 由 $(J/2)\hat{S}^+$ 实现，$4 \rightarrow 5$ 由 $J^z\hat{S}^z$ 实现，以补全相互作用项；而对已补全的相互作用，由单位算符 $\hat{I}$ 实现 $5 \rightarrow 5$。此外，所有串状态必须在最后一个格点右侧从状态 1 出发，并在第一个格点左侧终止于状态 5（即未来 MPO 的维数 $D_W$）。现在可以用下面这组取值为算符的矩阵来编码：

\begin{equation}
\hat{W}^{[i]} = \left[
\begin{array}{ccccc}
\hat{I} & 0 & 0 & 0 & 0 \\
\hat{S}^+ & 0 & 0 & 0 & 0 \\
\hat{S}^- & 0 & 0 & 0 & 0 \\
\hat{S}^z & 0 & 0 & 0 & 0 \\
-h\hat{S}^z & (J/2)\hat{S}^- & (J/2)\hat{S}^+ & J^z\hat{S}^z & \hat{I}
\end{array}
\right]
\end{equation}
而在第一个和最后一个格点上
\begin{equation}
\hat{W}^{[1]} = \left[
\begin{array}{ccccc}
-h\hat{S}^z & (J/2)\hat{S}^- & (J/2)\hat{S}^+ & J^z\hat{S}^z & \hat{I} 
\end{array}
\right]
\quad\quad 
\hat{W}^{[L]} = \left[
\begin{array}{c}
\hat{I} \\ \hat{S}^+ \\ \hat{S}^- \\ \hat{S}^z \\ -h\hat{S}^z 
\end{array}
\right] .
\end{equation}
事实上，简单的乘法即可看出哈密顿量是如何产生的。代入显式的算符表示，就得到 MPO 所需的 $W^{\sigma\sigma'}$ 矩阵。因此可以把哈密顿量精确地写成非常紧凑的 MPO 形式；\cite{Crosswhite08a} 给出了一套导致相同结果的类似规则。

对于作用范围更长的哈密顿量，还须引入更多的``中间状态''。让我们考虑一个只含 $\hat{S}^z \hat{S}^z$ 相互作用、但同时包含最近邻和次近邻耦合的模型，
\begin{equation}
\hat{H} = J_1 \sum_{i} \hat{S}^z_i \hat{S}^z_{i+1} + J_2 \sum_{i} \hat{S}^z_i \hat{S}^z_{i+2} .
\end{equation}
那么体算符将形如
\begin{equation}
\hat{W}^{[i]} = \left[
\begin{array}{cccc}
\hat{I} & 0 & 0 & 0 \\
\hat{S}^z & 0 & 0 & 0 \\
0 & \hat{I} & 0 & 0 \\
0 & J_1\hat{S}^z & J_2\hat{S}^z & \hat{I}
\end{array}
\right] .
\end{equation}
$J_1$ 相互作用可以像前面那样编码（按 $1 \rightarrow 2 \rightarrow 4$ 进行），而对次近邻相互作用，则必须在 2 与 4 之间插入额外一步，即中间状态 3，其中恰好插入一个单位算符（按 $1 \rightarrow 2 \rightarrow 3 \rightarrow 4$ 进行）。它仅仅起记账作用。类似地，可以构造作用范围更长的相互作用。除左上角和右下角外，$\hat{W}^{[i]}$ 的非零部分按构造都位于对角线以下。

看起来，对于作用范围更长的相互作用，随着所需中间状态越来越多（每单位相互作用范围、每个相互作用项各需增加一个状态），维数 $D_W$ 将迅速增长。虽然一般情形确实如此，但已知存在一些重要例外，可以紧凑得多地表述 \cite{Frowis10,Crosswhite08a}；例如考虑如下指数衰减的相互作用强度 $J(r)=J e^{-r/\xi} = J \lambda^r$，其中 $r>0$ 且 $\lambda = \exp(-1/\xi)$。相互作用项 $\sum_r J(r)\hat{S}^z_i \hat{S}^z_{i+r}$ 将被如下体算符所涵盖
\begin{equation}
\hat{W}^{[i]} = \left[
\begin{array}{ccc}
\hat{I} & 0 & 0  \\
\hat{S}^z & \lambda\hat{I} & 0  \\
0 & J\lambda \hat{S}^z & \hat{I}
\end{array}
\right] .
\end{equation}

但即使不出现这种简化，事实表明仍然可以找到维数相当小、精度损失适度的 MPO：或者把任意相互作用函数 $J(r)$ 近似为按上述方式编码的指数和\cite{VerstraetePirvu08,Crosswhite08a}，在 $\alpha_i, \lambda_i$ 中最小化 $L_2$ 距离 $\| J(r) - \sum_{i=1}^n \alpha_i \lambda_i^r \|$，其中 $n$ 由 $D_W$ 和愿意接受的精度损失决定；或者如 \cite{Frowis10}，在可行的情况下当然可以先构造精确的 MPO，再把 MPS 压缩技术加以改造，将其压缩到可接受的 $D_W$（及精度损失）。 

\subsection{将哈密顿量 MPO 作用于混合规范形式态}
\label{subsec:Hamiltonianmixedcanonical}
让我们考虑处于如下混合规范表示中的 $\ket{\psi}$，它与单格点 DMRG 表示相同，
\begin{equation}
\ket{\psi} = \sum_{\fat{\sigma}} A^{\sigma_1} \ldots A^{\sigma_{\ell-1}} \Psi^{\sigma_{\ell}} B^{\sigma_{\ell+1}} \ldots B^{\sigma_L} \ket{\fat{\sigma}} 
\end{equation}
或
\begin{equation}
\ket{\psi} = \sum_{a_{\ell-1},a_{\ell}} \ket{a_{\ell-1}}_A \Psi^{\sigma_{\ell}}_{a_{\ell-1}, a_{\ell}} \ket{a_{\ell}}_B .
\end{equation}
现在来看用 $\hat{H}$ 的 MPO 表示所得到的矩阵元 $\bra{a_{\ell-1} \sigma_{\ell} a_{\ell}} \hat{H}  \ket{a'_{\ell-1} \sigma'_{\ell} a'_{\ell}}$。通过两次插入恒等算符 $\hat{I} = \sum_{\fat{\sigma}} \ket{\fat{\sigma}}\bra{\fat{\sigma}}$，我们得到（带星号的求和不包括格点 $\ell$）
\begin{eqnarray*}
& & \bra{a_{\ell-1} \sigma_{\ell} a_{\ell}} \hat{H}  \ket{a'_{\ell-1} \sigma'_{\ell} a'_{\ell}} \\
&=& \sum_{\fat{\sigma}} \sum_{\fat{\sigma}'} W^{\sigma_1,\sigma'_1} \ldots W^{\sigma_L,\sigma'_L} 
\braket{a_{\ell-1} \sigma_{\ell} a_{\ell}}{\fat{\sigma}} \braket{\fat{\sigma}'}{a'_{\ell-1} \sigma'_{\ell} a'_{\ell}} \\
&=& \sum_{\fat{\sigma}*}\sum_{\fat{\sigma}'*} W^{\sigma_1,\sigma'_1} \ldots W^{\sigma_{\ell},\sigma'_{\ell}} \ldots W^{\sigma_L,\sigma'_L} \\
& & \braket{a_{\ell-1}}{\sigma_1 \ldots \sigma_{\ell-1}} 
\braket{a_{\ell}}{\sigma_{\ell+1} \ldots \sigma_L}  \braket{\sigma'_1 \ldots \sigma'_{\ell-1}}{a'_{\ell-1}}
\braket{\sigma'_{\ell+1} \ldots \sigma'_L}{a'_{\ell}} \\
&=& \sum_{\fat{\sigma}*}\sum_{\fat{\sigma}'*} W^{\sigma_1,\sigma'_1} \ldots W^{\sigma_{\ell},\sigma'_{\ell}} \ldots W^{\sigma_L,\sigma'_L} \\
& & (A^{\sigma_1} \ldots A^{\sigma_{\ell-1}})^*_{1,a_{\ell-1}} (B^{\sigma_{\ell+1}} \ldots B^{\sigma_L})^*_{a_{\ell},1} (A^{\sigma'_1} \ldots A^{\sigma'_{\ell-1}})_{1,a'_{\ell-1}} (B^{\sigma'_{\ell+1}} \ldots B^{\sigma'_L})_{a'_{\ell},1} \\
&=& \sum_{\{ a_i, b_i, a'_i\} } \left( \sum_{\sigma_1\sigma'_1} A^{\sigma_1*}_{1,a_1} W^{\sigma_1,\sigma'_1}_{1,b_1} A^{\sigma'_1}_{1,a'_1} \right) \left( \sum_{\sigma_2\sigma'_2} A^{\sigma_2*}_{a_1,a_2} W^{\sigma_2,\sigma'_2}_{b_1,b_2} A^{\sigma'_2}_{a'_1,a'_2} \right) \ldots 
\times W^{\sigma_{\ell},\sigma'_{\ell}}_{b_{\ell-1},b_{\ell}} \times \\
& &  \left( \sum_{\sigma_{\ell+1}\sigma'_{\ell+1}} B^{\sigma_{\ell+1}*}_{a_{\ell},a_{\ell+1}} W^{\sigma_{\ell+1},\sigma'_{\ell+1}}_{b_{\ell},b_{\ell+1}} B^{\sigma'_{\ell+1}}_{a'_{\ell},a'_{\ell+1}} \right) \ldots \left( \sum_{\sigma_L\sigma'_L} B^{\sigma_L*}_{a_{L-1},1} W^{\sigma_L,\sigma'_L}_{b_{L-1},1} B^{\sigma'_L}_{a'_{L-1},1} \right) .
\end{eqnarray*}   
MPO 表述的一切优美之处似乎都消失了，但图形表示能将其挽回（图~\ref{fig:mpodmrghamiltonian}）。从上面显式表达式的第二行或第三行最容易理解它：哈密顿量 MPO（以乘积基表达）投影到 A 块和 B 块的块态上，而这些块态在 $\fat{\sigma}$ 基中是有展开的。

\begin{figure}
\centering\includegraphics[scale=0.7]{PyMPS_MPOOnDMRGState.eps}
\caption{用 MPO/MPS 语言表示的 DMRG 表达式 $\bra{a_{\ell-1} \sigma_{\ell} a_{\ell}} \hat{H}  \ket{a'_{\ell-1} \sigma'_{\ell} a'_{\ell}}$。哈密顿量 MPO 与四个以 MPS 形式写出的块态展开缩并（两个左矢、两个右矢，其中两个在 A 块上、两个在 B 块上）。缩并后的网络解耦为 $L$、$W$ 和 $R$ 三部分，分别对应 A 块、B 块和中心格点。}
\label{fig:mpodmrghamiltonian}
\end{figure}

事实上，我们也可以把该表达式显而易见的三分结构编码为
\begin{equation}
\bra{a_{\ell-1} \sigma_{\ell} a_{\ell}} \hat{H}  \ket{a'_{\ell-1} \sigma'_{\ell} a'_{\ell}} 
= \sum_{b_{\ell-1}, b_{\ell}} L^{a_{\ell-1},a'_{\ell-1}}_{b_{\ell-1}}  W^{\sigma_{\ell},\sigma'_{\ell}}_{b_{\ell-1},b_{\ell}} R^{a_{\ell},a'_{\ell}}_{b_{\ell}} ,
\end{equation}
其中 $L$ 和 $R$ 包含图形网络缩并后的左部和右部：
\begin{eqnarray}
L^{a_{\ell-1},a'_{\ell-1}}_{b_{\ell-1}} &=& \sum_{\{ a_i, b_i, a'_i; i<\ell-1\} } \left( \sum_{\sigma_1\sigma'_1} A^{\sigma_1*}_{1,a_1} W^{\sigma_1,\sigma'_1}_{1,b_1} A^{\sigma'_1}_{1,a'_1} \right) \ldots \left( \sum_{\sigma_{\ell-1}\sigma'_{\ell-1}} A^{\sigma_{\ell-1} *}_{a_{\ell-2},a_{\ell-1}} W^{\sigma_{\ell-1},\sigma'_{\ell-1}}_{b_{\ell-2},b_{\ell-1}} A^{\sigma'_{\ell-1}}_{a'_{\ell-2},a'_{\ell-1}} \right) 
\label{eq:LdefineinMPO} \\
R^{a_{\ell},a'_{\ell}}_{b_{\ell}} &=& \sum_{\{ a_i, b_i, a'_i; i>\ell\} } 
\left( \sum_{\sigma_{\ell+1}\sigma'_{\ell+1}} B^{\sigma_{\ell+1}*}_{a_{\ell},a_{\ell+1}} W^{\sigma_{\ell+1},\sigma'_{\ell+1}}_{b_{\ell},b_{\ell+1}} B^{\sigma'_{\ell+1}}_{a'_{\ell},a'_{\ell+1}} \right) \ldots \left( \sum_{\sigma_L\sigma'_L} B^{\sigma_L*}_{a_{L-1},1} W^{\sigma_L,\sigma'_L}_{b_{L-1},1} B^{\sigma'_L}_{a'_{L-1},1} \right) 
\end{eqnarray}
现在，我们可以把 $\hat{H}$ 作用于处于混合规范形式或单格点 DMRG 表示的态 $\ket{\psi}$ 的结果写为
\begin{equation}
\hat{H}\ket{\psi} = \sum_{b_{\ell-1}, b_{\ell}} \sum_{a'_{\ell-1}, \sigma'_{\ell}, a'_{\ell}}
 L^{a_{\ell-1},a'_{\ell-1}}_{b_{\ell-1}}  W^{\sigma_{\ell},\sigma'_{\ell}}_{b_{\ell-1},b_{\ell}} R^{a_{\ell},a'_{\ell}}_{b_{\ell}} \Psi^{\sigma'_{\ell}}_{a'_{\ell-1},a'_{\ell}} \ket{a_{\ell-1}}_A \ket{\sigma_{\ell}} \ket{a_{\ell}}_B .
\end{equation}
我们马上就会讨论到，$\hat{H}\ket{\psi}$ 是迭代基态搜索中的关键操作。直接计算这个表达式慢得无法接受；可以在两个方面使其大幅加速：其一，$L$ 和 $R$ 可以迭代地构建，以最大限度地复用已有信息；这涉及网络缩并的最优安排。此外，$L$、$R$ 和 $W$ 对 $\ket{\psi}$ 的最终作用也可以安排得非常高效。

让我们首先考虑构建 $L$ 和 $R$。在实际应用中，我们绝不会执行它们背后完整的网络收缩，因为本着 DMRG 的精神，我们面对的是逐格点增大和缩小的块。不过，$L$ 和 $R$ 的构建方式是迭代的，恰好与块的生长和收缩直接对应。我将以 $L$ 为例加以说明，使用 $A$ 矩阵；左正规化只在某一点被显式地用于进一步简化，因此各公式都是一般性的。我们从考虑大小为 1 的块开始：将 $A^{[1]}$ 与 $A^{[1]\dagger}$ 同 $W^{[1]}$ 收缩。块的基表示则由下式给出
\begin{equation}
F^{[1]}_{a_1,b_1;a'_1} = \sum_{\sigma_1,\sigma_1',a_0,b_0,a'_0} W^{[1]\sigma_1\sigma'_1}_{b_0,b_1} (A^{[1]\sigma_1\dagger})_{a_1,a_0} F^{[0]}_{a_0,b_0,a'_0} A^{[1]\sigma'_1}_{a'_0,a'_1}
\end{equation}
其中我们引入了一个哑标量 $F^{[0]}_{a_0,b_0,a'_0}=1$，且 $a_0,b_0,a'_0$ 只能取值 1；这只是为了让第一步与后续各步保持一致。所得对象是一个张量 $F^{[1]}_{a_1,b_1,a'_1}$，对应于伸出的三条腿。

现在我们可以简单地继续在下一个格点上收缩 $A$、$A^\dagger$ 和 $W$，收缩更新为
\begin{equation}
F^{[i]}_{a_i,b_{i},a'_i} =  \sum_{\sigma_i,\sigma_i',a_{i-1},b_{i-1},a'_{i-1}} W^{[i]\sigma_i\sigma'_i}_{b_{i-1},b_{i}}  
(A^{[i]\sigma_i\dagger})_{a_i,a_{i-1}} F^{[i-1]}_{a_{i-1},b_{i-1},a'_{i-1}} A^{[i]\sigma'_i}_{a'_{i-1},a'_i} 
\end{equation}
其图示表示如图~\ref{fig:constructionEmatrices} 所示。
\begin{figure}
\centering\includegraphics[scale=0.7]{PyMPS_EMatrixUpdate.eps}
\caption{通过与 $A^{[i]*}$、$W^{[i]}$ 和 $A^{[i]}$ 收缩，从 $F^{[i-1]}$ 更新到 $F^{[i]}$。虽然在数学上把这三个新加入的张量视为一个对象是合理的，但在数值实践中，出于效率考虑，它们是被逐一收缩进网络的。}
\label{fig:constructionEmatrices}
\end{figure}

通过最优加括号方式，这一构建可以最有效地计算为
\begin{equation}
F^{[i]}_{a_i,b_{i},a'_i} = \sum_{\sigma_i,a_{i-1}}
(A^{[i]\sigma_i\dagger})_{a_i,a_{i-1}} 
\left( \sum_{\sigma_i',b_{i-1}}
W^{[i]\sigma_i\sigma'_i}_{b_{i-1},b_{i}}
\left( \sum_{a'_{i-1}} F^{[i-1]}_{a_{i-1},b_{i-1},a'_{i-1}} A^{[i]\sigma'_i}_{a'_{i-1},a'_i}  \right) \right) .
\end{equation}
这里，我们把三个新张量逐一收缩进网络，其运算量分别是最内层括号 $O(dD^3 D_W)$、其次 $O(d^2 D^2 D_W^2)$、最后 $O(d D^3 D_W)$。实际上，第二步运算在实践中更快，因为我们知道 $\hat{W}$ 中大多数算符干脆就是零；剩余的算符通常也具有简单的结构。在从左正规化矩阵构建 $L$ 且指标 $b_i=D_W$ 的情形下，如果我们按照上一节概述的规则构建 $\hat{H}$，还可以进一步加速：我们知道在这种情况下只有恒等算符向左作用，这意味着 $F^{[i]}_{a_i,D_W,a'_i} = \delta_{a_i,a'_i}$，从而同时简化了最内层括号和最外层运算。同样的想法也适用于右侧指标 $b_i=1$ 的情形，即从右正规化矩阵构建 $R$。

请注意，这一构建是针对 DMRG 所解释的表示更新的推广：一个典型情形是 $\hat{O}_i \hat{O}_j$ 的表示，其中两个算符分别局域地作用在格点 $i$ 和 $j$ 上。此时 MPO 处处维数为 $(1\times 1)$，且除格点 $i$ 和 $j$ 外，处处有 $W^{\sigma\sigma'} = \delta_{\sigma,\sigma'}$，而在这两个格点上则为 $W^{\sigma_i\sigma'_i} = O^{[i]\sigma_i,\sigma'_i}$，$j$ 处类似。向前推进收缩，$F^{[i-1]}$ 仍是标量 1。于是
\begin{equation}
F^{[i]} =  \sum_{\sigma_i,\sigma_i'} O^{[i]\sigma_i,\sigma'_i} A^{[i]\sigma_i\dagger} A^{[i]\sigma'_i}
\end{equation}
其中 $F^{[i]}$、$A^{[i]\sigma_i\dagger}$ 和 $A^{[i]\sigma'_i}$ 都是矩阵，且暗含乘积 $A^{[i]\sigma_i\dagger} A^{[i]\sigma'_i}$。$F^{[i]}$ 矩阵正是涵盖格点 1 到 $i$ 的块基中的算符表示。

更新到格点 $j-1$ 之前，各步都简化为
\begin{equation}
F^{[k]} =  \sum_{\sigma_k} A^{[k]\sigma_k\dagger} F^{[k-1]} A^{[k]\sigma_k} ,\quad (i<k<j)
\end{equation}
（暗含矩阵乘法），而在格点 $j$ 处我们再次得到一个非平凡的步骤，
\begin{equation}
F^{[j]} = \sum_{\sigma_j,\sigma'_j} O^{[j]\sigma_j,\sigma'_j} A^{[j]\sigma_j\dagger} F^{[j-1]} A^{[j]\sigma'_j},
\end{equation}
此后的更新照旧按前面各格点的方式进行。若把矩阵乘法显式写出，就会看到这正是 DMRG 算法中讨论过的构建。
  
最后，$\hat{H}\ket{\psi}$ 可以有利地按如下方式加括号：
\begin{equation}
\hat{H}\ket{\psi} = \sum_{b_{\ell-1}, a'_{\ell-1}} 
 L^{a_{\ell-1},a'_{\ell-1}}_{b_{\ell-1}} 
\left( \sum_{b_{\ell} \sigma'_{\ell}} 
W^{\sigma_{\ell},\sigma'_{\ell}}_{b_{\ell-1},b_{\ell}}
 \left( \sum_{a'_{\ell}} R^{a_{\ell},a'_{\ell}}_{b_{\ell}} \Psi^{\sigma'_{\ell}}_{a'_{\ell-1},a'_{\ell}} \right) \right) \ket{a_{\ell-1}}_A \ket{\sigma_{\ell}} \ket{a_{\ell}}_B ,
\end{equation}
其最坏标度行为为 $O(D^3)$。更精确地说，最内层运算为 $O(D^3 D_W d)$；下一层为 $O(D^2 D_W^2 d^2)$，再往后是一个代价为 $O(D^3 D_W^2 d^2)$ 的求和。跟踪 $W$ 的结构是有利的，即弄清对哪些 $(b_{\ell-1},b_{\ell})$ 取值组合它为零从而无需计算（通常占大多数），并在态处于混合规范形式时利用刚才讨论的关于 $L$ 和 $R$ 的简化。
 
\subsection{迭代基态搜索}

现在让我们转向算法本身。设 $\hat{H}$ 以 MPO 形式给出，并考虑一类具有预定矩阵维数的 MPS（只需设想一个由形状和大小符合要求的矩阵 $M^\sigma$ 构成的随机 MPS，暂且不假设其归一化）。为了在这一类态中找到对基态的最优逼近，我们必须找到使下式取极小值的 MPS $\ket{\psi}$：
\begin{equation}
E = \frac{\bra{\psi} \hat{H} \ket{\psi}}{\braket{\psi}{\psi}} .
\end{equation} 
事实证明，这可以转化为一个远比从某个随机态出发做虚时演化高效的基态算法。为了求解此问题，我们引入拉格朗日乘子 $\lambda$，并对
\begin{equation}
\bra{\psi} \hat{H} \ket{\psi} - \lambda {\braket{\psi}{\psi}} ;
\label{eq:Lagrange}
\end{equation}
求极值；最终，$\ket{\psi}$ 将是所求的基态，而 $\lambda$ 则是基态能量。表示式~(\ref{eq:Lagrange}) 的 MPS 网络如图~\ref{fig:extremizationexpression} 所示。
\begin{figure}
\centering\includegraphics[scale=0.7]{PyMPS_MPSEnergyMinimization.eps}
\caption{为求得用于寻找基态及其能量而需取极值的泛函所要收缩的网络。左边表示 $\bra{\psi} \hat{H} \ket{\psi}$ 项，右边表示平方范数 $\braket{\psi}{\psi}$。}
\label{fig:extremizationexpression}
\end{figure}

这一方法的问题在于，变量（矩阵元 $M^\sigma_{aa'}$）以乘积的形式出现，使其成为一个高度非线性的优化问题。但它同样可以迭代地进行，而这正是驱动 DMRG 的同一思想：保持除一个格点（$\ell$）之外所有格点上的矩阵不变，只把格点 $\ell$ 上的矩阵元 $M^{\sigma_\ell}_{a_{\ell-1}a_\ell}$ 当作变量。于是变量在式~(\ref{eq:Lagrange}) 中只以二次形式出现，对它求极值是一个良性的线性代数问题。这会降低能量，找到一个变分上更好的态，但当然不是最优的。接着继续变动另一格点上的矩阵元，寻找能量更低的态，如此遍历所有格点多轮，直到能量不再改善为止。

让我们先考虑重叠的计算，同时把选定的 $M^{\sigma_\ell}$ 保持显式写出。我们发现
\begin{equation}
\braket{\psi}{\psi} = \sum_{\sigma_\ell} \sum_{a_{\ell-1}a_{\ell}} \sum_{a'_{\ell-1}a'_{\ell}} \Psi^A_{a_{\ell-1},a'_{\ell-1}} M^{\sigma_\ell*}_{a_{\ell-1}, a_{\ell}} M^{\sigma_\ell}_{a'_{\ell-1}, a'_{\ell}} \Psi^B_{a_\ell, a_\ell'},
\end{equation}
其中
\begin{eqnarray}
\Psi^A_{a_{\ell-1},a'_{\ell-1}} &=& \sum_{\sigma_1,\ldots,\sigma_{\ell-1}} (M^{\sigma_{\ell-1}\dagger} \ldots M^{\sigma_1\dagger} M^{\sigma_1} \ldots M^{\sigma_{\ell-1}})_{a_{\ell-1},a'_{\ell-1}} \\
\Psi^B_{a_\ell, a'_\ell} &=& \sum_{\sigma_{\ell+1},\ldots,\sigma_L} (M^{\sigma_{\ell+1}} \ldots M^{\sigma_L} M^{\sigma_L\dagger} \ldots M^{\sigma_{\ell+1}\dagger})_{a'_\ell,a_\ell} .
\end{eqnarray}
如图形表示所特别清楚地显示的那样，在得到最后两个表达式时，适用于重叠的同一套聪明收缩规则同样适用；此外，如果我们在格点 $\ell$ 上从一个相邻格点移动到下一个相邻格点，它们可以迭代地更新，从而使计算代价最小。当格点 1 到 $\ell-1$ 为左正规化、格点 $\ell+1$ 到 $L$ 为右正规化时，归一化条件带来进一步的简化，即
\begin{equation} 
\Psi^A_{a_{\ell-1},a'_{\ell-1}} = \delta_{a_{\ell-1},a'_{\ell-1}} \quad\quad  \Psi^B_{a_\ell a'_\ell} = \delta_{a_\ell a'_\ell} .
\end{equation}

现在让我们考虑 $\hat{H}$ 以 MPO 语言表述时的 $\bra{\psi} \hat{H} \ket{\psi}$。结合上一节对 $\hat{H}\ket{\psi}$ 的分析，我们可以立即写出
\begin{equation}
\bra{\psi} \hat{H} \ket{\psi} = \sum_{\sigma_\ell,\sigma'_\ell} \sum_{a'_{\ell-1}a'_{\ell}}  \sum_{a_{\ell-1}a_{\ell}} \sum_{b_{\ell-1},b_\ell} L^{a_{\ell-1},a'_{\ell-1}}_{b_{\ell-1}} W^{\sigma_\ell,\sigma'_\ell}_{b_{\ell-1},b_\ell} R^{a_\ell,a'_\ell}_{b_\ell} M^{\sigma_\ell*}_{a_{\ell-1},a_\ell} M^{\sigma'_\ell}_{a'_{\ell-1},a'_\ell} 
\end{equation}
其中 $L$ 和 $R$ 如前面所定义；至于如何高效地求值这样的表达式，前文已经讨论过。

现在若对 $M^{\sigma_\ell*}_{a_{\ell-1},a_\ell}$ 取式~(\ref{eq:Lagrange}) 的极值，我们得到
\begin{equation}
\sum_{\sigma'_\ell}  \sum_{a'_{\ell-1}a'_{\ell}} \sum_{b_{\ell-1},b_\ell} L^{a_{\ell-1},a'_{\ell-1}}_{b_{\ell-1}} W^{\sigma_\ell,\sigma'_\ell}_{b_{\ell-1},b_\ell} R^{a_\ell,a'_\ell}_{b_\ell}  M^{\sigma'_\ell}_{a'_{\ell-1},a'_\ell} 
 - \lambda  \sum_{a'_{\ell-1}a'_{\ell}} \Psi^A_{a_{\ell-1},a'_{\ell-1}} \Psi^B_{a_\ell a'_\ell} M^{\sigma_\ell}_{a'_{\ell-1}, a'_{\ell}} = 0.
 \label{eq:vMPSeigenequation}
\end{equation}
这实际上是一个非常简单的本征值方程；如果我们通过重排指标引入矩阵 $H$ 和 $N$：$H_{(\sigma_\ell a_{\ell-1} a_\ell), (\sigma'_\ell a'_{\ell-1} a'_\ell)} = \sum_{b_{\ell-1},b_\ell} L^{a_{\ell-1},a'_{\ell-1}}_{b_{\ell-1}} W^{\sigma_\ell,\sigma'_\ell}_{b_{\ell-1},b_\ell} R^{a_\ell,a'_\ell}_{b_\ell} $ 以及
$N_{(\sigma_\ell a_{\ell-1} a_\ell), (\sigma'_\ell a'_{\ell-1} a'_\ell)} = \Psi^A_{a_{\ell-1},a'_{\ell-1}} \Psi^B_{a_\ell, a'_\ell} \delta_{\sigma_\ell,\sigma'_\ell}$，以及满足 $v_{\sigma_\ell a_{\ell-1} a_\ell} = M^{\sigma_\ell}_{a_{\ell-1}, a_{\ell}}$ 的矢量 $v$，我们就得到一个{\em 广义本征值问题}
矩阵维数为 $(dD^2 \times dD^2)$，
\begin{equation}
H v - \lambda N v = 0,
\end{equation}
其图示见图~\ref{fig:equationsystem}。
解出最低本征值 $\lambda_0$ 就得到 $v^0_{\sigma_\ell a_{\ell-1} a_\ell}$，再把它重排回 $M^{\sigma_\ell}_{a_{\ell-1}a_{\ell}}$，而 $\lambda_0$ 即为当前的基态能量估计值。

\begin{figure}
\centering\includegraphics[scale=0.7]{PyMPS_MPSEnergyEquationSystem.eps}
\caption{用于优化 $M^{\ \sigma_\ell}_{\ a_{\ell-1},a_\ell}$ 的广义本征值问题。未知矩阵在左、右两个网络中被圈出。}
\label{fig:equationsystem}
\end{figure}

有几点评注需要说明。

\begin{itemize}
\item 该问题是厄米的；从构造方式以及所用 MPO 的厄米性可以看出，$H$ 和 $N$ 都是厄米的。
\item 一般来说，$dD^2$ 对于精确对角化而言太大了，但由于我们只关心最低的本征值和本征态，一个瞄准谱端部的迭代本征求解器就足够了。典型的方法有 Lanczos 或 Jacobi-Davidson 大型稀疏矩阵求解器。这类方法的收敛速度最终取决于初始出发向量或猜测向量的质量。由于这个本征值问题是迭代求基态方法的一部分，当前的 $M^{\sigma_\ell}$ 就是一个有效的猜测，在接近收敛时会大大加快计算。
\item 当 $N$ 的条件数变差时，广义本征值问题在数值上可能非常苛刻。但对于开边界条件，只要保证态在格点 $\ell-1$ 及其之前是左规范的、从格点 $\ell+1$ 起是右规范的，这就不成问题。于是
对 $\Psi^A$ 和 $\Psi^B$ 的种种简化意味着 $N$ 就是单位矩阵 $I$。本征值问题随之简化为一个标准问题，
\begin{equation}
\sum_{\sigma'_\ell} \sum_{a'_{\ell-1}a'_{\ell}} \sum_{b_{\ell-1},b_\ell}
L^{a_{\ell-1},a'_{\ell-1}}_{b_{\ell-1}} W^{\sigma_\ell,\sigma'_\ell}_{b_{\ell-1},b_\ell} R^{a_\ell,a'_\ell}_{b_\ell}
M^{\sigma'_\ell}_{a'_{\ell-1}, a'_\ell} - \lambda  M^{\sigma_\ell}_{a_{\ell-1}, a_{\ell}} = 0 .
\end{equation}
即 $Hv-\lambda v=0$，如图~\ref{fig:normalizedequationsystem} 所示。各求和的计算将利用 $\hat{H}\ket{\psi}$ 的最优加括号方式来完成。
为实现这一简化，需让位置 $\ell$ 沿链从右到左再反向往复扫描，使得每次极小化之后只需执行一步左规范化或右规范化程序，就能维持最优的规范化配置。
\end{itemize}

\begin{figure}
\centering\includegraphics[scale=0.7]{PyMPS_MPSNormalizedEnergyEquationSystem.eps}
\caption{用于优化 $M^{\ \sigma_\ell}_{\ a_{\ell-1},a_\ell}$ 的标准本征值问题。未知矩阵在左侧网络中被圈出。}
\label{fig:normalizedequationsystem}
\end{figure}

于是最优算法按如下方式进行。
\begin{itemize}
\item 从对 $\ket{\psi}$ 的某个初始猜测出发，该态是右规范的，即仅由 $B$ 矩阵组成。
\item 迭代地计算从格点位置 $L-1$ 到 1 的所有 $R$ 表达式。
\item {\em 右扫描：}从格点 $\ell=1$ 到格点 $L-1$，按如下方式沿晶格向右扫描：用迭代本征求解器求解 $M^{\sigma_\ell}$ 的标准本征值问题，以其当前值作为出发点。得到解之后，通过 SVD（或 QR）把 $M^{\sigma_\ell}$ 左规范化为 $A^{\sigma_\ell}$，以维持所需的规范化结构。SVD 余下的矩阵被乘到右边的 $M^{\sigma_{\ell+1}}$ 上，作为下一格点本征求解器的初始猜测。通过再纳入一个格点迭代地构造 $L$ 表达式。前进一个格点，$\ell \rightarrow \ell+1$，然后重复。
\item {\em 左扫描：}从格点 $\ell=L$ 到格点 $2$，按如下方式沿晶格向左扫描：用迭代本征求解器求解 $M^{\sigma_\ell}$ 的标准本征值问题，以其当前值作为出发点。得到解之后，通过 SVD（或 QR）把 $M^{\sigma_\ell}$ 右规范化为 $B^{\sigma_\ell}$，以维持所需的规范化结构。SVD 余下的矩阵被乘到左边的 $M^{\sigma_{\ell-1}}$ 上，作为下一格点本征求解器的初始猜测。通过再纳入一个格点迭代地构造 $R$ 表达式。前进一个格点，$\ell \rightarrow \ell-1$，然后重复。
\item 重复右扫描和左扫描，直至达到收敛。只要能量收敛即认为达到收敛，但最好的检验是（利用 MPO）考察
$\bra{\psi} \hat{H}^2 \ket{\psi} - (\bra{\psi} \hat{H} \ket{\psi})^2$，看是否已经得到某个本征态；这个表达式应尽可能地接近 0。
\end{itemize}

如果我们根据各自的规范化情况把矩阵记为 $A$、$B$、$M$（$M$ 始终指当前正在处理的格点上的那个矩阵），并给它们加上下标 $i$ 以标记它们被本征求解器更新的次数，则该算法可形式化为
\begin{eqnarray*}
& & M_0 B_0 B_0 B_0 B_0 B_0 \\
&\stackrel{diag}{\rightarrow}& M_1 B_0 B_0 B_0 B_0 B_0
\stackrel{SVD}{\rightarrow} A_1 M_0 B_0 B_0 B_0 B_0 \\
&\stackrel{diag}{\rightarrow}& A_1 M_1 B_0 B_0 B_0 B_0
\stackrel{SVD}{\rightarrow} A_1 A_1 M_0 B_0 B_0 B_0 \\
&\stackrel{diag}{\rightarrow}& A_1 A_1 M_1 B_0 B_0 B_0
\stackrel{SVD}{\rightarrow} A_1 A_1 A_1 M_0 B_0 B_0 \\
& & \ldots \\
&\stackrel{diag}{\rightarrow}& A_1 A_1 A_1 A_1 M_1 B_0
\stackrel{SVD}{\rightarrow} A_1 A_1 A_1 A_1 A_1 M_0 \\
&\stackrel{diag}{\rightarrow}& A_1 A_1 A_1 A_1 A_1 M_1
\stackrel{SVD}{\rightarrow} A_1 A_1 A_1 A_1 M_1 B_1 \\
&\stackrel{diag}{\rightarrow}& A_1 A_1 A_1 A_1 M_2 B_1
\stackrel{SVD}{\rightarrow} A_1 A_1 A_1 M_1 B_2 B_1 \\
& & \ldots \\
&\stackrel{diag}{\rightarrow}& A_1 M_2 B_2 B_2 B_2 B_1
\stackrel{SVD}{\rightarrow} M_1 B_2 B_2 B_2 B_2 B_1 \\
\end{eqnarray*}
然后再一次从左向右移动，从一次对角化步骤开始。

在这一迭代过程中，能量只会下降，因为我们通过改变参数不断地改进。会出现两个问题：若从随机态出发，最初几次扫描中迭代本征求解器对 $M^{\sigma_\ell}$ 的猜测会非常糟糕，导致迭代次数很大、性能很差。此外，我们无法保证这一过程真正达到全局极小，而不是陷在一个非全局的极小之中。

解决第一个问题的一种办法是先用无限系统 DMRG 产生一个初始猜测；无限系统 DMRG 的一个最优 MPS 版本将在第~\ref{sec:infinite} 节讨论。这个初始猜测虽然可能离真正的解还很远，但通常会比随机的出发态好得多。此外，可以尝试平衡迭代次数（最初几次扫描中很高）：先从小的 $D$ 出发，在该拟设类中收敛，然后增大 $D$ 并在新矩阵元处补零，再次收敛，如此往复。当 $D$ 变大时，猜测态有望已经非常接近最终态，从而只需很少的迭代。然而结果表明，若以过小的 $D$ 起步，可能落入一个非全局极小，并且在增大 $D$ 之后也无法脱身。很一般地，正如在 DMRG 中那样，绝不要只对单个 $D$ 计算结果，而应在多次运行中逐步增大它，直到结果收敛（在 $D\rightarrow\infty$ 极限下结果保证是精确的）。

如果我们寻找的是低能激发态而非基态，会出现两种典型情形：
(i) 已知该激发态是希尔伯特空间按某个好量子数分解后另一扇区的基态。那么计算不过是在那个不同扇区中做基态计算。(ii) 该激发态是基态所在扇区中的第一、第二或更高阶激发。那么我们必须迭代地计算这些激发，并让所求的态相对已经找到的更低能态保持正交归一；这显然把该方法限制在少数低能激发上。需要修改算法的地方在迭代本征求解器之中；例如在 Lanczos 迭代里，新生成的下一个 Lanczos 态不仅要相对先前的 Lanczos 态正交归一，还要相对已经构造出的哈密顿量本征态正交归一。这是 Lanczos 算法的一个标准推广。

刚刚介绍的变分 MPS 算法相当容易陷入困境。究竟会怎样陷入，实际上与过程中初始态的选取方式有一定关系：假定像各向异性海森堡链的情形那样，存在某个量子对称性，即有与之对易的算符 $[\hat{H},\hat{O}]=0$，此处即总磁化强度算符 $\hat{M}=\sum_i \hat{S}^z_i$，它使磁化强度 $M$ 成为一个好量子数。于是初始态分为两类，取决于它们是不是 $\hat{M}$ 的本征态。若态是随机选取的，通常属于后者；若由无限系统 DMRG 或其 MPS 变体生成，则属于前者。

把希尔伯特空间按磁化强度的本征态分解，${\mathcal H} = \oplus_M {\mathcal H}_M$，我们可以把初始态写成
\begin{equation}
\ket{\psi} = \sum_M \psi_M \ket{M} \quad\quad \ket{M} \in {\mathcal H}_M .
\end{equation}
我们要找的基态是 $\ket{\psi_0} \in {\mathcal H}_{\tilde{M}}$；于是初始态要么具有任意的 $\psi_M$，要么在 $M \neq \tilde{M}$ 时 $\psi_M = 0$（这里假定我们不会遭遇在错误对称扇区中给出初始态这一灾难）。让我们假定在第一种情形下，扫描会消去来自错误对称扇区的贡献；如果做不到，那么由于错误的混入得以幸存，变分最优态无论如何都无法达到。由于通过例如 Lanczos 进行的迭代基态搜索是求最大本征值及相应本征态的幂方法 $\lim_{n\rightarrow\infty} \hat{H}^n \ket{\psi}$ 的优化版本，可以证明在全希尔伯特空间中，错误的对称扇区一定会被投影掉。而在我们的算法中，这种迭代投影是在一个高度受限的态空间中进行的，由于它按顺序逐个查看各波函数分量，效率可能不高。由于随机的出发态效率极低，我在这里无法报告太多实际经验。
无论如何，一旦我们进入一个确定的对称扇区，对任意 Schmidt 分解 $\ket{\psi} = \sum_{a_\ell} s_{a_\ell} \ket{a_\ell}_A \ket{a_\ell}_B$ 而言，每个态都会带有好量子数（不同量子数的叠加会立即与全局好量子数矛盾），即 $m_a^A$ 和 $m_a^B$ 满足 $m_a^A + m_a^B = \tilde{M}$，这里我对指标做了简化。以 $m_a^B$ 为例，它们会分布在某个范围内，比如说 1 个态带磁化强度 $m$，3 个态带磁化强度 $m'$，5 个态带磁化强度 $m''$，依此类推。正如我接下来要展示的，这一分布在后续扫描中保持{\em 不变}。这意味着，如果它不对应于变分最优态所给出的分布，就永远无法到达那个态。在随机态方案中，人们或许指望其他总磁化强度被缓慢消去能把我们``抚顺''到正确的分布上，但这没有任何保证；在无限系统方案中，人们只能指望这种热身方案立刻产生正确的分布，而这相当不可能发生。

这一分布保持不变的原因，可以从对 $M^{\sigma_{\ell}}_{a_{\ell-1},a_{\ell}}$ 做 SVD 以执行一步（例如）左规范化中看出来：把矩阵 $M^{\sigma_{\ell}}$ 重新整理成某个 $\Psi$ 并做 SVD，至多给出 $D$ 个非零奇异值；$V^\dagger$ 中的右奇异向量正是 $\Psi^\dagger \Psi$ 的本征向量，而后者是块对角的，因为态 $\ket{a_{\ell}}_B$ 带有好量子数。因此，右奇异向量（本征向量）编码的是同一量子数的块内部的基变换，于是具有给定量子数的态的数目保持不变；又由于 Schmidt 分解所要求的量子数匹配，系统另一部分中具有该量子数的态的数目也保持不变。

人们已经提出了多种摆脱这一潜在陷阱的办法。第一种是修改算法以同时考虑两个格点，就像常规（两格点）DMRG 那样；我们将在下一节讨论其 MPS 实现。这一做法虽然更慢（大约慢一个 $d$ 因子），但它提供了稍稍扩大的拟设空间以及随之而来的截断，使算法在基态搜索中更不易陷入局部能量极小。特别地，两格点算法扩大了的拟设空间允许通过截断对量子数分布进行重组。一旦这一步收敛，就可以切换到单格点算法，正如 Takasaki 等人 \cite{Takasaki99} 所建议的那样，尽管这样做是否严格地导向最优结果并不清楚\cite{McCulloch08}。

好得多的是，White \cite{White05} 提出了一种能相当可靠地防止陷入并保证重组的程序。它对于可靠的单格点 DMRG（或变分 MPS，我们将证明二者等价）算法至关重要，并使之成为该方法的最新最高水平形式。它始于这样一个观察：子系 A 的量子数会因哈密顿量中把 A 与系统其余部分连接起来的那些部分所引起的量子涨落而改变。因此我们必须更仔细地考察 $\hat{H}\ket{\psi}$ 的结构。

考虑单格点算法中{\em 能量优化之后}的 $\psi^{\sigma_{\ell}}_{a_{\ell-1},a_{\ell}}$。我们也可以写成
\begin{equation}
\hat{H} = \sum_{b_{\ell}} \hat{H}^{A\bullet}_{b_{\ell}} \hat{H}^B_{b_{\ell}}
\end{equation}
其中
\begin{equation}
\hat{H}^{A\bullet}_{b_{\ell}} = \sum_{\fat{\sigma}, \fat{\sigma}'\in A\bullet} ( W^{\sigma_1,\sigma'_1} \ldots W^{\sigma_{\ell},\sigma'_{\ell}} )_{b_{\ell}} \quad
\hat{H}^B_{b_{\ell}} = \sum_{\fat{\sigma}, \fat{\sigma}'\in B} ( W^{\sigma_{\ell+1},\sigma'_{\ell+1}} \ldots W^{\sigma_L,\sigma'_L} )_{b_{\ell}} ,
\end{equation}
使得这个求和中共有 $D_W$ 项。如果从块态的角度思考，我们想知道在 $\hat{H}^{A\bullet}_{b_{\ell}}$ 的作用下，A$\bullet$ 上能到达哪些新的态。
把这一作用的结果投影到 A$\bullet$B 基上，它就是
\begin{eqnarray}
(\hat{H}^{A\bullet}_{b_{\ell}} \ket{\psi})^{a_{\ell-1}, \sigma_{\ell},a_{\ell}} &=&
\sum_{\sigma'_{\ell}} \sum_{a'_{\ell-1},b_{\ell-1}} \sum_{\{ a_i, b_i, a'_i; i<\ell-1\} } \left( \sum_{\sigma_1\sigma'_1} A^{\sigma_1*}_{1,a_1} W^{\sigma_1,\sigma'_1}_{1,b_1} A^{\sigma'_1}_{1,a'_1} \right) \ldots \times \nonumber \\
& & \left( \sum_{\sigma_{\ell-1}\sigma'_{\ell-1}} A^{\sigma_{\ell-1} *}_{a_{\ell-2},a_{\ell-1}} W^{\sigma_{\ell-1},\sigma'_{\ell-1}}_{b_{\ell-2},b_{\ell-1}} A^{\sigma'_{\ell-1}}_{a'_{\ell-2},a'_{\ell-1}} \right)  W^{\sigma_{\ell},\sigma'_{\ell}}_{b_{\ell-1},b_{\ell}} \Psi^{\sigma'_{\ell}}_{a'_{\ell-1},a_{\ell}}
\end{eqnarray}
而这正是
\begin{equation}
(\hat{H}^{A\bullet}_{b_{\ell}} \ket{\psi})^{a_{\ell-1}, \sigma_{\ell},a_{\ell}} =
\sum_{\sigma'_{\ell}} \sum_{a'_{\ell-1},b_{\ell-1}}
L^{a_{\ell-1},a'_{\ell-1}}_{b_{\ell-1}} W^{\sigma_{\ell},\sigma'_{\ell}}_{b_{\ell-1},b_{\ell}} \Psi^{\sigma'_{\ell}}_{a'_{\ell-1},a_{\ell}}
\end{equation}
其中用了来自式~(\ref{eq:LdefineinMPO}) 的 $L^{a_{\ell-1},a'_{\ell-1}}_{b_{\ell-1}}$，如图~\ref{fig:Whiteimprovement} 所直观显示。这表明实际的计算代价非常低，因为最复杂的部分我们已经完成了。

\begin{figure}
\centering\includegraphics[scale=0.7]{PyMPS_WhitePredictionObject.eps}
\caption{$(\hat{H}^{A\bullet}_{b_\ell}\ket{\psi})^{a_{\ell-1},\sigma_{\ell},a_{\ell}}$ 的图形表示：除一个 $W$ 张量和一个 $\Psi$ 张量之外，为得到 $L^{a_{\ell-1},a'_{\ell-1}}_{b_{\ell-1}}$ 所需的全部缩并都已经计算过了。}
\label{fig:Whiteimprovement}
\end{figure}

现在我们想把 $\hat{H}^{A\bullet}_{b_{\ell}}$ 生成的那些态纳入为 A$\bullet$ 寻找好基的过程之中。这里，DMRG 提供了瞄准多个态的可能性。常规算法接下来会计算 $\hat{\rho}^{A\bullet} = \tr_B \ket{\psi}\bra{\psi}$ 或 $\rho_{(a_{\ell-1} \sigma_{\ell}),(a'_{\ell-1}  \sigma'_{\ell})} = \sum_{a_{\ell}} \Psi^{\sigma_{\ell}}_{a_{\ell-1},a_{\ell}}\Psi^{\sigma'_{\ell}\dagger}_{a_{\ell},a'_{\ell-1}}$，求出本征值（奇异值的平方）、本征向量（左奇异向量），然后截断等等。但我们可以考察一个修改过的密度矩阵，它也把新的项考虑进来：
\begin{equation}
\hat{\rho}^{A\bullet} = \tr_B \ket{\psi}\bra{\psi} + \alpha \sum_{b_{\ell}} \tr_B \hat{H}^{A\bullet}_{b_{\ell}} \ket{\psi} \bra{\psi} \hat{H}^{A\bullet}_{b_{\ell}} ,
\end{equation}
其中 $\alpha$ 是一个小数，比如 $10^{-4}$，它给常规算法可能遗漏的贡献赋予一点权重。付出的代价是：谱的末端会有少数几个来自 $\ket{\psi}\bra{\psi}$ 的权重极小的态被舍弃。经过多次扫描之后，$\alpha$ 会被缓慢地趋于零。

接着对新的密度矩阵进行对角化，并按最大的 $D$ 个本征值作截断，得到一个列正交归一的 $(dD \times D)$ 矩阵，$U_{(a_{\ell-1}\sigma_{\ell}), a_{\ell}} \rightarrow A^{\sigma_{\ell}}_{a_{\ell-1}, a_{\ell}}$，然后我们在下一个格点上继续；由于我们只利用修正后密度矩阵本征值的相对次序，而不作他用，因此它们之和不等于 1 也无关紧要。为给大型稀疏矩阵求解器预测下一个 $\Psi^{\sigma_{\ell+1}}$，我们使用下一节推导的 DMRG 预测公式，
\begin{equation}
\Psi^{\sigma_{\ell+1}}_{a_{\ell}a_{\ell+1}} = 
\sum_{a_{\ell-1}\sigma_{\ell} a_{\ell}'} 
A^{\sigma_{\ell} \dagger}_{a_{\ell},a_{\ell-1}} \Psi^{\sigma_{\ell}}_{a_{\ell-1}a_{\ell}'}  B^{\sigma_{\ell+1}}_{a_{\ell}', a_{\ell+1}} .
\end{equation}
除此之外，其余一切保持不变。对于一种通常收敛快得多、也远不容易陷入非最优结果的算法而言，额外的数值开销与编程工作量是微不足道的。

\subsection{用 MPS 语言表述的传统 DMRG：微妙的差别}   
\label{subsec:conventionalDMRGinMPS}
前一节的做法与传统 DMRG 有何关联？本质性的答案是：对于开边界条件（OBC），MPS 方法与有限系统 DMRG 完全相同，只是前提是我们每次移动一个格点而不是两个，
亦即\ 考虑``单格点''而非``双格点''DMRG：此时考虑的是块-格点-块构型 A$\bullet$B，而不是块-格点-格点-块构型 A$\bullet\bullet$B。

让我们先回顾一下这些算法的关键步骤，假定扫描方向向右：(i) 给定某个构型 A$\bullet$B（或相应的构型 $AAAMBB$）和哈密顿量 $\hat{H}$，用大型稀疏矩阵本征值求解器寻找最优的 $\psi_{a_{\ell-1} \sigma_\ell a_\ell}$（DMRG 中）或 $M^{\sigma_\ell}_{a_{\ell-1}, a_\ell}$（MPS 中）以求得基态；A$\bullet\bullet$B 情形类似。(ii) 得到基态后，MPS 导出一组左规范化的 $A$-矩阵，而 DMRG 则找到新的块态，其结构可以用左规范化的 $A$-矩阵编码。(iii) 所有算法都切换到新的 A$\bullet$B、A$\bullet\bullet$B 或 $AAAAMB$ 构型，其中活动中心向右移动一个格点，并为下一个基态的计算提供初猜，然后回到步骤 (i)。

步骤 (i)：如果使用相同的态构型和相同的哈密顿量，结果必须相同。由于 DMRG 从左右两侧分别生长块 A 和块 B，而每个块生长步骤 A$\bullet \rightarrow$A 都可以用 $A$-矩阵编码，$\bullet$B $\rightarrow$ B 同理，我们可以断言 A 上的所有矩阵都是左规范化的，B 上的所有矩阵都是右规范化的，于是双格点 DMRG 态取如下形式
\begin{equation}
\ket{\psi} = \sum_{a_{\ell-1} \sigma_\ell \sigma_{\ell+1} a_{\ell+1}} \Psi^{\sigma_\ell\sigma_{\ell+1}}_{a_{\ell-1}, a_{\ell+1}} \ket{a_{\ell-1}}_A \ket{\sigma_\ell} \ket{\sigma_{\ell+1}} \ket{a_{\ell+1}}_B = \sum_{\fat{\sigma}} A^{\sigma_1} \ldots A^{\sigma_{\ell-1}} 
\Psi^{\sigma_\ell \sigma_{\ell+1}} B^{\sigma_{\ell+2}} \ldots B^{\sigma_L} \ket{\fat{\sigma}},
\end{equation}
其中单格点 DMRG 态只需作显而易见的替换 $\Psi^{\sigma_\ell\sigma_{\ell+1}} \rightarrow \Psi^{\sigma_\ell}$。这与变分 MPS 方法中的混合规范形式态完全一致，我们面对的是同样的态结构。 

尚需证明两者的哈密顿量也相同。严格地说，情况并非如此：
我们刚才使用的 $\hat{H}$ 的 MPO 表示显然是精确的。另一方面，DMRG 中 $\hat{H}$ 的表示包含一系列约化基变换，因此本质上是不精确的。这样看来，两种表示似乎互不相干，且 MPO 一方因精确而占优。但更仔细的分析表明，在计算 MPS 与 DMRG 基态搜索中出现的期望值 $\bra{\psi} \hat{H} \ket{\psi}$ 的层面上，两种表示给出完全相同的结果（对于更高的矩，例如 $\bra{\psi} \hat{H}^2 \ket{\psi}$，两者并不相同，此时 MPO 表示以一定数值代价换取明显更高的精度，见下文）。

DMRG 哈密顿量与 MPO 哈密顿量都包含精确哈密顿量的全部项。正如我们在把哈密顿量 MPO 作用于混合规范形式态时已经看到的（第~\ref{subsec:Hamiltonianmixedcanonical} 节），大型稀疏本征值问题即式~(\ref{eq:vMPSeigenequation}) 中出现的 $L$ 与 $R$ 对象的计算，无非就是 DMRG 中直到当前 A$\bullet$B 构型为止的那一串约化基变换。因此，对于 $\bra{\psi} \hat{H} \ket{\psi}$（但一般也仅限于此量！），两种方法完全相同。

此外，本征值问题中出现的运算 $\hat{H}\ket{\psi}$，其操作计数并不比相应的 DMRG 流程更差。为看清这一点，让我们聚焦于各向异性最近邻海森堡链的那个示例 MPO。考虑对 $b_{\ell-1}, b_{\ell}$ 的双重求和时，效率上似乎存在差别。从 $W$-矩阵的结构可以清楚地看出，$D_W^2$ 个（示例中为 25 个）矩阵元中大多数是零，因此我们可以强烈地限制求和范围。但即便如此，计数结果仍是：把场设为 0，对海森堡哈密顿量而言，DMRG 框架中有 8 项贡献：严格作用于 A 和 B 的哈密顿量部分 $\hat{H}_A$ 与 $\hat{H}_B$ 各一项，每个块三项，对应于连接块与格点的三种算符组合。这些项在另一个块中都是对角的，故总共有 8 个代价为 $O(D^3)$ 的操作。在 MPO 计算中，双矩阵-矩阵乘法让人天真地估计为 16 个代价 $O(D^3)$ 的操作，对应 $W^{\sigma_{\ell},\sigma'_{\ell}}$ 的 8 个非零矩阵元。但我们随后可以利用以下规则：若 $b_{\ell-1}=d_W$，则左边没有操作，$L^{a_{\ell-1}, a'_{\ell-1}}_{d_W}= \delta_{a_{\ell-1}, a'_{\ell-1}}$，于是一个操作被省去；类似地，若 $b_{\ell}=1$，则右边没有操作，$R^{a_{\ell},a'_{\ell}}_{b_{\ell}}=\delta_{a_{\ell},a'_{\ell}}$，又省去一个操作。审视 $W$ 的结构可知，所有非零矩阵元都满足这两个条件之一，于是计数减半至 8 个操作。只有更远程的相互作用不符合这一图景，但它们在 DMRG 中同样要付出 $2O(D^3)$ 的代价。

步骤 (ii)：能量极小化之后，变分 MPS 与 DMRG 产生（相同的）$M^{\sigma_\ell}$ 与 $\Psi^{\sigma_\ell}$，或 $\Psi^{\sigma_\ell \sigma_{\ell+1}}$。两种方法随后对这个结果的处理看起来不同，实际上却是一回事：在变分 MPS 中，做一次 SVD 以保证左规范化，然后向右移动一个格点，继续在下一个格点上极小化。在 DMRG 中，则先做密度矩阵分析以确定新的（截断后的）块基。但如果执行相应的 SVD，就会发现非零奇异值（因而非零密度矩阵本征值）的个数受矩阵维数 $D$ 的限制：
\begin{equation}
\Psi^{\sigma_\ell}_{a_{\ell-1},a_ \ell} \rightarrow \Psi_{(a_{\ell-1} \sigma_ \ell), a_ \ell} = \sum_{k=1}^{\min(dD,d)=D} A^{\sigma_ \ell}_{a_{\ell-1},k} S_{kk} (V^\dagger)_{k, a_ \ell} .
\end{equation}
因此，这里没有发生任何截断，我们只是在做一次酉变换，为新的更大的块 A 获得正交归一的态（由于 SVD 与密度矩阵对角化之间的联系，这恰好就是 MPS 中的左规范化）。两种表述的作用完全相同；既然没有截断发生，用瘦 QR 分解也同样可行。

另一方面，在双格点 DMRG 中，同一步骤变为
\begin{equation}
\Psi^{\sigma_ \ell\sigma_{\ell +1}}_{a_{\ell-1},a_{\ell +1}} \rightarrow \Psi_{(a_{\ell-1} \sigma_ \ell), (\sigma_{\ell +1} a_{\ell +1})} = \sum_{k=1}^{\min(dD,dD)=dD} A^{\sigma_ \ell}_{a_{\ell-1},k} S_{k,k} (V^\dagger)_{k,(\sigma_{\ell +1} a_{\ell +1})} .
\end{equation}
但新块中只能保留 $D$ 个态，因此必须发生截断！这就是变分 MPS 与单格点 DMRG 一方、双格点 DMRG 另一方之间{\em 唯一}的差别。

步骤 (iii)：在 DMRG 中，一次迭代完成后，自由格点向右移动一位，导致块 A 生长、块 B 收缩。在这一步上，所有方法再次一致：在变分 MPS 中，B 的收缩只不过体现在态由一串 $B$-矩阵构成、而最左边的那个已经脱落；A 的生长则由类似的一串矩阵给出，只是新加进了一个 $A$-矩阵。自由格点上的矩阵在所有方法中都是待定的，因此关于它的规范化无话可说。

如前所述，基态能量的极小化是一个代价高昂的大型稀疏矩阵问题。由于这些方法是迭代的，一个好的初猜是可取的。DMRG 为此提供了某种``态预测''\cite{White96}。事实上，人们发现预测的结果恰恰就是在变分 MPS 语言中自然而然得到的东西，而无需为寻找态预测耗费智力上的努力。

设在单格点 DMRG 中我们刚刚优化了 $\Psi^{\sigma_\ell}$，导出了新的 $A^{\sigma_ \ell}$。那么在 MPS 语言中，下一个态是 $\Psi^{\sigma_{\ell +1}}= SV^{\dagger} B^{\sigma_{\ell +1}}$，其中 $S$ 与 $V^\dagger$ 来自 SVD。在 DMRG 语言中，我们取
\begin{equation}
\ket{\psi} = \sum_{a_{\ell-1}\sigma_ \ell a_{\ell }'} \Psi^{\sigma_ \ell}_{a_{\ell-1}, a_{\ell }} \ket{a_{\ell-1}}_A \ket{\sigma_ \ell} \ket{a_{\ell}'}_B
\end{equation}
并两次插入近似恒等算符 $I= \sum \ket{a_ \ell}_A\phantom\langle_A\bra{a_ \ell}$ 与
$I= \sum \ket{\sigma_{\ell +1}}\ket{a_{\ell+1}}_B \phantom\langle_B\bra{a_{\ell +1}} \bra{\sigma_{\ell+1}}$。用 $A$ 与 $B$ 矩阵表示矩阵元后，该态现在读作
\begin{equation}
\ket{\psi} = \sum_{a_\ell\sigma_{\ell +1} a_{\ell +1}} \left( \sum_{a_{\ell-1}\sigma_ \ell  a_{\ell}'} 
A^{\sigma_ \ell \dagger}_{a_ \ell,a_{\ell-1}} \Psi^{\sigma_ \ell}_{a_{\ell-1}, a_{\ell}'}  B^{\sigma_{\ell +1}}_{a_{\ell}' a_{\ell +1}} \right) 
\ket{a_{\ell}}_A \ket{\sigma_{\ell +1}} \ket{a_{\ell +1}}_B
\end{equation}
于是预测公式为
\begin{equation}
\Psi^{\sigma_{\ell +1}}_{a_{\ell}, a_{\ell+1}} = 
\sum_{a_{\ell-1}\sigma_ \ell a_{\ell}'} 
A^{\sigma_ \ell \dagger}_{a_ \ell,a_{\ell-1}} \Psi^{\sigma_ \ell}_{a_{\ell-1}, a_{\ell}'}  B^{\sigma_{\ell +1}}_{a_{\ell}' a_{\ell +1}} .
\end{equation}
但这正是下一个本征值问题的 MPS 拟设，因为 $ A^{\sigma_ \ell \dagger} \Psi^{\sigma_ \ell}  B^{\sigma_{\ell +1}} =
A^{\sigma_ \ell \dagger} A^{\sigma_ \ell } S V^{\dagger} B^{\sigma_{\ell +1}}$，而这又正是
$S V^{\dagger} B^{\sigma_{\ell +1}}$，因为在拟设中对 $\sigma_ \ell $ 求和且左规范化成立。双格点 DMRG 可依此类推，留作读者的练习。

这澄清了变分 MPS、单格点 DMRG（二者相同）与双格点 DMRG（不同）之间的关系。但必须指出，信息存储方式或更隐式或更显式的差别，即使在算法严格说来完全相同的情况下也会带来差异——预测在一种表述中轻而易举而在另一种中却非如此，这一事实已经给了我们一个例子。此外还有更多。

(i) 在 DMRG 中，表示态的等效基与表示哈密顿量或其他算符的等效基是捆绑在一起的。这就是为什么当我们同时考虑若干不同态（如基态与第一激发态）时，会出现诸如多态瞄准这样的概念。此时人们考虑混合约化密度算符
\begin{equation}
\hat{\rho}_A = \sum_{i} \alpha_i \tr_B \ket{\psi_i}\bra{\psi_i}
\end{equation}
其中 $\ket{\psi_i}$ 为目标态，且 $0 < \alpha_i \leq 1$，$\sum_i \alpha_i =1$，以便为所有感兴趣的态提供一个共同的基集合。当然，在给定数值资源的条件下，这只能以一定的精度损失为代价、且只对少数几个态可行。在 MPO/MPS 表述中，若愿意为每个态重新计算收缩，则态的基只在精确完备基的层面才被捆绑在一起。因此，一旦涉及多个态，MPO/MPS 表述相对于传统 DMRG 语言便发挥出其全部潜力。

(ii) MPO/MPS 表述更优越的另一个情形，尽管数值代价更高，是计算 $\bra{\psi} \hat{H}^2 \ket{\psi}$，例如在估计基态被求得的精度如何时它很有用。在 MPO 形式体系中，通过收缩图~\ref{fig:calculationH2} 所示的网络，可以把它精确算出（当然要对 $\ket{\psi}$ 作内在的近似）。对程序员而言，最经济的做法当然是计算 $\hat{H}\ket{\psi}$ 再取范数，这两个操作在现阶段都是现成的。这样做的操作代价是：MPO 的作用需 $O(LD^2D_W^2 d^2)$，范数计算需 $O(LD^3 D_W^3 d)$。后者非常昂贵，因此更高效的做法是像计算 $\bra{\psi} \hat{H} \ket{\psi}$ 时那样采用迭代式的构造。我要作一个重要说明：维数 $D_W^2$ 只是 $\hat{H}^2$ 的最坏情形\cite{Frowis10}：把平方展开并为表达式引入规则，可以得到更高效的 MPO，其最优性可以通过数值方法检验——做一次 SVD 压缩，查看奇异值中哪些为零即可。对我们的各向异性海森堡哈密顿量，$\hat{H}^2 $ 只需 $D_W=9$ 而不是 25。对更高次幂，收益更为可观，且可以通过数值方法获得：构造 $\hat{H}^n$ 的显式 MPO 后作压缩，只丢弃奇异值中的那些零。

\begin{figure}
\centering\includegraphics[scale=0.7]{PyMPS_MPOCalculationHsquare.eps}
\caption{$\hat{H}^2$ 期望值的``精确''计算：哈密顿量 MPO 被重复两次，底部夹 $\ket{\psi}$，顶部夹 $\bra{\psi}$。}
\label{fig:calculationH2}
\end{figure}
 
在 DMRG 计算中，得到的将是序列 $\hat{H}(\hat{H} \ket{\psi})$，位于 DMRG 块-格点基中，如图~\ref{fig:DMRGcalculationH2} 所示。关键在于，在第二次施加 $\hat{H}$ 之前，会发生一次到约化块基上的投影，它不是恒等变换，从而丢失信息。

\begin{figure}
\centering\includegraphics[scale=0.7]{PyMPS_DMRGCalculationHsquare.eps}
\caption{$\hat{H}^2$ 期望值的 DMRG 计算：在步骤 (1) 中，哈密顿量 MPO 在 DMRG 基中被施加于 $\ket{\psi}$ 一次，即结果被投影到约化基上，得到某个 $\Phi^{\, \sigma_{\ell}}$。它随后在步骤 (2) 中第二次施加 $\hat{H}$ 时取代 $\Psi^{\, \sigma_{\ell}}$。最终，结果再被投影到块基上。}
\label{fig:DMRGcalculationH2}
\end{figure}

MPS 与 DMRG 的比较在算法上意味着什么？首先，传统 DMRG 的截断误差——它已被证明是衡量结果质量的高度可靠的工具——不过是有些反常的双格点设置造成的人为产物。在变分 MPS 或单格点 DMRG 中，它必须由别的判据取代，例如能量的方差。其次，虽然所有方法在``寻求给定类型拟设所能达到的最低能量''这一意义上都是变分的，但双格点 DMRG 的拟设逐格点而变（因为拟设中的 $\Psi^{\sigma\sigma}$ 反常），而单格点 DMRG 的拟设自始至终保持不变，在概念上更为优雅。当然，这以可能陷入局部极小为代价，这也提醒我们：数学上最美的未必是最实用的。
 
\section{用 MPS 处理时间演化（实时的与虚时的）}
计算 $\eul^{-\imag \hat{H}t}$ 或 $\eul^{-\beta\hat{H}}$ 这类算符对量子态的作用，无论对于量子态的实时演化还是对于量子统计力学，都是量子力学中的核心问题；$\beta$ 可以解释为虚时。MPS 最吸引人的特性之一，就是这类实时或虚时演化可以被非常整洁而高效地编码。这一点对纯态和混合态（在有限温度下十分重要）都成立。在下文中，我将聚焦于基于演化算符 Trotter 分解的时间演化\cite{Vidal03a,Vidal04b,Daley04,WhiteFeiguin04,VerstraeteRipoll04}：先解释 Trotter 分解以及纯态情形的算法结构，再讲 Trotter 分解的 MPO 表示。在此之后，我将讨论模拟混合态动力学所需的改动。
\subsection{传统时间演化：纯态}

\subsubsection{时间演化的 Trotter 分解}
设 $\hat{H}$ 只包含最近邻相互作用，即 $\hat{H} = \sum_i \hat{h}_i$，其中 $\hat{h}_i$ 包含格点 $i$ 与 $i+1$ 之间的相互作用。于是我们可以把时间离散化为 $t = N \tau$（$\tau\rightarrow 0$，$N \rightarrow \infty$），并且（在最朴素的方案中）做一阶 Trotter 分解：
\begin{equation}
\eul^{-\imag \hat{H}\tau} = \eul^{-\imag  \hat{h}_1 \tau}\eul^{-\imag  \hat{h}_2 \tau}\eul^{-\imag  \hat{h}_3 \tau} \ldots  \eul^{-\imag  \hat{h}_{L-3} \tau}\eul^{-\imag  \hat{h}_{L-2} \tau}\eul^{-\imag  \hat{h}_{L-1} \tau} + O(\tau^2),
\end{equation}
它包含一个因键哈密顿量互不对易而产生的误差，一般有 $[ \hat{h}_i,\hat{h}_{i+1}] \neq 0$；更高阶的分解将在本节后面讨论。所有作用在奇数键上的时间演化（$\eul^{-\imag \hat{H}_{{\rm odd}} \tau}$）彼此对易，偶数键上的（$\eul^{-\imag \hat{H}_{{\rm even}} \tau}$）也彼此对易，因此可以同时进行。于是我们要寻找一个在奇数键上执行无穷小时间步的 MPO，以及另一个在偶数键上做同样事情的 MPO。

既然任何算符都保证可以表示为 MPO，让我们暂且假定确实能够高效地构造这些无穷小时间步的表示（显式构造见下一节）。我们将会看到，无穷小时间步 MPO 的最大键维数是 $d^2$，因为 $\eul^{-\imag \hat{h} \tau}$ 的维数是 $(d^2 \times d^2)$。因此，施加无穷小时间步 MPO 会使键维数从 $D$ 增大到 $d^2 D$。反复施加无穷小时间演化 MPO 会导致矩阵维数指数式增长，所以每个时间步之后都必须作截断。

由此得到的时间演化算法形式非常简单：从 $\ket{\psi(t=0)}$ 出发，重复以下步骤：
\begin{itemize}
\item 将奇数键的 MPO 作用于 $\ket{\psi(t)}$。
\item 将偶数键的 MPO 作用于 $\eul^{-\imag \hat{H}_{{\rm odd}}\tau}\ket{\psi(t)}$。
\item 将 MPS $\ket{\psi(t+\tau)}=\eul^{-\imag \hat{H}_{{\rm even}}\tau}\eul^{-\imag \hat{H}_{{\rm odd}}\tau}\ket{\psi(t)}$ 从维数 $d^2 D$ 压缩到 $D$，并监控误差。显然，也可以允许某个压缩误差（态间距离）$\epsilon$ 并选取随时间变化的 $D$：它通常会随时间强烈增长，从而限制可达的时间尺度。仿照基态计算的做法，所有结果都应在 $D\rightarrow\infty$ 或 $\epsilon\rightarrow 0$ 下外推。
\end{itemize}

在每个时间步之后，我们可以按标准方式计算可观测量，$\langle O(t) \rangle = \bra{\psi(t)} \hat{O} \ket{\psi(t)}$。但我们还能做得更多：我们可以如下计算含时关联子
\begin{equation}
\langle \hat{O}(t) \hat{P} \rangle = \bra{\psi} \eul^{+i\hat{H}t} \hat{O} \eul^{-\imag \hat{H}t} \hat{P} \ket{\psi}
= \bra{\psi(t)} \hat{O} \ket{\phi(t)}
\end{equation}
其中 $\ket{\psi(t)} = \eul^{-\imag \hat{H}t}\ket{\psi}$ 以及 $\ket{\phi(t)} = \eul^{-\imag \hat{H}t} \hat{P} \ket{\psi}$。例如取 $\hat{O} = \hat{S}^z_i$ 与 $\hat{P} = \hat{S}^z_j$，我们便可以计算 
$\langle \hat{S}^z_i(t) \hat{S}^z_j \rangle$，并经双重傅里叶变换得到结构函数
\begin{equation}
S^{zz}(k,\omega) \propto \int dt \sum_{n} \langle \hat{S}^z_i(t) \hat{S}^z_{i+n} \rangle \eul^{\imag kn}\eul^{-\imag \omega t},
\end{equation}
为使公式简洁，这里我假设了晶格的平移不变性与无限延展。

对上述算法的一个简单改进是采用二阶 Trotter 分解
\begin{equation}
\eul^{-\imag  \hat{H}\tau} = \eul^{-\imag  \hat{H}_{{\rm odd}} \tau/2}\eul^{-\imag  \hat{H}_{{\rm even}} \tau}\eul^{-\imag  \hat{H}_{{\rm odd}} \tau/2} + O(\tau^3),
\end{equation}
这样每个时间步的误差又降低了一个 $\tau$ 的量级。如果我们在每个时间步之后不做计算，就可以把半步归并在一起，在不比一阶 Trotter 分解多花任何代价的情况下进行工作。

一种源于量子蒙特卡罗、非常流行的四阶 Trotter 分解实现由 Suzuki\cite{Suzuki76,Suzuki91} 给出的如下公式：
\begin{equation}
\eul^{-\imag  \hat{H}\tau} = \hat{U}(\tau_1) \hat{U}(\tau_2) \hat{U}(\tau_3) \hat{U}(\tau_2) \hat{U}(\tau_1),  
\end{equation}
其中
\begin{equation}
\hat{U} (\tau_i) = \eul^{-\imag  \hat{H}_{{\rm odd}} \tau_i/2}\eul^{-\imag  \hat{H}_{{\rm even}} \tau_i}\eul^{-\imag  \hat{H}_{{\rm odd}}
\tau_i/2}
\end{equation}
以及
\begin{equation}
\tau_1 = \tau_2 = \frac{1}{4-4^{1/3}} \tau  \quad\quad \tau_3 = \tau -2 \tau_1 - 2\tau_2 .
\end{equation}
采用对称性较低的公式\cite{McLachlan95}，可以在相近的代价下得到更小的误差。至此算法的阐述即告完成（误差分析将在讲完其他时间演化方法之后给出），现在我们需要构造各个 MPO。

\subsubsection{纯态时间演化的 MPO}
让我们考虑链上所有奇数键的 Trotter 步：
\begin{equation}
\eul^{-\imag  \hat{h}_1 \tau} \otimes \eul^{-\imag  \hat{h}_3 \tau} \otimes \ldots \otimes \eul^{-\imag  \hat{h}_{L-1} \tau} \ket{\psi} ;
\label{eq:oddevolution}
\end{equation}
每个像 $\eul^{-\imag  \hat{h}_1 \tau}$ 这样的键演化算符都具有形式 $\sum_{\sigma_1 \sigma_2, \sigma'_1 \sigma'_2} O^{\sigma_1\sigma_2,\sigma'_1\sigma'_2} \ket{\sigma_1\sigma_2}\bra{\sigma'_1\sigma'_2}$。无论在图示还是显式的数学表示中，都可以明显看出这个算符破坏了 MPS 形式（图~\ref{fig:trotterstep}）。

\begin{figure}
\centering\includegraphics[width=250pt]{PyMPS_TrotterStep.eps}
\caption{一个 Trotter 步：在所有奇数键上进行（无穷小的）键时间演化。这使两个格点合并，乍看之下 MPS 的简单乘积形式被破坏，但时间演化可以转写为 MPO。由于时间演化是因式化的，各 MPO 在所有偶数键上的维数为 1（细线）。}
\label{fig:trotterstep}
\end{figure}

因此，若能把 $O^{\sigma_1\sigma_2,\sigma'_1\sigma'_2}$ 写成包含张量积 $O^{\sigma_1, \sigma'_1} \otimes O^{\sigma_2, \sigma'_2}$ 的某种形式，以保持 MPS 形式，将是理想的。为此，我们把任意态分解为 MPS 的流程加以改造，用于算符（每个格点有两个指标）。这一做法之所以可行，是因为指标数目极少。将 $O$ 重排以归组局域指标，然后进行奇异值分解：
\begin{eqnarray*}
O^{\sigma_1 \sigma_2, \sigma'_1 \sigma'_2} &=& P_{(\sigma_1 \sigma'_1), (\sigma_2 \sigma'_2)} \\
&=& \sum_k U_{\sigma_1 \sigma'_1,k} S_{k,k} (V^\dagger)_{k, (\sigma_2 \sigma'_2)} \\
&=& \sum_k U^{\sigma_1 \sigma'_1}_{k} \overline{U}^{\sigma_2 \sigma'_2}_{k} \\
&=& \sum_k U^{\sigma_1 \sigma'_1}_{1,k} \overline{U}^{\sigma_2 \sigma'_2}_{k,1}
\end{eqnarray*}
其中 $U^{\sigma_1 \sigma'_1}_{k} =  U_{(\sigma_1 \sigma'_1),k}\sqrt{S_{k,k}}$ 以及
$\overline{U}^{\sigma_2 \sigma'_2}_{k} = \sqrt{S_{k,k}} (V^\dagger)_{k, (\sigma_2 \sigma'_2)}$。在推导的最后一步中，我们引入了一个取值为 1 的哑指标，以得到 MPO 矩阵的形式。指标 $k$ 最大可取到 $d^2$，这就给出 MPO 的键维数 $D_W$。


表示式~(\ref{eq:oddevolution}) 中算符的 MPO，即 $U^{\sigma_1\sigma'_1} \overline{U}^{\sigma_2\sigma'_2}U^{\sigma_3\sigma'_3} \overline{U}^{\sigma_4\sigma'_4}\ldots$，与原先的各幺正算符一样每隔一根键便因式化。若某个格点不参与任何键演化，我们只需给它分配作为 $(1\times 1)$-矩阵的恒等幺正算符：$I^{\sigma,\sigma'}_{1,1}=\delta_{\sigma,\sigma'}$。于是全局 MPO 可以由局域 MPO 平凡地构造出来。在所有奇数键上做时间演化的 MPO 形如
$U\overline{U}U\overline{U}U\overline{U} \ldots$，而偶数键时间步则为 $IU\overline{U}U\overline{U}U\overline{U}\ldots I$（图~\ref{fig:oddevenTrotter}）。

\begin{figure}
\centering\includegraphics[width=250pt]{PyMPS_OddEvenTrotter.eps}
\caption{完整的一阶 Trotter 时间步（奇数键与偶数键）。粗线与细线分别对应 MPO 键上的维数 1 与 $>1$。顶行中第一和最后格点上的 MPO 是平凡的标量恒等式 1。}
\label{fig:oddevenTrotter}
\end{figure}

\subsection{常规时间演化：混合态}

\subsubsection{混合态的纯化}
有限温度计算可以基于对任意混合量子态的纯化\cite{VerstraeteRipoll04}来进行：若我们考虑物理空间 P 中由正交归一态构成的混合态，则可以把它理解为对 PQ 上的一个纯态的 Schmidt 分解作偏迹的结果，其中 Q 是辅助空间：
\begin{equation}
\hat{\rho}_P = \sum_{a=1}^r s_a^2 \ket{a}_P \bra{a}_P \rightarrow \ket{\psi} = \sum_{a=1}^r s_a \ket{a}_P \ket{a}_Q \quad \quad \hat{\rho}_P = \tr_Q \ket{\psi}\bra{\psi} .
\end{equation}
辅助态空间可以简单地取为原空间的一个副本，于是链上的有限温度密度算符可以表示为梯子上的纯态（见图~\ref{fig:finitetemperaturesketch}）。

为了计算热密度算符 $\hat{\rho}_\beta= Z(\beta)^{-1} \eul^{-\beta \hat{H}}$，$Z(\beta) = \tr_P \eul^{-\beta \hat{H}}$，我们写
\begin{equation}
\hat{\rho}_\beta = Z(\beta)^{-1} \eul^{-\beta \hat{H}} = Z(\beta)^{-1} \eul^{-\beta \hat{H}/2} \cdot \hat{I} \cdot \eul^{-\beta \hat{H}/2} .
\end{equation}
恒等算符 $\hat{I}$ 不过就是 $Z(0) \hat{\rho}_0$，即无穷大温度下的密度算符乘以无穷大温度的配分函数。假设我们已知 $\hat{\rho}_0$ 作为 MPS 的纯化，即 $\ket{\psi_{\beta=0}}$。那么
\begin{equation}
\hat{\rho}_\beta = (Z(0)/Z(\beta)) \eul^{-\beta \hat{H}/2} \cdot \tr_Q \ket{\psi_0}\bra{\psi_0} \cdot \eul^{-\beta \hat{H}/2} = (Z(0)/Z(\beta)) \tr_Q \eul^{-\beta \hat{H}/2}\ket{\psi_0} \bra{\psi_0}\eul^{-\beta \hat{H}/2} .
\end{equation}
由于哈密顿量不作用于 Q，对 Q 的迹可以提出来。但这一结果意味着我们必须做一次虚时演化
\begin{equation}
\ket{\psi_\beta} = \eul^{-\beta \hat{H}/2}\ket{\psi_0} .
\end{equation}
期望值由
\begin{equation}
\langle \hat{O} \rangle_\beta = \tr_P \hat{O}\hat{\rho}_\beta =
(Z(0)/Z(\beta)) \tr_P \hat{O} \tr_Q \ket{\psi_\beta} \bra{\psi_\beta}
= (Z(0)/Z(\beta)) \bra{\psi_\beta} \hat{O} \ket{\psi_\beta} .
\end{equation} 
$(Z(0)/Z(\beta)) = (d^L/Z(\beta))$ 看似难以求得，但它可以从恒等算符的期望值轻易得到，
\begin{equation}
1 = \langle \hat{I} \rangle_\beta = \tr_P \hat{\rho}_\beta =
(Z(0)/Z(\beta)) \tr_P \tr_Q \ket{\psi_\beta} \bra{\psi_\beta}
= (Z(0)/Z(\beta)) \braket{\psi_\beta}{\psi_\beta} ,
\end{equation} 
因此 $Z(\beta)/Z(0)= \braket{\psi_\beta}{\psi_\beta}$，或者，与标准量子力学完全一致地，
\begin{equation}
\langle \hat{O} \rangle_\beta = \frac{ \bra{\psi_\beta} \hat{O}\ket{\psi_\beta}}{\braket{\psi_\beta}{\psi_\beta}} .
\end{equation}
但这恰好把我们带回到已知道如何计算的表达式。我们要做的只是找出 $\ket{\psi_0}$，进行虚时演化直至 $-\imag\beta/2$，然后像对纯态那样计算期望值。我们甚至可以让纯化后的态 $\ket{\psi_\beta}$ 继续经受实时演化，以处理有限 $T$ 下的时间依赖性。

我们也可以很简单地处理热力学，因为 $Z(\beta)/Z(0)$ 由 $\ket{\psi_\beta}$ 的范数的平方给出，且 $Z(0)=d^L$。这意味着，只要在所有温度下保持纯化态归一化，并随着温度下降、$\beta$ 增大而累积归一化因子，就可以得到 $Z(\beta)$。由 $Z(\beta)$ 可得 $F(\beta)=-\beta^{-1} \ln Z(\beta)$。同时，$U(\beta) = \langle \hat{H} \rangle_\beta = \bra{\psi_\beta} \hat{H} \ket{\psi_\beta}$。而这又给出 $S(\beta) = \beta(U(\beta)-F(\beta))$。其余热力学量可类似地得到。

\begin{figure}
\centering\includegraphics[scale=0.7]{PyMPS_FiniteTemperatureSketch.eps}
\caption{有限温度模拟的示意表示：不再用一条链，而是搭建一个梯子，其中包含该链的一个完全相同的副本。物理格点标号为奇数，辅助格点标号为偶数。物理腿与辅助腿上的对应格点由最大纠缠态相连。为达到逆温度 $\beta$，进行虚时演化直至``时间'' $-\imag\beta/2$。}
\label{fig:finitetemperaturesketch}
\end{figure}

无限温度混合态的纯化是一个维数为 1 的简单 MPS，因为它会因子化（若我们把梯子的一个横档取作一个大格点）：
\begin{equation}
\hat{\rho}_0 = \frac{1}{d^L} \hat{I} = \left(\frac{1}{d} \hat{I}\right)^{\otimes L} .
\end{equation}
由于 $\hat{\rho}_0$ 是因子化的，我们现在可以把每个物理格点上的局域混合态纯化为横档 $i$ 上的一个纯态，得到 $\ket{\psi_{i0}}$，于是 
$\ket{\psi_0} = \ket{\psi_{1,0}}\ket{\psi_{2,0}}\ket{\psi_{3,0}} \ldots$，是一个乘积态，即维数为 1 的 MPS。若考虑梯子的某个横档 $i$，其物理格点 $2i-1$ 与辅助格点 $2i$ 上分别有态 $\ket{\sigma}_P$ 与 $\ket{\sigma}_Q$，则可以如下纯化：
\begin{equation}
\frac{1}{d} \hat{I} = \sum_{\sigma} \frac{1}{d} \phantom\langle_P\ket{\sigma}\bra{\sigma}_P = \tr_Q \left[ \left( \sum_{\sigma} \frac{1}{\sqrt{d}} \ket{\sigma}_P \ket{\sigma}_Q \right) \left( \sum_{\sigma} \frac{1}{\sqrt{d}} \bra{\sigma}_P \bra{\sigma}_Q  \right) \right] .
\end{equation}
因此纯化由一个最大纠缠态给出（纠缠熵为 $\log_2 d$），
\begin{equation}
\ket{\psi_{i0}} = \sum_{\sigma} \frac{1}{\sqrt{d}} \ket{\sigma}_P \ket{\sigma}_Q .
\end{equation}
容易看出，可以对 P 和 Q 各自分别作局域幺正变换而保持该结构不变。例如，对自旋-1/2 链的纯化而言，用单态作为局域纯化是有利的，
\begin{equation}
\ket{\psi_{i,0}} = \frac{1}{\sqrt{2}} [ \ket{\uparrow_P\downarrow_Q} - \ket{\downarrow_P\uparrow_Q} ],
\end{equation}
以防程序懂得利用好量子数：这种态可以同时保持总 $S=0$ 与 $S^z=0$。此时，四个 $A$ 矩阵为
\begin{equation}
A^{\uparrow_P\uparrow_Q} =0 \quad 
A^{\uparrow_P\downarrow_Q} =1/\sqrt{2} \quad 
A^{\downarrow_P\uparrow_Q} =-1/\sqrt{2} \quad 
A^{\downarrow_P\downarrow_Q} =0  ,
\end{equation}
于是 $\beta=0$ 的纯化初态 $\ket{\psi_0}$ 现在由梯子上单态键的乘积给出。实际上，对重塑后的矩阵 $A_{\sigma,\sigma'}$ 作 SVD 使我们可以引入真正局域于格点的 $A$ 矩阵，它们在奇数格点上维数为 $(1\times 2)$，在偶数格点上为 $(2\times 1)$：
\begin{equation}
A^{\uparrow_{2i-1}} = [ 1\ \ 0] \quad
A^{\downarrow_{2i-1}} = [ 0\ \ -1] \quad
A^{\uparrow_{2i}} = [ 0\ \ 1/\sqrt{2}]^T \quad
A^{\downarrow_{2i}} = [ 1/\sqrt{2} \ \ 0]^T \quad .
\end{equation}
为了应用纯态时间演化算法，剩下的就是找出 MPO。
 
\subsubsection{混合态演化的 MPO}

\begin{figure}
\centering\includegraphics[width=250pt]{PyMPS_LadderSetup.eps}
\caption{链上混合态的时间演化可以看作梯子上纯态的时间演化，其中物理格点位于第一条腿上，而额外的辅助格点位于第二条腿上。这个梯子被映射为一条链。由于时间演化只作用于物理态，就产生了次近邻相互作用。}
\label{fig:laddersetup}
\end{figure}

混合态模拟中出现的梯子可以映射为一条链（图~\ref{fig:laddersetup}），其中物理哈密顿量只作用于奇数格点 $1,3,5,\ldots$，而辅助格点是偶数的 $2,4,6,\ldots$。于是唯一非平庸的时间演化连接 $(1,3), (3,5), (5,7)$。处理这类更长程的相互作用有几种办法：一种是显式地构造这种长程相互作用，另一种是使用所谓的 swap 门，把它约化为最近邻相互作用。

对``更长程''相互作用 $(1,3)$ 的直接 MPO 描述必然要在格点 $2$ 上放一个非平庸的张量，而格点 $4$ 是惰性的。$(3,5)$ 同样如此：$4$ 上有一个非平庸张量，而格点 $6$ 是惰性的。这提示了一种 Trotter 分解：``奇''步骤取 $(1,2,3,4), (5,6,7,8),\ldots$，``偶''步骤取 $(3,4,5,6),(7,8,9,10),\ldots$。

于是格点 1 至 4 上的四格点演化算符为
\begin{equation}
O^{(\sigma_1\sigma_2\sigma_3\sigma_4), (\sigma'_1\sigma'_2\sigma'_3\sigma'_4)} =
O^{(\sigma_1\sigma_2\sigma_3), (\sigma'_1\sigma'_2\sigma'_3)}\cdot \delta_{\sigma_4,\sigma'_4}
\end{equation}
而且只需对包含实际物理时间演化的两格点幺正算符稍作修改，就能构造出三格点幺正算符：
\begin{equation}
O^{(\sigma_1\sigma_2\sigma_3), (\sigma'_1\sigma'_2\sigma'_3)} = O^{(\sigma_1\sigma_3), (\sigma'_1\sigma'_3)} \cdot \delta_{\sigma_2,\sigma'_2} .
\end{equation}
接着对这个三格点幺正算符作两次 SVD。为了记号方便，我们先把指标下移并逐格点重排。然后通过重塑及随后的 SVD 逐步迭代地把 $\sigma_3,\sigma'_3$、$\sigma_2,\sigma'_2$ 和 $\sigma_1,\sigma'_1$ 分离出来： 
\begin{eqnarray*}
& & O_{(\sigma_1\sigma_2\sigma_3), (\sigma'_1\sigma'_2\sigma'_3)} \\
&=& P_{(\sigma_1 \sigma'_1\sigma_2\sigma'_2),(\sigma_3\sigma'_3)} \\
&=& \sum_{k_2} U_{(\sigma_1 \sigma'_1\sigma_2\sigma'_2),k_2} S^{[2]}_{k_2,k_2} (V^\dagger_{23})_{k_2,(\sigma_3\sigma'_3)} \\
&=&  \sum_{k_2} U_{(\sigma_1 \sigma'_1),(\sigma_2\sigma'_2 k_2)} S^{[2]}_{k_2,k_2} (V^\dagger_{23})_{k_2,(\sigma_3\sigma'_3)} \\
&=& \sum_{k_1,k_2} U_{(\sigma_1 \sigma'_1),k_1} S^{[1]}_{k_1,k_1} (V^\dagger_{12})_{k_1, (\sigma_2\sigma'_2 k_2)} S^{[2]}_{k_2,k_2} (V^\dagger_{23})_{k_2,(\sigma_3\sigma'_3)} \\
&=& \sum_{k_1,k_2} W^{\sigma_1 \sigma'_1}_{1,k_1} W^{\sigma_2 \sigma'_2}_{k_1,k_2} W^{\sigma_3 \sigma'_3}_{k_2,1} 
\end{eqnarray*}
其中，通过引入哑指标以及格点 4 上的惰性张量：
\begin{eqnarray}
W^{\sigma_1 \sigma'_1}_{1,k_1} &=& U_{(\sigma_1 \sigma'_1),k_1} \sqrt{S^{[1]}_{k_1,k_1}} \\
W^{\sigma_2 \sigma'_2}_{k_1,k_2} &=& \sqrt{S^{[1]}_{k_1,k_1}} (V^\dagger_{12})_{k_1, (\sigma_2\sigma'_2 k_2)} \sqrt{S^{[2]}_{k_2,k_2}} \\
W^{\sigma_3 \sigma'_3}_{k_2,1} &=& \sqrt{S^{[2]}_{k_2,k_2}} (V^\dagger_{23})_{k_2,(\sigma_3\sigma'_3)} \\
W^{\sigma_4 \sigma'_4}_{1,1} &=& \delta_{\sigma_4,\sigma'_4}
\end{eqnarray}
由此可以像纯态时间演化那样构造整条链的 MPO。我们所做的无非是对 4 个格点上的 MPO 作迭代分解。同样，这仍然是可处理的，因为只涉及 4 个格点。 

显然，直接写出作用于全部四条键 $(1,2)$、$(1,3)$、$(2,4)$ 和 $(3,4)$ 的演化算符并对其施行类似的一系列 SVD，是一步简单的推广，这样就可以考虑真实梯子上的纯态时间演化。 

实现超越直接邻居的相互作用，另一种途径是使用 {\em swap 门}\cite{Stoudenmire10}。以真实梯子为例，其四条键上有相互作用，其中两条 [$(1,3)$ 和 $(2,4)$] 是次近邻相互作用。但如果交换格点 $2 \leftrightarrow 3$ 上的态，它们就变成最近邻相互作用。梯子上的时间演化于是可以如下进行：(i) 演化键 $(1,2)$ 和 $(3,4)$；(ii) 交换格点 2 和 3 上的态；(iii) 用``旧''键 $(1,3)$ 和 $(2,4)$ 的演化算符演化``新''键 $(1,2)$ 和 $(3,4)$；(iv) 再次交换格点 2 和 3 上的态。在其他情形下，需要找到更巧妙的方案，最好能生成一个只在最近邻之间进行交换的序列。格点 $i$ 与 $j$ 之间的 swap 算符简单地由下式给出
\begin{equation}
\hat{S}_{ij} = \sum_{\sigma_i\sigma'_i \sigma_j\sigma'_j} S^{\sigma_i\sigma_j \sigma'_i \sigma'_j} \ket{\sigma_i\sigma_j}\bra{\sigma'_i\sigma'_j} \quad\quad 
S^{\sigma_i\sigma_j \sigma'_i \sigma'_j} = \delta_{\sigma_i,\sigma'_j} \delta_{\sigma_j,\sigma'_i},
\end{equation}
它是幺正的且为其自身的逆。它交换 MPS 中两个格点的物理指标；对于最近邻之间的交换，很容易恢复 MPS 的原有形式：设 MPS 是左规范的；作用在格点 $i$ 与 $i+1$ 上的幺正算符只影响这两个格点。特别地，块态 $\ket{a_{k<i}}_A$ 与 $\ket{a_{i+1}}_A$ 的正交归一性不受影响。
如果我们引入矩阵 $M_{(a_{i-1}\sigma_{i+1}),(\sigma_i a_{i+1})} = \sum_{a_i} A^{[i]\sigma_{i+1}}_{a_{i-1},a_i} A^{[i+1]\sigma_i}_{a_i,a_{i+1}}$，就可以构成 $\tilde{M}_{(a_{i-1}\sigma_i),(\sigma_{i+1}a_{i+1})}=M_{(a_{i-1}\sigma_{i+1}),(\sigma_i a_{i+1})}$ 并作一次 SVD，其中 $U_{(a_{i-1}\sigma_i),a_i}$ 给出新的左规范 $\tilde{A}^{\sigma_i}_{a_{i-1},a_i}$，而
$S_{a_i,a_i} (V^\dagger)_{a_i,(\sigma_{i+1}a_{i+1})}$ 给出新的左规范 $\tilde{A}^{\sigma_{i+1}}_{a_{i},a_{i+1}}$。后者是左规范的，这一点由 $\tilde{A}^{\sigma_i}$ 的左规范性以及保持不变的正交归一性 $\ket{a_{i+1}}_A$ 得出。

在结束这篇关于超越最近邻相互作用的展望之前，我要指出使用 MPO 还允许其他 Trotter 分解，例如把海森堡哈密顿量分解为依赖于 $x$、$y$ 和 $z$ 的各个部分，这对长程相互作用很有用\cite{VerstraetePirvu08}。 
 

\subsection{tDMRG、TEBD 与 MPS 时间演化的比较：微小的差别}
在基于 MPS 的时间演化（tMPS）发展出来之前不久，另有两个算法被提出用于模拟一维量子链的实时动力学：时间演化块抽取（time-evolving block decimation, TEBD）\cite{Vidal03a,Vidal04b} 与实时或含时 DMRG（tDMRG）\cite{Daley04,WhiteFeiguin04}。这两个算法同样以 MPS 为基础，但细看之下与 tMPS 有所不同。在进入这些细节之前，我要强调：它们全都基于对 MPS 作时间演化的思想，该思想最早在 \cite{Vidal03a,Vidal04b} 中提出，因此它们都只是同一主题下的次要变体。tDMRG 与 TEBD 在数学上是等价的，即在精确算术下应给出相同的结果；但在数值上它们是明显不同的算法，各自都执行着对方方法中没有对应物的操作，各有优劣。我们先讨论 tDMRG，因为它的语言更接近 tMPS，然后再讨论 TEBD，看看这些微小差别有多重要（或者多么不重要）。

\subsubsection{含时 DMRG（tDMRG）}
把整个晶格上的全局时间演化分解为键上无穷小时间演化的 Trotter 序列，对这里讨论的三种算法都是一样的。
因此我们聚焦于格点 $\ell+1$ 与 $\ell+2$ 上的一个无穷小时间演化 $\eul^{-\imag  \hat{h}_{\ell+1} \tau}$。 
用局域基表示的演化算符为
\begin{equation}
U^{(\sigma_{\ell+1} \sigma_{\ell+2}), (\sigma'_{\ell+1} \sigma'_{\ell+2})} = \bra{\sigma_{\ell+1} \sigma_{\ell+2}} \eul^{-\imag  \hat{h}_{\ell+1} \tau} \ket{\sigma'_{\ell+1} \sigma'_{\ell+2}} .
\end{equation}
当前态 $\psi$ 用两格点 DMRG 记法、以左规范和右规范矩阵表示为
\begin{equation}
\ket{\psi} = \sum_{\fat{\sigma}} A^{\sigma_1} \ldots A^{\sigma_\ell} \Psi^{\sigma_{\ell+1}\sigma_{\ell+2}} B^{\sigma_{\ell+3}} \ldots B^{\sigma_L} \ket{\fat{\sigma}} .
\end{equation}
时间演化把 $\Psi^{\sigma_{\ell+1}\sigma_{\ell+2}}$ 变为
\begin{equation}
\Phi^{\sigma_{\ell+1}\sigma_{\ell+2}}_{a_\ell, a_{\ell+2}} = \sum_{\sigma'_{\ell+1} \sigma'_{\ell+2}}
U^{(\sigma_{\ell+1} \sigma_{\ell+2}), (\sigma'_{\ell+1} \sigma'_{\ell+2})}  \Psi^{\sigma'_{\ell+1}\sigma'_{\ell+2}}_{a_\ell, a_{\ell+2}} . 
\end{equation} 
它与 $A$、$B$ 矩阵一起定义了一个合法的 DMRG 态，我们称之为 $\ket{\phi}$。为了继续推进，即把 $\eul^{-\imag  \hat{h}_{\ell+3} \tau}$ 作用于下一对格点，我们必须把态化成如下形式 
\begin{equation}
\ket{\phi} = \sum_{\fat{\sigma}} A^{\sigma_1} \ldots A^{\sigma_{\ell+2}} \Phi^{\sigma_{\ell+3}\sigma_{\ell+4}} B^{\sigma_{\ell+5}} \ldots B^{\sigma_L} \ket{\fat{\sigma}} .
\end{equation}
改变只可能涉及格点 $\ell+1$ 至 $\ell+4$：前两个格点是因为演化算符的作用，后两个格点是因为要化成 DMRG 形式。我们先在格点 $\ell+1$ 和 $\ell+2$ 上生成新的 $A$ 矩阵：把 $\Phi^{\sigma_{\ell+1}\sigma_{\ell+2}}_{a_\ell, a_{\ell+2}} = \Phi_{(a_\ell \sigma_{\ell+1}), (\sigma_{\ell+2} a_{\ell+2})}$ 重塑后作 SVD（DMRG 传统上是通过密度矩阵分析以及移动格点时的 DMRG 预言来做的，结果相同）：
\begin{equation}
\Phi_{(a_\ell \sigma_{\ell+1}),(\sigma_{\ell+2} a_{\ell+2})} = \sum_{a_{\ell+1}} U_{(a_\ell \sigma_{\ell+1}),a_{\ell+1}} S_{a_{\ell+1},a_{\ell+1}} (V^\dagger)_{a_{\ell+1}, (\sigma_{\ell+2} a_{\ell+2})}.
\end{equation}
$U$ 可以立即重塑为一个合法的 $A$ 矩阵，但其列维数高达 $dD$，必须截断到 $D$，同时保持对维数 $dD$ 的 MPS 的最佳逼近。答案一如既往：只保留最大的 $D$ 个奇异值，并相应地收缩矩阵 $U$、$S$、$V^\dagger$。方法的近似就在这里（在明显的 Trotter 误差之外）。这一步做完后，我们重塑为
\begin{equation}
\sum_{a_{\ell+1}} A^{\sigma_{\ell+1}}_{a_\ell, a_{\ell+1}}  S_{a_{\ell+1},a_{\ell+1}} (V^\dagger)^{\sigma_{\ell+2}}_{a_{\ell+1},  a_{\ell+2}}
\end{equation}
并构成移动一个格点后的 $\Phi$ 为
\begin{equation}
\Phi^{\sigma_{\ell+2}\sigma_{\ell+3}}_{a_{\ell+1},a_{\ell+3}} = \sum_{a_{\ell+2}} S_{a_{\ell+1},a_{\ell+1}} (V^\dagger)^{\sigma_{\ell+2}}_{a_{\ell+1},  a_{\ell+2}} B^{\sigma_{\ell+3}}_{a_{\ell+2},  a_{\ell+3}} .
\end{equation}
但我们还须再移动一个格点，做法是把 $\Phi^{\sigma_{\ell+2}\sigma_{\ell+3}}_{a_{\ell+1},a_{\ell+3}}$ 重塑为 $\Phi_{(a_{\ell+1}\sigma_{\ell+2}),(\sigma_{\ell+3} a_{\ell+3})}$，像前面那样作一次 SVD，从 $dD$ 个奇异值中保留对应最大的 $D$ 个的态，重塑，记下 $A^{\sigma_{\ell+2}}$，并构成移动两个格点后的 $\Phi$ 为
\begin{equation}
\Phi^{\sigma_{\ell+3}\sigma_{\ell+4}}_{a_{\ell+2},a_{\ell+4}} = \sum_{a_{\ell+3}} S_{a_{\ell+2},a_{\ell+2}} (V^\dagger)^{\sigma_{\ell+3}}_{a_{\ell+2},  a_{\ell+3}} B^{\sigma_{\ell+4}}_{a_{\ell+3},  a_{\ell+4}} .
\end{equation}
第二次 SVD 及其把奇异值截断到 $D$ 并不损失更多信息，因为非零奇异值至多有 $D$ 个，尽管形式上本可以有 $dD$ 个。原因在于，在格点 $\ell+1$ 和 $\ell+2$ 上作时间演化之前，跨越键 $\ell+2$ 的 Schmidt 秩至多为 $D$（这是 MPS 构造的结果）。而若两个态由一个只作用在 A 部或 B 部上的幺正变换相联系，则它们的 Schmidt 秩相同。但无穷小时间演化正是作用在 A 部上的幺正变换。

现在我们可以按照 tMPS 的思路继续进行下一个无穷小局域时间演化步骤。   
\subsubsection{时间演化块消去（TEBD）}
这里，我们假设已将 $\ket{\psi}$ 写成 $\Gamma\Lambda$ 记法，
\begin{equation}
\ket{\psi} = \sum_{\fat{\sigma}} 
\Gamma^{\sigma_1} \Lambda^{[1]} \Gamma^{\sigma_2} \Lambda^{[2]} \ldots
\Gamma^{\sigma_{\ell}} \Lambda^{[\ell]}  \Gamma^{\sigma_{\ell+1}} \Lambda^{[\ell+1]}
\Gamma^{\sigma_{\ell+2}} \Lambda^{[\ell+2]}  \Gamma^{\sigma_{\ell+3}} \Lambda^{[\ell+3]} \ldots
\Gamma^{\sigma_L}   
\ket{\fat{\sigma}}.
\label{eq:TEBDstart}
\end{equation}
这个态可以立即与两格点 DMRG 记法联系起来。特别地，
\begin{equation}
\Psi^{\sigma_{\ell+1} \sigma_{\ell+2}} = \Lambda^{[\ell]}  \Gamma^{\sigma_{\ell+1}} \Lambda^{[\ell+1]}
\Gamma^{\sigma_{\ell+2}} \Lambda^{[\ell+2]}  .
\end{equation}
这与 DMRG 的 $\Psi$ 完全相同，演化算符 $U^{(\sigma_{\ell+1} \sigma_{\ell+2}), (\sigma'_{\ell+1} \sigma'_{\ell+2})}$ 也一样，因而  
\begin{equation}
\Phi^{\sigma_{\ell+1}\sigma_{\ell+2}}_{a_\ell, a_{\ell+2}} = \sum_{\sigma'_{\ell+1} \sigma'_{\ell+2}}
U^{(\sigma_{\ell+1} \sigma_{\ell+2}), (\sigma'_{\ell+1} \sigma'_{\ell+2})}  \Psi^{\sigma_{\ell+1}'\sigma_{\ell+2}'}_{a_\ell, a_{\ell+2}} . 
\end{equation} 
为了继续推进，必须在两个活动格点上恢复 $\Gamma\Lambda$ 记法。与 tDMRG 完全类似，通过 SVD 得到
\begin{equation}
\Phi_{(a_\ell \sigma_{\ell+1}), (\sigma_{\ell+2} a_{\ell+2})} = \sum_{a_{\ell+1}} U_{(a_\ell \sigma_{\ell+1}),a_{\ell+1}} \Lambda^{[\ell+1]}_{a_{\ell+1},a_{\ell+1}} (V^\dagger)_{a_{\ell+1}, (\sigma_{\ell+2} a_{\ell+2})} ,
\end{equation}
所缺的是 $\Lambda^{[\ell]}$ 和 $\Lambda^{[\ell+2]}$。于是我们写成（把 $U$ 和 $V^\dagger$ 重塑形状并省略 $a$ 指标）
\begin{equation}
\Phi^{\sigma_{\ell+1}\sigma_{\ell+2}} = \Lambda^{[\ell]} (\Lambda^{[\ell]})^{-1} U^{\sigma_{\ell+1}} 
 \Lambda^{[\ell+1]} V^{\sigma_{\ell+2}\dagger} ( \Lambda^{[\ell+2]})^{-1}  \Lambda^{[\ell+2]} .
\end{equation}
现在，与 tDMRG 中一样，$ \Lambda^{[\ell+1]}$ 中有多达 $dD$ 个奇异值，我们将其截断到最大的 $D$ 个，同样如 tDMRG 那样，也相应地截断相邻的矩阵。现在我们引入
\begin{equation}
\Gamma^{\sigma_{\ell+1}}_{a_\ell,a_{\ell+1}} =  (\Lambda^{[\ell]})^{-1}_{a_\ell,a_\ell} U^{\sigma_{\ell+1}}_{a_\ell,a_{\ell+1}} \quad 
\Gamma^{\sigma_{\ell+2}}_{a_{\ell+1},a_{\ell+2}} =  V^{\sigma_{\ell+2}\dagger}_{a_{\ell+1},a_{\ell+2}} ( \Lambda^{[\ell+2]})^{-1}_{a_{\ell+2},a_{\ell+2}} 
\end{equation}
并得到
\begin{equation}
\Phi^{\sigma_{\ell+1}\sigma_{\ell+2}} = \Lambda^{[\ell]}  \Gamma^{\sigma_{\ell+1}} 
 \Lambda^{[\ell+1]} \Gamma^{\sigma_{\ell+2}}  \Lambda^{[\ell+2]} ,
\end{equation}
回到了规范形式。为了考虑下一个键上的时间演化，我们无需做任何 SVD，只需把
\begin{equation}
\Psi^{\sigma_{\ell+3} \sigma_{\ell+4}} = \Lambda^{[\ell+2]}  \Gamma^{\sigma_{\ell+3}} \Lambda^{[\ell+3]}
\Gamma^{\sigma_{\ell+4}} \Lambda^{[\ell+4]}  
\end{equation}
组合起来然后继续。与 tDMRG 中一样，这一步不损失任何信息，只是在这里更为显式。

在高精度计算中当 $D$ 变得非常大时，奇异值往往会非常小，而如前所述，除以它们是数值不稳定性的一处来源。在热力学极限 iTEBD 方法（我们稍后讨论）的背景下，Hastings 提出了一种优雅的变通办法，其数值代价极低\cite{Hastings09}，但它也容易推广到有限系统 TEBD。假设我们从一个以表示式 (\ref{eq:TEBDstart}) 给出的态出发。然后我们把所有配对组合成右规范的 $B$ 矩阵，
\begin{equation}
B^{\sigma_i} = \Gamma^{\sigma_i} \Lambda^{[i]} ,
\end{equation}
但记住各个 $\Lambda^{[i]}$ 以备后用。然后我们构造
\begin{equation}
\overline{\Psi}^{\sigma_{\ell+1} \sigma_{\ell+2}} = B^{\sigma_{\ell+1}} B^{\sigma_{\ell+2}} =  \Gamma^{\sigma_{\ell+1}} \Lambda^{[\ell+1]}
\Gamma^{\sigma_{\ell+2}} \Lambda^{[\ell+2]} ,
\end{equation}
于是 $\Psi^{\sigma_{\ell+1} \sigma_{\ell+2}} = \Lambda^{[\ell]} \overline{\Psi}^{\sigma_{\ell+1} \sigma_{\ell+2}}$。我们对 $\overline{\Psi}^{\sigma_{\ell+1} \sigma_{\ell+2}}$ 施行时间演化，得到
\begin{equation}
\overline{\Phi}^{\sigma_{\ell+1}\sigma_{\ell+2}}_{a_\ell, a_{\ell+2}} = \sum_{\sigma'_{\ell+1} \sigma'_{\ell+2}}
U^{(\sigma_{\ell+1} \sigma_{\ell+2}), (\sigma'_{\ell+1} \sigma'_{\ell+2})}  \overline{\Psi}^{\sigma_{\ell+1}'\sigma_{\ell+2}'}_{a_\ell, a_{\ell+2}} .
\end{equation} 
然后 $\Phi^{\sigma_{\ell+1} \sigma_{\ell+2}} = \Lambda^{[\ell]} \overline{\Phi}^{\sigma_{\ell+1} \sigma_{\ell+2}}$。如前所述，我们对 $\Phi^{\sigma_{\ell+1} \sigma_{\ell+2}}$ 做 SVD，得到 
\begin{equation}
\Phi_{(a_\ell \sigma_{\ell+1}), (\sigma_{\ell+2} a_{\ell+2})} = \sum_{a_{\ell+1}} U_{(a_\ell \sigma_{\ell+1}),a_{\ell+1}} \Lambda^{[\ell+1]}_{a_{\ell+1},a_{\ell+1}} (V^\dagger)_{a_{\ell+1}, (\sigma_{\ell+2} a_{\ell+2})} = A^{\sigma_{\ell+1}} \Lambda^{[\ell+1]} B^{\sigma_{\ell+2}} .
\end{equation}
截断到最大的 $D$ 个奇异值后，我们就得到了新的 $\Lambda^{[\ell+1]}$（保留下来供后续使用）以及新的 $B^{\sigma_{\ell+2}}$。新的 $B^{\sigma_{\ell+1}}$
由下式给出
\begin{equation}
B^{\sigma_{\ell+1}} = \sum_{\sigma_{\ell+2}}  \overline{\Phi}^{\sigma_{\ell+1}\sigma_{\ell+2}}
B^{\sigma_{\ell+2}\dagger},
\end{equation}
因而只需一个简单的矩阵乘法；避免了除法。对于最后一个等式，我们利用 $B^{\sigma_{\ell+2}}$ 的右规范性，即 $\sum_{\sigma_{\ell+2}} \Phi^{\sigma_{\ell+1}\sigma_{\ell+2}} B^{\sigma_{\ell+2}\dagger} =  A^{\sigma_{\ell+1}} \Lambda^{[\ell+1]}$。同时，$B^{\sigma_{\ell+1}} = \Gamma^{\sigma_{\ell+1}} \Lambda^{[\ell+1]} = (\Lambda^{[\ell]})^{-1} A^{\sigma_{\ell+1}}  \Lambda^{[\ell+1]}$。把这两个恒等式与 $\Phi^{\sigma_{\ell+1} \sigma_{\ell+2}} = \Lambda^{[\ell]} \overline{\Phi}^{\sigma_{\ell+1} \sigma_{\ell+2}}$ 结合起来即得结果。

\subsubsection{算法比较}
将 TEBD 与 tDMRG 逐步比较，立刻就能看出两种方法在数学上完全等价。tDMRG 中的第二次 SVD 只不过是移动左规范矩阵与右规范矩阵之间的分界，而在 TEBD 中这只需对 $\Gamma$ 与 $\Lambda$ 重新加括号即可简单实现。尽管如此，两者仍有差别：每演化一个键，tDMRG 要做两次代价高昂的 SVD 分解（或等价的密度矩阵分析），而 TEBD 只做一次。另一方面，TEBD 会遇到除以可能非常小的奇异值的情况，这是潜在数值不精确性的一个重要来源；不过这些可以\cite{Hastings09}以低数值代价消除。从数值的角度看，tDMRG 并非先出现的 TEBD 的简单翻译，而是有其自身强弱之处的独立算法。

两种方法共享一个核心特征：时间演化与截断交织在一起——每次键演化之后都进行一次 SVD 截断。相比之下，tMPS 先演化所有键，然后通过 SVD 或迭代方式把矩阵维数从 $d^2 D \rightarrow D$ 压缩，从而截断整个态。 

tMPS 是更整洁的方法，而且还可以证明它更为精确。事实上，对于实时演化，它之于 tDMRG 或 TEBD，恰如迭代变分压缩之于 SVD 压缩，这意味着当态的改变很小时（例如时间步长非常小），差异会减小，因为截断之间的相互依赖变得不那么严重，只剩下非常温和的截断。上述关系确实存在，这一点可以通过不用变分、而仅用 SVD 来压缩一个 tMPS 态看出： 

取 $\ket{\psi}$ 为右规范的，在第一个键上做一步 tDMRG/TEBD 演化，或在所有奇数键上做 tMPS 步骤。截断现在要用 SVD 来进行，我的论断是：SVD 察觉不到这两个截然不同的时间演化后的态之间的差别。在第一个键本身上，所有方法产生相同的结构，但在所有其他格点上二者不同。这就是说，
\begin{equation}
\ket{\phi}_{{\rm tDMRG}} = \sum_{\fat{\sigma}} \Psi^{\sigma_1\sigma_2} B^{\sigma_3} B^{\sigma_4} \ldots \ket{\fat{\sigma}}  
\end{equation}
等于 
\begin{equation}
\ket{\phi}_{{\rm tMPS}} = \sum_{\fat{\sigma}} \Psi^{\sigma_1\sigma_2} \overline{B}^{\sigma_3} \overline{B}^{\sigma_4} \ldots \ket{\fat{\sigma}}  ,
\end{equation}
其中 $\overline{B}$ 矩阵来自与时间演化键算符的收缩。只要集合 $\{ B \}$ 与 $\{ \overline{B} \}$ 都生成正交归一的态集合，对第一个键的 SVD 对两个态就是等价的。情况确实如此：$B$ 矩阵按定义就生成这样的集合，而 $\overline{B}$ 矩阵生成的态与第一组正交归一态之间只差一个幺正变换（实时演化！）。这一方法等价的观察对链上更远的键同样成立。

因此，三种算法的差别只有在变分压缩的层面才显现出来。
 
\subsection{我们能走多远？} 
在这一节中，我描述了纯态与混合态时间演化的基本算法。误差有两个来源。其一是 Trotter 分解：对于 $n$ 阶分解，每个时间步 $\tau$ 产生 $O(\tau^{n+1})$ 的误差。由于共有 $t/\tau$ 个时间步，误差最终为 $O(\tau^n t)$，即随时间线性增长。这意味着它随时间的增长只是适度的，并且可以通过更小的时间步和/或更高阶的分解来压低。这一点为当前所有方法所共有\cite{Vidal03a,Vidal04b,Daley04,WhiteFeiguin04,VerstraeteRipoll04}。事实上，还有其他计算矩阵指数的方法，例如 Krylov 方法\cite{Schmitteckert04}，或如 \cite{Feiguin05} 中的前瞻（lookahead）步骤，它们能进一步减小这一误差。无论如何，从长远看它并不十分令人担心。 

另一方面，还有每个时间步之后对 MPS 膨胀起来的键维数做截断所带来的误差。这个误差是严重的；早已有工作表明它会导致随时间指数式爆增的误差\cite{Gobert05}。然而截断误差只是症状，并非根本问题：真正的原因是——遵照 Lieb-Robertson 定理——对于量子态的非平衡演化，纠缠 $S$ 可以随时间线性增长：$S(t) \leq S(0) + c t$，其中 $c$ 是与晶格中激发传播速度有关的某个常数\cite{Calabrese06}。对于许多量子淬火，这一线性界实际上会达到，此时哈密顿量的某个参数被突然改变，使得总能量发生广延量级的变化。
无论从 $D \sim 2^S$ 还是从严格的分析\cite{Osborne06}都可得出：在这种情形下，矩阵维数必须随时间指数增长，$D(t) \sim 2^t$；或者说，在矩阵维数固定时精度将指数式恶化。 

尽管如此，在许多情形下矩阵尺寸的增长足够缓慢，数值资源仍足以观察所关心的含时现象：含时 DMRG 从那时起已被广泛使用，并被发现为强关联一维系统的非平衡行为开辟了全新的视角（仅举几例：\cite{WhiteFeiguin04,Feiguin05a,Feiguin05,Gobert05,Kollath05,Trebst06,Cramer08,Barthel09a,Kollath06}）。
 
\section{含时模拟：扩展适用范围}
\label{sec:beyondtime}

时间演化——按历史先后次序，无论它是用 TEBD、tDMRG 还是 tMPS 实现的——从根本上受限于所能达到的时间。深层原因在于非平衡量子态中纠缠随时间（至多）线性累积，若要维持给定精度，这会转化为键维数 $D(t)$（至多）随时间指数式增长。我们会以指数的速度撞上“时间之墙”。能否把它推向更远的未来？有限温度计算也面临类似问题。虽然它们未必涉及动力学，但作为虚时演化来看，它们提出了自己的问题：静态或热力学计算中 $T\rightarrow 0$ 极限的问题，以及动力学方面——尤其是高温下的动力学——的问题。

我们尚未讨论的第二个问题涉及{\em 耗散}时间演化：我们关注的不再是纯哈密顿动力学中那样的封闭量子系统，而是开放量子系统。如果动力学是马尔可夫的（即环境没有记忆，这在许多情形下是一个极不平凡的假设），那么最一般的动力学由 Lindblad 方程给出。虽然形式上容易证明用 MPS 可以相当容易地模拟它，但在实际操作中这在数值上相当棘手，因此非常需要有更简单的方案。 

在本节中，我们将首先考虑通过评估时间演化张量网络的不同方案来扩展模拟时间范围的尝试。正如我们在波函数重叠的计算这类简单例子中已经看到的，收缩的次序可能极大地改变计算量。第二步，我们将考察一种预测方法：它拾取数值原始数据，并在数据满足某种数学形式的前提下，非常成功地把它们外推一个数量级，以有限温度的情形为例。第三步，我们将考察一种截然不同的有限温度模拟方法。最后，我将处理耗散动力学问题，并展示该领域取得的出色进展。

\subsection{收缩的次序：光锥、海森堡绘景与横向折叠}
让我们考虑如下含时期望值的计算
\begin{equation}
\bra{\psi(t)} \hat{O} \hat{P} \ket{\psi(t)} = \bra{\psi} \eul^{+\imag\hat{H}t} \hat{O} \hat{P} \eul^{-\imag\hat{H}t} \ket{\psi} .
\end{equation}
从 $\ket{\psi}$ 出发，将其演化到时刻 $t$，得到 $\ket{\psi(t)}$。然后如前所述，把两个算符夹在 $ \bra{\psi(t)}$ 与 $\ket{\psi(t)}$ 之间来计算期望值。但我们也可以把这一过程表示为对一个二维张量网络的（近似）收缩，如图~\ref{fig:2Dcontractiontime} 所示，随后沿时间方向逐行向内收缩。

\begin{figure}
\centering\includegraphics[scale=0.7]{PyMPS_MPSContractionTime2D.eps}
\caption{用于计算含时期望值的二维张量网络：顶行与底行分别表示 $\bra{\psi}$ 和 $\ket{\psi}$。两列 MPO 阵列（用括号标出）以 Trotter 化形式表示 $\eul^{+\imag\hat{H}t}$ 和 $\eul^{-\imag\hat{H}t}$；我不区分不同的局域 MPO，例如出现在最左列与最右列某些格点上的恒等算符。在中央行，我们放置恒等算符（白色方块）和待求值的算符。虚线标示沿时间方向的收缩积累。}
\label{fig:2Dcontractiontime}
\end{figure}

设 $t=N\Delta t$，每个 Trotter 步有 $n$ 个 MPO（例如一阶时为 2 个），我们便得到一个 $L \times (2nN+3)$ 个格点的点阵，即宽度为 $L$、奇数高度为 $2nN+3$。若把位于第 $i$ 行第 $j$ 列（如矩阵那样）格点上的张量记为 $T^{[i,j]}$，并把指标标记为上 $u$、下 $d$、左 $l$、右 $r$，仿照 MPO 把 $T^{[i,j]}$ 写成带指标 $T^{[i,j]u,d}_{l,r}$ 的形式，那么例如对 $i=1$（左矢态所在处）可以确认：
\begin{equation}
T^{[1,1] 1,d}_{1,r} = A^{[1]d*}_{1,r} \quad
T^{[1,j] 1,d}_{l,r} = A^{[j]d*}_{l,r} \quad
T^{[1,L] 1,d}_{l,1} = A^{[L]d*}_{l,1}
\end{equation}
其中 $1<j<L$。类似地，对 $i=2nN+3$（右矢态所在处）：
\begin{equation}
T^{[2nN+3,1] u,1}_{1,r} = A^{[1]u}_{1,r} \quad
T^{[2nN+3,j] u,1}_{l,r} = A^{[j]u}_{l,r} \quad
T^{[2nN+3,L] u,1}_{l,1} = A^{[L]u}_{l,1}
\end{equation}
而在第 $i=nN+2$ 行（算符所在处）：
\begin{equation}
T^{[nN+2,j] u,d}_{1,1} = \left\{ \begin{array}{cl} \hat{O}^{u,d} & {\rm on\ operator\ location\ } j \\
\delta_{u,d} & {\rm else} \end{array} \right\} .
\end{equation}
在这一行中，网络在水平方向上是若干标量的乘积，因此指标为 $(1,1)$。在所有其余各行上，张量 $T^{[i,j]}$ 由局域 MPO 给出：在满足 $nN+2 < i < 2nN+3$ 的所有行上 $T^{[i,j]u,d}_{l,r}= W^{[\alpha]u,d}_{l,r}$（类型 $\alpha$ 取决于所选的分解），而
$T^{[i,j]u,d}_{l,r}= W^{[\alpha]d,u*}_{l,r}$ 出现在满足 $1 < i < nN+2$ 的所有行上，后者对应于左矢的时间演化。

\subsubsection{时间演化中的光锥}
把左矢与右矢的时间演化放在一起考虑，实际上能带来重要的简化。在图~\ref{fig:lightcone1} 中，我恢复了一阶 Trotter 分解中键演化的交替模式，并用白色方块明确标出了单位算符的位置。我们想计算 $\langle \hat{O}(t) \rangle$，其中算符位于格点 2 上。让我们看最后几个 Trotter 步（第 5 行和第 7 行）。第 5 行有若干演化算符 $\eul^{+\imag \hat{h} \Delta t}$，第 7 行有相应的算符  $\eul^{-\imag \hat{h} \Delta t}$。这意味着它们相互抵消，可以用单位算符替换，但第 2、3 列除外，因为那里中间插有 $\hat{O}$。反过来再看第 4 行和第 8 行，此时第 5 至 8 列中出现了相互抵消的演化算符；其余那些不抵消，因为它们夹着非单位算符。就这样，我们一路向左矢和右矢态推进，直到再无抵消可能为止。所得的张量网络在很大程度上表现为：把键（行）维数大于 1 的复杂张量替换为键维数为 1 的单位张量，这意味着收缩变得平庸，无需压缩。剩下的是一个由演化算符构成的算法性 {\em 光锥}，它们能“看到”待求值的非平庸算符 $\hat{O}$ 的存在。注意，不要把这个算法光锥与物理光锥混为一谈：若令 $\Delta t\rightarrow 0$，则对任意固定时间 $t$，算法光锥都会变得无限宽。在物理上，Lieb-Robinson 定理指出，在宽度为 $x=2ct$ 的“光锥”之外（其中 $c$ 是某个依赖于问题的“速度”），关联呈指数快速衰减，这与相对论性光锥所强加的硬截断不同。物理光锥以及关联的这种特殊衰减，是上一节的 MPS/tDMRG/TEBD 算法由 Hastings\cite{Hastings09,Hastings08} 发展出的非常有趣的算法推广的基础，此处不再展开。

\begin{figure}
\centering\includegraphics[scale=0.25]{PyMPS_MPSContractionLightCone.eps}$\quad\quad$\includegraphics[scale=0.25]{PyMPS_MPSContractionLightCone2.eps}
\vskip 1cm
\centering\includegraphics[scale=0.25]{PyMPS_MPSContractionLightCone3.eps}$\quad\quad$\includegraphics[scale=0.25]{PyMPS_MPSContractionLightCone4.eps}
$\quad\quad$\includegraphics[scale=0.25]{PyMPS_MPSContractionLightCone5.eps}

\caption{光锥：各图应从左到右阅读，先上行、后下行。左上图中，一个算符夹在两个经时间演化的态之间；计入 Trotter 演化后，恒等算符用白色方块标出。从中心向上、下两边推进时，只要键演化算符不夹着该算符或（因而）幸存的键算符，它们就相互抵消为单位算符。结果便是一个由演化算符构成的光锥，四周围绕着在数值上平庸的单位算符。}
\label{fig:lightcone1}
\end{figure}

虽然这一结构在我们考察例如 $n$ 点关联子时会变得更复杂，但我们可以获得巨大的算法节省，代价是：对算符的不同位置，必须考虑新的网络。

若要在不同时刻（例如 $\langle \hat{O} (t_1) \rangle$、$\langle \hat{O} (t_2) \rangle$ 等）进行计算，代价是什么？这是一个很自然的问题，因为我们可能关心例如某个局域密度的时间演化。如果不用光锥，那么我们只是向内逐行计算收缩，算出某个平均值，取回存储的直到算符所在那一行的收缩结果，再加若干 Trotter 步，再收缩，再算某个平均值，如此往复。这正是我们一直以来的做法。当然，在时刻 $t_2 > t_1$ 作用的算符生成的光锥不同于且大于在时刻 $t_1$ 作用的算符生成的光锥。但如果哈密顿量不依赖于时间，较大的光锥在其尖端包含了较小的那个。因此，把时间演化反过来、由未来向过去推进，在数值上是有意义的。而这恰恰就是切换到海森堡绘景。

\subsubsection{海森堡绘景}
从数学上看，切换到海森堡绘景不过是重新加括号：
\begin{equation}
\langle \hat{O}(t) \rangle = \bra{\psi(t)} \hat{O} \ket{\psi(t)} = \bra{\psi} \eul^{+\imag \hat{H} t} \hat{O} \eul^{-\imag \hat{H} t} \ket{\psi} = \bra{\psi} \hat{O}(t) \ket{\psi},
\end{equation}
其中我们引入了含时算符
\begin{equation}
\hat{O}(t) =  \eul^{+\imag \hat{H} t} \hat{O} \eul^{-\imag \hat{H} t} .
\end{equation}
如果把其中出现的时间演化 Trotter 化，我们便恰好得到上一节的光锥结构，只是它尚未与左矢和右矢一起收缩；见图~\ref{fig:Heisenbergtime}。

\begin{figure}
\centering\includegraphics[scale=0.5]{PyMPS_MPSHeisenbergTime.eps}
\caption{海森堡绘景中算符 $\hat{O}$ 的时间演化，转换到 MPS/MPO 表示。算符周围每一组对称的层对应一个时间步（更确切地说，是一个 Trotter 时间步的一部分）。随着点阵中越来越大的区域受到 $\hat{O}$ 的影响，光锥不断变宽。恒等算符（白色方块）是出于对称性考虑而插入的，并无实际作用。}
\label{fig:Heisenbergtime}
\end{figure}

这使得可以在海森堡绘景中建立时间演化 \cite{Prosen09,Hartmann09}。具体而言，人们通过 MPO-MPO 乘法构造一个在空间上不断增长的 MPO，例如在计算 $\hat{H}^2$ 的作用时就会遇到这种乘法。设时间演化算符当前的 MPO 由形如 $O^{\sigma_i \sigma'_i}_{a_{i-1},a_i}$、键维数为 $D$ 的局域 MPO 组成（在其确实需要考虑的格点 $i$ 上），而时间演化 MPO（对 $\eul^{-\imag \hat{h} t}$ 而言）为
$W^{\sigma_i \sigma'_i}_{b_{i-1},b_i}$、键维数为 $D_W$，那么经过（部分）时间步之后，该算符为
\begin{equation}
O^{\sigma_i, \sigma'_i}_{(b_{i-1},a_{i-1},c_{i-1}), (b_i,a_i,c_i)}
= \sum_{\sigma^{''}_i} W^{\sigma^{''}_i,\sigma_i *}_{b_{i-1},b_i} \left( \sum_{\sigma^{'''}_i}  O^{\sigma^{''}_i \sigma^{'''}_i}_{a_{i-1},a_i} W^{\sigma^{'''}_i,\sigma'_i}_{c_{i-1},c_i} \right) . 
\end{equation}
键维数为 $DD_W^2$。随后按前文对 MPO 的说明，把这个算符压缩回键维数 $D$，基本上就是使用压缩 MPS 的方法。

这就建立了一种“常规”时间演化：经受时间演化的不再是一个态，而是一个以 MPO 形式表示的算符，它在每个时间步之后被压缩到可控的矩阵维数。我们基本上可以复用该算法的大部分内容。

这种表述有哪些潜在的优点与缺点？首先，算法光锥带来的节省被立即纳入其中。其次，我们可以指望截断更小：虽然海森堡绘景与薛定谔绘景中被收缩的网络完全相同，但薛定谔绘景中的{\em 截断}并未考虑算符，因而针对性较弱——可以设想，对于局域密度这类“简单”算符，态演化中的许多精细结构其实并不真正需要，而演化算符本身会告诉我们哪些信息是专门为这个算符所需要的。

相应的缺点当然是，对不同的算符需要重新进行计算；而在薛定谔绘景中，只要存储了时间演化后的波函数，就可以随时对任意算符求值。当然，这里相应的优点是：只要存储了时间演化后的算符，就可以随时对不同的态求值。

就写作之时而言，对于简单算符，可达到的时间跨度似乎确实能大幅延长，但我很难许诺某种经验法则。

\subsubsection{横向收缩与折叠}

当然，态随时间演化而逐步建立的做法符合我们对世界的直觉，但没有什么能阻止我们在空间方向上收缩同一网络，即逐列收缩；既然收缩次序对效率的影响可能相当大，这也许会有帮助。为了复用现有程序，可以简单地把当前网络逆时针旋转 90 度，便得到宽度为 $(2nN+3)$、高度为 $L$ 的点阵。如果我们继续按先竖后横、并按矩阵的行列逻辑来标记张量 $T^{[i,j]}$，那么新点阵中的张量为
\begin{equation}
\overline{T}^{[i,j] u,d}_{l,r} = T^{[L+1-j,i] r,l}_{u,d} ,
\end{equation}
如图~\ref{fig:latticerotation} 所示。然后我们可以再次逐行收缩。

事实证明，简单的旋转（或横向收缩）并不能延长可达到的时间尺度。正是通过一个额外的{\em 折叠}步骤，时间尺度才可能被大幅延长 \cite{Banuls09}。折叠沿新的“时间”（即真实空间）轴进行，并把新的“空间”（即真实时间）区域的范围减半（见图~\ref{fig:latticefolding}）。于是我们拥有的不再是格点 1 到 $2nN+3$，而是双重格点 1 到 $nN+2$：双重格点 1 包含旧格点 1 与 $2nM+3$，双重格点 2 包含旧格点 2 与 $2nN+2$；一般地，$i$ 包含旧格点 $i$ 与 $2nN+4-i$，直到数对 $(nN+1, nN+3)$。格点 $nN+2$ 是特殊的：由于我们折叠的是奇数个格点，会有一个格点保持单独，这就是格点 $nN+2$，它对应于含有在时刻 $t$ 待求值算符的那一行。在所有其余格点上，我们把对应于“相同”时间步——一个时间上向前、一个向后——的张量相互折叠在一起。

预期是，这种对向前与向后时间步的折叠会导致纠缠积累中的相消，从而能够达到更长的时间（$D$ 的增长没有那么快）。

为简化记号，我们定义 $L' = 2nN+3 \equiv 2\ell +1$。
如果把折叠的下端读作一个 MPS，那么折叠后的态同样以一个 MPS 开头，其矩阵构造为
\begin{equation}
M^{f[i]\Sigma_i}_{a^f_{i-1},a^f_i} = M^{f[i]\sigma_i, \sigma_{L'+1-i}}_{(a_{i-1},a_{L'+1-i}),(a_i,a_{L'-i})}
= \overline{T}^{[L,i]\sigma_i}_{a_{i-1},a_i} \overline{T}^{[L,L'+1-i]\sigma_{L'+1-i}}_{a_{L'-i},a_{L'+1-i}} 
\end{equation}  
对所有 $1 \leq i \leq \ell$ 成立，而
\begin{equation}
M^{f[\ell+1]\Sigma_{\ell+1}}_{a^f_{\ell},1} = M^{f[\ell+1]\sigma_{\ell+1}}_{(a_{\ell},a_{\ell+1}),1} 
=\overline{T}^{[L,\ell+1]\sigma_{\ell+1}}_{a_{\ell},a_{\ell+1}}.
\end{equation}  
我们在每个格点上（除格点 $\ell+1$ 外，那里维数仍为 $d$）定义了 $d^2$ 个“胖”局域态 $\ket{\Sigma_i} = \ket{\sigma_i}\ket{\sigma_{L'+1-i}}$（在编程时当然也可以引入一个虚设的格点）。类似地，我们为折叠后的 MPO 构造新的张量。

\begin{figure}
\centering\includegraphics[scale=0.5]{PyMPS_LatticeRotation.eps}
\caption{旋转点阵以复用代码：设采用时空标记 $[i,j]$，其中 $i$ 为时间、$j$ 为空间（旋转后指虚构的“时间”与“空间”），则张量指标 $u$、$d$ 与 $l$、$r$ 按图示交换位置。}
\label{fig:latticerotation}
\end{figure}
 
\begin{figure}
\centering\includegraphics[scale=0.5]{PyMPS_LatticeFolding.eps}
\caption{旋转后的点阵此时沿``空间''（原来的时间）轴折叠，该轴有 $2nN+3 \equiv L^{'}$ 个格点。对应同一时刻的各条线彼此重叠（用椭圆标示）；在末时刻对算符求值的那条线在弯折处保持单独。}
\label{fig:latticefolding}
\end{figure}

如果我们假设原点阵以及作用在其上的哈密顿量是平移不变的（至少对平移偶数个格点是如此），就可以方便地用转移算符写出这些收缩。若把折叠点阵第一行（第一个``时间''切片）上的态记为 $\ket{\psi_L}$（对应原点阵的格点 $L$），把最下面一行（最后一个``时间''切片）上的态记为 $\bra{\psi_1}$，则（为简单起见取 $i$ 为奇数）
\begin{equation}
\langle \hat{O}_i \rangle = \frac{\bra{\psi_L} E^{(i-3)/2} E_O E^{(L-i-1)/2} \ket{\psi_1}}{ \bra{\psi_L}  E^{(L-2)/2} \ket{\psi_1}} .
\end{equation}
这里我们引入了作用在长为 $\ell$、宽为 2 的条带上的转移算符 $E$ 与 $E_O$，如图~\ref{fig:bigtransferoperator} 所示（为便于表示，画的是未折叠、未旋转的形式）。$E_O$ 由 $E$ 得来：把格点 $\ell+1$ 处的单位算符换成 $\hat{O}$。

\begin{figure}
\centering\includegraphics[scale=0.5]{PyMPS_BigTransferOperator.eps}
\caption{转移算符 $E$（虚线矩形；注意作用方向）在时间演化网络的未折叠、未旋转表示中的示意。除对算符求值的格点被相应修改外，它通过平移两个格点在整个点阵中重复出现。}
\label{fig:bigtransferoperator}
\end{figure}

对于空间上有限的点阵，这可以通过迭代收缩与压缩来求值，但也可以取热力学极限。设 $E$ 有本征矢分解
\begin{equation}
E = \sum_{i} \lambda_i \ket{i} \bra{i} .
\end{equation}
注意 $E$ 不是厄米的，因此 $\ket{i}$ 与 $\bra{i}$ 不是互为伴随，而是各自不同的右本征矢与左本征矢。由二者的双正交归一性，
\begin{equation}
\lim_{L\rightarrow\infty} E^L = \lim_{L\rightarrow\infty} 
\sum_i \lambda_i^L \ket{i} \bra{i} = \lambda_0^L \ket{R}\bra{L},
\end{equation}
其中我假定最大本征值 $\lambda_0$ 非简并（通常正是如此），并把相应的右、左本征矢的记号改为 $\ket{R}$ 与 $\bra{L}$。

于是在热力学极限下得到期望值
\begin{equation}
\langle \hat{O}_i \rangle =  \frac{\bra{L} E_O \ket{R}}{\lambda_0} ,
\end{equation}
其中 $\lambda_0 = \bra{L} E \ket{R}$。两点关联子则由下式给出
\begin{equation}
\langle \hat{O}_i \hat{P}_j \rangle =  \frac{\bra{L} E_O E^r E_P \ket{R}}{\lambda_0^{r+2}},
\end{equation}
其中 $r$ 是格点 $i$ 与 $j$ 之间转移算符的个数。

为了对这类表达式求值，我们显然需要 $\lambda_0$、$\ket{R}$ 与 $\bra{L}$。由于 $E$ 的矩阵元是显式已知的，我们同样容易构造其转置，从而把问题全部化为求两个右本征矢。既然要找的是最大本征值，幂法（把 $E$ 反复作用到一个猜测矢量上）即可奏效。但也可以把 $E$ 看作某个非厄米算符的 MPO 表示，复用迭代的基态搜索技术，只需两处修改：求最低本征值改为求最高本征值，且常规 Lanczos 算法必须换成非厄米方法，即双正交 Lanczos 算法或 Arnoldi 方法。

虽然编码比标准时间演化更为复杂，但可达到的时间尺度大幅延长，3 到 5 倍似乎很容易实现。
  
\subsection{线性预测与谱函数}
谱函数 $S(k,\omega)$ 是多体物理中最重要的理论与实验量之一。虽然在 $T=0$ 下有非常精确的直接计算方法 \cite{Hallberg95,Ramasesha97,Kuhner99,Jeckelmann02}，但也有一种间接方法，由 \cite{WhiteFeiguin04} 开创：先计算实时实空间关联子如 $\langle \hat{S}^+_i(t) \hat{S}^-_j(0) \rangle$，再作到动量空间与频率空间的双重傅里叶变换。该方法的优点是可以无缝推广到有限 $T$，但受限于可达到的长度与时间尺度。

其中时间的限制严重得多，原因在于纠缠随时间迅速增长。可达到的时间尺度大多十分有限，以至于直接做傅里叶变换会产生强烈的混叠，或者必须对原始数据加窗而使谱信息被明显抹平。不过，这一限制可以借助线性预测技术在极低的数值代价下规避，无论 $T=0$\cite{Pereira08,White08} 还是 $T>0$\cite{Barthel09b}；该技术延长了可达到的 $t$，从而大大细化频域中的结果。

对于由 DMRG 得到的、在等间隔时刻 $t_n=n\*\Delta t$（最大时间 $t_\obs:= N \Delta t$）处的复数据时间序列 $x_0, x_1, \ldots, x_n, \ldots, x_N$，我们对 $x_{N+1}, x_{N+2}, \ldots$ 作预测。
对于 $t=t_\obs$ 之后的数据点，线性预测作如下拟设
\begin{equation}\label{eq:predictionAnsatz}
\tilde{x}_n = - \sum_{i=1}^p a_i x_{n-i} .
\end{equation}
时间步 $n$ 处的（预测）值 $\tilde{x}_n$ 被假定为前面 $p$ 个值 $\{x_{n-1},\dots,x_{n-p}\}$ 的线性组合。$a_i$ 一旦由已知数据确定，就被用来计算所有 $n>N$ 的 $x_n$（的近似）。

系数 $a_i$ 的确定方法是：在已知数据的子区间 $t_n \in (t_\obs-t_\fit,t_\obs]$（对应集合 $\Omega = \{ n\, |\, t_\obs - t_\fit < n\Delta t \leq t_\obs \}$）上最小化预测的最小二乘误差，即最简单的做法是最小化 $E\equiv\sum_{n\in\Omega} |\tilde x_n-x_n|^2$。$t_\fit=t_\obs/2$ 常是稳健的选择，既能减小短时间的影响，又有足够多的数据点。对 $a_i$ 最小化 $E$ 得到线性方程组
\begin{equation}
\label{eq:linPred:stationaryError}
R \vec{a} = -\vec{r},
\end{equation}
其中 $R$ 与 $\vec{r}$ 是自关联
$R_{ji} = \sum_{n\in\Omega} x^*_{n-j} x_{n-i}$ 与 $r_{j} = \sum_{n\in\Omega} x^*_{n-j} x_n$。
式~(\ref{eq:linPred:stationaryError}) 的解为 $\vec{a}=-R^{-1}\vec{r}$。

人们或许会问，为什么朝无限时间的外推能够以这种方式实现。
如下文所示，线性预测生成的是振荡项与指数衰减（或增长）项的叠加，这正是多体物理中自然出现的一类时间依赖性：
典型形式为 $G(k,\omega) = (\omega - \epsilon_k - \Sigma(k,\omega))^{-1}$ 的格林函数在时间-动量表示中由极点主导；例如对于位于 $\omega=\omega_1 -\imag \eta_1$、留数为 $c_1$ 的单个单极点，格林函数将为
$G(k,t) = c_1 \eul^{-\imag \omega_1 t - \eta_1 t}$，而对更复杂的极点结构，它也类似地是这类项的叠加。通常只有少数极点起作用，线性预测的拟设很适合我们所关心的响应量的典型性质。当这样的拟设不成立时，该方法多半不适用。

为看出预测所生成时间序列的特殊形式，
我们引入矢量 $\vec{x}_n := [x_{n}, \ldots, x_{n-p+1}]^T$，使得 (\ref{eq:predictionAnsatz}) 取如下形式
\begin{equation}
\tilde{\vec{x}}_{n+1} = A \* \vec{x}_n,	
\end{equation}
 其中
\begin{equation}\label{eq:predictionMatrix}
A\equiv
\left[
\begin{array}{ccccc}
 -a_1&-a_2&-a_3&\cdots&-a_p\\
  1  & 0  & 0  &\cdots& 0  \\
  0  & 1  & 0  &\cdots& 0  \\
\vdots&\ddots&\ddots&\ddots&\vdots\\
  0   &\cdots& 0  & 1  & 0
\end{array}  
\right] ,
\end{equation}
其中 $a_i$ 就是上面求得的矢量 $\vec{a}$ 的元素。
预测因此对应于把 $A$ 的幂作用到初始矢量 $\vec{x}_N$ 上。对 $A$ 作（非厄米的）本征矢分解，本征值为 $\alpha_i$，得到
\begin{equation}
\tilde x_{N+m} = [A^m \* \vec{x}_N]_1=\sum_{i=1}^p c_i \alpha_i^m,
\end{equation}
其中系数 $c_i$ 由 $\vec{x}_N$ 与 $A$ 的本征矢确定。本征值 $\alpha_i$ 编码了物理共振频率与阻尼。二者的联系为
$\alpha_i = \eul^{\imag  \omega_i \Delta t - \eta_i \Delta t}$。
可能出现 $|\alpha_i|\geq 1$ 的伪本征值，但可以按 \cite{Barthel09b} 所述加以处理。

在 $T=0$ 下，临界一维系统的时间关联子呈幂律衰减。此时用指数衰减的叠加来模拟这些幂律 \cite{Pereira08}。在有限温度下，时间关联子 $S(k,t)$ 在长时间后典型地呈指数衰减（源于热致展宽），这使线性预测特别适合这种情形。这也接近典型的实验情形，例如一维磁链的非弹性中子散射。

作为例子，我们考虑无外场、$J^z=0$（$XY$ 链）与 $J^z=1$ 的海森堡反铁磁体。前一情形允许精确的解析解。结果表明，预测可以把时间序列 $S(k,t)$ 延长一个数量级以上而精度无明显损失。在频率空间中，这对应于极高精度的谱线形（图 \ref{fig:xychain}）。

\begin{figure}
\centering\includegraphics[width=250pt]{PyMPS_TransIsing.eps}
\caption{线与点分别表示 $XY$ 链谱函数 $S^{+-}(k,\omega)$ 的线形在温度 $\beta=10$ 与 $\beta=50$（宽与窄线形）、（自左向右）$k=\pi/8$、$k=\pi/3 $、$k=3\pi/4 $ 处的精确解析解与数值解。虚线是如果不作预测、而是在傅里叶变换前对原始数据加窗时，从 $\beta=10$ 模拟中能最优提取出的线形。改编自 \cite{Barthel09b}。}
\label{fig:xychain}
\end{figure}

由于 $XY$ 链的色散关系只是一条简单的磁振子线，其自能结构非常简单，因此预测方法很容易应用。作为更有挑战性的例子，我们考虑各向同性 $S=1/2$ 链的自旋子连续谱；见图~\ref{fig:USHAFM}。在零温极限下，结果与 Bethe 拟设结果符合极好（残余差异难以归因：这里的 Bethe 拟设只能近似求值 \cite{Caux06}）。在有限温度下，不同精度的模拟表明结果已完全收敛且本质精确。这让我们有理由期待该方法将成为例如模拟中子散射实验结果的有力工具。

\begin{figure}
\centering\includegraphics[height=200pt]{PyMPS_HAFM.eps}
\caption{各向同性海森堡链在两个动量、四个不同温度下的谱函数。在 $T=0$ 时，Bethe 拟设（B.A.）与数值结果符合极好。改编自 \cite{Barthel09b}。}
\label{fig:USHAFM}
\end{figure}

\subsection{最小纠缠典型热态}
为计算静态与动态性质而模拟热密度算符的效果非常好。理论上，这一方法不存在任何限制。但在实践中会遇到若干局限。一方面，在极低温度 $T\rightarrow 0$ 下模拟变得困难。在此极限下，物理系统 P 上的混合态将演化趋向于基态上的纯态投影算符 $\hat{\rho}^P_\infty = \ket{\psi_\infty}_P \phantom{\,}_P\bra{\psi_\infty}$（这里我用 $\ket{\psi_\infty}_P$ 表示物理哈密顿量 $\hat{H}$ 的基态，与纯化一节一样指 $\beta=\infty$）。但在此极限下 P 不再与辅助系统 Q 纠缠，我们模拟的是两个纯态的乘积：为简单起见，设基态能量为零。现在考虑辅助系统的约化密度算符。在不计无关紧要的范数的意义下，
\begin{equation}
\hat{\rho}^Q_{\beta\rightarrow\infty} = \lim_{\beta\rightarrow\infty} \tr_P \eul^{-\beta\hat{H}/2} \ket{\psi_0} \bra{\psi_0} \eul^{-\beta\hat{H}/2} = \phantom{\,}_P\braket{\psi_\infty}{\psi_0} \braket{\psi_0}{\psi_\infty}_P ,
\end{equation}
因为在 $\beta\rightarrow\infty$ 的极限下，迹化简为基态的贡献，
$\tr_P \rightarrow \phantom{\,}_P\bra{\psi_\infty} \cdot \ket{\psi_\infty}_P$。利用密度算符的 $\beta=0$ 纯化 $\ket{\psi_0} = d^{-L/2} \sum_{\fat{\sigma}} \ket{\fat{\sigma}}_P \ket{\fat{\sigma}}_Q$，最后一个表达式（同样不计无关紧要的范数）化简为
\begin{equation}
\hat{\rho}^Q_{\beta\rightarrow\infty} = \sum_{\fat{\sigma}\fat{\sigma}'} \psi^{\fat{\sigma}*}_\infty \ket{\fat{\sigma}}_Q \phantom{\,}_Q \bra{\fat{\sigma}'} \psi^{\fat{\sigma}'}_\infty =
\ket{\psi_\infty}_Q \phantom{\,}_Q \bra{\psi_\infty} ,
\end{equation}
其中 $\psi^{\fat{\sigma}}_\infty$ 是基态的展开系数。因此，零温纯化是一个乘积态
\begin{equation}
\ket{\psi_\infty} = \ket{\psi_\infty}_P \ket{\psi_\infty}_Q,
\end{equation}
其中后一个态正是在辅助态空间上定义的物理基态。
假设用矩阵维数 $D$ 就足以足够精确地描述该态，那么这个乘积将需要维数为 $D^2$ 的矩阵来描述。这实际上意味着，对所研究的问题，我们算法的标度是特征矩阵维数的六次方而非三次方。另一方面——这正是预测方法试图解决的问题——我们在高温 $T\rightarrow\infty$ 时尤其会遇到长时间模拟的困难：许多态以相近但不完全相同的权重作出贡献，而 MPS 在编码如此众多的贡献态方面效率不高。

正如 White \cite{Stoudenmire10,White09} 所指出的，可以完全避开纯化方法，转而在一组精心挑选的热态——即所谓{\em 最小纠缠典型热态}（METTS）——上作抽样。已有工作表明，这一方法能大大缓解第一个局限，而对第二个局限目前所知尚不多。

热平均由下式给出
\begin{equation}
\langle \hat{A} \rangle = \frac{1}{Z} \tr \eul^{-\beta \hat{H}} \hat{A} = \frac{1}{Z} \sum_n \eul^{-\beta E_n} \bra{n} \hat{A} \ket{n} ,
\end{equation}
其中与所有教科书一样，我选用了热密度算符的能量表象。正如 Schr\"{o}dinger 许多年前就已指出的，这在数学上是正确的，却不合物理实际，因为有限温度下的真实系统通常并不处于本征态的统计混合之中——本征态在与环境耦合时是极为脆弱的。但取迹所用基的选择是任意的，也可以写成
\begin{equation}
\langle A \rangle = \frac{1}{Z} \sum_i \bra{i} \eul^{-\beta\hat{H}/2} \hat{A} \eul^{-\beta\hat{H}/2} \ket{i}
= \frac{1}{Z} \sum_i P(i) \bra{\phi(i)} \hat{A} \ket{\phi(i)}
\end{equation}
其中 $\{ \ket{i} \}$ 是任意一组正交归一基，而 $\ket{\phi(i)} = P(i)^{-1/2} \eul^{-\beta\hat{H}/2} \ket{i}$。由 $P(i) = \bra{i} \eul^{-\beta \hat{H}} \ket{i}$ 可知 $\ket{\phi_{i}}$ 是归一化的。容易看出 $\sum_i P(i) = Z$，故 $P(i)/Z$ 是概率。于是，可以以概率 $P(i)/Z$ 对 $\ket{\phi(i)}$ 进行{\em 抽样}，并对 $\bra{\phi(i)} \hat{A} \ket{\phi(i)}$ 求平均，从而对 $\langle \hat{A}\rangle$ 作统计估计。

要把这一切变成实际可行的算法，还有几个问题。既然我们并不知道那个复杂的概率分布，怎样才能正确抽样？能否选出一组态，使得平均值收敛最快？考虑到虚时演化将是算法的一部分，能否找到低纠缠的基 $\{ \ket{i} \}$，使得时间演化可以用适中的 $D$ 完成，即运行得很快？

为解决这些问题，White 选取由乘积态构成的计算基作为正交归一基，
\begin{equation}
\ket{i} = \ket{i_1}\ket{i_2}\ket{i_3} \ldots \ket{i_L} ,
\end{equation}
即经典（无纠缠）乘积态（CPS），它们可以精确地用 $D=1$ 表示。相应的态
\begin{equation}
\ket{\phi(i)} = \frac{1}{\sqrt{P(i)}} \eul^{-\beta \hat{H}/2} \ket{i}
\end{equation}
就是所谓的{\em 最小纠缠典型热态}（METTS）：我们的希望是，虽然虚时演化会因哈密顿量的作用而引入纠缠，但可以合理地预期，最终的纠缠会低于对已经纠缠的态作类似演化所得的纠缠。虽然从严格的数学意义上讲这并不完全成立，但从实用的角度看似乎确实如此！与纯化相比，这将是快得多的计算，尤其是因为纯化的因子化问题不会出现。

为了以正确的概率分布（我们无法算出它）对 $\ket{\phi(i)}$ 抽样，人们使用与蒙特卡罗相同的技巧，生成一个由态构成的马尔可夫链，
$\ket{i_1} \rightarrow \ket{i_2} \rightarrow \ket{i_3} \rightarrow \ldots$，使其复现正确的概率分布。从这一分布我们可以生成 $\ket{\phi(i_1)}$、$\ket{\phi(i_2)}$、$\ket{\phi(i_3)}$、$\ldots$，用以计算平均值。

算法运行如下：从一个随机的 CPS $\ket{i}$ 出发，重复以下三个步骤，直到结果的统计性足够好：
\begin{itemize}  
\item 通过虚时演化计算 $\eul^{-\beta\hat{H}/2} \ket{i}$ 并将态归一化（其模方为 $P(i)$，但算法中用不到它）。
\item 按形如 $\bra{\phi(i)} \hat{A} \ket{\phi(i)}$ 的方式计算所需的量，以供平均。
\item 通过量子测量把态 $\ket{\phi(i)}$ 坍缩为一个新的 CPS $\ket{i'}$，其概率为 $p(i\rightarrow i') = |\braket{i'}{\phi(i)}|^2$，然后从这个新态重新开始。
\end{itemize}

按照 \cite{Stoudenmire10}，我们来确认这种做法确实给出正确的抽样。由于 $\ket{\phi(i)}$ 与 $\ket{i}$ 服从相同的分布，只需证明后者被正确抽样。设前一个 CPS $\ket{i}$ 已按正确概率 $P(i)/Z$ 被抽取，问此时以多大的概率坍缩到某个 $\ket{j}$，可以发现
\begin{equation}
\sum_i \frac{P(i)}{Z} p(i\rightarrow j) = \sum_i \frac{P(i)}{Z} | \braket{j}{\phi(i)} |^2 = \sum_i
\frac{1}{Z} | \bra{j} \eul^{-\beta \hat{H}/2} \ket{i}|^2 = \frac{1}{Z} \bra{j} \eul^{-\beta\hat{H}} \ket{j} = \frac{P(j)}{Z}.
\end{equation}
这表明所期望的分布是更新过程的不动点。因此这一做法是有效的；不过，与蒙特卡罗中一样，明智的做法是舍弃最初的若干数据点，以消除初始 CPS 带来的偏差。事实证明——在舍弃最初几个 METTS 以消除初选的影响之后——只需对一百个左右的态求平均，就能以很高的精度计算局域静态量（磁化强度、键能）。

虽然我们已经知道如何做虚时演化，但还须讨论坍缩步骤。事实证明，利用 MPS 的结构，可以使算法这一部分比虚时演化快得多。

对每个格点 $i$，我们选取一个{\em 任意}的 $d$ 维正交归一基 $\{ \ket{\tilde{\sigma}_i} \}$，以区别于计算基 $\{ \ket{\sigma_i} \}$。由此可以构造投影算符 $\hat{P}^{\tilde{\sigma}_i} = \ket{\tilde{\sigma}_i} \bra{\tilde{\sigma}_i}$，局域坍缩到态 $\ket{\tilde{\sigma}_i}$ 的概率按标准量子力学给出为
$p_{\tilde{\sigma}_i} =\bra{\psi} \hat{P}^{\tilde{\sigma}_i} \ket{\psi}$。若把 $\ket{\psi}$ 坍缩为 CPS
$\ket{\psi'} = \ket{ \tilde{\sigma}_1} \ket{ \tilde{\sigma}_2} \ldots \ket{ \tilde{\sigma}_L}$，则概率为  $\bra{\psi} \hat{P}^{\tilde{\sigma}_1} \ldots \hat{P}^{\tilde{\sigma}_L} \ket{\psi} = | \braket{\psi'}{\psi} |^2$，正如算法所要求的。单格点坍缩之后，波函数变为
\begin{equation}
\ket{\psi} \rightarrow  p_{\tilde{\sigma}_i}^{-1/2} \hat{P}^{\tilde{\sigma}_i}\ket{\psi} ,
\end{equation}
其中前置因子保证坍缩后态的正确归一化，与初等量子力学中一样。举例来说，对于沿任意轴 ${\mathbf n}$ 测量的 $S=\frac{1}{2}$ 自旋，投影算符为
\begin{equation}
\hat{P}^{\uparrow_n,\downarrow_n} = \frac{1}{2} \pm {\mathbf n} \cdot \hat{{\mathbf S}}_i .
\end{equation}

正如 \cite{Stoudenmire10} 所指出的，如果利用 MPS 与 CPS 的两个特点，对所有格点依次进行局域测量和坍缩可以非常高效地完成：(i) 局域期望值若位于显式格点上（用 DMRG 的语言），或位于混合规范形式中左规范与右规范格点之间的格点上（用 MPS 的语言），就可以非常高效地求值；(ii) 坍缩之后，所得的态是所有已坍缩格点上的局域态与未坍缩的其余部分构成的乘积态。

设 $\ket{\psi}$ 是右规范的（放宽一点：格点 1 上的归一化无关紧要）。于是 $p_{\tilde{\sigma}_1} =\bra{\psi} \hat{P}^{\tilde{\sigma}_1} \ket{\psi}$ 的计算变得平凡，因为格点 2 到 $L$ 上期望值网络的收缩恰好给出 $\delta_{a_1,a'_1}$。于是
\begin{equation}
\bra{\psi} \hat{P}^{\tilde{\sigma}_i} \ket{\psi} = \sum_{a_1,\sigma_1,\sigma'_1} B^{\sigma_1 *}_{a_1} \bra{\sigma_1} \hat{P}^{\tilde{\sigma}_1} \ket{\sigma'_1} B^{\sigma'_1}_{a_1} = 
\sum_{a_1} \left[ \sum_{\sigma_1} B^{\sigma_1*}_{a_1} \braket{\sigma_1}{\tilde{\sigma}_1} \right]
\left[ \sum_{\sigma'_1} B^{\sigma'_1}_{a_1} \braket{\tilde{\sigma}_1}{\sigma'_1} \right] .
\label{eq:probsinmetts}
\end{equation}
这个表达式看起来非常特殊地依赖于第一个格点（因为开边界），但正如我们将看到的，其实不然！

算出格点 1 上坍缩的概率之后，按刚才生成的分布随机抽取一个具体的坍缩 $\ket{\tilde{\sigma}_1}$。坍缩后的态将形如 $\ket{\tilde{\sigma}_1} \ket{\psi_{{\rm rest}}}$，因而是一个乘积态。因此，格点 1 上的新矩阵（我们记作 $A^{\sigma_1}$）必定都是标量，即 $D=1$ 的矩阵。由
$\ket{\tilde{\sigma}} = \sum_{\sigma} \ket{\sigma} \braket{\sigma}{\tilde{\sigma}} = \sum_\sigma \ket{\sigma} A^{\sigma}$ 可知它们为
\begin{equation}
A^{\sigma_1}_{1,1} = \braket{\sigma_1}{\tilde{\sigma}_1} ;
\end{equation}
容易看出左规范自动得到满足，因此用字母 $A$ 标记。但 $A^{\sigma_1}$ 维数的这一改变意味着 $B^{\sigma_2}$ 也必须改变，即
\begin{equation}
B^{\sigma_2}_{a_1,a_2} \rightarrow M^{\sigma_2}_{\tilde{\sigma}_1,a_2} =
p_{\tilde{\sigma}_1}^{-1/2} \sum_{\sigma_1,a_1} \braket{\tilde{\sigma}_1}{\sigma_1} B^{\sigma_1}_{a_1} B^{\sigma_2}_{a_1 a_2} .
\end{equation}
由于标记 $\tilde{\sigma}_1$ 取确定值，它只是一个哑指标，$M^{\sigma_2}$ 的行维数也只是 1，与第一个格点上的矩阵一样。因此，式~(\ref{eq:probsinmetts}) 可以推广到所有格点，而代价最高的步骤是 $B^{\sigma_2}$ 的更新，其标度为 $D^2 d^2$，而不是像时间演化那样的 $D^3$。

为看清这一代换，我们把 $\hat{P}^{\tilde{\sigma}_1}$ 写成 MPO：$W^{\sigma_1 \sigma'_1} = \braket{\sigma_1}{\tilde{\sigma}_1} \braket{\tilde{\sigma}_1}{\sigma'_1}$。于是坍缩后的 $\ket{\psi}$ 为
\begin{eqnarray*}
& & p_{\tilde{\sigma}_1}^{-1/2} \sum_{\fat{\sigma}} \sum_{\sigma'_1} \braket{\sigma_1}{\tilde{\sigma}_1} \braket{\tilde{\sigma}_1}{\sigma'_1} B^{\sigma'_1} B^{\sigma_2} \ldots \ket{\fat{\sigma}} = \\
& & \sum_{\fat{\sigma}} \sum_{a_2} A^{\sigma_1}_{1,1} \left( \sum_{\sigma'_1 a_1}  p_{\tilde{\sigma}_1}^{-1/2}  \braket{\tilde{\sigma}_1}{\sigma'_1} B^{\sigma'_1}_{1,a_1} B^{\sigma_2}_{a_1,a_2} \right) (B^{\sigma_3} \ldots B^{\sigma_L})_{a_2,1} \ket{\fat{\sigma}} ,
\end{eqnarray*}
由此即得上述代换。

再作几点评论。每个格点上的测量基都可以随机选取；为了获得马尔可夫链较短的自关联``时间''（即高质量的抽样），这样做当然很好，但比总是坍缩到同一个基要昂贵得多，而后者却会产生遍历性问题。建议在两个基之间交替切换，使得对每个基而言，其投影算符在另一个基中都是最大程度混合的（例如，沿 $x$ 轴与 $z$ 轴（或 $y$ 轴）交替测量自旋）。这样自关联时间可以降到 5 步左右 \cite{Stoudenmire10}。至于统计误差的估计，在最简单的情形下，只需对长度大于自关联时间的分箱求平均，再考察这些分箱平均值的统计分布，即可得到误差棒。

仍有一些耐人寻味的问题，既涉及这一算法的潜力，也涉及其理论基础：对于结构函数所需的更长程的关联子，它表现如何？动力学量不难得到，因为弱纠缠的 METTS 的时间演化代价并不高——但对少数几个``典型"态求平均的高效率是否仍会保持？
\subsection{耗散动力学：量子跳跃}
当物理系统 A 与某个环境 B 相耦合、从而使 A 成为{\em 开放}量子系统时，就会出现耗散（即非哈密顿的）动力学。这是一个内容非常丰富的物理学领域，这里只回顾几条在此有用的核心结果。系统密度算符的时间演化总可以写成 Kraus 表示
\begin{equation}
\hat{\rho}_A (t) = \sum_j \hat{E}_A^j(t) \hat{\rho}_A (0) \hat{E}_A^{j\dagger}(t),
\end{equation}
其中 Kraus 算符满足条件
\begin{equation}
\sum_j \hat{E}_A^{j\dagger}(t) \hat{E}_A^j(t) = \hat{I}_A .
\end{equation}
如果动力学是无记忆的（马尔可夫型），它只依赖于无穷小时间之前的密度算符，这时可以导出一个主方程，即 Lindblad 方程。在 $\diff t\rightarrow 0$ 的极限下，环境以 $p_0 \rightarrow 1$ 的概率保持不变，而以正比于 $\diff t$ 的概率发生改变（量子跳跃）。若把 Kraus 算符 $\hat{E}_A^0(t)$ 与"无改变"对应起来，一个合理的按时间标度的拟设是
\begin{equation}
\hat{E}_A^0(\diff t) = \hat{I}_A + O(\diff t) \quad\quad \hat{E}_A^j(\diff t) = \sqrt{\diff t} \hat{L}_A^j  \quad (j>0)
\end{equation}
或更精确地
\begin{equation}
\hat{E}_A^0(\diff t) = \hat{I}_A + (\hat{K}_A - \imag \hat{H}_A) \diff t ,
\end{equation}
其中 $\hat{K}_A$ 和 $\hat{H}_A$ 是两个厄米算符。Kraus 算符的归一化条件给出
\begin{equation}
\hat{K}_A = -\frac{1}{2} \sum_{j>0} \hat{L}_A^{j\dagger} \hat{L}_A^j .
\end{equation}
由这些拟设可以从 Kraus 演化公式导出一个微分方程，即 Lindblad 方程
\begin{equation}
\frac{\diff \hat{\rho}}{\diff t} = - \imag [\hat{H},\hat{\rho}] + \sum_{j>0} \left( \hat{L}^j \hat{\rho} \hat{L}^{j\dagger} - \frac{1}{2} \{ \hat{L}^{j\dagger} \hat{L}^j, \hat{\rho} \} \right) ,
\end{equation} 
其中我略去了指标 A。的确，在没有量子跳跃时（只有 $j=0$），就回到 von Neumann 方程。以非厄米性为代价，此方程可以化简。若引入 $\hat{H}_{{\rm eff}} := \hat{H} + \imag \hat{K}$，则最后一项消失，得到
\begin{equation}
\frac{\diff \hat{\rho}}{\diff t} = - \imag [\hat{H}_{{\rm eff}}\hat{\rho}-\hat{\rho}\hat{H}_{{\rm eff}}^\dagger] + \sum_{j>0} \hat{L}^j \hat{\rho} \hat{L}^{j\dagger} .
\label{eq:Lindblad2}
\end{equation} 

在 MPS 形式框架下模拟 Lindblad 方程是相当容易的，特别是利用 MPO \cite{VerstraeteRipoll04}，也可以采用超算符形式 \cite{VidalZwolak04}。这一做法的问题在于，与一个态在哈密顿量下的演化相比，它在数值上代价更高。\cite{Daley09} 提出了一种非常有吸引力的替代方案，它可以最大限度地复用现有的纯态程序，将纯态时间演化与量子轨迹方法结合起来。

量子轨迹方法在量子光学中已得到广泛应用\cite{Plenio98}。量子轨迹方法不直接使用 Lindblad 方程（后者同时考虑了通过 $\hat{\rho}(0)$ 体现的初始态的概率分布以及所有可能的量子跳跃序列），而是对初始态的分布进行抽样，并对每个抽样态做纯态时间演化，其中随机的量子跳跃在随机时刻发生。这些选取方式使得如果对物理量在这族随时间演化的态的分布上取平均，就能重新得到 Lindblad 方程的结果。让我们略去对初始态的抽样（假定初始态总是相同的），转而关注概念上更困难的对量子跳跃的平均。

算法随后在时间区间 $[0,T]$ 内（其中 $T$ 是模拟的终止时刻）按如下方式生成 ${\mathcal N}$ 条量子轨迹：
\begin{itemize}
\item 生成一个初始态 $\ket{\psi(0)}$；根据物理问题的不同，它要么正确地抽样 $t=0$ 时刻的密度算符，要么干脆总是取同一个态。
\item 在 $[0,1]$ 中选取一个均匀分布的随机数 $p_0$。
\item 用我们的某种纯态时间演化方法，让 $\ket{\psi(0)}$ 在 $\hat{H}_{{\rm eff}}$ 作用下做时间演化。由于有效哈密顿量是非厄米的，态的范数会随时间减小。在时刻 $t_1$ 停止时间演化，该时刻由
$\braket{\psi(t_1)}{\psi(t_1)} = p_0$ 定义；这就是第一次量子跳跃发生的时刻。注意若 $T<t_1$，则模拟在 $T$ 处停止，我们得到一条没有跳跃的轨迹，并对末态做归一化。
\item 为了在 $t_1$ 时刻实施量子跳跃，我们计算
\begin{equation}
\tilde{p}_j = \bra{\psi(t_1)} \hat{L}^{j\dagger} \hat{L}^j \ket{\psi(t_1)} \quad\quad p_j = \frac{\tilde{p}_j}{\sum_{j>0} \tilde{p}_j} \quad (j>0)
\end{equation}
并按照归一化的概率分布 $\{ p_j \}$ 选取一个 $j$。
\item 我们实施这一跳跃并将态归一化，
\begin{equation}
\ket{\psi(t_1^+)} = \frac{\hat{L}^j \ket{\psi(t_1)}}{\| \hat{L}^j \ket{\psi(t_1)} \|} .
\end{equation}
\item 此后，我们继续寻找新的 $p_0$，由它出发，$\ket{\psi(t_1^+)}$ 在 $\hat{H}_{{\rm eff}}$ 作用下的时间演化给出第二次量子跳跃的位置 $t_2$，如此继续，直到超过 $T$。
\end{itemize}
现在把直到时刻 $T$ 的物理量对已生成的 ${\mathcal N}$ 条量子轨迹取平均。如果所有态在任何时刻都被归一化，就会得到正确的概率；但算法中情形并非如此，因此必须把例如时刻 $t$ 的范数考虑进去。显然，需要对 ${\mathcal N} \rightarrow \infty$ 的收敛性做仔细分析，不过看起来对于少量的跳跃算符，甚至几百条轨迹就可能给出高度可靠的结果 \cite{Daley09}。

关于这种抽样能重现 Lindblad 方程动力学的观察，属于量子轨迹方面的标准文献内容（原文此处拼写为 trajactories）。证明可以分两步进行，这里只作简述。第一步，考虑固定的时间步长 $\diff t$，计算该时间区间内"无跳跃"与"发生跳跃 $j$"的概率（$p_j = \diff t \bra{\psi(t)} \hat{L}^{j\dagger} \hat{L}^j \ket{\psi(t)}$，$p_0 = 1 - \sum_j p_j$）。然后按照生成的分布，要么在 $\hat{H}_{{\rm eff}}$ 作用下演化 $\diff t$ 并归一化，要么实施跳跃并归一化。可以证明这样能重现 Lindblad 方程。第二步，证明按此方式生成的量子跳跃分布与我们在算法中使用的分布是相同的。

\section{DMRG 与 NRG}
\subsection{Wilson 的数值重正化群（NRG）与 MPS}
Wilson 的数值重正化群（NRG）\cite{Wilson75,Bulla08,Krishnamurthy80} 源于解释为何含少量磁性杂质的金属表现出非单调电阻率行为的尝试。
人们发现，一个足够好的极小模型由
\begin{equation}
\hat{H}_A = \sum_{k\sigma} \epsilon_k \hat{c}^\dagger_{k\sigma} \hat{c}^{\phantom\dagger}_{k\sigma} + \sum_{k\sigma} V_k (\hat{f}^\dagger_\sigma \hat{c}^{\phantom\dagger}_{k\sigma} + {\rm h.c.}) + U (\hat{n}_{f\uparrow}-1/2) (\hat{n}_{f\downarrow}-1/2).
\end{equation}
这个单杂质 Anderson 模型包含一个最多可被两个电子占据的杂质格点（算符 $\hat{f}^\dagger_\sigma$），具有在位排斥 $U$，并通过某种杂化函数 $V_k$ 与能散为 $\epsilon_k$ 的导带（算符 $\hat{c}^\dagger_{k\sigma}$）耦合。

为了使问题易于处理，人们从动量表象改用能量表象，假定只有低能各向同性的 $s$ 波散射是重要的，并引入对数离散化：能带被表示为一系列能带片段，其能量宽度在靠近费米能量 $\epsilon_F$ 处按指数方式减小。这对应于如下观察：量子杂质物理的决定性特征，即在局域杂质谱函数中费米能量处出现的极窄共振峰，与费米能量附近的态以指数强度的方式联系在一起。然而，要在技术层面上让 NRG 真正运转起来，对数离散化同样是必需的！

经过进一步的变换（对此我请读者参阅 \cite{Wilson75,Bulla08}），Anderson 哈密顿量最终被映射为一条除第一个格点外无相互作用的半无限链：
\begin{equation} 
\hat{H}= U (\hat{n}_{f\uparrow}-1/2) (\hat{n}_{f\downarrow}-1/2) + t_{-1} \sum_\sigma (\hat{f}^\dagger_\sigma \hat{d}_{0\sigma} + {\rm h.c.}) + \sum_{\sigma,n=0}^\infty t_n (\hat{d}^\dagger_{n\sigma} \hat{d}_{n+1,\sigma} + {\rm h.c.}),
\end{equation}
其中 $\hat{d}_{\sigma}$ 是费米子算符。关键之处在于 $t_n$ 指数衰减，$t_n \sim \Lambda^{-n}$，其中 $\Lambda$ 是对数离散化中能带的收缩因子，通常取 1.5 到 2 量级的值。这显然是一个我们的方法可以处理的模型，例如基态搜索——由于跃迁按指数衰减，我们无须考虑一条真正无限的链。

NRG 建立在如下观察之上：指数衰减导致能量尺度的分离——假定我们已知部分链到某一长度为止的谱，那么由于链深处指数小的能量尺度，其余所有格点只会对它做出指数小的修正。求基态（以及更一般地求低能激发谱）现在通过迭代精确对角化来实现：假定我们对链的某个左端部分有 $D$ 维的有效本征空间。那么更大一号的链的态空间维数是 $dD=4D$；为避免指数增长，我们必须截断回 $D$ 个态。NRG 的处方是对该系统对角化并保留 $D$ 个最低的本征态。从仍可精确对角化的很短的链出发，这一过程对最低能态的分辨按指数方式变好，其正当性由能量尺度分离保证：在某一步保留哪些态的决定，即便事后来看也不会被急剧改变，因为链上所有更远的格点都在小得多的能量上相互作用。所得到的各个能量尺度（链长）下的本征谱，随后可用来提取 RG 流信息，或计算杂质问题的热力学量或动力学量。

既然 MPS 的构建单元 $A^\sigma$ 可以被解读为在把块增长一个格点时对一次状态抽取步骤进行编码，而且与具体的状态抽取方案无关，那么立刻可以看出，NRG 与 DMRG 一样，也可以被视为在 MPS 上运作\cite{Weichselbaum09}。这闭环了一段历史：事实上，正是把 NRG 朴素地应用于海森堡模型和 Hubbard 模型时失败的分析，催生了 DMRG 的发展。一个 NRG 态看起来形如
\begin{equation}
\ket{a_\ell} = \sum_{\sigma_1,\ldots,\sigma_\ell} (A^{\sigma_1} \ldots  A^{\sigma_\ell})_{a_\ell} \ket{\sigma_1\ldots\sigma_\ell}.
\end{equation}
在每个长度 $\ell$ 处，我们都得到一个由 $D$ 个态构成的谱。

既然 DMRG 是在 MPS 拟设空间上做变分的，那么可以合理地预期，相对于 NRG 方法至少总能有一些改进。事实的确如此 \cite{Weichselbaum09}；下一节中，我将讨论一些已经牢固定型的改进，以及其他一些更具推测性的改进，即尚缺乏在相关复杂问题上进行基准测试的改进。

\subsection{超越数值重正化群}
事实上，考虑 NRG 的 MPS 表述即便不借助与 DMRG 这类变分方法的联系也有帮助，例如对某些求和规则的严格执行 \cite{Weichselbaum07,Peters06}，但这超出了本综述的范围。

然而我们还能做更多的事：把 NRG 生成的最终 MPS 构造施以类似 DMRG 的扫描。这会在一定程度上改进基态的质量，但首要的是，高能（短链）处的截断过程将获悉低能处的截断，反之亦然。与 NRG 不同，现在各能量尺度之间存在反馈。在这个意义上，把 NRG 用于杂质问题，其概念上的推进类似于无限系统 DMRG 为变分的有限系统 DMRG 所提供的热身过程。

在对数离散化下，能量尺度的分离足够大，以致这一效应是次要的；对于一个聚焦于 Abrikosov-Kondo-Suhl 共振的简单单杂质问题，最终的改进非常有限，因为 NRG 本来就是为最优地描述这一特征而设计的。要点在于，由于反馈的存在，能量尺度分离现在可以完全放弃，因而对数离散化也可以放弃，我们可以在物理上合适之处选择对能带更细粒度的分辨。这可能找到多方面的应用。

在一个应用中，人们对处于外场中的杂质问题做了 MPS 上的变分计算。外场导致共振峰分裂为两个依赖自旋的峰，分别移到费米能量之上和之下。在图 \ref{fig:NRGDMRG} 中我们考虑其中一个峰，使用了三种技术：NRG、一种解析方法\cite{Rosch03,Garst05} 以及变分 MPS（DMRG）计算。由于对数离散化，NRG 聚焦于 $\epsilon_F$，完全看不到这个依赖外场的峰。放松对数离散化，并在预期峰位附近、$\epsilon_F$ 之外提供足够细的能量区间，移动后的共振就可以被清晰地分辨出来，甚至与解析结果符合得非常好。

这项工作的第二个有趣应用，是在动力学平均场理论（DMFT）的框架中以之取代 NRG 作为杂质求解器 \cite{Georges96,Nishimoto04,Garcia04,Raas04,Raas05}。在那种情形下，需要费米能量处金属共振之外的信息，因此改进其他能量尺度上的谱分辨将是非常值得的。

\begin{figure}
\centering\includegraphics[height=200pt]{PyMPS_Spectral.eps}
\caption{外场中 Kondo 模型的杂质谱函数。细竖线标示能量区间：在某一能量以下，对数离散化转为线性。变分 MPS 计算（带与不带去卷积）揭示了 NRG 错过的一个峰，其峰位与 Rosch {\em et al.} \cite{Rosch03} 的计算加上微扰求得的峰移符合得非常好。取自 \cite{Weichselbaum09}。}
\label{fig:NRGDMRG}
\end{figure}

由于这条半无限链除第一个格点外没有相互作用，我们可以考虑把它展开为一条自旋 $1/2$ 的无限链：杂质位于中心，自旋向上或自旋向下费米子的有无对应两种自旋态，链的左半部分对应自旋向上的费米子，右半部分对应自旋向下的费米子 \cite{Raas04}。相近的能量现在不再被归在一起，但在类 DMRG 的做法中这已不再要紧！通过中心杂质才发生相互作用的自旋本质上可能并不纠缠，这一直觉得到了实际计算的证实。这一点很重要，因为它意味着我们不必为矩阵维数的增大付出沉重代价。恰恰相反：如果在 NRG 做法中我们看到的本质上是两条彼此平行的无耦合自旋链，这就意味着，若自旋链的维数为 $D$，则相应的 MPS 维数为 $O(D^2)$。因此我们可以预期，一个态数为 $D$ 的 NRG 计算可以被一个态数为 $O(\sqrt{D})$ 的更快的 DMRG 计算所取代。

除了这一加速之外，当杂质与多个能带耦合时（此时 NRG 呈指数复杂性\cite{Weichselbaum09}），当然也可以应用展开。中心格点当然保持不变，而它的数值处理可能变得极其昂贵，因此必须为这个格点设计新的策略。这方面还有大量工作要做，但围绕这些想法的首批有趣的后续研究已经出现\cite{Saberi08,Holzner10}。

\section{无限尺寸算法}
\label{sec:infinite}

\subsection{用 MPS 语言重新表述无限系统 DMRG}

在对有限系统算法的广泛讨论之后，现在让我们完全用 MPS 语言重新表述无限系统 DMRG。为顾及格点的迭代插入，方便起见对局域态稍作不同的标记；我们把各态称为 $\ket{\sigma_1^A} \ket{\sigma_2^A} \ldots \ket{\sigma_\ell^A}\ket{\sigma_\ell^B} \ldots \ket{\sigma_2^B}\ket{\sigma_1^B}$。此外，给矩阵 $A$ 和 $B$ 各配两个标记将会非常有用，
因为矩阵与其生成时所在格点之间的联系将会消失。

从大小为 0 的块出发（即一条 2 格点的链），基态波函数为 $\ket{\psi_1}=\sum_{\sigma_1^A \sigma_1^B} \Psi^{\sigma_1^A \sigma_1^B} \ket{\sigma_1^A} \ket{\sigma_1^B}$。把 $\Psi^{\sigma_1^A\sigma_1^B}$ 读作矩阵 $(\Psi^1)_{\sigma_1^A,\sigma_1^B}$，对它做奇异值分解得  $\Psi^1=U_1 \Lambda^{[1]} V^\dagger_1$。由此我们读出
\begin{equation}
A^{[1]\sigma_1^A}_{1,a_1} = (U_1)_{\sigma_1^A,a_1} \quad\quad B^{[1]\sigma_1^B}_{a_1,1} = (V^\dagger_1)_{a_1,\sigma_1^B} .
\end{equation} 
$A$ 和 $B$ 从 $U$ 和 $V^\dagger$ 继承左规范与右规范的性质，态取如下形式
\begin{equation}
\ket{\psi_1} = \sum_{\sigma_1^A \sigma_1^B} A^{[1]\sigma_1^A} \Lambda^{[1]} B^{[1]\sigma_1^B} \ket{\sigma_1^A \sigma_1^B} .
\end{equation}
现在插入两个格点并对 $\hat{H}$ 的能量求极小，我们得到
\begin{equation}
\ket{\psi_2} = \sum_{\sigma_1^A \sigma_2^A \sigma_2^B \sigma_1^B} A^{[1]\sigma_1^A} 
\Psi^{\sigma_2^A \sigma_2^B} 
B^{[1]\sigma_1^B} \ket{\sigma_1^A \sigma_2^A \sigma_2^B \sigma_1^B},
\end{equation}
其中每个 $\Psi^{\sigma_2^A \sigma_2^B}$ 都是维数与 $A$ 和 $B$ 相匹配的矩阵，隐含地假定矩阵乘积 $A\Psi B$。把这组 $\Psi$ 矩阵重塑为一个矩阵，
\begin{equation}
(\Psi_2)_{(a_1^A \sigma_2^A),(a_1^B \sigma_2^B)} =  (\Psi^{\sigma_2^A \sigma_2^B})_{a_1^A,a_1^B},
\end{equation}
SVD 给出 $\Psi_2 = U_2 \Lambda^{[2]} V^\dagger_2$，由此我们可以构造
\begin{equation}
A^{[2]\sigma_2^A}_{a_1^A,a_2^A} = U_{(a_1^A \sigma_2^A), a_2^A} \quad\quad  
B^{[2]\sigma_2^B}_{a_2^B,a_1^B} = V^\dagger_{(a_1^B \sigma_2^B), a_2^B}
\end{equation}
于是
\begin{equation}
\ket{\psi_2} = \sum_{\sigma_1^A \sigma_2^A \sigma_2^B \sigma_1^B} 
A^{[1]\sigma_1^A}A^{[2]\sigma_2^A} \Lambda^{[2]} B^{[2]\sigma_2^B}B^{[1]\sigma_1^B} 
\ket{\sigma_1^A \sigma_2^A \sigma_2^B \sigma_1^B}.
\end{equation}
在第 $\ell$ 步，波函数将成为
\begin{equation}
\ket{\psi_\ell} = \sum_{\sigma_1^A \ldots \sigma_\ell^A \sigma_\ell^B \ldots \sigma_1^B} 
A^{[1]\sigma_1^A} \ldots A^{[\ell]\sigma_\ell^A} \Lambda^{[\ell]} B^{[\ell]\sigma_\ell^B} \ldots B^{[1]\sigma_1^B}  
\ket{\sigma_1^A \ldots \sigma_\ell^A \sigma_\ell^B \ldots \sigma_1^B}
\end{equation}
其形状如图~\ref{fig:infiniteDMRG_AB} 所示。

\begin{figure}
\centering\includegraphics[scale=0.7]{PyMPS_InfiniteDMRG_MPS.eps}
\caption{无限系统 DMRG 在第 $\ell$ 步生成的 MPS 结构：一串左规范的 $A$、一串右规范的 $B$，由对角奇异值矩阵 $\Lambda^{[\ell]}$ 连接起来。注意从结构上看，中央单元并不重复。}
\label{fig:infiniteDMRG_AB}
\end{figure}

当然，一旦矩阵维数超过 $D$，我们在每一步都会丢弃最小的那些奇异值及其相应的奇异向量，这其实就是最初表述中基于密度矩阵的截断。
在每一步（新链长）我们都可以把该长度下的 $\hat{H}$ 写成一个 MPO 并做能量最小化。其他算符也能像有限尺寸情形那样找到类似的表示。 

让我简要过一遍这一算法在 $\Gamma\Lambda$ 记法下的重新表述。第一步我们只是把 $A$ 和 $B$ 改名为 $\Gamma$，与有限系统情形中边界格点的转换方式相一致：
\begin{equation}
\ket{\psi_1} =  \sum_{\sigma_1^A \sigma_1^B} \Gamma^{[1]\sigma_1^A} \Lambda^{[1]} \Gamma^{[1]\sigma_1^B} \ket{\sigma_1^A \sigma_1^B} .
\end{equation}
然后我们在下式中最小化 $\Psi^{\sigma_2^A \sigma_2^B}$：
\begin{equation}
\ket{\psi_2} = \sum_{\sigma_1^A \sigma_2^A \sigma_2^B \sigma_1^B} \Gamma^{[1]\sigma_1^A} 
\Psi^{\sigma_2^A \sigma_2^B} 
\Gamma^{[1]\sigma_1^B} \ket{\sigma_1^A \sigma_2^A \sigma_2^B \sigma_1^B},
\end{equation}
并像前面一样把它分解为 $A^{[2]\sigma_2^A} \Lambda^{[2]} B^{[2]\sigma_2^B}$。现在我们定义（由于标记方式的原因，$B$ 矩阵与有限系统设置相比略有变化）
\begin{equation}
\Lambda^{[1]}_a \Gamma^{\sigma_2^A}_{ab} = A^{[2]\sigma_2^A}_{ab} \quad\quad
\Gamma^{\sigma_2^B}_{ab} \Lambda^{[1]}_b = B^{[2]\sigma_2^B}_{ab}
\end{equation}
从而得到
\begin{equation}
\ket{\psi_2} =  \sum_{\sigma_1^A \sigma_2^A \sigma_1^B \sigma_2^B} \Gamma^{[1]\sigma_1^A} \Lambda^{[1]} \Gamma^{[2]\sigma_2^A} \Lambda^{[2]} \Gamma^{[2]\sigma_2^B} \Lambda^{[1]} \Gamma^{[1]\sigma_1^B}   \ket{\sigma_1^A \sigma_2^A \sigma_2^B \sigma_1^B} ,
\end{equation}
如图~\ref{fig:infiniteDMRG_GL} 所示。 

我们现在可以问两个问题：(i) 在 DMRG 中，用 Lanczos 这类迭代求解器求出基态是最耗时的部分。能否通过提供一个好的初始猜测来实现加速？在有限系统 DMRG 中，MPS 表述自动给出了 White 的预测方法；而对无限系统 DMRG 的加速尝试过去也有人做过，但成败参半\cite{Schollwoeck98,Sun02}。(ii) 我们能否利用链中心的信息，在热力学极限下构建一个（周期至多为 2 的）平移不变态，求出它的基态或让它做时间演化？

两个问题的答案都是肯定的，而这建立在识别出态中两格点重复结构的基础上。  与 $A,B$ 记法相反——那里即使在热力学极限下中央单元也不自我重复——在 $\Gamma\Lambda$ 记法中很容易读出一个两格点重复单元，即
\begin{equation}
\Lambda^{[\ell-1]} \Gamma^{[\ell]\sigma_\ell^A} \Lambda^{[\ell]} \Gamma^{[\ell]\sigma_\ell^B} .
\end{equation}
利用转换规则可以把它翻译成 $A$、$B$ 的语言：
\begin{equation}
A^{[\ell]\sigma_\ell^A} \Lambda^{[\ell]} B^{[\ell]\sigma_\ell^B} (\Lambda^{[\ell-1]})^{-1} .
\end{equation}
这个结果也可以直接从 $A$、$B$ 记法得到，但论证要比在 $\Gamma\Lambda$ 记法中更复杂。当然应当理解：重复这些态片段并不会生成本来从中取出它们的那个态；这里只是断言，在热力学极限 $\ell\rightarrow\infty$ 下，当所有格点都彼此等价时，这种重复可以相当接近。无论如何，它们是对态的样貌的一种有依据的猜测！ 

\begin{figure}
\centering\includegraphics[scale=0.7]{PyMPS_InfiniteDMRG_TEBD.eps}
\caption{在 $\Gamma\Lambda$ 记法下无限系统 DMRG 第 $\ell$ 步生成的 MPS 结构：$\Gamma$ 与 $\Lambda$ 矩阵交替出现，并给出（重复的）两格点砌块的一种可能认定。}
\label{fig:infiniteDMRG_GL}
\end{figure}

我们现在要让这个态片段派上多重用场：先用于由无限系统 DMRG 生成的有限系统，再用于热力学极限态，两者都涵盖基态搜索与时间演化。在前一种情形，它将为下一个量子态提供一个极其高效的猜测；在此态上计算观测量则与其他有限系统中的做法完全相同。
在后一种情形，可以构建基态与时间演化两种算法（iDMRG 与 iTEBD），但二者都需要对热力学极限下态的正交归一性问题做（相同的）分析。

\subsection{无限系统 DMRG 中的态预测} 
识别出态的``单元胞''，使我们能够为无限系统 DMRG 定义一个好的初始猜测\cite{McCulloch08}，它避免了前人遇到的所有问题，并且即使对短链也带来戏剧性的加速——尽管在这些短链上，``链中心可以代表热力学极限态的物理''这一隐含假设肯定是错的：为使链增长，我们只需插入一个单元胞，即便对短链来说``态不过是这些单元胞的重复''这一想法并未得到很好的验证——但即便如此，也要远胜于随机猜测。从
\begin{equation}
\ket{\psi_{\ell}} =  \sum_{\fat{\sigma}} A^{[1]\sigma_1^A} \ldots A^{[\ell-1]\sigma_{\ell-1}^A} (A^{[\ell]\sigma_\ell^A} \Lambda^{[\ell]} B^{[\ell]\sigma_\ell^B} [\Lambda^{[\ell-1]}]^{-1}) \Lambda^{[\ell-1]} B^{[\ell-1]\sigma_{\ell-1}^B} \ldots B^{[1]\sigma_1^B} \ket{\fat{\sigma}},
\end{equation} 
其中重复单元已被括出，该猜测便为
\begin{eqnarray}
\ket{\psi_{\ell+1}^{{\rm guess}}} &=& \sum_{\fat{\sigma}} A^{[1]\sigma_1^A} \ldots  A^{[\ell-1]\sigma_{\ell-1}^A} \times \nonumber \\
& & ( A^{[\ell]\sigma_\ell^A} \Lambda^{[\ell]} B^{[\ell]\sigma_{\ell+1}^A} [\Lambda^{[\ell-1]}]^{-1} ) (A^{[\ell]\sigma_{\ell+1}^B} \Lambda^{[\ell]} B^{[\ell]\sigma_\ell^B}  
[\Lambda^{[\ell-1]}]^{-1} ) \times \nonumber \\
& &  \Lambda^{[\ell-1]} B^{[\ell-1]\sigma_{\ell-1}^B} \ldots B^{[1]\sigma_1^B} \ket{\fat{\sigma}} 
\end{eqnarray}
或者展开乘积，得 
\begin{equation}
\ket{\psi_{\ell+1}^{{\rm guess}}} = \sum_{\fat{\sigma}} A^{[1]\sigma_1^A} \ldots A^{[\ell]\sigma_\ell^A} \Lambda^{[\ell]} B^{[\ell]\sigma_{\ell+1}^A} [\Lambda^{[\ell-1]}]^{-1} A^{[\ell]\sigma_{\ell+1}^B} \Lambda^{[\ell]} B^{[\ell]\sigma_\ell^B} B^{[\ell-1]\sigma_{\ell-1}^B} \ldots B^{[1]\sigma_1^B} \ket{\fat{\sigma}} .
\end{equation}
在这个拟设中，我们现在可以把对 $\Psi^{\sigma_{\ell+1}^A\sigma_{\ell+1}^B}$ 的一个猜测认定为
\begin{equation}
\Psi^{\sigma_{\ell+1}^A\sigma_{\ell+1}^B}_{{\rm guess}} =
\Lambda^{[\ell]} B^{[\ell]\sigma_{\ell+1}^A} [\Lambda^{[\ell-1]}]^{-1} A^{[\ell]\sigma_{\ell+1}^B} \Lambda^{[\ell]} .
\label{eq:guess2sites}
\end{equation}
从这个拟设出发，接下来就可以在无限系统 DMRG 框架内迭代地求出使能量最小的 $\Psi^{\sigma_{\ell+1}^A\sigma_{\ell+1}^B}$，并由它生成 $A^{[\ell+1]\sigma_{\ell+1}^A}$、$\Lambda^{[\ell+1]}$ 和 $B^{[\ell+1]\sigma_{\ell+1}^B}$。

另外，这个拟设也可以化成更优雅的形式。目前，$B$ 矩阵出现在晶格的 A 侧，反之亦然。但我们可以利用把 MPS 化为规范形式的能力，用 SVD 把 $\Lambda^{[\ell]} B^{[\ell]\sigma_{\ell+1}^A}$ 规范化为
$A^{[\ell+1]\sigma_{\ell+1}^A} \Lambda_R^{[\ell]}$，其中 $A$ 由 $U$ 得出，$\Lambda$ 由 $DV^\dagger$ 得出，方法如前所述（$\Lambda_a B^\sigma_{ab} = M_{a\sigma,b} = \sum_k U_{a\sigma,k} D_k (V^\dagger)_{k,b} = A^\sigma_{ak} \Lambda_{kb}$）。类似地，我们从右侧对  $A^{[\ell]\sigma_{\ell+1}^B} \Lambda^{[\ell]}$ 做规范化，得到 $\Lambda_L^{[\ell]} B^{[\ell+1]\sigma_{\ell+1}^B}$，其中 $B$ 来自 $V^\dagger$。于是我们有拟设
\begin{equation}
\ket{\psi_{\ell+1}^{{\rm guess}}} = \sum_{\fat{\sigma}} A^{[1]\sigma_1} \ldots  A^{[\ell+1]\sigma_{\ell+1}} \Lambda^{[\ell+1]}_{{\rm guess}} B^{[\ell+1]\sigma_{\ell+1}}  \ldots B^{[1]\sigma_1} \ket{\fat{\sigma}} ,
\end{equation}
其中 
\begin{equation}
\Lambda^{[\ell+1]}_{{\rm guess}} =  \Lambda_R^{[\ell]} [\Lambda^{[\ell-1]}]^{-1} \Lambda^{[\ell]}_L .
\end{equation}
从这个拟设出发，接下来就可以迭代地求出使能量最小的 $\Lambda^{[\ell+1]}$，只需对单格点变分 MPS 的最小化部分稍作修改。一般而言，所得结果不会具有由 SVD 产生的对角形式，因为 $\Lambda_R$ 和 $\Lambda_L$ 本来就不是对角的。但对它做一次 SVD 会得到两个幺正矩阵，可以把它们吸收进相邻的 $A$ 和 $B$ 而不改变其归一化性质，使得最终的 $\Lambda^{[\ell+1]}$ 是对角的。在这种形式下，算法可表示为图~\ref{fig:infinitebuildup}。

\begin{figure}
\centering\includegraphics[scale=0.7]{PyMPS_InfiniteSystemDMRG.eps}
\caption{包含预测步骤的无限系统算法的概括表示。每一步引入两个新格点，用上一次迭代算出的一个拟设（白色），从而得到一个新的奇异值分解（黑色）。}
\label{fig:infinitebuildup}
\end{figure}


正如 McCulloch 所证明的\cite{McCulloch08}，这一预测使无限系统 DMRG 大幅提速，并与有限系统 DMRG 的预测算法相得益彰：预测态与计算所得态之间的重叠常常接近于 1，误差约在 $10^{-10}$ 的量级！

\subsection{iDMRG：热力学极限下的变分基态搜索}
利用前几节的思想，现在可以非常简单地把无限系统 DMRG 变成一个高效而精确的算法，称为 {\em iDMRG}，指的是 DMRG 的热力学极限版本：
\begin{itemize}
\item 搭建一个无限系统 DMRG 算法。
\item 把上一节的预测步骤加进最小化算法。
\item 运行修改后的算法直至达到收敛。波函数
是否收敛到热力学极限，可以通过考察关系式 式~(\ref{eq:rhoarecursion}) 来判断
\begin{equation}
\sum_{\sigma_\ell} A^{[\ell]\sigma_\ell} \rho_A^{[\ell]} A^{[\ell]\sigma_\ell\dagger} = \rho_A^{[\ell-1]} ,
\end{equation}
其中 $\rho_A^{[\ell]} = \Lambda^{[\ell]}\Lambda^{[\ell]\dagger}$ 和 $\rho_A^{[\ell-1]} = \Lambda^{[\ell-1]}\Lambda^{[\ell-1]\dagger}$ 是系统左半部分的约化密度算符。注意，这一关系在同一个有限系统内的约化密度算符之间总是成立，但这里的算符来自长度为 $2(\ell-1)$ 和 $2\ell$ 的两个不同系统，因此预期这一关系只会在 $\ell\rightarrow\infty$ 时作为不动点关系成立。按照与为更大系统构造拟设时相同的论证，我们可以把 $A^{[\ell]\sigma_\ell}\Lambda^{[\ell]}$ 变换为 $\Lambda^{[\ell]}_L B^{[\ell+1]\sigma_\ell}$。于是，利用右归一化，不动点关系的左端可简化为 $\Lambda^{[\ell]}_L\Lambda^{[\ell]\dagger}_L \equiv \hat{\rho}^{[\ell]}_L$，关系式变为 $\hat{\rho}^{[\ell]}_L = \rho_A^{[\ell-1]}$。如果这一关系以高精度成立，就表明已到达热力学不动点。  
衡量两个密度算符彼此接近程度的一种方式由保真度给出 \cite{McCulloch08}
\begin{equation}
F(\hat{\rho}^{[\ell]}_L, \hat{\rho}_A^{[\ell-1]}) = \tr \sqrt{\sqrt{\hat{\rho}^{[\ell]}_L}\hat{\rho}_A^{[\ell-1]} \sqrt{\hat{\rho}^{[\ell]}_L}} .
\end{equation}
代入定义并利用迹的循环性质，可以证明 $F(\hat{\rho}^{[\ell]}_L, \hat{\rho}^{[\ell-1]}_A) = \sum_i s_i$，其中 $s_i$ 是  $\Lambda^{[\ell]}_L \Lambda^{[\ell-1]\dagger}$ 的奇异值。 
\end{itemize}
当然，这个算法可以随时停下，但那样我们得到的只是有限系统结果，它肯定可以通过有限系统 DMRG 加以改进。上面给出的收敛判据确实使我们能够得到热力学极限态，我们可以把它形式上写为
\begin{equation}
\ket{\psi} = \sum_{\fat{\sigma}} \ldots A^{[\ell]\sigma_i}\Lambda^{[\ell]} B^{[\ell]\sigma_{i+1}} [\Lambda^{[\ell-1]}]^{-1}A^{[\ell]\sigma_{i+2}}\Lambda^{[\ell]} B^{[\ell]\sigma_{i+3}} [\Lambda^{[\ell-1]}]^{-1}A^{[\ell]\sigma_{i+4}}\Lambda^{[\ell]} B^{[\ell]\sigma_{i+5}} [\Lambda^{[\ell-1]}]^{-1} \ldots 
 \ket{\fat{\sigma}} ,
\end{equation}
其中我们把 $\ell$ 取为收敛判据得到满足的那次迭代步。现在的问题是如何计算期望值。显然，我们无法像前面那样写出一个有限的网络收缩；这个网络将是无限大的，因此无法天真地收缩。收缩之所以可行，只有在我们能把收缩的数目约减为有限多个时才行。对于有限系统网络，我们看到左正交归一与右正交归一使我们能够消去大部分收缩：对于格点 $i$ 和 $j$ 上的观测量，只需考虑这两个格点上以及二者之间的收缩。然而，热力学极限态没有任何理由必须满足归一化判据；事实上，通常它并不满足。因此我们需要一个正交归一化手续。此后，期望值就可以像对具有左、右正交归一性的有限晶格那样来计算。由于这一手续对下一个算法也很重要，且在概念上稍微深入一些，我把它留到单独的一节。 

此处应当指出，iDMRG 可以与更早的、以 PWFRG（乘积波函数重正化群，product wave function renormalization group）为名的算法途径联系起来 \cite{NishinoOkunishi95,Hieida97,Ueda06}，它们已经包含了上述思想的一部分；iDMRG 则把这些思想推进到其自然的完成形态 \cite{Ueda08,Ueda10}。

\subsection{iTEBD：热力学极限中的时间演化}
在本节中，我改用 $\Gamma\Lambda$ 记法，尽管这些公式可以很容易地转写成 $A,B$ 形式。利用我们的态片段 $\Lambda^{[\ell-1]} \Gamma^{[\ell]\sigma_\ell^A} \Lambda^{[\ell]} \Gamma^{[\ell]\sigma_\ell^B}$，我们可以搭建一条无限链
\begin{equation}
\ket{\psi} = \sum_{\fat{\sigma}} \ldots (\Lambda^A \Gamma^{A \sigma_i} \Lambda^B \Gamma^{B \sigma_{i+1}})(\Lambda^A \Gamma^{A \sigma_{i+2}} \Lambda^B \Gamma^{B \sigma_{i+3}}) \ldots \ket{\fat{\sigma}} ,
\end{equation}
与上一节中一样，其中 $\Gamma^{A\sigma}= \Gamma^{[\ell]\sigma_\ell^A}$，
$\Gamma^{B\sigma}= \Gamma^{[\ell]\sigma_\ell^B}$、$\Lambda^A=\Lambda^{[\ell-1]}$ 且
$\Lambda^B=\Lambda^{[\ell]}$。该片段可以是已收敛的基态计算的结果，也可以来自某个能够精确构造的简单起始态。

现在我们可以按 Trotter 形式写下时间演化：先对所有奇数键（我将其称为 AB）施加一个无穷小时间步，然后再对所有偶数键（我将其称为 BA）施加。键演化算符与 tMPS/tDMRG/TEBD 情形中的完全相同，我将把它们记为 $U_{AB} = \sum_{\sigma_A\sigma_B\sigma'_A\sigma'_B} U_{AB}^ {\sigma_A\sigma_B, \sigma'_A\sigma'_B} \ket{\sigma_A\sigma_B} \bra{\sigma'_A\sigma'_B}$，类似地还有 $U_{BA}$。

正如我们已经看到的 [参见式~(\ref{eq:guess2sites})]，一个完整的两格点片段由五个矩阵的乘积构成，
$\Lambda^A \Gamma^{A \sigma_A} \Lambda^B \Gamma^{B \sigma_{B}} \Lambda^A$。于是在键 AB 上的时间演化给出一组矩阵
\begin{equation}
M^{\sigma_A\sigma_{B}} = \sum_{\sigma'_A\sigma'_{B}} U_{AB}^ {\sigma_A\sigma_{B}, \sigma'_A\sigma'_{B}}  \Lambda^A \Gamma^{A \sigma'_A} \Lambda^B \Gamma^{B \sigma'_{B}} \Lambda^A .
\end{equation}
经过如今已是标准做法的重整形，我们通过 SVD 得到
\begin{equation}
M^{\sigma_A\sigma_{B}}=A^{\sigma_A} \Lambda^B B^{\sigma_{B}} ,
\end{equation}
其中新的 $\Lambda^B$（相应地还有 $A^{\sigma_A}$ 和 $B^{\sigma_{B}}$）像在 tDMRG 或 TEBD 中那样被截断，用以替换旧的。利用 $\Lambda^A$（仍取自上一轮迭代），我们可以定义新的 $\Gamma^{A \sigma_A}$ 和 $\Gamma^{B \sigma_{B}} $（通过 $A^{\sigma_A} = \Lambda^A \Gamma^{A \sigma_A}$ 及
$B^{\sigma_{B}} = \Gamma^{B \sigma_{B}} \Lambda^A$）。这就定义了一个新的“单胞”。

如果把它写下来并在后面再接上一个，我们可以读出两个 AB 单胞中央处的键 BA，即
$ \Lambda^B \Gamma^{B \sigma_{B}} \Lambda^A \Gamma^{A \sigma_{A}} \Lambda^B$，对它作时间演化得到
\begin{equation}
N^{\sigma_{B}\sigma_{A}} = \sum_{\sigma'_{B}\sigma'_{A}} U_{BA}^ {\sigma_{B}\sigma_{A}, \sigma'_{B}\sigma'_{A}}  \Lambda^B \Gamma^{B \sigma'_{B}} \Lambda^A \Gamma^{A \sigma'_{A}} \Lambda^B .
\end{equation}
重整形并作 SVD 给出
\begin{equation}
N^{\sigma_{B}\sigma_{A}}=A^{\sigma_{B}} \Lambda^A B^{\sigma_{A}} ,
\end{equation}
其中 $\Lambda^A$（相应地还有其他矩阵）再次被截断并替换旧的。利用 $\Lambda^B$（仍取自上一轮迭代），我们可以定义新的 $\Gamma^{B \sigma_{B}}$ 和 $\Gamma^{A \sigma_{A}} $（通过 $B^{\sigma_{A}} =   \Gamma^{A \sigma_{A}} \Lambda^B$ 及 $A^{\sigma_{B}} =\Lambda^B \Gamma^{B \sigma_{B}}$）。除以小奇异值带来的麻烦可以用针对 TEBD 已经讨论过的简单修改来避免 \cite{Hastings09}。

因此，通过依次施加无穷小的键演化算符，我们可以在热力学极限下建立实时或虚时演化。这一算法被称为 {\em iTEBD}，因为它给出了 TEBD 的无限尺寸推广 \cite{Vidal07}。

正交归一的问题再次出现 \cite{Orus08}。假设初始态满足正交归一判据。纯实时演化会生成一串作用于该态的幺正算符，这保持正交归一性质。但每个时间步之后不可避免的截断会破坏这一性质，哪怕截断可能很小。为了把这一方法变成可行的算法，我们必须像 iDMRG 那样，在每一步之后处理热力学极限下的正交化问题。

在此我想提到，McCulloch\cite{McCulloch08} 已经证明，只要把中心键上的极小化替换为中心键上的时间演化，就可以把 iDMRG 变成 iTEBD，这相对原始 iTEBD 算法有一些概念上的优点。

最后用几句话谈谈外推，以此结束本节。在有限系统 DMRG（或 MPS）中，惯常做法是对每个有限长度 $L$ 先在 $D$ 上外推结果以使精度最大化，然后再把这些结果在 $L$ 上外推到热力学极限。而在这里，我们直接在热力学极限下工作（假设 iDMRG 已被带到收敛），剩下的只是对 $D$ 的外推。有趣的是，在临界点处，这种对 $D$ 的外推与纠缠熵标度及临界指数有着深刻而有用的联系 \cite{Tagliacozzo08}。实际上，有限的矩阵维数在临界系统中引入了一个有限的关联长度，这与有限系统 DMRG 的情形颇为相似：在那里，MPS 特有的指数叠加能够恰当地模拟幂律的长度标度也随 $D$ 的某个幂次标度 \cite{Andersson99}。

\subsection{热力学极限态的正交化}

在有限系统上做 iDMRG 时，$A$ 矩阵和 $B$ 矩阵保持左规范和右规范；这意味着左块态与右块态彼此正交，如前所示。我们将稍微不按惯例地把具有这种性质的态称为 {\em 正交}态。正如上一节所见，我们可以使用片段
\begin{equation}
A^{[\ell]\sigma_\ell^A}\Lambda^{[\ell]} B^{[\ell]\sigma_\ell^B} (\Lambda^{[\ell-1]})^{-1}
\end{equation}
我们可以重复它来搭建一条无限长链，
\begin{equation}
\ket{\psi} = \sum_{\fat{\sigma}} \ldots A^{\sigma_i}\Lambda^{[\ell]} B^{\sigma_{i+1}} (\Lambda^{[\ell-1]})^{-1}A^{\sigma_{i+2}}\Lambda^{[\ell]} B^{\sigma_{i+3}} (\Lambda^{[\ell-1]})^{-1}A^{\sigma_{i+4}}\Lambda^{[\ell]} B^{\sigma_{i+5}} (\Lambda^{[\ell-1]})^{-1} \ldots \ket{\fat{\sigma}} ,
\end{equation}
这里我简化了 $A$、$B$ 的记号。这些态的问题在于，对任意分成两块的二分，左块和右块上的态一般来说{\em 并不正交}：如果我们按上述方式把 $\Lambda^{[\ell]} B$ 变换为 $\tilde{A} \Lambda_R^{[\ell]}$，链将变成
\begin{equation}
\ket{\psi} = \sum_{\fat{\sigma}} \ldots A^{\sigma_1}\tilde{A}^{\sigma_{2}} P
A^{\sigma_{3}} \tilde{A}^{\sigma_{4}} P A^{\sigma_{5}}\tilde{A}^{\sigma_{6}} P \ldots \ket{\fat{\sigma}} ,
\end{equation}
其中 $P= \Lambda_R^{[\ell]} (\Lambda^{[\ell-1]})^{-1}$。如果把 $P$ 吸收到它左边的 $\tilde{A}$ 中，即 $\tilde{A}P \rightarrow \overline{A}$，归一化条件就变成
\begin{equation}
\sum_\sigma \overline{A}^{\sigma\dagger} \overline{A}^{\sigma} = P^\dagger \sum_\sigma \tilde{A}^{\sigma\dagger} \tilde{A}^\sigma P = P^\dagger P.
\end{equation}
然而一般而言 $P^\dagger P \neq I$。这不仅出现在 $\ell$ 较小、我们离无限系统不动点还很远的情形；只要丢弃态权重有限——在 DMRG 计算中通常就是如此，即使它非常小——在不动点处也是如此。

正如 Orus 和 Vidal\cite{Orus08} 所指出的（这里沿用 \cite{McCulloch08} 的表述），检测正交归一的一个条件是检验单位算符在两个块态 $\ket{a}$、$\ket{a'}$ 之间的期望值是否为 $\delta_{a,a'}$（见 图~\ref{fig:infiniteorthogonalization}）。考虑像有限系统那样的期望值收缩，并假设我们已从左边的 $-\infty$ 一路收缩到格点 $0$。结果将是一个矩阵状的对象 $C_{a'a}$，对应于开放的腿。现在考虑操作 $E_L(C)$，它把收缩再向前推进两步，即跨过格点 1 和 2。这个{\em 转移算符}为 [参见第~\ref{subsubsec:transferoperator}~节]
\begin{equation}
E_L(C) = \sum_{\sigma_1\sigma_2} P^\dagger \tilde{A}^{\sigma_2\dagger} A^{\sigma_1\dagger} C A^{\sigma_1} \tilde{A}^{\sigma_2} P.
\end{equation}
对于正交归一的态，我们希望 $E_L(I) = I$，这正是单位算符对正交归一块态产生的期望值矩阵。然而，利用左规范条件，我们实际得到的是 $E_L(I)=P^\dagger P$。当系统延伸至无穷时，$E_L(I) = I$ 必须与最大本征值相联系；整个态的可归一化意味着最大本征值必须为 1。$E_L$ 的“二次”形式意味着与之相应的本征矩阵 $V_L$ 是厄米的且非负。借助本征值分解或奇异值分解，可以把它分解为 $V_L = X^\dagger X$，其中 $X$ 可逆。我们可以在每个 $P$ 之后插入 $X^{-1}X$，使单胞变成 $X A^{\sigma_1} \tilde{A}^{\sigma_2} P X^{-1}$，而新的转移算符为
\begin{equation}
E'_L(C) = \sum_{\sigma_1\sigma_2} X^{-1\dagger}P^\dagger \tilde{A}^{\sigma_2\dagger} A^{\sigma_1\dagger} X^\dagger C X A^{\sigma_1} \tilde{A}^{\sigma_2} P X^{-1}.
\end{equation}
于是
\begin{equation}
E'_L(I) = X^{-1\dagger} V_L X^{-1} = I
\end{equation}
这由 $V_L$ 相对于 $E_L$ 的本征矩阵性质得出。
（如果 $E_L$ 的最大本征值碰巧不为 1，则须对 $X$ 作适当缩放。）

把 $P$ 的定义代入 $X A^{\sigma_1}\tilde{A}^{\sigma_2}P X^{-1}$，撤销到 $\tilde{A}$ 的变换，并作变换 $A^{\sigma_1} \Lambda^{[\ell]} \rightarrow \Lambda^{[\ell]}_L \tilde{B}^{\sigma_1}$，单胞变成 $X \Lambda^{[\ell]}_L \tilde{B}^{\sigma_1} B^{\sigma_2} (\Lambda^{[\ell-1]})^{-1} X^{-1}$。把单胞的起点移位后，它变成 $(\Lambda^{[\ell-1]})^{-1} X^{-1} X \Lambda^{[\ell]}_L \tilde{B}^{\sigma_1} B^{\sigma_2} = Q \tilde{B}^{\sigma_1} B^{\sigma_2}$，其中 $Q = (\Lambda^{[\ell-1]})^{-1} X^{-1} X \Lambda_L^{[\ell]} = (\Lambda^{[\ell-1]})^{-1} \Lambda_L^{[\ell]}$，与 $X$ 无关。从右边计算收缩得到转移算符
\begin{equation}
E_R(C) = \sum_{\sigma_1\sigma_2} Q \tilde{B}^{\sigma_1} B^{\sigma_2} C B^{\sigma_2\dagger} \tilde{B}^{\sigma_1\dagger} Q^\dagger .
\end{equation}
与前面相同的本征值/本征矩阵论证给出主导本征矩阵 $V_R = YY^\dagger$（$Y$ 可逆）以及单胞 $Y^{-1}Q \tilde{B}^{\sigma_1} B^{\sigma_2} Y$。
这又给出 $E'_R(C)$，且 $E'_R(I)=I$。

如果把 $Q$ 的定义代入单胞，从 $\tilde{B}^{\sigma_1}$ 回到 $A^{\sigma_1}$，并把 $Q$ 显式写出，单胞即为
$Y^{-1}  (\Lambda^{[\ell-1]})^{-1} X^{-1} X A^{\sigma_1} \Lambda^{[\ell]} B^{\sigma_2} Y$，移动其原点后我们得到
\begin{equation}
X A^{\sigma_1} \Lambda^{[\ell]} B^{\sigma_2} Y Y^{-1}  (\Lambda^{[\ell-1]})^{-1} X^{-1} ,
\end{equation}
通过令
$A^{\sigma_1} \leftarrow XA^{\sigma_1}$、$B^{\sigma_2} \leftarrow B^{\sigma_2}Y$ 及 $\Lambda^{[\ell-1]} \leftarrow X\Lambda^{[\ell-1]}Y$，可以把它带回到单胞的原始形式，但现在{\em 恰当地保证了左规范和右规范。}

更精确地说，新的单胞给出 $E_L(I)=I$ 和 $E_R(I)=I$。但要注意，$E_L$ 和 $E_R$ 是由略有移位的单胞构造的：如图所示，$E_L$ 用的是 $A^{\sigma_1} \Lambda^{[\ell]} B^{\sigma_2} (\Lambda^{[\ell-1]})^{-1}$，而 $E_R$ 用的是 $(\Lambda^{[\ell-1]})^{-1} A^{\sigma_1} \Lambda^{[\ell]} B^{\sigma_2}$。我们可以把第一个和第二个单胞合并成左规范和右规范的两格点矩阵 $A^{\sigma_1\sigma_2}$ 与 $B^{\sigma_1\sigma_2}$。然后可以用标准方式把它们分解为左规范和右规范的矩阵，得到 $A^{[1]\sigma_1}A^{[2]\sigma_2}$ 和 $B^{[1]\sigma_1}B^{[2]\sigma_2}$。注意，这些矩阵当然不同于我们最初所拥有的那些。

\begin{figure}
\centering\includegraphics[scale=0.7]{PyMPS_InfiniteSystemOrthogonalization.eps}
\caption{若矩阵得到恰当的归一化，热力学极限拟设会同时生成左规范条件和右规范条件。注意，两个转移算符定义在热力学极限态的两个不同的两格点单胞上。}
\label{fig:infiniteorthogonalization}
\end{figure}

热力学极限态现在可以写为
\begin{equation}
\ket{\psi} = \sum_{\fat{\sigma}} \ldots A^{[1]\sigma_1} A^{[2]\sigma_2} A^{[1]\sigma_3} A^{[2]\sigma_4} \ldots \ket{\fat{\sigma}}
\end{equation}
或者类似地用 $B^{[1,2]}$ 写出，分别对应左规范和右规范形式。对计算期望值而言特别有意义的是混合规范形式，即左边为 $A$ 矩阵、右边为 $B$ 矩阵。如果考察其背后的两个单胞，就会看到在边界处，两者都想并入同一个 $(\Lambda^{[\ell-1]})^{-1}$ 来生成 $A$ 和 $B$ 矩阵。这个问题可以通过在出问题的 $(\Lambda^{[\ell-1]})^{-1}$ 之后插入恒等式 $I = \Lambda^{[\ell-1]} (\Lambda^{[\ell-1]})^{-1}$ 来解决。于是我们可以立即写下混合规范形式
\begin{equation}
\ket{\psi} = \sum_{\fat{\sigma}} \ldots A^{[1]\sigma_1} A^{[2]\sigma_2} A^{[1]\sigma_3} A^{[2]\sigma_4} \Lambda^{[\ell-1]}
B^{[1]\sigma_5} B^{[2]\sigma_6} B^{[1]\sigma_7} B^{[2]\sigma_8}
\ldots \ket{\fat{\sigma}} .
\label{eq:infinitemixedcanonical}
\end{equation}
\subsubsection{热力学极限中期望值的计算}
在有限系统中，我们已经看到如何通过传递一个对象 $C^{[i]}$ 来计算期望值：它以哑标量 $C^{[0]}=1$ 出发，借助转移算符 $E^{[i]}_O$ 从 $C^{[i]}$ 传递到 $C^{[i+1]}$，其中 $hat{O}$ 是期望值结构中局域作用的算符（大多是恒等算符，除非——比如说——在考察两点关联子时作用在两个格点上）。在恒等算符的情形，对左规范矩阵而言，从左到右的转移算符把恒等算符映射为恒等算符；类似地，若由右规范矩阵构成，从右到左的转移算符也把恒等算符映射为恒等算符。

上一节中已经在热力学极限态及其两个两格点转移算符 $E_L$ 和 $E_R$ 中发现了同样的结构。这使得旧有结果可以直接移植过来。假设我们想要计算 $\bra{\psi} \hat{O}_1 \hat{O}_i \ket{\psi}$；那么我们把该态化成式~(\ref{eq:infinitemixedcanonical}) 的混合规范形式，即 $A$ 矩阵一直排到格点 $i$ 或 $i+1$（取决于奇偶结构），然后是 $\Lambda^{[\ell-1]} $，再往后是直到无穷的 $B$ 矩阵。从左边收缩掉所有直到 0 的 $A$ 矩阵、从右边收缩掉所有 $B$ 矩阵之后，我们就得到如图~\ref{fig:finitecontractionremainder} 所示的剩余有限网络。
这个期望值随后可以像有限网络那样求值；如果我们使用 $C^{[i]}$ 与转移算符的记法，则从 $C^{[0]}=I$ 出发。区别在于，在最后 $\Lambda^{[\ell-1]}$ 会出现在收缩之中，而最终约化为标量则由收尾的 $\delta_{aa'}$ 线完成（它可以被理解为求迹）：假设最后一个格点是 $i$，那么最终的期望值在最后一步由下式给出
\begin{equation}
\bra{\psi} O^1 \otimes \ldots O^i \ket{\psi} = \tr  \Lambda^{[\ell-1]\dagger} C^{[i]} \Lambda^{[\ell-1]} = 
\tr   \Lambda^{[\ell-1]} \Lambda^{[\ell-1]\dagger} C^{[i]} = \tr \rho_A C^{[i]} ,
\end{equation}
其中我们用到了 $\Lambda$ 矩阵与规范形式下约化密度算符之间的关系。

\begin{figure}
\centering\includegraphics[scale=0.7]{PyMPS_FiniteContractionRemainder.eps}
\caption{利用左归一化和右归一化，无限收缩的求值可以约化为一个有限收缩，从而可以像处理任何有限 MPS 那样处理它。}
\label{fig:finitecontractionremainder}
\end{figure}
    
这一计算可以很容易地约化到两个态的重叠的情形，后者不过是单位算符在这两个态之间的矩阵元。假设两个态都处于平移不变形式，例如分别使用左归一化矩阵 $A^{[1]}$ 和 $A^{[2]}$（以及 $\tilde{A}^{[1]}$、$\tilde{A}^{[2]}$）。现在我们用 $E_L$ 把一个无限重叠计算向右推进两个格点（比如 1 和 2）：如果当前的重叠矩阵是 $C$，则它按如下方式推进：
\begin{equation}
E_L (C) = \sum_{\sigma_1\sigma_2} \tilde{A}^{[2]\sigma_2\dagger} \tilde{A}^{[1]\sigma_1\dagger} C A^{[1[\sigma_1} A^{[2]\sigma_2} .
\end{equation}  
如果把 $C$ 按 $E_L$ 的本征矩阵分解，那么在热力学极限下只有最大本征值的贡献会存活下来。对于一个正交归一的态，它与自身的重叠给出主导本征对 $C=I$ 和 $\lambda=1$。两个态的重叠中较小的 $\lambda$ 可以被解读为每格点重叠；当然，在热力学极限下这两个态彼此是正交的（重叠 $\lim_{L\rightarrow\infty} \lambda^L = 0$）。这种每格点的热力学重叠（或保真度）可以非常好地用来探测量子相变：只需把参数稍有不同的两个哈密顿量的基态作重叠\cite{McCulloch08}。
 
\section{结论：其他值得关注的课题}
在遍历了这份长长的课题清单——聚焦于基本的算法构件、基态搜索、热力学极限算法以及零温和有限温下丰富的实时与虚时方法——之后，可以看到 DMRG 与基于 MPS 的算法的可能性还远未穷尽。有几个我未曾论及、但想简要提及的课题（同样是一份不完全的清单）是：转移矩阵 DMRG 方法（参考文献见引言部分）、带周期边界条件的 DMRG 与 MPS，以及作为最新进展的连续空间 MPS\cite{Verstraete10}——它表现为相干态的一种漂亮推广，应当能在场论中找到有趣的应用。对于 {\em 周期边界条件}，已有的结果已经相当多，所以这里只做一个简短的概述。在 DMRG 框架内处理 PBC 的做法，是在一条开边界链的格点 1 与 $L$ 之间引入一个长程相互作用（参见例如 \cite{Schmitteckert96,Rapsch99,Meden03a,Meden03b}；然而人们一致发现其精度标度比开边界条件差得多。其深层原因（粗略地说）在于：在环上 A 与 B 之间的界面加倍，纠缠也随之加倍；考虑到纠缠与 MPS 维数之间的指数关系，这意味着资源必须从 $D$ 增至最多 $D^2$，也就是说，为达到相近的精度，算法需要平方的时间（有时被称为 $D^6$ 标度，这是相对于开边界条件的 $D$ 而言的）。物理上恰当的周期边界条件 MPS 拟设由式 (\ref{eq:MPSforPBC}) 给出；所需的 $D$ 与 OBC 大致相同，但在其上重新运行变分基态搜索算法的标度是 $D^5$（因为在格点 1 和 $L$ 上矢量代替矩阵的简化不再出现）\cite{VerstraetePorras04}。同时，从广义本征值问题到标准本征值问题的简化也不再发生，这可能导致病态的条件数。PBC 的 MPS 表示有一个很好的特点：可以生成动量的本征态。对于 $k = n (2\pi/L)$ 以及一个（非平移不变的）MPS $\ket{\psi} = \sum_{{\fat{\sigma}}} \tr (A^{[1]\sigma_1} \ldots A^{[L]\sigma_L}) \ket{\fat{\sigma}}$，下面的态就是动量 $k$ 的平移不变本征态：\cite{Porras06}
\begin{equation}
\ket{\psi_k} = \sum_{n=0}^{L-1} \eul^{\imag k n}  \tr (A^{[1]\sigma_{1+n}} \ldots A^{[L]\sigma_{L+n}}) \ket{\fat{\sigma}} .
\end{equation}
最近已有一些改进 $D^5$ 标度的有趣方案被提出 \cite{Pippan10,Pirvu10}，这是一个仍在持续受到关注的领域。文献 \cite{VerstraeteMurg08} 相当详尽地讨论了这一课题。

我想可以这样作结：虽然 MPS 的基本框架如今已经非常成熟，DMRG 也已成长为强关联量子体系可用的最强大数值方法之一，但即使在一维体系这一发展完善的领域中，本文介绍的许多算法仍有进一步改进的余地，从而把新的应用带入我们的视野。事实上相当令人惊讶的是，对于这里介绍的不少方法（其他方法也一样），人们对它们在真实问题中的细致行为知之甚少，而分析这些行为或许会带来有趣的进一步想法。此外，已完成的应用与可完成的应用之比，对未来的激动人心的研究而言似乎十分有利。

\section{致谢}
我要感谢 Aspen Center for Physics、意大利的里雅斯特的 International Center for Theoretical Physics 以及柏林的 Wissenschaftskolleg（Institute for Advanced Study）的热情接待，本工作的主要部分就是在这些地方完成的。兼具激发思考的讨论与潜心钻研的空间，正是物理学家梦寐以求的环境。当然还要感谢许多与我讨论过这些课题、并使我受益良多的人：Mari-Carmen Banuls, Thomas Barthel, Ignacio Cirac, Andrew Daley, Jan von Delft, Jens Eisert, Adrian Feiguin, Juanjo Garcia-Ripoll, Fabian Heidrich-Meisner, Corinna Kollath, Ian McCulloch, Valentin Murg, Volker Nebendahl, Tomotoshi Nishino, Reinhard Noack, Tobias Osborne, Lode Pollet, Matthias Troyer, Frank Verstraete, Guifre Vidal, Andreas Weichselbaum, Steve White 以及 Tao Xiang，此处仅举其要者。特别感谢 Matt Hastings 向我指出他的一篇重要论文，并感谢 Stefan Depenbrock 对本综述一个接近定稿的版本所作的极为仔细的阅读（所有残留的错误都是在那之后加入的）。

\begin{thebibliography}{999}

\bibitem{Bloch08}
I.~Bloch, J.~Dalibard and W.~Zwerger, Rev.\ Mod.\ Phys.\ {\bf 80}, 885 (2008).

\bibitem{Bethe31}
H.~Bethe, Z.\ Phys.\ A {\bf 71}, 205 (1931).

\bibitem{Lieb68}
E.~H.~Lieb and F.~Y.~Wu, Phys.\ Rev.\ Lett.\ {\bf 20}, 1445 (1968).

\bibitem{Haldane83}
F.~D.~M.~Haldane, Phys.\ Rev.\ Lett.\ {\bf 50}, 1153 (1983). 

\bibitem{White92}
S.~R. White, Phys.\ Rev.\ Lett.\ {\bf 69}, 2863 (1992). 

\bibitem{White93}
S.~R. White, Phys.\ Rev.\ B {\bf 48}, 10345 (1993). 

\bibitem{Schollwoeck05} U.~Schollw\"{o}ck, Rev.\ Mod.\ Phys. {\bf 77}, 259 (2005).

\bibitem{Hallberg06} K.~Hallberg, Adv.\ Phys.\ {\bf 55}, 477 (2006).

\bibitem{Hallberg95}
K.~Hallberg, Phys.\ Rev.\ B {\bf 52}, 9827 (1995).

\bibitem{Ramasesha97}
S.~Ramasesha, S.~K.~Pati, H.~R.~Krishnamurthy, Z.~Shuai and J.~L.~Br{\'e}das, Synth.\ Met.\ {\bf 85}, 1019 (1997).

\bibitem{Kuhner99}
T.~D. K\"{u}hner and S.~R. White, Phys.\ Rev.\ B {\bf 60}, 335 (1999).

\bibitem{Jeckelmann02}
E.~Jeckelmann, Phys.\ Rev.\ B {\bf 66}, 045114 (2002).
 
\bibitem{Nishino95}
T.~Nishino, J. Phys.\ Soc.\ Jpn.\ {\bf 64}, 3598 (1995).

\bibitem{Wang97}
X.~Q. Wang and T.~Xiang, Phys.\ Rev.\ B {\bf 56}, 5061 (1997).

\bibitem{Shibata97}
N.~Shibata, J. Phys.\ Soc.\ Jpn. {\bf 66}, 2221 (1997).

\bibitem{Sirker05}
J.~Sirker and A.~Kl\"{u}mper, Phys.\ Rev.\ B {\bf 71}, 241101 (2005).

\bibitem{Hieida98}
Y.~Hieida, J. Phys.\ Soc.\ Jpn. {\bf 67}, 369 (1998).

\bibitem{Carlon99}
E.~Carlon, M.~Henkel and U.~Schollw\"{o}ck, Eur.\ J. Phys.\ B {\bf 12}, 99 (1999).

\bibitem{Carlon01}
E.~Carlon, M.~Henkel and U.~Schollw\"{o}ck, Phys.\ Rev.\ E {\bf 63}, 036101 (2001).

\bibitem{Henkel01}
M.~Henkel and U.~Schollw\"{o}ck, J.\ Phys.\ A {\bf 34}, 3333 (2001).

\bibitem{White96}
S.~R.~White, Phys.\ Rev.\ Lett. {\bf 77}, 3633 (1996).

\bibitem{White98a} 
S.~R.~White and D.~J.~Scalapino, Phys.\ Rev.\ Lett.\ {\bf 80}, 1272 (1998).

\bibitem{White98b} 
S.~R.~White and D.~J.~Scalapino, Phys.\ Rev.\ Lett.\ {\bf 81}, 3227 (1998).

\bibitem{White00} 
S.~R.~White and D.~J.~Scalapino, Phys.\ Rev.\ B {\bf 61}, 6320 (2000).

\bibitem{White03} 
S.~R.~White and D.~J.~Scalapino, Phys.\ Rev.\ Lett.\ {\bf 91}, 136403 (2003).

\bibitem{Chernyshev03}
A.~L.~Chernyshev, S.~R.~White and A.~H.~Castro-Neto, Phys.\ Rev.\ B {\bf 65}, 214527 (2003).

\bibitem{Chernyshev05}
A.~L.~Chernyshev, A.~H.~Castro-Neto and S.~R.~White, Phys.\ Rev.\ Lett.\ {\bf 94}, 036407 (2005).

\bibitem{White07}
S.~R.~White and A.~L.~Chernyshev, Phys.\ Rev.\ Lett.\ {\bf 99}, 127004 (2007).

\bibitem{Vidal03}
G.~Vidal, J.~I.~Latorre, E.~Rico and A.~Kitaev, Phys.\ Rev.\ Lett.\ {\bf 90}, 227902 (2003).

\bibitem{Latorre04}
J.~I.~Latorre, E.~Rico and G.~Vidal, Quant.\ Inf.\ Comp. {\bf 4}, 48 (2004).

\bibitem{Eisert10}
J.~Eisert, M.~Cramer and M.~B.~Plenio, Rev.\ Mod.\ Phys.\ {\bf 82}, 277 (2010).

\bibitem{Baxter68}
R. J. Baxter, J.\ Math.\ Phys.\ {\bf 9}, 650 (1968).

\bibitem{Affleck87}
I. Affleck, T. Kennedy, E.~H. Lieb and H. Tasaki, Phys.\ Rev.\ Lett. {\bf 59},799 (1987).

\bibitem{Affleck88}
I. Affleck, T. Kennedy, E.~H. Lieb and H. Tasaki, Comm.\ Math.\ Phys. {\bf 115}, 477 (1988).

\bibitem{Fannes89}
M. Fannes, B. Nachtergaele and R.~F. Werner, Europhys.\ Lett. {\bf 10}, 633 (1989).

\bibitem{Fannes92}
M. Fannes, B. Nachtergaele and R.~F. Werner,
Comm.\ Math.\ Phys. {\bf 144}, 443 (1992).

\bibitem{Klumper93}
A. Kl\"{u}mper, A. Schadschneider and J. Zittartz,
Europhys.\ Lett. {\bf 24}, 293 (1993).

\bibitem{Kolezhuk96a}
A.~K. Kolezhuk, R. Roth and U. Schollw\"{o}ck, Phys.\ Rev.\ Lett. {\bf 77}, 5142 (1996).

\bibitem{Kolezhuk96b}
A.~K. Kolezhuk, R. Roth and U. Schollw\"{o}ck, Phys.\ Rev.\ B {\bf 55}, 8928 (1997).

\bibitem{Kolezhuk98}
A.~K. Kolezhuk and H.~J. Mikeska, Phys.\ Rev.\ Lett. {\bf 80}, 2709 (1998).

\bibitem{Kolezhuk97}
A.~K. Kolezhuk, H.~J. Mikeska and S. Yamamoto, Phys.\ Rev.\ B {\bf 55}, R3336 (1997).

\bibitem{Kolezhuk99}
A.~K. Kolezhuk, H.~J. Mikeska, K. Maisinger and U. Schollw\"{o}ck, Phys.\ Rev.\ B {\bf 59}, 13565 (1997).

\bibitem{Ostlund95}
S. Ostlund and S. Rommer, Phys.\ Rev.\ Lett. {\bf 75}, 3537 (1995).

\bibitem{Dukelsky98}
J. Dukelsky. M.A. Mart\'{\i}n-Delgado, T. Nishino and G. Sierra, Europhys.\ Lett. {\bf 43}, 457 (1998).

\bibitem{Takasaki99}
H. Takasaki, T. Hikihara and T. Nishino, J. Phys.\ Soc.\ Jpn. {\bf 68}, 1537 (1999).

\bibitem{White05}
S. R. White, Phys.\ Rev.\ B {\bf 72}, 180403 (2005).

\bibitem{Vidal03a}
G. Vidal, Phys.\ Rev.\ Lett.\ {\bf 91}, 147902 (2003).

\bibitem{Vidal04b}
G. Vidal, Phys.\ Rev.\ Lett.\ {\bf 93}, 040502  (2004).

\bibitem{Daley04}
A.~J. Daley, C. Kollath, U. Schollw{\"o}ck, and G. Vidal, J. Stat. Mech.: 
Theor. Exp.,  P04005  (2004).

\bibitem{WhiteFeiguin04}
S.~R. White and A.~E. Feiguin, Phys. Rev. Lett {\bf 93},  076401  (2004).

\bibitem{VerstraeteRipoll04}
F. Verstraete, J.~J. Garc{\'i}a-Ripoll, and J.~I. Cirac, Phys.\ Rev.\ Lett. {\bf
  93},  207204  (2004).

\bibitem{Prosen09}
T. Prosen and M. Znidaric, J.\ Stat.\ Mech.\ P02035 (2009).

\bibitem{Hartmann09}
M.~J.~Hartmann, J.~Prior, S.~R.~Clark and M.~B.~Plenio, 
Phys.\ Rev.\ Lett. {\bf 102}, 057202 (2009).

\bibitem{Banuls09}
M.-C.~Banuls, M.~B.~Hastings, F.~Verstraete and J.~I.~Cirac, 
Phys.\ Rev.\ Lett. {\bf 102}, 240603 (2009).

\bibitem{Feiguin05a}
A.~E.~Feiguin and S.~R.~White, 
Phys.\ Rev.\ B {\bf 72}, 220401 (2005).

\bibitem{VerstraetePorras04}
F. Verstraete, D. Porras, and J.~I. Cirac, Phys. Rev. Lett {\bf 93},  227205
  (2004).

\bibitem{Pippan10}
P.~Pippan, S.~R. White and H.-G. Evertz, Phys.\ Rev.\ B {\bf 81}, 081103(R) (2010).

\bibitem{Pirvu10} B. Pirvu, F. Verstraete and G. Vidal, arXiV:1005.5195 (2010).

\bibitem{Wilson75}
K.~G. Wilson, Rev.\ Mod.\ Phys.\ {\bf 47}, 773 (1975).

\bibitem{Bulla08}
R. Bulla, T.~A. Costi and T. Pruschke, Rev.\ Mod.\ Phys.\ {\bf 80}, 395 (2008).

\bibitem{Krishnamurthy80}
H.~R. Krishna-murthy, J.~W. Wilkins and K.~G. Wilson, Phys.\ Rev.\ B {\bf 21}, 1003 (1980).

\bibitem{Weichselbaum09}
A. Weichselbaum, F. Verstraete, U. Schollw\"{o}ck, J.I. Cirac and J. von Delft, Phys.\ Rev.\ B {\bf 80}, 165117 (2009).

\bibitem{Vidal07}
G. Vidal, Phys.\ Rev.\ Lett.\ {\bf 98}, 070201 (2007).

\bibitem{Orus08}
R. Orus and G. Vidal, Phys.\ Rev.\ B {\bf 78}, 155117 (2008).

\bibitem{McCulloch08}
I. P. McCulloch, arXiv:0804.2509.

\bibitem{Verstraete10}
F.~Verstraete and J.~I.~Cirac, Phys.\ Rev.\ Lett.\ {\bf 104}, 190405 (2010).

\bibitem{VerstraeteCirac04}
F.~Verstraete and J.~I. Cirac, arXiv:cond-mat/0407066 (2004).

\bibitem{Nishino01}
T.~Nishino, Y.~Hieida, K.~Okunishi, N.~Maeshima, Y.~Akutsu and A.~Gendiar, Prog.\ Theor.\ Phys.\ {\bf 105}, 409 (2001).

\bibitem{VerstraeteMurg08}
F. Verstraete, V. Murg and J. I. Cirac, Adv.\ Phys.\ {\bf 57}, 143 (2008).

\bibitem{McCulloch07}
I.P. McCulloch, J. Stat. Mech.: Theor. Exp. P10014 (2007).

\bibitem{Peschel98}
I. Peschel, K. Hallberg, X. Wang and M. Kaulke (Eds.), {\em Density Matrix
Renormalization: a New Numerical Method}, Lecture Notes in Physics 528, Springer, New York (1998).

\bibitem{VerstraetePirvu08}
B. Pirvu, V. Murg, J.I. Cirac and F. Verstraete, New J.\ Phys.\ {\bf 12}, 025012 (2010).

\bibitem{Crosswhite08}
G. M. Crosswhite and D. Bacon, Phys.\ Rev.\ A {\bf 78}, 012356 (2008).

\bibitem{Frowis10}
F. Fr\"{o}wis, V. Nebendahl and W. D\"{u}r, Phys.\ Rev.\ A {\bf 81}, 062337 (2010).

\bibitem{WhiteNoack92} 
S.~R. White and R.~M. Noack, Phys.\ Rev.\ Lett. {\bf 68}, 3487 (1992).

\bibitem{SierraNishino97} G.~Sierra and T.~Nishino, Nucl.\ Phys.\ B {\bf 495}, 505 (1997).

\bibitem{McCulloch00} I.~P.~McCulloch and M.~Gulacsi, Aust.\ J.\ Phys.\ {\bf 53}, 597 (2000).

\bibitem{McCulloch01} I.~P.~McCulloch and M.~Gulacsi, Philos.\ Mag.\ Lett.\ {\bf 81}, 447 (2001).

\bibitem{McCulloch02} I.~P.~McCulloch and M.~Gulacsi, Europhys.\ Lett.\ {\bf 57}, 852 (2002).

\bibitem{McCulloch08a} I.~P.~McCulloch, R.~Kube, M.~Kurz, A.~Kleine, U.~Schollw\"{o}ck and A.~K.~Kolezhuk, Phys.\ Rev.\ B {\bf 77}, 094404 (2008).

\bibitem{Smerat09} S.~Smerat, I.~P.~McCulloch, H.~Schoeller and U.~Schollw\"{o}ck,
Phys.\ Rev.\ B {\bf 79}, 235107 (2009).

\bibitem{Langer09} S.~Langer, F.~Heidrich-Meisner, J.~Gemmer, I.~P.~McCulloch and U.~Schollw\"{o}ck,
Phys.\ Rev.\ B {\bf 79}, 214409 (2009).

\bibitem{Holzner09} A.~Holzner, I.~P.~McCulloch,  U.~Schollw\"{o}ck, J.~von Delft and F.~Heidrich-Meisner,
Phys.\ Rev.\ B {\bf 80}, 205114 (2009).

\bibitem{Schollwoeck03}
U. Schollw\"{o}ck, S. Chakravarty, J.~O. Fjaerestad, J. Marston and M. Troyer,
Phys.\ Rev.\ Lett. {\bf 90}, 186401 (2003).

\bibitem{Peschel99}
I.~Peschel, M.~Kaulke and \"{O}.~Legeza, Ann.\ Phys.\ (Leipzig) {\bf 8}, 153 (1999).

\bibitem{Okunishi99}
K.~Okunishi, Y.~Hieida and Y.~Akutsu, Phys.\ Rev.\ E {\bf 59}, R6227 (1999).

\bibitem{Chung00}
M.~C.~Chung and I.~Peschel, Phys.\ Rev.\ B {\bf 62}, 4191 (2000).

\bibitem{Chung01}
M.~C.~Chung and I.~Peschel, Phys.\ Rev.\ B {\bf 64}, 064412 (2001).

\bibitem{Chan02}
G.~K.~L.~Chan, P.~W.~Ayers and I.~E.~S.~Croot, J.\ Stat.\ Phys.\ {\bf 109}, 289 (2002). 

\bibitem{Bekenstein73}
J. D. Bekenstein, Phys.\ Rev.\ D {\bf 7}, 2333 (1973).

\bibitem{Srednicki93}
M. Srednicki, Phys.\ Rev.\ Lett. {\bf 71}, 666 (1993).

\bibitem{Callan94}
C. G. Callan and F. Wilczek, Phys.\ Lett.\ B {\bf 333}, 55 (1994).

\bibitem{Plenio05}
M. B. Plenio, J. Eisert, J. Dreissig, and M. Cramer, Phys.\ Rev.\ Lett. {\bf 94}, 
060503 (2005).

\bibitem{Barthel06}
T. Barthel, M.C. Chung and U. Schollw\"{o}ck, Phys.\ Rev.\ A {\bf 74}, 022329 (2006).

\bibitem{Gioev06}
D. Gioev and I. Klich, Phys.\ Rev.\ Lett.\ {\bf 96}, 100503 (2006).

\bibitem{Schuch08}
N.~Schuch, M.~M.~Wolf, F.~Verstraete and J.~I.~Cirac,
Phys.\ Rev.\ Lett.\ {\bf 100}, 030504 (2008).

\bibitem{Verstraete06}
F.~Verstraete and J.~I.~Cirac, Phys.\ Rev.\ B {\bf 73}, 094423 (2006). 

\bibitem{Weichselbaum05}
A.~Weichselbaum, private communication (2005).

\bibitem{Garel96}
U.~Schollw\"{o}ck, Th.~Jolicoeur and Th.~Garel, Phys.\ Rev.\ B {\bf 53}, 3304 (1996).

\bibitem{Stoudenmire10}
E. M. Stoudenmire and S. R. White, New J.\ Phys. {\bf 12}, 055026 (2010).

\bibitem{Crosswhite08a}
G. M. Crosswhite, A. C. Doherty and G. Vidal, Phys.\ Rev.\ B {\bf 78}, 035116 (2008).

\bibitem{Suzuki76} M. Suzuki, Prog.\ Theor.\ Phys.\ {\bf 56}, 1454 (1976).

\bibitem{Suzuki91} M. Suzuki, J. Math.\ Phys.\ {\bf 32}, 400 (1991).

\bibitem{McLachlan95} R. I. McLachlan, SIAM J.\ Sci.\ Comp.\ {\bf 16}, 151 (1995).

\bibitem{Hastings09} M.~Hastings, J.\ Math.\ Phys.\ {\bf 50}, 095207 (2009).

\bibitem{Schmitteckert04}
P. Schmitteckert, Phys.\ Rev.\ B {\bf 70}, 121302 (2004).

\bibitem{Feiguin05}
A.~E. Feiguin and S.~R. White, Phys.\ Rev.\ B {\bf 72}, 020404 (2005).

\bibitem{Gobert05}
D. Gobert, C. Kollath, U. Schollw{\"o}ck, and G. Sch{\"u}tz, Phys. Rev. E {\bf
  71},  036102  (2005).

\bibitem{Calabrese06}
P. Calabrese and J. Cardy, J. Stat.\ Mech.: Theor.\ Exp. P04010 (2006).

\bibitem{Osborne06}
T.~Osborne, Phys.\ Rev.\ Lett.\ {\bf 97}, 157202 (2006).

\bibitem{Kollath05}
C. Kollath, U. Schollw\"{o}ck and W. Zwerger, Phys.\ Rev.\ Lett. {\bf 95}, 176401 (2005).

\bibitem{Trebst06}
S. Trebst, U. Schollw\"{o}ck, M. Troyer and P. Zoller, Phys.\ Rev.\ Lett. {\bf 96}, 250402 (2006).

\bibitem{Cramer08}
M. Cramer, A. Flesch, I. McCulloch, U. Schollw\"{o}ck and J. Eisert, Phys.\ Rev.\ Lett. {\bf 101}, 063001 (2008).

\bibitem{Barthel09a}
T. Barthel, C. Kasztelan, I.P. McCulloch and U. Schollw\"{o}ck, Phys.\ Rev.\ A {\bf 79}, 053627 (2009).

\bibitem{Kollath06}
C. Kollath, A. Iucci, T. Giamarchi, W. Hofstetter and U. Schollw\"{o}ck, Phys.\ Rev.\ Lett. {\bf 97}, 050402 (2006).

\bibitem{Hastings08}
M. Hastings, Phys.\ Rev.\ B {\bf 77}, 144302 (2008).

\bibitem{Pereira08}
R.~G. Pereira, S.~R. White and I. Affleck,
Phys.\ Rev.\ Lett. {\bf 100}, 027206 (2008).

\bibitem{White08}
S.~R. White and I. Affleck,
Phys.\ Rev.\ B {\bf 77}, 134437 (2008).

\bibitem{Barthel09b}
T. Barthel, U. Schollw\"{o}ck and S.~R. White, Phys.\ Rev.\ B {\bf 79}, 245101 (2009).

\bibitem{Caux06}
J.-S. Caux and R. Hagemans, J. Stat.\ Mech.: Theor.\ Exp. P12013 (2006).

\bibitem{White09}
S.~R. White,
Phys.\ Rev.\ Lett. {\bf 102}, 190601 (2009).

\bibitem{Saberi08}
H.~Saberi, A.~Weichselbaum and J.~von Delft, Phys.\ Rev.\ B {\bf 78}, 035124 (2008).

\bibitem{Holzner10}
A.~Holzner, A.~Weichselbaum and J.~von Delft, Phys.\ Rev.\ B {\bf 81}, 125126 (2010).

\bibitem{VidalZwolak04}
M. Zwolak and G. Vidal,  Phys.\ Rev.\ Lett.\ {\bf 93}, 207205 (2004).

\bibitem{Daley09}
A. J. Daley, J. M. Taylor, S. Diehl, M. Baranov and P. Zoller, Phys.\ Rev.\ Lett.\ {\bf 102}, 040402 (2009).

\bibitem{Plenio98}
M. B. Plenio and P. L. Knight, Rev.\ Mod.\ Phys.\ {\bf 70}, 101 (1998).

\bibitem{Weichselbaum07} 
A. Weichselbaum and J. von Delft, Phys.\ Rev.\ Lett.\ {\bf 99}, 076402 (2007).

\bibitem{Peters06} R. Peters, T. Pruschke and F. B. Anders, Phys.\ Rev.\ B {\bf 74}, 245114 (2006).

\bibitem{Rosch03}
A. Rosch, T.A. Costi, J. Paaske and P. W\"{o}lfle, Phys.\ Rev.\ B {\bf 68}, 014430 (2003).

\bibitem{Garst05} 
M. Garst, P. W\"{o}lfle, L. Borda, J. von Delft and L. Glazman, Phys.\ Rev.\ B {\bf 72}, 205125 (2005).

\bibitem{Georges96} A. Georges, G. Kotliar, W. Krauth and M. J. Rozenberg, Rev.\ Mod.\ Phys.\ {\bf 68}, 13 (1996).

\bibitem{Nishimoto04} S. Nishimoto, F. Gebhard and E. Jeckelmann, J.\ Phys.: Cond.\ Mat. {\bf 16}, 7063 (2004).

\bibitem{Garcia04} D. J. Garcia, K. Hallberg and M. J. Rozenberg, Phys.\ Rev.\ Lett.\ {\bf 93}, 246403 (2004). 

\bibitem{Raas04}
C. Raas, G.~S. Uhrig and F.~B. Anders, Phys.\ Rev. B {\bf 69}, 041102(R) (2004).

\bibitem{Raas05}
C. Raas and G.~S. Uhrig, Eur.\ Phys.\ J. B {\bf 45}, 293 (2005).

\bibitem{Schollwoeck98} 
U. Schollw\"{o}ck, Phys.\ Rev.\ B {\bf 58}, 8194 (1998).

\bibitem{Sun02}
L. Q. Sun, J. Zang, S. J. Qin and Y. Lei, Phys.\ Rev.\ B {\bf 65}, 132412 (2002).

\bibitem{NishinoOkunishi95} T.~Nishino and K.~Okunishi, J.\ Phys.\ Soc.\ Jpn.\ {\bf 64}, 4084 (1995).

\bibitem{Hieida97} Y.~Hieida, K.~Okunishi and Y.~Akutsu, Phys.\ Lett.\ A {\bf 233}, 464 (1997).

\bibitem{Ueda06} K.~Ueda, T.~Nishino, K.~Okunishi, Y.~Hieida, R.~Derian and A.~Gendiar, J.\ Phys.\ Soc.\ Jpn.\ {\bf 75}, 014003 (2006).

\bibitem{Ueda08} H.~Ueda, T.~Nishino, K.~Kusakabe, J.\ Phys.\ Soc.\ Jpn.\ {\bf 77}, 114002 (2008).

\bibitem{Ueda10} H.~Ueda, A.~Gendiar and T.~Nishino, J.\ Phys.\ Soc.\ Jpn.\ {\bf 79}, 044001 (2010).

\bibitem{Tagliacozzo08} L. Tagliacozzo, T. R. de Oliveira, S. Iblisdir and J. I. Latorre, Phys.\ Rev.\ B {\bf 78}, 024410 (2008).

\bibitem{Andersson99} M. Andersson, M. Boman and S. \"{O}stlund, Phys.\ Rev.\ B {\bf 59}, 10493 (1999).

\bibitem{Schmitteckert96} P. Schmitteckert and U. Eckern, Phys.\ Rev.\ B {\bf 53}, 15397 (1996).

\bibitem{Rapsch99} S. Rapsch, U. Schollw\"{o}ck and W. Zwerger, Europhys.\ Lett.\ {\bf 46}, 559 (1999).

\bibitem{Meden03a} V. Meden and U. Schollw\"{o}ck, Phys.\ Rev.\ B {\bf 67}, 035106 (2003).

\bibitem{Meden03b} V. Meden and U. Schollw\"{o}ck, Phys.\ Rev.\ B {\bf 67}, 193303 (2003).

\bibitem{Porras06} D. Porras, F. Verstraete and J. I. Cirac, Phys.\ Rev.\ B {\bf 73}, 014410 (2006).

\end{thebibliography}

\end{document}


